% Options for packages loaded elsewhere
% Options for packages loaded elsewhere
\PassOptionsToPackage{unicode}{hyperref}
\PassOptionsToPackage{hyphens}{url}
%
\documentclass[
  letterpaper,
]{book}
\usepackage{xcolor}
\usepackage[margin=2cm]{geometry}
\usepackage{amsmath,amssymb}
\setcounter{secnumdepth}{5}
\usepackage{iftex}
\ifPDFTeX
  \usepackage[T1]{fontenc}
  \usepackage[utf8]{inputenc}
  \usepackage{textcomp} % provide euro and other symbols
\else % if luatex or xetex
  \usepackage{unicode-math} % this also loads fontspec
  \defaultfontfeatures{Scale=MatchLowercase}
  \defaultfontfeatures[\rmfamily]{Ligatures=TeX,Scale=1}
\fi
\usepackage{lmodern}
\ifPDFTeX\else
  % xetex/luatex font selection
\fi
% Use upquote if available, for straight quotes in verbatim environments
\IfFileExists{upquote.sty}{\usepackage{upquote}}{}
\IfFileExists{microtype.sty}{% use microtype if available
  \usepackage[]{microtype}
  \UseMicrotypeSet[protrusion]{basicmath} % disable protrusion for tt fonts
}{}
\makeatletter
\@ifundefined{KOMAClassName}{% if non-KOMA class
  \IfFileExists{parskip.sty}{%
    \usepackage{parskip}
  }{% else
    \setlength{\parindent}{0pt}
    \setlength{\parskip}{6pt plus 2pt minus 1pt}}
}{% if KOMA class
  \KOMAoptions{parskip=half}}
\makeatother
% Make \paragraph and \subparagraph free-standing
\makeatletter
\ifx\paragraph\undefined\else
  \let\oldparagraph\paragraph
  \renewcommand{\paragraph}{
    \@ifstar
      \xxxParagraphStar
      \xxxParagraphNoStar
  }
  \newcommand{\xxxParagraphStar}[1]{\oldparagraph*{#1}\mbox{}}
  \newcommand{\xxxParagraphNoStar}[1]{\oldparagraph{#1}\mbox{}}
\fi
\ifx\subparagraph\undefined\else
  \let\oldsubparagraph\subparagraph
  \renewcommand{\subparagraph}{
    \@ifstar
      \xxxSubParagraphStar
      \xxxSubParagraphNoStar
  }
  \newcommand{\xxxSubParagraphStar}[1]{\oldsubparagraph*{#1}\mbox{}}
  \newcommand{\xxxSubParagraphNoStar}[1]{\oldsubparagraph{#1}\mbox{}}
\fi
\makeatother


\usepackage{longtable,booktabs,array}
\usepackage{calc} % for calculating minipage widths
% Correct order of tables after \paragraph or \subparagraph
\usepackage{etoolbox}
\makeatletter
\patchcmd\longtable{\par}{\if@noskipsec\mbox{}\fi\par}{}{}
\makeatother
% Allow footnotes in longtable head/foot
\IfFileExists{footnotehyper.sty}{\usepackage{footnotehyper}}{\usepackage{footnote}}
\makesavenoteenv{longtable}
\usepackage{graphicx}
\makeatletter
\newsavebox\pandoc@box
\newcommand*\pandocbounded[1]{% scales image to fit in text height/width
  \sbox\pandoc@box{#1}%
  \Gscale@div\@tempa{\textheight}{\dimexpr\ht\pandoc@box+\dp\pandoc@box\relax}%
  \Gscale@div\@tempb{\linewidth}{\wd\pandoc@box}%
  \ifdim\@tempb\p@<\@tempa\p@\let\@tempa\@tempb\fi% select the smaller of both
  \ifdim\@tempa\p@<\p@\scalebox{\@tempa}{\usebox\pandoc@box}%
  \else\usebox{\pandoc@box}%
  \fi%
}
% Set default figure placement to htbp
\def\fps@figure{htbp}
\makeatother





\setlength{\emergencystretch}{3em} % prevent overfull lines

\providecommand{\tightlist}{%
  \setlength{\itemsep}{0pt}\setlength{\parskip}{0pt}}



 


% --- Soporte UTF-8 y español ---
\usepackage{fontspec}
\setmainfont{Times New Roman}   % Fuente principal
\usepackage{polyglossia}
\setdefaultlanguage{spanish}   % Para acentos y ñ

% --- Capítulos sin "Chapter XX:" ---
\usepackage{titlesec}
% 0pt evita que LaTeX genere número basado en el título
\titleformat{\chapter}[hang]{\huge\bfseries}{}{0pt}{}

% --- Centrar figuras automáticamente ---
\usepackage[justification=centering]{caption}
\usepackage{float}
\restylefloat{figure}

% --- Tablas profesionales ---
\usepackage{booktabs}
\usepackage{longtable}

% --- Matemáticas ---
\usepackage{amsmath, amssymb}

% --- Hipervínculos PDF ---
\usepackage{hyperref}icables en PDF
\makeatletter
\@ifpackageloaded{tcolorbox}{}{\usepackage[skins,breakable]{tcolorbox}}
\@ifpackageloaded{fontawesome5}{}{\usepackage{fontawesome5}}
\definecolor{quarto-callout-color}{HTML}{909090}
\definecolor{quarto-callout-note-color}{HTML}{0758E5}
\definecolor{quarto-callout-important-color}{HTML}{CC1914}
\definecolor{quarto-callout-warning-color}{HTML}{EB9113}
\definecolor{quarto-callout-tip-color}{HTML}{00A047}
\definecolor{quarto-callout-caution-color}{HTML}{FC5300}
\definecolor{quarto-callout-color-frame}{HTML}{acacac}
\definecolor{quarto-callout-note-color-frame}{HTML}{4582ec}
\definecolor{quarto-callout-important-color-frame}{HTML}{d9534f}
\definecolor{quarto-callout-warning-color-frame}{HTML}{f0ad4e}
\definecolor{quarto-callout-tip-color-frame}{HTML}{02b875}
\definecolor{quarto-callout-caution-color-frame}{HTML}{fd7e14}
\makeatother
\makeatletter
\@ifpackageloaded{bookmark}{}{\usepackage{bookmark}}
\makeatother
\makeatletter
\@ifpackageloaded{caption}{}{\usepackage{caption}}
\AtBeginDocument{%
\ifdefined\contentsname
  \renewcommand*\contentsname{Table of contents}
\else
  \newcommand\contentsname{Table of contents}
\fi
\ifdefined\listfigurename
  \renewcommand*\listfigurename{List of Figures}
\else
  \newcommand\listfigurename{List of Figures}
\fi
\ifdefined\listtablename
  \renewcommand*\listtablename{List of Tables}
\else
  \newcommand\listtablename{List of Tables}
\fi
\ifdefined\figurename
  \renewcommand*\figurename{Figure}
\else
  \newcommand\figurename{Figure}
\fi
\ifdefined\tablename
  \renewcommand*\tablename{Table}
\else
  \newcommand\tablename{Table}
\fi
}
\@ifpackageloaded{float}{}{\usepackage{float}}
\floatstyle{ruled}
\@ifundefined{c@chapter}{\newfloat{codelisting}{h}{lop}}{\newfloat{codelisting}{h}{lop}[chapter]}
\floatname{codelisting}{Listing}
\newcommand*\listoflistings{\listof{codelisting}{List of Listings}}
\usepackage{amsthm}
\theoremstyle{plain}
\newtheorem{theorem}{Theorem}[chapter]
\theoremstyle{remark}
\AtBeginDocument{\renewcommand*{\proofname}{Proof}}
\newtheorem*{remark}{Remark}
\newtheorem*{solution}{Solution}
\newtheorem{refremark}{Remark}[chapter]
\newtheorem{refsolution}{Solution}[chapter]
\makeatother
\makeatletter
\makeatother
\makeatletter
\@ifpackageloaded{caption}{}{\usepackage{caption}}
\@ifpackageloaded{subcaption}{}{\usepackage{subcaption}}
\makeatother
\usepackage{bookmark}
\IfFileExists{xurl.sty}{\usepackage{xurl}}{} % add URL line breaks if available
\urlstyle{same}
\hypersetup{
  pdftitle={Notas de teoría - Cálculo I},
  hidelinks,
  pdfcreator={LaTeX via pandoc}}


\title{Notas de teoría - Cálculo I}
\author{}
\date{}
\begin{document}
\frontmatter
\maketitle

\renewcommand*\contentsname{Table of contents}
{
\setcounter{tocdepth}{2}
\tableofcontents
}

\mainmatter
\bookmarksetup{startatroot}

\chapter{Cálculo I}\label{cuxe1lculo-i}

\bookmarksetup{startatroot}

\chapter{Presentación}\label{presentaciuxf3n}

Este material ha sido elaborado como apoyo para el cursado de la
asignatura \textbf{Cálculo I}, con el objetivo de facilitar el estudio y
la comprensión de los principales conceptos del análisis matemático.

A lo largo de los distintos capítulos se desarrollan, de manera
progresiva y sistemática, los temas fundamentales del curso: funciones,
límites, continuidad, derivadas e integrales. Cada uno de estos
contenidos constituye una herramienta esencial para el estudio de
fenómenos propios de la matemática, la física, la ingeniería y otras
ciencias aplicadas.

El enfoque adoptado combina el desarrollo teórico con ejemplos resueltos
paso a paso, buscando fortalecer tanto la comprensión conceptual como la
capacidad de resolución de problemas. Se prioriza una presentación
clara, ordenada y gradual, que permita al estudiante construir los
conocimientos progresivamente.

Este libro está pensado como un complemento al trabajo en clase y al
estudio personal, y pretende servir como una guía de consulta permanente
a lo largo del cursado.

Se recomienda al lector trabajar activamente con los ejemplos y
ejercicios, y utilizar este material como apoyo para afianzar los
contenidos vistos en cada unidad.

\bookmarksetup{startatroot}

\chapter{1. Funciones}\label{funciones}

\section{1.1. Definición y ejemplos}\label{seccion_1.1}

\begin{tcolorbox}[enhanced jigsaw, opacityback=0, left=2mm, opacitybacktitle=0.6, toprule=.15mm, bottomrule=.15mm, colbacktitle=quarto-callout-note-color!10!white, titlerule=0mm, bottomtitle=1mm, title=\textcolor{quarto-callout-note-color}{\faInfo}\hspace{0.5em}{Definición (Función)}, toptitle=1mm, colframe=quarto-callout-note-color-frame, colback=white, breakable, rightrule=.15mm, arc=.35mm, leftrule=.75mm, coltitle=black]

Una \textbf{función} \(f\) es una regla que asigna a cada elemento \(x\)
de un conjunto \(D\) exactamente un elemento de un conjunto \(E\),
denotado \(f(x)\). Esta asignación podemos indicarla mediante la
expresión \(f\colon D\rightarrow E\).

\end{tcolorbox}

El conjunto \(D\) se llama \textbf{dominio} de la función. Usualmente
escribiremos \(D=\text{Dom}(f)\) o \(D=D_f\). El \textbf{rango} o
\textbf{conjunto imagen} de \(f\) consiste en todos los valores \(f(x)\)
conforme \(x\) varía en su dominio, es decir

\[\operatorname{Im}(f)=\{y: y=f(x) \text{ para algún }x\in D\}\subseteq E.\]

Un símbolo que representa un número arbitrario en el dominio de \(f\) se
llama \textbf{variable independiente}, mientras que uno en el rango de
\(f\) se llama \textbf{variable dependiente}.

El número \(f(x)\) es el valor de la función \(f\) en \(x\), y lo leemos
``\(f\) de \(x\)''. En general trabajaremos con funciones para las que
\(D\) y \(E\) son subconjuntos de los números reales. Éstas reciben el
nombre de \textbf{funciones escalares}.

Podemos pensar una función como si fuese una máquina: si \(x\) está en
el dominio de \(f\) entonces ingresa en la máquina, que la acepta como
una entrada. La máquina produce como salida \(f(x)\), según la regla que
define a la función.

\begin{center}
\includegraphics[width=0.4\linewidth,height=\textheight,keepaspectratio]{index_files/mediabag/imagenes_cap1/1.1-funcion-maquina.pdf}
\end{center}

Así, podemos concebir al dominio como el conjunto de todas las entradas
posibles \(x\) y a la imagen como el conjunto de todas las salidas
posibles \(f(x)\).

Otra forma de representar una función es mediante un diagrama de flechas
o \textbf{diagrama de Venn}. Cada flecha relaciona un elemento en \(D\)
con su único valor asociado mediante \(f\) en el conjunto \(E\).

\begin{center}
\includegraphics[width=0.4\linewidth,height=\textheight,keepaspectratio]{index_files/mediabag/imagenes_cap1/1.1-funcion-venn.pdf}
\end{center}

El método más común para visualizar una función es su \textbf{gráfico}
(o gráfica). Si \(f\) es una función con dominio \(D\) su gráfico es el
conjunto de pares ordenados \[\text{Gráf}(f)=\{(x,f(x)): x\in D\}.\]

En otras palabras, el gráfico está formado por los pares \((x,y)\) del
plano que verifican que \(y=f(x)\), con \(x\) en el dominio de \(f\).

\begin{center}
\includegraphics[width=0.5\linewidth,height=\textheight,keepaspectratio]{index_files/mediabag/imagenes_cap1/1.1-grafico-funcion.pdf}
\end{center}

El gráfico da un esbozo útil del comportamiento o de la ``historia de
vida'' de una función. En cada \(x\) del dominio, el valor
correspondiente a \(f(x)\) se puede leer como la altura del gráfico
sobre \(x\).

\begin{center}
\includegraphics[width=0.6\linewidth,height=\textheight,keepaspectratio]{index_files/mediabag/imagenes_cap1/1.1-dominio-rango.pdf}
\end{center}

Ejemplo

Dado el siguiente gráfico

\begin{center}
\includegraphics[width=0.4\linewidth,height=\textheight,keepaspectratio]{index_files/mediabag/imagenes_cap1/1.1-Ejemplo1.pdf}
\end{center}

\begin{itemize}
\tightlist
\item
  Encontrar el valor de \(f(1)\) y \(f(2)\).
\item
  ¿Cuáles son el dominio y el rango de \(f\)?
\end{itemize}

\begin{tcolorbox}[enhanced jigsaw, opacityback=0, left=2mm, opacitybacktitle=0.6, toprule=.15mm, bottomrule=.15mm, colbacktitle=quarto-callout-tip-color!10!white, titlerule=0mm, bottomtitle=1mm, title=\textcolor{quarto-callout-tip-color}{\faLightbulb}\hspace{0.5em}{Solución}, toptitle=1mm, colframe=quarto-callout-tip-color-frame, colback=white, breakable, rightrule=.15mm, arc=.35mm, leftrule=.75mm, coltitle=black]

Observando el gráfico de la función podemos ver que \(f(1)=0\) y que
\(f(2)=-1\). El dominio de \(f\) es \(D=[0,3]\) y el rango corresponde
al intervalo \([-1,2]\).

\end{tcolorbox}

Ejemplo

Trazar la gráfica y encontrar el dominio y rango de cada función.

\begin{itemize}
\tightlist
\item
  \(f(x)=2x-1\).
\item
  \(g(x)=x^2\).
\end{itemize}

\begin{tcolorbox}[enhanced jigsaw, opacityback=0, left=2mm, opacitybacktitle=0.6, toprule=.15mm, bottomrule=.15mm, colbacktitle=quarto-callout-tip-color!10!white, titlerule=0mm, bottomtitle=1mm, title=\textcolor{quarto-callout-tip-color}{\faLightbulb}\hspace{0.5em}{Solución}, toptitle=1mm, colframe=quarto-callout-tip-color-frame, colback=white, breakable, rightrule=.15mm, arc=.35mm, leftrule=.75mm, coltitle=black]

Observemos que \(f\) es una función y lineal y \(g\) una cuadrática, por
lo que sus gráficos corresponden a una recta y una parábola,
respectivamente.

\includegraphics[width=0.55\linewidth,height=\textheight,keepaspectratio]{index_files/mediabag/imagenes_cap1/1.1-ejemplo-recta-parabola.pdf}

En relación a la función \(f\), tenemos que
\(D=\text{Im}(f)=\mathbb{R}\). Con respecto a \(g\), tenemos que
\(D=\mathbb{R}\) mientras que
\(\text{Im}(g)=\mathbb{R}_0^+=[0,\infty)\), dado que \(x^2\geq 0\) para
todo número real \(x\).

\end{tcolorbox}

Ejemplo

Un recipiente rectangular para almacenamiento, con su parte superior
abierta, tiene un volumen de 10 \(m^3\). La longitud de su base es el
doble de su ancho. Si el material de la base cuesta 10 dólares por
\(m^2\) y el de los lados 6 dólares por \(m^2\), expresar el costo de
fabricación en función del ancho de la base.

\begin{tcolorbox}[enhanced jigsaw, opacityback=0, left=2mm, opacitybacktitle=0.6, toprule=.15mm, bottomrule=.15mm, colbacktitle=quarto-callout-tip-color!10!white, titlerule=0mm, bottomtitle=1mm, title=\textcolor{quarto-callout-tip-color}{\faLightbulb}\hspace{0.5em}{Solución}, toptitle=1mm, colframe=quarto-callout-tip-color-frame, colback=white, breakable, rightrule=.15mm, arc=.35mm, leftrule=.75mm, coltitle=black]

\includegraphics[width=0.4\linewidth,height=\textheight,keepaspectratio]{index_files/mediabag/imagenes_cap1/1.1-ejemplo-caja.pdf}

Si denotamos con \(w\) al ancho de la base, entonces su longitud es
\(2w\). Si además \(h\) representa la altura de la caja, entonces
podemos calcular el volumen como \[ V= (2w) \cdot w\cdot h=2w^2h.\]

Dado que el volumen es de 10 \(m^3\), tenemos que
\(h=\displaystyle \frac{5}{w^2}\). Con esta relación podemos expresar el
área de la base \(A_b\) de la caja y su área lateral \(A_\ell\) en
función de \(w\).

\[A_b=(2w)\cdot w=2w^2 \quad \text{ y }\quad A_\ell=2\cdot (2w)\cdot h+ 2\cdot w\cdot h=6wh=\frac{30}{w}.\]

Ahora calculamos los costos \(C_b\) y \(C_\ell\) para la base y el
lateral, respectivamente.

\[C_b=10\, \frac{\text{dólares}}{m^2}\cdot (2w^2)\, m^2= 20w^2\, \text{dólares}.\]

\[C_\ell=6\, \frac{\text{dólares}}{m^2}\cdot \frac{30}{w}\, m^2= \frac{180}{w}\, \text{dólares}.\]

Finalmente, el costo total es
\(C=C_b+C_\ell=20w^2+\displaystyle\frac{180}{w}\), expresado en dólares.

\end{tcolorbox}

Ejemplo

Encuentre el dominio de cada función.

\begin{itemize}
\tightlist
\item
  \(f(x)=\sqrt{x+2}\)
\item
  \(g(x)=\displaystyle \frac{1}{x^2-x}\)
\end{itemize}

\begin{tcolorbox}[enhanced jigsaw, opacityback=0, left=2mm, opacitybacktitle=0.6, toprule=.15mm, bottomrule=.15mm, colbacktitle=quarto-callout-tip-color!10!white, titlerule=0mm, bottomtitle=1mm, title=\textcolor{quarto-callout-tip-color}{\faLightbulb}\hspace{0.5em}{Solución}, toptitle=1mm, colframe=quarto-callout-tip-color-frame, colback=white, breakable, rightrule=.15mm, arc=.35mm, leftrule=.75mm, coltitle=black]

Para determinar el dominio de \(f\), observemos que la función está bien
definida siempre y cuando \(x+2\geq 0\), o equivalentemente
\(x\geq -2\). Por lo tanto, \(D=[-2,\infty)\).

En el caso de \(g\), recordemos que la división por cero no está
definida. Ya que \(x^2-x=x(x-1)\), entonces \(x\) no puede tomar el
valor 0 ni el valor 1. Por lo tanto, \(D=\mathbb{R}\setminus \{0,1\}\).

\end{tcolorbox}

\subsection{Relación entre curvas y gráfico de
funciones}\label{relaciuxf3n-entre-curvas-y-gruxe1fico-de-funciones}

La gráfica de una función resulta ser una curva en el plano \(xy\). Pero
dada cualquier curva plana, ¿cómo saber si se trata del gráfico de una
función? Para responder esta pregunta, utilizamos la siguiente
propiedad.

Prueba de la línea vertical

Una curva en el plano \(xy\) es la gráfica de una función de \(x\) si y
sólo si ninguna línea vertical se interseca con la curva más de una vez.

Si cada recta vertical \(x=a\) corta a la curva sólo una vez, en el
punto \((a,b)\), entonces definimos la función \(f\) de modo tal que
\(f(a)=b\).

Sin embargo, si hay algún valor \(a\) para el cual \(x=a\) corta a la
curva en más de un punto, ésta no puede representar a una función pues
no podemos asignar dos valores distintos \(b\) y \(c\) a \(x=a\).

\begin{center}
\includegraphics[width=0.6\linewidth,height=\textheight,keepaspectratio]{index_files/mediabag/imagenes_cap1/1.1-lineavertical.pdf}
\end{center}

Por ejemplo, si consideramos \(x=y^2-2\), ésta no es la gráfica de una
función de \(x\), según la prueba de la línea vertical.

\begin{center}
\includegraphics[width=0.4\linewidth,height=\textheight,keepaspectratio]{index_files/mediabag/imagenes_cap1/1.1-parabola-horizontal.pdf}
\end{center}

Sin embargo, si despejamos la variable \(y\) obtenemos
\[|y|=\sqrt{x+2}\quad\text{ o bien }\quad y=\pm\sqrt{x+2},\]

y cada una de estas ecuaciones se corresponde con las dos mitades de la
parábola de arriba. Cada una de estas curvas obtenidas resulta ser el
gráfico de una función de \(x\).

\begin{center}
\includegraphics[width=0.6\linewidth,height=\textheight,keepaspectratio]{index_files/mediabag/imagenes_cap1/1.1-parabola-mitades.pdf}
\end{center}

\subsection{Funciones seccionalmente
definidas}\label{funciones-seccionalmente-definidas}

Estas funciones se definen con diferentes reglas en diferentes partes de
su dominio.

Ejemplo

Sea \(f\) la función definida por \[
f(x)=
\begin{cases}
1-x & \text{ si } x\leq 1,\\
x^2 & \text{ si } x>1.
\end{cases}
\]

Calcular \(f(0)\), \(f(1)\), \(f(2)\) y trazar la gráfica.

\begin{tcolorbox}[enhanced jigsaw, opacityback=0, left=2mm, opacitybacktitle=0.6, toprule=.15mm, bottomrule=.15mm, colbacktitle=quarto-callout-tip-color!10!white, titlerule=0mm, bottomtitle=1mm, title=\textcolor{quarto-callout-tip-color}{\faLightbulb}\hspace{0.5em}{Solución}, toptitle=1mm, colframe=quarto-callout-tip-color-frame, colback=white, breakable, rightrule=.15mm, arc=.35mm, leftrule=.75mm, coltitle=black]

Para calcular \(f(0)\) observemos que debemos usar la expresión \(1-x\),
ya que \(0\leq 1\). Entonces resulta \(f(0)=1-0=1\). De la misma manera,
\(f(1)=1-1=0\) y \(f(2)=2^2=4\), ya que \(2>1\).

La gráfica de la función resulta

\includegraphics[width=0.4\linewidth,height=\textheight,keepaspectratio]{index_files/mediabag/imagenes_cap1/1.1-ejemplo-funcion-partes.pdf}

\end{tcolorbox}

Un ejemplo muy común de función seccionalmente definida y que
utilizaremos con frecuencia es la función \textbf{valor absoluto}. El
valor absoluto de un número \(x\) se define como

\phantomsection\label{form-valorabs}
\[
|x|=
\begin{cases}
x & \text{si } x\geq 0,\\
-x & \text{si } x<0.
\end{cases}
\]

El valor absoluto de \(x\) puede interpretarse como la
\textbf{distancia} desde \(x\) hasta \(0\) en la recta numérica. De la
definición podemos ver que \(|x|\geq 0\) para todo \(x\).

Ejemplo

Trazar la gráfica de \(y=|x|\).

\begin{tcolorbox}[enhanced jigsaw, opacityback=0, left=2mm, opacitybacktitle=0.6, toprule=.15mm, bottomrule=.15mm, colbacktitle=quarto-callout-tip-color!10!white, titlerule=0mm, bottomtitle=1mm, title=\textcolor{quarto-callout-tip-color}{\faLightbulb}\hspace{0.5em}{Solución}, toptitle=1mm, colframe=quarto-callout-tip-color-frame, colback=white, breakable, rightrule=.15mm, arc=.35mm, leftrule=.75mm, coltitle=black]

Usando la definición de valor absoluto, obtenemos la gráfica al dibujar
la recta \(y=x\) sólamente para \(x\geq 0\) e \(y=-x\) para cuando
\(x<0\).

\includegraphics[width=0.4\linewidth,height=\textheight,keepaspectratio]{index_files/mediabag/imagenes_cap1/1.1-ejemplo-valor-absoluto.pdf}

\end{tcolorbox}

Ejemplo

Encontrar una fórmula explícita para la función dada en el siguiente
gráfico.

\begin{center}
\includegraphics[width=0.4\linewidth,height=\textheight,keepaspectratio]{index_files/mediabag/imagenes_cap1/1.1-ejemplo-funcion-tramos.pdf}
\end{center}

\begin{tcolorbox}[enhanced jigsaw, opacityback=0, left=2mm, opacitybacktitle=0.6, toprule=.15mm, bottomrule=.15mm, colbacktitle=quarto-callout-tip-color!10!white, titlerule=0mm, bottomtitle=1mm, title=\textcolor{quarto-callout-tip-color}{\faLightbulb}\hspace{0.5em}{Solución}, toptitle=1mm, colframe=quarto-callout-tip-color-frame, colback=white, breakable, rightrule=.15mm, arc=.35mm, leftrule=.75mm, coltitle=black]

Notemos que el gráfico de la función está compuesto por tres tramos de
recta. El primero de ellos es el segmento que une los puntos \((0,0)\) y
\((1,1)\). Utilizando la \hyperref[recta-dos-puntos]{ecuación de la
recta que pasa por dos puntos} \((x_1,y_1)\) y \((x_2,y_2)\) obtenemos

\phantomsection\label{recta-dos-puntos}
\begin{quote}
\textbf{📐 Ecuación de la recta por dos puntos}

\[
y-y_1=\frac{y_2-y_1}{x_2-x_1}(x-x_1)
\]
\end{quote}

\[
y-0=\frac{1-0}{1-0}(x-0) \quad \text{ o, equivalentemente }\quad  y=x. 
\]

Repitiendo el mismo proceso para el segundo tramo de recta, que conecta
los puntos \((1,1)\) y \((2,0)\) resulta

\[
y-1=\frac{0-1}{2-1}(x-1) \quad \text{ o, equivalentemente }\quad  y=-x+2. 
\]

Finalmente, en el tercer tramo de recta la función es constante y nula,
con lo cual su ecuación es simplemente \(y=0\).

Resumiendo todo lo realizado, escribimos a \(f\) como una función
definida por partes \[
f(x)=\begin{cases}
x & \text{si } 0\leq x\leq 1,\\
-x+2 & \text{si } 1<x\leq 2,\\
0 & \text{si } x>2.
\end{cases}
\]

\end{tcolorbox}

\subsection{Simetría}\label{simetruxeda}

\begin{tcolorbox}[enhanced jigsaw, opacityback=0, left=2mm, opacitybacktitle=0.6, toprule=.15mm, bottomrule=.15mm, colbacktitle=quarto-callout-note-color!10!white, titlerule=0mm, bottomtitle=1mm, title=\textcolor{quarto-callout-note-color}{\faInfo}\hspace{0.5em}{Definición (Función par)}, toptitle=1mm, colframe=quarto-callout-note-color-frame, colback=white, breakable, rightrule=.15mm, arc=.35mm, leftrule=.75mm, coltitle=black]

Decimos que una función \(f\) es \textbf{par} si satisface
\[f(x)=f(-x)\]

para todo \(x\) en el dominio de \(f\).

\end{tcolorbox}

Geométricamente, las funciones pares son simétricas respecto del eje
\(y\). Si conocemos el gráfico de \(f\) para \(x\geq 0\), basta con
reflejarlo sobre el eje \(y\) para obtener el gráfico completo.

\begin{center}
\includegraphics[width=0.5\linewidth,height=\textheight,keepaspectratio]{index_files/mediabag/imagenes_cap1/1.1-funcion-par.pdf}
\end{center}

Por ejemplo, la función \(f(x)=x^2\) es par, pues
\[f(-x)=(-x)^2=x^2=f(x),\]

para todo \(x\).

\begin{tcolorbox}[enhanced jigsaw, opacityback=0, left=2mm, opacitybacktitle=0.6, toprule=.15mm, bottomrule=.15mm, colbacktitle=quarto-callout-note-color!10!white, titlerule=0mm, bottomtitle=1mm, title=\textcolor{quarto-callout-note-color}{\faInfo}\hspace{0.5em}{Definición (Función impar)}, toptitle=1mm, colframe=quarto-callout-note-color-frame, colback=white, breakable, rightrule=.15mm, arc=.35mm, leftrule=.75mm, coltitle=black]

Decimos que una función \(f\) es \textbf{impar} si satisface
\[f(x)=-f(-x)\]

para todo \(x\) en el dominio de \(f\).

\end{tcolorbox}

La gráfica de las funciones impares son simétricas respecto del origen
de coordenadas. Si conocemos el gráfico de \(f\) para \(x\geq 0\), basta
con rotarlo 180 grados respecto al origen para conocer el gráfico
completo.

\begin{center}
\includegraphics[width=0.5\linewidth,height=\textheight,keepaspectratio]{index_files/mediabag/imagenes_cap1/1.1-funcion-impar.pdf}
\end{center}

Como ejemplo de función impar podemos mencionar \(f(x)=x^3\). En efecto,
\[f(-x)=(-x)^3=-x^3=-f(x),\] para todo \(x\).

Ejemplo

Determinar si cada una de las siguientes funciones es par, impar, o
ninguna.

\begin{itemize}
\tightlist
\item
  \(f(x)=x^3+x\).
\item
  \(g(x)=1-x^4\).
\item
  \(h(x)=2x-x^2\).
\end{itemize}

\begin{tcolorbox}[enhanced jigsaw, opacityback=0, left=2mm, opacitybacktitle=0.6, toprule=.15mm, bottomrule=.15mm, colbacktitle=quarto-callout-tip-color!10!white, titlerule=0mm, bottomtitle=1mm, title=\textcolor{quarto-callout-tip-color}{\faLightbulb}\hspace{0.5em}{Solución}, toptitle=1mm, colframe=quarto-callout-tip-color-frame, colback=white, breakable, rightrule=.15mm, arc=.35mm, leftrule=.75mm, coltitle=black]

Comencemos por la función \(f\). Notemos que
\[f(-x)=(-x)^3+(-x)=-x^3-x=-(x^3+x)=-f(x)\]

para todo \(x\in\mathbb{R}\). Entonces \(f\) es impar.

En el caso de \(g\) obtenemos \[g(-x)=1-(-x)^4=1-x^4=g(x)\] para todo
\(x\in\mathbb{R}\). En consecuencia, \(g\) es par.

Por último \[h(-x)=2(-x)-(-x)^2=-2x-x^2.\] Ésta última expresión no
coincide con \(h(x)\) ni con \(-h(x)\), por lo tanto \(h\) no es par ni
impar.

\end{tcolorbox}

\subsection{Funciones crecientes y
decrecientes}\label{funciones-crecientes-y-decrecientes}

En la siguiente figura vemos que el gráfico de \(f\) sube desde \(A\)
hasta \(B\), desciende de \(B\) a \(C\) y luego vuelve a subir desde
\(C\) hasta \(D\).

\begin{center}
\includegraphics[width=0.6\linewidth,height=\textheight,keepaspectratio]{index_files/mediabag/imagenes_cap1/1.1-creciente-decreciente.pdf}
\end{center}

En este caso decimos que \(f\) está \textbf{creciendo} sobre el
intervalo \([a,b]\), \textbf{decreciendo} sobre \([b,c]\) y luego
\textbf{crece} de nuevo en \([c,d]\).

\begin{tcolorbox}[enhanced jigsaw, opacityback=0, left=2mm, opacitybacktitle=0.6, toprule=.15mm, bottomrule=.15mm, colbacktitle=quarto-callout-note-color!10!white, titlerule=0mm, bottomtitle=1mm, title=\textcolor{quarto-callout-note-color}{\faInfo}\hspace{0.5em}{Definición (Función creciente y decreciente)}, toptitle=1mm, colframe=quarto-callout-note-color-frame, colback=white, breakable, rightrule=.15mm, arc=.35mm, leftrule=.75mm, coltitle=black]

Decimos que la función \(f\) es \textbf{creciente} sobre un intervalo
\(I\) si
\[f(x_1)<f(x_2) \quad\text{ siempre que }\quad x_1<x_2 \quad \text{ en }\quad  I.\]

Análogamente, decimos que \(f\) es \textbf{decreciente} sobre el
intervalo \(I\) si

\[f(x_1)>f(x_2) \quad\text{ siempre que }\quad x_1<x_2 \quad \text{ en }\quad  I.\]

\end{tcolorbox}

Observar que para que una función sea creciente o decreciente en un
intervalo, las correspondientes desigualdades de arriba tienen que valer
\textbf{para todo par de números} \(x_1,x_2\) en el intervalo.

\hyperref[seccion_1.1]{↑ Volver al inicio de la sección}

\section{1.2. Modelos matemáticos: un catálogo de funciones
básicas}\label{seccion_1.2}

\subsection{Modelos lineales}\label{modelos-lineales}

Una \textbf{función lineal} es de la forma

\[
f(x)=mx+b
\]

donde \(m\) es la \textbf{pendiente} y \(b\) la \textbf{ordenada al
origen}.

\begin{center}
\includegraphics[width=0.3\linewidth,height=\textheight,keepaspectratio]{index_files/mediabag/imagenes_cap1/1.2-modelo-lineal.pdf}
\end{center}

El gráfico de estas funciones corresponde a una recta en el plano y
crecen en una proporción constante (es decir, \(y\) se incrementa de
igual forma cuando \(x\) se incrementa a intervalos constantes).

Ejemplo

A medida que el aire seco se mueve hacia arriba, se expande y se enfría.
Supongamos que la función que modela la temperatura en función de la
altura es lineal.

\begin{enumerate}
\def\labelenumi{\arabic{enumi}.}
\tightlist
\item
  Si la temperatura del suelo es de 20°C y a una altura de 1 km es de
  10°C, expresar la temperatura \(T\) (en °C) como una función de la
  altura \(h\) (en kilómetros).
\item
  Trazar la gráfica de la función del inciso (1). ¿Qué representa la
  pendiente?
\item
  ¿Cuál es la temperatura a una altura de 2.5 km?
\end{enumerate}

\begin{tcolorbox}[enhanced jigsaw, opacityback=0, left=2mm, opacitybacktitle=0.6, toprule=.15mm, bottomrule=.15mm, colbacktitle=quarto-callout-tip-color!10!white, titlerule=0mm, bottomtitle=1mm, title=\textcolor{quarto-callout-tip-color}{\faLightbulb}\hspace{0.5em}{Solución}, toptitle=1mm, colframe=quarto-callout-tip-color-frame, colback=white, breakable, rightrule=.15mm, arc=.35mm, leftrule=.75mm, coltitle=black]

\begin{enumerate}
\def\labelenumi{\arabic{enumi}.}
\item
  Queremos expresar la temperatura como una función lineal de la altura.
  Entonces escribimos \(T(h)=mh+b\), donde \(m\) y \(b\) son los
  coeficientes que queremos determinar. Como la temperatura al nivel del
  suelo es de 20°C, tenemos que \(T(0)=20\), es decir

  \[
  20=m\cdot 0+b=b,
  \]

  de donde \(b=20\). Por otra parte \(T(1)=10\), con lo cual

  \[
  10=m\cdot 1+b=m+20,
  \]

  con lo que deducimos que \(m=-10\). Finalmente nuestra función \(T\)
  resulta

  \[
  T(h)=-10h+20
  \]

  y nos da la temperatura expresada en °C de la altura \(h\) dada en
  kilómetros.
\item
  El gráfico de la función \(T\) es el siguiente

  \includegraphics[width=0.5\linewidth,height=\textheight,keepaspectratio]{index_files/mediabag/imagenes_cap1/1.2-ejemplo-temperatura.pdf}

  En este caso la pendiente representa la relación de cambio de la
  temperatura con respecto a la altura. Dado que \(m=-10\), esto
  significa que por cada kilómetro que se ascienda en altura, la
  temperatura desciende 10°C.
\item
  Si \(h=2.5\), entonces \(T(2.5)=-10(2.5)+20=-5\).
\end{enumerate}

\end{tcolorbox}

\subsection{Polinomios}\label{polinomios}

A una función \(P\) se la llama \textbf{polinomio} si es de la forma

\[p(x)=a_nx^n+a_{n-1}x^{n-1}+\dots+a_2x^2+a_1x+a_0=\sum_{i=0}^n a_ix^i,\]

donde \(n\) es un entero no negativo y los números \(a_i\) son
constantes llamadas \textbf{coeficientes} del polinomio, para
\(1\leq i\leq n\). A \(a_n\) se lo llama \textbf{coeficiente principal}
y a \(a_0\) \textbf{término independiente} de \(p\).

El dominio de cualquier polinomio es \(\mathbb{R}\). Si \(a_n\neq 0\),
decimos que el \textbf{grado} de \(p\) es \(n\). Por ejemplo,
\[p(x)=2x^6-x^4+\frac{2}{3}x^3+\sqrt{2}\] es un polinomio de grado 6.

\begin{itemize}
\item
  Los polinomios de grado 1 son de la forma \(p(x)=mx+b\), es decir, son
  \textbf{funciones lineales}.
\item
  Si \(p\) tiene grado 2, entonces \[p(x)=ax^2+bx+c\] y se llama
  \textbf{función cuadrática}. Su gráfica es una parábola que abre hacia
  arriba cuando \(a>0\) y hacia abajo cuando \(a<0\).
\end{itemize}

\begin{center}
\includegraphics[width=0.6\linewidth,height=\textheight,keepaspectratio]{index_files/mediabag/imagenes_cap1/1.2-modelo-cuadratico.pdf}
\end{center}

\begin{itemize}
\tightlist
\item
  Un polinomio de grado 3 tiene la forma
  \[p(x)=ax^3+bx^2+cx+d,\quad a\neq 0\] y recibe el nombre de
  \textbf{función cúbica}. Algunos ejemplos son
\end{itemize}

\begin{center}
\includegraphics[width=0.6\linewidth,height=\textheight,keepaspectratio]{index_files/mediabag/imagenes_cap1/1.2-modelo-cubico.pdf}
\end{center}

\subsection{Funciones de potencia}\label{funciones-de-potencia}

Una función de la forma

\[f(x)=x^a\]

donde \(a\) es una constante, recibe el nombre de \textbf{función
potencia}. Consideramos diferentes casos:

\begin{itemize}
\item
  \(a=n\), con \(n\) entero positivo. Observemos algunos ejemplos, que
  incluyen algunos modelos ya estudiados.

  \begin{center}
  \includegraphics[width=0.6\linewidth,height=\textheight,keepaspectratio]{index_files/mediabag/imagenes_cap1/1.2-modelo-potencia-natural.pdf}
  \end{center}

  \begin{center}
  \includegraphics[width=0.6\linewidth,height=\textheight,keepaspectratio]{index_files/mediabag/imagenes_cap1/1.2-modelo-potencia-natural2.pdf}
  \end{center}

  La forma general de la gráfica de \(f(x)=x^n\) depende de si \(n\) es
  par o impar, obteniendo respectivamente una gráfica simétrica respecto
  del eje \(y\) o del origen de coordenadas.

  Conforme aumenta \(n\), la gráfica se hace más aplastada cerca del
  origen y se vuelve más pronunciada cuando \(|x|\geq 1\).

  \begin{center}
  \includegraphics[width=0.7\linewidth,height=\textheight,keepaspectratio]{index_files/mediabag/imagenes_cap1/1.2-comparacion-potencias.pdf}
  \end{center}
\item
  \(a=1/n\), donde \(n\) es un entero positivo. La función
  \(f(x)=x^{1/n}=\sqrt[n]{x}\) se denomina \textbf{función raíz}. Cuando
  \(n\) es par, el dominio de la función es \([0,\infty)\). Para \(n\)
  impar, el dominio es \(\mathbb{R}\) y la gráfica es simétrica respecto
  del origen de coordenadas.
\end{itemize}

\begin{center}
\includegraphics[width=0.6\linewidth,height=\textheight,keepaspectratio]{index_files/mediabag/imagenes_cap1/1.2-funciones-raiz.pdf}
\end{center}

\begin{itemize}
\tightlist
\item
  \(a=-1\). La función \(\displaystyle f(x)=x^{-1}= \frac{1}{x}\) se
  denomina \textbf{función recíproca}. Se trata de una hipérbola que
  tiene a los ejes coordenados como asíntotas. El dominio de esta
  función es \(D=\mathbb{R}\setminus \{0\}\).
\end{itemize}

\begin{center}
\includegraphics[width=0.4\linewidth,height=\textheight,keepaspectratio]{index_files/mediabag/imagenes_cap1/1.2-funcion-reciproca.pdf}
\end{center}

\subsection{Funciones racionales}\label{funciones-racionales}

Una función \textbf{racional} \(f\) es un cociente de dos polinomios

\[
f(x)=\frac{P(x)}{Q(x)}.
\]

Dado que no es posible dividir por cero, el dominio de \(f\) consiste de
todos los valores de \(x\) tales que \(Q(x)\neq 0\).

Por ejemplo, las funciones \[
f(x)=\frac{1}{x},\quad g(x)=\frac{2x^4-x^2+1}{x^2-4}
\]

son racionales.

\subsection{Funciones algebraicas}\label{funciones-algebraicas}

Son aquellas funciones que pueden construirse utilizando
\textbf{operaciones algebraicas} (suma, resta, multiplicación, división,
raíces). En particular, cualquier función racional es también
algebraica. Otros ejemplos que podemos mencionar son \[
f(x)=\sqrt{x^2+1}, \quad g(x)=\frac{x^4-16x^2}{x+\sqrt{x}}+(x-2)\sqrt[3]{x+1}.
\]

\subsection{Funciones
trigonométricas}\label{funciones-trigonomuxe9tricas}

Por convención, siempre utilizaremos el argumento de estas funciones
medido en \textbf{radianes} (recordemos que \(180^{\circ}=\pi\)
radianes). Veamos cómo lucen los gráficos para las funciones
\textbf{seno} y \textbf{coseno}.

\begin{center}
\includegraphics[width=0.6\linewidth,height=\textheight,keepaspectratio]{index_files/mediabag/imagenes_cap1/1.2-seno-coseno.pdf}
\end{center}

El dominio de estas funciones es \(D=\mathbb{R}\) y el rango es el
intervalo \([-1,1]\). Es decir, tenemos que \[
-1\leq \sin x\leq 1 \quad \text { y }\quad  -1\leq \cos x\leq 1
\] o, en términos de valor absoluto \[
|\sin x|\leq 1 \quad \text{ y }\quad |\cos x|\leq 1.
\]

Además \[
\sin x = 0 \quad\text{ si y sólo si }\quad x=n\pi, \,\,n\in\mathbb{Z}
\]

y

\[
\cos x = 0 \quad\text{ si y sólo si }\quad x=(2n+1)\pi/2, \,\,n\in\mathbb{Z}.
\]

Estas dos funciones son periódicas de período \(2\pi\). Esto significa
que \[
\sin(x+2\pi)=\sin x\quad \text{ y }\quad \cos(x+2\pi)=\cos x
\] para todo valor de \(x\).

La función \textbf{tangente} está dada por
\(\displaystyle\tan x=\frac{\sin x}{\cos x}.\)

\begin{center}
\includegraphics[width=0.6\linewidth,height=\textheight,keepaspectratio]{index_files/mediabag/imagenes_cap1/1.2-tangente.pdf}
\end{center}

Observar que esta función no está definida cuando \(\cos x=0\), es
decir, cuando \(x=(2n+1)\pi/2\), con \(n\in\mathbb{Z}\). Su rango es
todo \(\mathbb{R}\), es decir, toma todos los valores reales.

La función tangente es periódica de período \(\pi\). Esto quiere decir
que \[
\tan(x+\pi)=\tan x
\] para todo \(x\) en el dominio.

\subsection{Funciones exponenciales}\label{funciones-exponenciales}

Son de la forma

\[
f(x)=a^x
\]

donde \(a\) es un número \textbf{positivo}. A continuación vemos las
gráficas para \(a=2\) y \(a=1/2\), respectivamente.

\begin{center}
\includegraphics[width=0.6\linewidth,height=\textheight,keepaspectratio]{index_files/mediabag/imagenes_cap1/1.2-exponencial.pdf}
\end{center}

En ambos casos el dominio es \(\mathbb{R}\) y el rango el intervalo
\((0,\infty)\). Observemos además que estas funciones son crecientes
cuando \(a>1\) y decrecientes cuando \(0<a<1\).

\subsection{Funciones logarítmicas}\label{funciones-logaruxedtmicas}

Son de la forma

\[
f(x)=\log_a x
\]

donde \(a>0\) y \(a\neq 1\). Estas son las inversas de las funciones
exponenciales. Su dominio es \((0,\infty)\) y su rango es
\(\mathbb{R}\).

\begin{center}
\includegraphics[width=0.6\linewidth,height=\textheight,keepaspectratio]{index_files/mediabag/imagenes_cap1/1.2-logaritmos.pdf}
\end{center}

\subsection{Funciones trascendentes}\label{funciones-trascendentes}

Estas funciones no son algebraicas. Incluyen a las trigonométricas y sus
inversas, exponenciales, logarítmicas, entre otras.

Ejemplo

Clasificar las siguientes funciones como algunos de los tipos analizados
anteriormente.

\begin{itemize}
\tightlist
\item
  \(\displaystyle f(x)=5^x\phantom{\frac{x^5}{x^3}}\)
\item
  \(g(x)=x^5\)
\end{itemize}

\begin{itemize}
\tightlist
\item
  \(\displaystyle h(x)=\frac{1+x}{1-\sqrt{x}}\)
\item
  \(u(t)=1-t+5t^4\)
\end{itemize}

\begin{tcolorbox}[enhanced jigsaw, opacityback=0, left=2mm, opacitybacktitle=0.6, toprule=.15mm, bottomrule=.15mm, colbacktitle=quarto-callout-tip-color!10!white, titlerule=0mm, bottomtitle=1mm, title=\textcolor{quarto-callout-tip-color}{\faLightbulb}\hspace{0.5em}{Solución}, toptitle=1mm, colframe=quarto-callout-tip-color-frame, colback=white, breakable, rightrule=.15mm, arc=.35mm, leftrule=.75mm, coltitle=black]

Las funciones \(g\) y \(u\) son polinómicas. Por otro lado, \(f\) es una
función exponencial de base \(a=5\) y \(h\) es una función algebraica.

\end{tcolorbox}

\hyperref[seccion_1.2]{↑ Volver al inicio de la sección}

\section{1.3: Funciones nuevas a partir de funciones
antiguas}\label{seccion_1.3}

En esta sección veremos cómo obtener el gráfico de ciertas funciones más
generales que las aprendidas en la sección previa. Veamos algunos tipos
de transformaciones que pueden aplicarse a la gráfica de una función y
qué efecto tienen.

\subsection{Traslaciones}\label{traslaciones}

Consideremos \(c>0\) y partimos del gráfico de \(y=f(x)\). Entonces

\begin{longtable}[]{@{}ll@{}}
\toprule\noalign{}
Para obtener el gráfico de & Se desplaza el de \(y=f(x)\) \\
\midrule\noalign{}
\endhead
\bottomrule\noalign{}
\endlastfoot
\(y=f(x)+c\) & \(c\) unidades hacia arriba \\
\(y=f(x)-c\) & \(c\) unidades hacia abajo \\
\(y=f(x-c)\) & \(c\) unidades a la derecha \\
\(y=f(x+c)\) & \(c\) unidades a la izquierda \\
\end{longtable}

\begin{center}
\includegraphics[width=0.6\linewidth,height=\textheight,keepaspectratio]{index_files/mediabag/imagenes_cap1/1.3-traslaciones.pdf}
\end{center}

\subsection{Alargamientos y reflexione verticales y
horizontales}\label{alargamientos-y-reflexione-verticales-y-horizontales}

Consideremos ahora \(c>1\) y partimos del gráfico de \(y=f(x)\).
Entonces

\begin{longtable}[]{@{}ll@{}}
\toprule\noalign{}
Para obtener el gráfico de & El de \(y=f(x)\) debe \\
\midrule\noalign{}
\endhead
\bottomrule\noalign{}
\endlastfoot
\(y=cf(x)\) & alargarse verticalmente en un factor de \(c\) \\
\(y=(1/c)f(x)\) & comprimirse verticalmente en un factor de \(c\) \\
\(y=f(cx)\) & \(c\) comprimirse horizontalmente en un factor \(c\) \\
\(y=f(x/c)\) & \(c\) alargarse horizontalmente en un factor \(c\) \\
\(y=-f(x)\) & \(c\) reflejarse respecto del eje \(x\) \\
\(y=f(-x)\) & \(c\) reflejarse respecto del eje \(y\) \\
\end{longtable}

\begin{center}
\includegraphics[width=0.6\linewidth,height=\textheight,keepaspectratio]{index_files/mediabag/imagenes_cap1/1.3-alargamientos-compresiones.pdf}
\end{center}

Ilustremos estas transformaciones con la función \(y=\cos x\). Dibujamos
la gráfica junto a la de las funciones \(y=c\cos x\), para \(c=2\) y
\(c=1/2\).

\begin{center}
\includegraphics[width=0.6\linewidth,height=\textheight,keepaspectratio]{index_files/mediabag/imagenes_cap1/1.3-transformaciones-coseno.pdf}
\end{center}

Dibujamos ahora \(y=\cos(cx)\), también para \(c=2\) y \(c=1/2\).

\begin{center}
\includegraphics[width=0.6\linewidth,height=\textheight,keepaspectratio]{index_files/mediabag/imagenes_cap1/1.3-transformaciones-coseno2.pdf}
\end{center}

Ejemplo

Partiendo del gráfico de \(y=f(x)=\sqrt{x}\) usar transformaciones para
dibujar el de las funciones

\begin{itemize}
\item
  \(f_1(x)=\sqrt{x}-2\)
\item
  \(f_2(x)=\sqrt{x-2}\)
\item
  \(f_3(x)=-\sqrt{x}\)
\item
  \(f_4(x)=2\sqrt{x}\)
\item
  \(f_5(x)=\sqrt{-x}\)
\end{itemize}

\begin{tcolorbox}[enhanced jigsaw, opacityback=0, left=2mm, opacitybacktitle=0.6, toprule=.15mm, bottomrule=.15mm, colbacktitle=quarto-callout-tip-color!10!white, titlerule=0mm, bottomtitle=1mm, title=\textcolor{quarto-callout-tip-color}{\faLightbulb}\hspace{0.5em}{Solución}, toptitle=1mm, colframe=quarto-callout-tip-color-frame, colback=white, breakable, rightrule=.15mm, arc=.35mm, leftrule=.75mm, coltitle=black]

El gráfico de \(f_1\) se obtiene trasladando el de \(f\) dos unidades
verticalmente hacia abajo.

\includegraphics[width=0.4\linewidth,height=\textheight,keepaspectratio]{index_files/mediabag/imagenes_cap1/1.3-ejemplo-traslacion-raiz.pdf}

Si trasladamos el gráfico de \(f\) dos unidades hacia la derecha,
obtenemos el gráfico de \(f_2\).

\includegraphics[width=0.4\linewidth,height=\textheight,keepaspectratio]{index_files/mediabag/imagenes_cap1/1.3-ejemplo-traslacion-raiz2.pdf}

Para el caso de \(f_3\) basta con reflejar el gráfico de \(f\) con
respecto al eje \(x\).

\includegraphics[width=0.4\linewidth,height=\textheight,keepaspectratio]{index_files/mediabag/imagenes_cap1/1.3-ejemplo-traslacion-raiz3.pdf}

El gráfico de \(f_4\) resulta de alargar el de \(f\) de forma vertical
en un factor de 2.

\includegraphics[width=0.35\linewidth,height=\textheight,keepaspectratio]{index_files/mediabag/imagenes_cap1/1.3-ejemplo-traslacion-raiz4.pdf}

Por último, reflejando el gráfico de \(f\) respecto al eje \(y\)
obtenemos el de \(f_5\).

\includegraphics[width=0.5\linewidth,height=\textheight,keepaspectratio]{index_files/mediabag/imagenes_cap1/1.3-ejemplo-traslacion-raiz5.pdf}

\end{tcolorbox}

Ejemplo

Esbozar el gráfico de la función \(f(x)=x^2+6x+10\).

\begin{tcolorbox}[enhanced jigsaw, opacityback=0, left=2mm, opacitybacktitle=0.6, toprule=.15mm, bottomrule=.15mm, colbacktitle=quarto-callout-tip-color!10!white, titlerule=0mm, bottomtitle=1mm, title=\textcolor{quarto-callout-tip-color}{\faLightbulb}\hspace{0.5em}{Solución}, toptitle=1mm, colframe=quarto-callout-tip-color-frame, colback=white, breakable, rightrule=.15mm, arc=.35mm, leftrule=.75mm, coltitle=black]

Podemos graficar sencillamente a \(f\) si completamos cuadrados. En
efecto, observemos que

\[
x^2+6x+10=(x^2+6x+9)-9+10=(x+3)^2+1.
\]

Ahora, partiendo del gráfico de \(y=x^2\) podemos obtener el de \(f\)
realizando un desplazamiento horizontal tres unidades para la izquierda,
seguido de un desplazamiento vertical de una unidad hacia arriba.

\includegraphics[width=0.5\linewidth,height=\textheight,keepaspectratio]{index_files/mediabag/imagenes_cap1/1.3-ejemplo-parabola.pdf}

\end{tcolorbox}

Ejemplo

Trazar el gráfico de las funciones

\begin{itemize}
\tightlist
\item
  \(y=\operatorname{sen}(2x)\)
\item
  \(y=1-\operatorname{sen} x\)
\end{itemize}

\begin{tcolorbox}[enhanced jigsaw, opacityback=0, left=2mm, opacitybacktitle=0.6, toprule=.15mm, bottomrule=.15mm, colbacktitle=quarto-callout-tip-color!10!white, titlerule=0mm, bottomtitle=1mm, title=\textcolor{quarto-callout-tip-color}{\faLightbulb}\hspace{0.5em}{Solución}, toptitle=1mm, colframe=quarto-callout-tip-color-frame, colback=white, breakable, rightrule=.15mm, arc=.35mm, leftrule=.75mm, coltitle=black]

En este caso tomamos de base el gráfico de la función seno, es decir,
\(y=\operatorname{sen} x\). Para la primera, debemos realizar una
compresión horizontal en un factor de 2 y el gráfico resulta

\includegraphics[width=0.6\linewidth,height=\textheight,keepaspectratio]{index_files/mediabag/imagenes_cap1/1.3-ejemplo-seno1.pdf}

Para la segunda función debemos realizar dos movimientos: primero
reflejamos \(y=\operatorname{sen} x\) con respecto al eje \(x\) y luego
trasladamos el gráfico una unidad verticalmente hacia arriba.

\includegraphics[width=0.6\linewidth,height=\textheight,keepaspectratio]{index_files/mediabag/imagenes_cap1/1.3-ejemplo-seno2.pdf}

\end{tcolorbox}

Ejemplo

Dibujar la función \(y=|x^2-1|\).

\begin{tcolorbox}[enhanced jigsaw, opacityback=0, left=2mm, opacitybacktitle=0.6, toprule=.15mm, bottomrule=.15mm, colbacktitle=quarto-callout-tip-color!10!white, titlerule=0mm, bottomtitle=1mm, title=\textcolor{quarto-callout-tip-color}{\faLightbulb}\hspace{0.5em}{Solución}, toptitle=1mm, colframe=quarto-callout-tip-color-frame, colback=white, breakable, rightrule=.15mm, arc=.35mm, leftrule=.75mm, coltitle=black]

Si aplicamos directamente la \hyperref[form-valorabs]{definición de
valor absoluto} , obtenemos \[
|x^2-1|=\begin{cases}
x^2-1 & \text{si } x^2-1\geq  0,\\
-(x^2-1) & \text{si } x^2-1<0.
\end{cases}
\]

Observemos a continuación que
\[ x^2-1\geq 0 \quad \Longleftrightarrow \quad x^2\geq 1 \quad \Longleftrightarrow \quad |x|\geq 1 \quad \Longleftrightarrow \quad  x\geq 1 \quad \text{ o }\quad x\leq -1,
\]

y también

\[ x^2-1< 0 \quad \Longleftrightarrow \quad x^2< 1 \quad \Longleftrightarrow \quad |x|< 1 \quad \Longleftrightarrow \quad  -1<x<1.
\]

Con esto, podemos reescribir la función de la siguiente manera \[
|x^2-1|=\begin{cases}
x^2-1 & \text{si } \quad x\geq 1 \quad \text{ o }\quad x\leq -1,\\
-x^2+1 & \text{si } \quad  -1<x<1.
\end{cases}
\]

Por lo tanto, el gráfico es

\includegraphics[width=0.4\linewidth,height=\textheight,keepaspectratio]{index_files/mediabag/imagenes_cap1/1.3-valor-absoluto-parabola.pdf}

\end{tcolorbox}

\begin{tcolorbox}[enhanced jigsaw, opacityback=0, left=2mm, opacitybacktitle=0.6, toprule=.15mm, bottomrule=.15mm, colbacktitle=quarto-callout-important-color!10!white, titlerule=0mm, bottomtitle=1mm, title=\textcolor{quarto-callout-important-color}{\faExclamation}\hspace{0.5em}{Observación}, toptitle=1mm, colframe=quarto-callout-important-color-frame, colback=white, breakable, rightrule=.15mm, arc=.35mm, leftrule=.75mm, coltitle=black]

El procedimiento utilizado en el ejemplo anterior se puede generalizar
de forma sencilla a cualquier función. En efecto, consideremos una
función \(f\) y supongamos que su gráfico está dado por

\includegraphics[width=0.4\linewidth,height=\textheight,keepaspectratio]{index_files/mediabag/imagenes_cap1/1.3-valor-absoluto-funcion1.pdf}

Si queremos dibujar \(y=|f(x)|\) podemos usar la
\hyperref[form-valorabs]{definición de valor absoluto} para escribir \[
|f(x)|=\begin{cases}
f(x) & \text{si } f(x)\geq 0,\\
-f(x) & \text{si }   f(x)<0.
\end{cases}
\]

Entonces el gráfico de \(y=|f(x)|\) coincide con el de \(f\) cuando los
valores de la función son no negativos (esto es, cuando la curva está
sobre o por encima del eje \(x\)) y en caso contrario se refleja
respecto de este eje. Así, resulta

\includegraphics[width=0.4\linewidth,height=\textheight,keepaspectratio]{index_files/mediabag/imagenes_cap1/1.3-valor-absoluto-funcion2.pdf}

\end{tcolorbox}

\subsection{Combinación de funciones}\label{combinaciuxf3n-de-funciones}

Dadas dos funciones \(f\) y \(g\), podemos formar otras nuevas
realizando operaciones entre ellas.

\begin{tcolorbox}[enhanced jigsaw, opacityback=0, left=2mm, opacitybacktitle=0.6, toprule=.15mm, bottomrule=.15mm, colbacktitle=quarto-callout-note-color!10!white, titlerule=0mm, bottomtitle=1mm, title=\textcolor{quarto-callout-note-color}{\faInfo}\hspace{0.5em}{Definición (Suma y resta de funciones)}, toptitle=1mm, colframe=quarto-callout-note-color-frame, colback=white, breakable, rightrule=.15mm, arc=.35mm, leftrule=.75mm, coltitle=black]

La \textbf{suma} y \textbf{resta} se definen como \[
(f+g)(x)=f(x)+g(x)\quad \text{ y }\quad (f-g)(x)=f(x)-g(x).
\]

Si \(A\) y \(B\) son los respectivos dominios de \(f\) y de \(g\), el
dominio de las funciones de arriba es \(D=A\cap B\).

\end{tcolorbox}

Por ejemplo, si \(f(x)=\sqrt{x}\) y \(g(x)=\sqrt{2-x}\), tenemos que
\(D_f=[0,\infty)\) y \(D_g=(-\infty,2]\). De esta manera, \[
(f+g)(x)=f(x)+g(x)=\sqrt{x}+\sqrt{2-x}
\] tiene como dominio \(D_f\cap D_g=[0,2]\).

\begin{tcolorbox}[enhanced jigsaw, opacityback=0, left=2mm, opacitybacktitle=0.6, toprule=.15mm, bottomrule=.15mm, colbacktitle=quarto-callout-note-color!10!white, titlerule=0mm, bottomtitle=1mm, title=\textcolor{quarto-callout-note-color}{\faInfo}\hspace{0.5em}{Definición (Producto y cociente de funciones)}, toptitle=1mm, colframe=quarto-callout-note-color-frame, colback=white, breakable, rightrule=.15mm, arc=.35mm, leftrule=.75mm, coltitle=black]

El \textbf{producto} y el \textbf{cociente} se definen de manera análoga
\[
(fg)(x)=f(x)g(x)\quad\text{ y }\quad \left(\frac{f}{g}\right)(x)=\frac{f(x)}{g(x)}.
\]

El dominio de \(fg\) es \(A\cap B\) y el de \(f/g\) es
\((A\cap B)\setminus\{x\in B: g(x)=0\}\).

\end{tcolorbox}

Por ejemplo, si \(f(x)=x^2\) y \(g(x)=x-1\), entonces la función
cociente es \[
\left(\frac{f}{g}\right)(x)=\frac{f(x)}{g(x)}=\frac{x^2}{x-1},
\] y su dominio es \(\mathbb{R}\setminus \{1\}\).

Existe otra manera de combinar dos funciones. Dadas \(f\) y \(g\), si
\(x\) está en el dominio de \(g\) y su imagen por \(g\), \(g(x)\) está
en el dominio de \(f\), entonces podemos calcular \(f(g(x))\). El
resultado es una nueva función \(h(x)=f(g(x))\) llamada
\textbf{composición} de \(f\) y \(g\), y denotada por \(f\circ g\).

\begin{center}
\includegraphics[width=0.6\linewidth,height=\textheight,keepaspectratio]{index_files/mediabag/imagenes_cap1/1.3-esquema-composicion.pdf}
\end{center}

El dominio de \(f\circ g\) es el conjunto de todos los \(x\) en el
dominio de \(g\) tales que \(g(x)\) está en el dominio de \(f\).

Ejemplo

Si \(f(x)=x^2\) y \(g(x)=x-3\), hallar \(f\circ g\) y \(g\circ f\).

\begin{tcolorbox}[enhanced jigsaw, opacityback=0, left=2mm, opacitybacktitle=0.6, toprule=.15mm, bottomrule=.15mm, colbacktitle=quarto-callout-tip-color!10!white, titlerule=0mm, bottomtitle=1mm, title=\textcolor{quarto-callout-tip-color}{\faLightbulb}\hspace{0.5em}{Solución}, toptitle=1mm, colframe=quarto-callout-tip-color-frame, colback=white, breakable, rightrule=.15mm, arc=.35mm, leftrule=.75mm, coltitle=black]

Por un lado tenemos que \[
(f\circ g)(x)=f(g(x))=f(x-3)=(x-3)^2
\]

mientras que

\[
(g\circ f)(x)=g(f(x))=g\left(x^2\right)=x^2-3.
\]

\end{tcolorbox}

\begin{tcolorbox}[enhanced jigsaw, opacityback=0, left=2mm, opacitybacktitle=0.6, toprule=.15mm, bottomrule=.15mm, colbacktitle=quarto-callout-important-color!10!white, titlerule=0mm, bottomtitle=1mm, title=\textcolor{quarto-callout-important-color}{\faExclamation}\hspace{0.5em}{Importante}, toptitle=1mm, colframe=quarto-callout-important-color-frame, colback=white, breakable, rightrule=.15mm, arc=.35mm, leftrule=.75mm, coltitle=black]

En general no es cierto que \(f\circ g=g\circ f\). La notación
\(f\circ g\) indica que primero se aplica \(g\) y luego \(f\) (y se lee
``\(g\) compuesta con \(f\)'').

\end{tcolorbox}

Ejemplo

Sean \(f(x)=\sqrt{x}\) y \(g(x)=\sqrt{2-x}\). Hallar las siguientes
funciones y sus respectivos dominios

\begin{itemize}
\tightlist
\item
  \(f\circ g\).
\item
  \(g\circ f\).
\end{itemize}

\begin{itemize}
\tightlist
\item
  \(f\circ f\).
\item
  \(g\circ g\).
\end{itemize}

\begin{tcolorbox}[enhanced jigsaw, opacityback=0, left=2mm, opacitybacktitle=0.6, toprule=.15mm, bottomrule=.15mm, colbacktitle=quarto-callout-tip-color!10!white, titlerule=0mm, bottomtitle=1mm, title=\textcolor{quarto-callout-tip-color}{\faLightbulb}\hspace{0.5em}{Solución}, toptitle=1mm, colframe=quarto-callout-tip-color-frame, colback=white, breakable, rightrule=.15mm, arc=.35mm, leftrule=.75mm, coltitle=black]

Comencemos calculando \(f\circ g\). Tenemos que \[
(f\circ g)(x)=f(g(x))=f\left(\sqrt{2-x}\right)=\sqrt{\sqrt{2-x}}=\sqrt[4]{2-x}.
\]

El dominio de \(f\circ g\) consiste en todos los números \(x\) tales que
\(2-x\geq 0\), es decir, \(D_{f\circ g}=(-\infty,2]\).

Ahora calculamos \(g\circ f\), esto es \[
(g\circ f)(x)=g(f(x))=g\left(\sqrt{x}\right)=\sqrt{2-\sqrt{x}}.
\]

Para que \(g\circ f\) esté bien definida es preciso que \(x\geq 0\) y
también que \(2-\sqrt{x}\geq 0\). Con esto resulta
\(D_{g\circ f}=[0,4]\).

Para el caso de \(f\circ f\) resulta

\[
(f\circ f)(x)=f(f(x))=f\left(\sqrt{x}\right)=\sqrt{\sqrt{x}}=\sqrt[4]{x},
\] cuyo dominio es \(D_{f\circ f}=[0,+\infty)\).

Finalmente, para \(g\circ g\) escribimos \[
(g\circ g)(x)=g(g(x))=g\left(\sqrt{2-x}\right)=\sqrt{2-\sqrt{2-x}}.
\] Esta función está bien definida siempre que se cumpla \(2-x\geq 0\) y
\(2-\sqrt{2-x}\geq 0\), lo que nos da \(D_{g\circ g}=[-2,2]\).

\end{tcolorbox}

Es posible componer más de dos funciones. Por ejemplo, si \(f\), \(g\) y
\(h\) son tres funciones entonces \[
(f\circ g\circ h)(x)=f(g(h(x))),
\] siempre que esta expresión pueda calcularse.

Ejemplo

Encontrar \(f\circ g\circ h\) siendo

\[
f(x)=\frac{x}{x+1},\quad g(x)=x^{10},\quad \text{ y }\quad h(x)=x+3.
\]

\begin{tcolorbox}[enhanced jigsaw, opacityback=0, left=2mm, opacitybacktitle=0.6, toprule=.15mm, bottomrule=.15mm, colbacktitle=quarto-callout-tip-color!10!white, titlerule=0mm, bottomtitle=1mm, title=\textcolor{quarto-callout-tip-color}{\faLightbulb}\hspace{0.5em}{Solución}, toptitle=1mm, colframe=quarto-callout-tip-color-frame, colback=white, breakable, rightrule=.15mm, arc=.35mm, leftrule=.75mm, coltitle=black]

Procedemos como en el ejemplo anterior, obteniendo \[
f(g(h(x)))=f(g(x+3))=f\left((x+3)^{10}\right)=\frac{(x+3)^{10}}{(x+3)^{10}+1}.
\]

\end{tcolorbox}

Ejemplo

Dada la función \(F(x)=\cos^2(x+9)\), encontrar funciones \(f\), \(g\) y
\(h\) de modo que \(F=f\circ g\circ h\).

\begin{tcolorbox}[enhanced jigsaw, opacityback=0, left=2mm, opacitybacktitle=0.6, toprule=.15mm, bottomrule=.15mm, colbacktitle=quarto-callout-tip-color!10!white, titlerule=0mm, bottomtitle=1mm, title=\textcolor{quarto-callout-tip-color}{\faLightbulb}\hspace{0.5em}{Solución}, toptitle=1mm, colframe=quarto-callout-tip-color-frame, colback=white, breakable, rightrule=.15mm, arc=.35mm, leftrule=.75mm, coltitle=black]

Notemos que para calcular \(F(x)\) primero tomamos el número \(x\) y
sumamos 9, a este resultado le calculamos el coseno y finalmente
elevamos al cuadrado. Teniendo en cuenta este orden de operaciones,
podemos definir \(h(x)=x+9\), \(g(x)=\cos x\) y \(f(x)=x^2\). Con esta
elección tenemos que \(F(x)=f(g(h(x)))\).

\end{tcolorbox}

\hyperref[seccion_1.3]{↑ Volver al inicio de la sección}

\section{1.5: Funciones exponenciales}\label{seccion_1.5}

Hemos visto que una función exponencial es de la forma

\[
f(x)=a^x,
\]

donde \(a\) es una constante positiva. Recordemos que

\begin{itemize}
\item
  Si \(x=n\), con \(n\in\mathbb{N}\), entonces \[
  a^x=a^n=\underbrace{a\times a\times \dots \times a}_{n \text{ factores}}.
  \]
\item
  Si \(x=0\), \(a^0=1\), y si \(x=-n\) con \(n\in\mathbb{N}\) entonces
\end{itemize}

\[
a^{-n}=\frac{1}{a^n}.
\]

\begin{itemize}
\tightlist
\item
  Si \(x=\displaystyle \frac{p}{q}\), donde \(p,q\in \mathbb{N}\),
  entonces
\end{itemize}

\[
a^x=a^{p/q}=\sqrt[q]{a^p}=\left(\sqrt[q]{a}\right)^p.
\]

¿Qué significa \(a^x\) cuando \(x\) es un número \textbf{irracional}?

En la siguiente figura vemos un esbozo del gráfico de \(y=a^x\) para
\(x\) racional.

\begin{center}
\includegraphics[width=0.4\linewidth,height=\textheight,keepaspectratio]{index_files/mediabag/imagenes_cap1/1.5-exponencial-racional.pdf}
\end{center}

Los ``huecos'' de esta gráfica corresponden a los valores \(f(x)\)
cuando \(x\) es un número irracional.

Podemos pensar en completar estos espacios definiendo a \(f\) como una
función creciente. Por ejemplo, si tomamos \(a=2\) y queremos conocer el
valor de \(2^{\sqrt{3}}\), observamos primero que

\[
1\text{.}7<\sqrt{3}<1\text{.}8
\]

y como la función es creciente, debe ocurrir que

\[2^{1\text{.}7}<2^{\sqrt{3}}<2^{1\text{.}8}.\]

Como tanto \(1\text{.}7\) como \(1\text{.}8\) son racionales, podemos
calcular \(2^{1\text{.}7}\) y \(2^{1\text{.}8}\), obteniendo una
aproximación del número \(2^{\sqrt{3}}\). Si se desea conocer con más
precisión, notemos que

\[
1\text{.}73<\sqrt{3}<1\text{.}74 \quad\Longrightarrow\quad 2^{1\text{.}73}<2^{\sqrt{3}}<2^{1\text{.}74}
\]

\[
1\text{.}732<\sqrt{3}<1\text{.}733 \quad\Longrightarrow\quad 2^{1\text{.}732}<2^{\sqrt{3}}<2^{1\text{.}733}
\]

\[
1\text{.}73205<\sqrt{3}<1\text{.}73206 \quad\Longrightarrow\quad 2^{1\text{.}73205}<2^{\sqrt{3}}<2^{1\text{.}73206}
\]

Es posible demostrar que existe un único número mayor que todos los
valores

\[
2^{1\text{.}73}, 2^{1\text{.}732}, 2^{1\text{.}73205}, \dots
\]

y menor que todos los números

\[
2^{1\text{.}74}, 2^{1\text{.}733}, 2^{1\text{.}73206}, \dots
\]

Entonces definimos \(2^{\sqrt{3}}\) como este número. Usando el proceso
de aproximación anterior podemos obtener la siguiente estimación
(correcta hasta seis cifras decimales)

\[
2^{\sqrt{3}}\approx 3\text{.}321997
\]

Este ejemplo puede repetirse con cualquier otro número irracional \(x\)
y, de esta manera, completar la gráfica punteada anterior.

\begin{center}
\includegraphics[width=0.6\linewidth,height=\textheight,keepaspectratio]{index_files/mediabag/imagenes_cap1/1.5-exponenciales.pdf}
\end{center}

Observemos que, si \(f(x)=a^{x}\) entonces

\[
\left(\frac{1}{a}\right)^x=\frac{1}{a^x}=a^{-x}=f(-x),
\]

con lo que la gráfica de \(y=(1/a)^x\) puede obtenerse de la de \(f\)
\textbf{reflejándola con respecto al eje} \(y\).

Proposición (Leyes de los exponentes)

Sean \(a,b\) positivos y \(x,y\) números reales. Entonces

\begin{enumerate}
\def\labelenumi{\arabic{enumi}.}
\tightlist
\item
  \(a^{x+y}=a^x\,a^y\).
\item
  \(a^{x-y}=\displaystyle \frac{a^x}{a^y}\).
\end{enumerate}

\begin{enumerate}
\def\labelenumi{\arabic{enumi}.}
\setcounter{enumi}{2}
\tightlist
\item
  \(\left(a^x\right)^y=a^{xy}\).
\item
  \((ab)^x=a^x\,b^x\).
\end{enumerate}

Ejemplo

Trazar la gráfica de \(y=3-2^x\). Determinar su dominio e imagen.

\begin{tcolorbox}[enhanced jigsaw, opacityback=0, left=2mm, opacitybacktitle=0.6, toprule=.15mm, bottomrule=.15mm, colbacktitle=quarto-callout-tip-color!10!white, titlerule=0mm, bottomtitle=1mm, title=\textcolor{quarto-callout-tip-color}{\faLightbulb}\hspace{0.5em}{Solución}, toptitle=1mm, colframe=quarto-callout-tip-color-frame, colback=white, breakable, rightrule=.15mm, arc=.35mm, leftrule=.75mm, coltitle=black]

Comencemos con el gráfico de \(y=f(x)=2^x\).

\includegraphics[width=0.35\linewidth,height=\textheight,keepaspectratio]{index_files/mediabag/imagenes_cap1/1.5-ejemplo-exponencial.pdf}

Para obtener el de la función deseada, reflejamos éste respecto del eje
\(x\) y luego realizamos una traslación vertical, tres unidades hacia
arriba. Así, resulta

\includegraphics[width=0.35\linewidth,height=\textheight,keepaspectratio]{index_files/mediabag/imagenes_cap1/1.5-ejemplo-exponencial2.pdf}

Como la expresión \(2^x\) está definida para cualquier número real,
tenemos que el dominio es \(D_f=\mathbb{R}\). Por otro lado, \(2^x>0\)
implica que \(-2^x<0\), y si sumamos 3 a ambos miembros de la
desigualdad concluimos que \(3-2^x<3\), es decir, \(y<3\). Con esto
vemos que \(\operatorname{Im}(f)=(-\infty,3)\).

\end{tcolorbox}

\subsection{Aplicación: modelos
poblacionales}\label{aplicaciuxf3n-modelos-poblacionales}

Consideremos una población de bacterias en un medio nutriente homogéneo.
Al realizar mediciones periódicas, se observa que la población se
duplica cada hora.

Si inicialmente se cuenta con 1000 individuos, ¿qué función permite
modelar la población en cada instante de tiempo?

Si \(p(t)\) denota la cantidad de individuos a tiempo \(t\), medido en
horas, entonces observemos que

\[
\begin{aligned}
p(0)&=1000\\
p(1)&=2p(0)=2^1\times 1000\\
p(2)&=2p(1)=2\times(2\times 1000)=2^2\times 1000\\
p(3)&=2p(2)=2\times(2^2\times 1000)=2^3\times 1000
\end{aligned}
\]

continuando este patrón, podemos ver que

\[
p(t)=1000\times 2^t.
\]

\subsection{\texorpdfstring{El número
\(e\)}{El número e}}\label{el-nuxfamero-e}

Hay una base \(a\) en la función exponencial que será especialmente
útil. La elección de la base influye en la forma en que la gráfica de
\(y=a^x\) cruza al eje \(y\).

La \textbf{recta tangente} en el punto (0,1) (que luego definiremos más
precisamente, pero que podemos pensar como la recta que sólo toca al
gráfico en ese punto) para \(y=2^x\) es \(m\approx 0\text{.}7\) y para
\(y=3^x\) es \(m\approx 1\text{.}1\).

\begin{center}
\includegraphics[width=0.6\linewidth,height=\textheight,keepaspectratio]{index_files/mediabag/imagenes_cap1/1.5-pendientes-exponenciales.pdf}
\end{center}

\begin{tcolorbox}[enhanced jigsaw, opacityback=0, left=2mm, opacitybacktitle=0.6, toprule=.15mm, bottomrule=.15mm, colbacktitle=quarto-callout-note-color!10!white, titlerule=0mm, bottomtitle=1mm, title=\textcolor{quarto-callout-note-color}{\faInfo}\hspace{0.5em}{Definición (El número \(e\))}, toptitle=1mm, colframe=quarto-callout-note-color-frame, colback=white, breakable, rightrule=.15mm, arc=.35mm, leftrule=.75mm, coltitle=black]

El número \(e\) es aquella base que se elige de modo que la recta
tangente al gráfico de \(y=e^x\) en el punto \((0,1)\) tenga pendiente
\(m=1\).

\end{tcolorbox}

El gráfico de \(y=e^x\) se encuentra entre los de \(y=2^x\) e \(y=3^x\).
El valor exacto hasta cinco cifras decimales de \(e\) es

\[
e\approx 2\text{.}71828
\]

\begin{center}
\includegraphics[width=0.4\linewidth,height=\textheight,keepaspectratio]{index_files/mediabag/imagenes_cap1/1.5-exponencial-2-e-3.pdf}
\end{center}

Ejemplo

Dibujar la función \(y=\displaystyle\frac{1}{2}e^{-x}-1\). Determinar su
dominio y rango.

\begin{tcolorbox}[enhanced jigsaw, opacityback=0, left=2mm, opacitybacktitle=0.6, toprule=.15mm, bottomrule=.15mm, colbacktitle=quarto-callout-tip-color!10!white, titlerule=0mm, bottomtitle=1mm, title=\textcolor{quarto-callout-tip-color}{\faLightbulb}\hspace{0.5em}{Solución}, toptitle=1mm, colframe=quarto-callout-tip-color-frame, colback=white, breakable, rightrule=.15mm, arc=.35mm, leftrule=.75mm, coltitle=black]

Para graficar la función dada podemos partir de \(y=e^x\). Primero
reflejamos la curva con respecto al eje \(y\) obteniendo el de
\(y=e^{-x}\), y seguidamente realizamos un alargamiento horizontal en un
factor 2 y finalmente un desplazamiento vertical de una unidad hacia
abajo. De esta manera, el gráfico resultante es

\includegraphics[width=0.4\linewidth,height=\textheight,keepaspectratio]{index_files/mediabag/imagenes_cap1/1.5-ejemplo-exponenciales.pdf}

Tenemos que el dominio es \(\mathbb{R}\) y, como \(e^{-x}>0\), al
multiplicar por \(1/2\) y restar \(1\) obtenemos
\(\displaystyle \frac{1}{2}e^{-x}-1>-1\), con lo cual el conjunto imagen
es \((-1,+\infty)\).

\end{tcolorbox}

\hyperref[seccion_1.5]{↑ Volver al inicio de la sección}

\section{1.6: Funciones inversas y logarítmicas}\label{seccion_1.6}

\begin{tcolorbox}[enhanced jigsaw, opacityback=0, left=2mm, opacitybacktitle=0.6, toprule=.15mm, bottomrule=.15mm, colbacktitle=quarto-callout-note-color!10!white, titlerule=0mm, bottomtitle=1mm, title=\textcolor{quarto-callout-note-color}{\faInfo}\hspace{0.5em}{Definición (Función uno a uno)}, toptitle=1mm, colframe=quarto-callout-note-color-frame, colback=white, breakable, rightrule=.15mm, arc=.35mm, leftrule=.75mm, coltitle=black]

Una función \(f\) se llama \textbf{uno a uno} o \textbf{inyectiva} si
nunca toma el mismo valor dos veces. Es decir, si se cumple que

\begin{equation}\phantomsection\label{eq-inyectiva_no}{
x_1\neq x_2 \Longrightarrow f(x_1)\neq f(x_2)
}\end{equation}

para todos \(x_1,x_2\in D_f\). O, equivalentemente, si

\begin{equation}\phantomsection\label{eq-inyectiva_si}{
f(x_1)=f(x_2) \Longrightarrow x_1=x_2
}\end{equation}

para todos \(x_1,x_2\in D_f\).

\end{tcolorbox}

En el siguiente esquema vemos dos funciones: \(f\) es uno a uno, pero
\(g\) no lo es.

\begin{center}
\includegraphics[width=0.4\linewidth,height=\textheight,keepaspectratio]{index_files/mediabag/imagenes_cap1/1.6-venns.pdf}
\end{center}

Geométricamente, si \(f\) es una función uno a uno, entonces ninguna
recta horizontal puede cortar al gráfico más de una vez. Si hay al menos
una recta horizontal que corta al gráfico de \(f\) más de una vez,
entonces \(f\) no es uno a uno.

\begin{center}
\includegraphics[width=0.4\linewidth,height=\textheight,keepaspectratio]{index_files/mediabag/imagenes_cap1/1.6-funcion-no-unouno.pdf}
\end{center}

Prueba de la línea horizontal

Una función es uno a uno si y sólo si ninguna recta horizontal corta a
su gráfico más de una vez.

Ejemplo

¿Es \(f(x)=x^3\) una función uno a uno?

\begin{tcolorbox}[enhanced jigsaw, opacityback=0, left=2mm, opacitybacktitle=0.6, toprule=.15mm, bottomrule=.15mm, colbacktitle=quarto-callout-tip-color!10!white, titlerule=0mm, bottomtitle=1mm, title=\textcolor{quarto-callout-tip-color}{\faLightbulb}\hspace{0.5em}{Solución}, toptitle=1mm, colframe=quarto-callout-tip-color-frame, colback=white, breakable, rightrule=.15mm, arc=.35mm, leftrule=.75mm, coltitle=black]

Si dibujamos el gráfico de \(f\), la prueba de la línea horizontal
sugiere que \(f\) es uno a uno.

\includegraphics[width=0.3\linewidth,height=\textheight,keepaspectratio]{index_files/mediabag/imagenes_cap1/1.6-ejemplo-cubica.pdf}

Ahora vamos a utilizar la ecuación (\ref{eq-inyectiva_si}) para
demostrarlo analíticamente. Sean \(x_1\) y \(x_2\) en \(\mathbb{R}\) de
modo que \(f(x_1)=f(x_2)\). Esto significa que \(x_1^3=x_2^3\), y
aplicando raíz cúbica a ambos miembros obtenemos que \(x_1=x_2\). Es
decir, la igualdad \(f(x_1)=f(x_2)\) implica que \(x_1=x_2\). En
consecuencia, \(f\) resulta ser uno a uno.

\end{tcolorbox}

Ejemplo

¿Es \(f(x)=x^2\) una función uno a uno?

\begin{tcolorbox}[enhanced jigsaw, opacityback=0, left=2mm, opacitybacktitle=0.6, toprule=.15mm, bottomrule=.15mm, colbacktitle=quarto-callout-tip-color!10!white, titlerule=0mm, bottomtitle=1mm, title=\textcolor{quarto-callout-tip-color}{\faLightbulb}\hspace{0.5em}{Solución}, toptitle=1mm, colframe=quarto-callout-tip-color-frame, colback=white, breakable, rightrule=.15mm, arc=.35mm, leftrule=.75mm, coltitle=black]

Si graficamos \(f(x)=x^2\), la prueba de la línea horizontal indica que
\(f\) no es uno a uno, es decir, existen puntos distintos del dominio
donde \(f\) toma el mismo valor.

\includegraphics[width=0.3\linewidth,height=\textheight,keepaspectratio]{index_files/mediabag/imagenes_cap1/1.6-ejemplo-cuadratica.pdf}

Para comprobar esto analíticamente, veamos que la ecuación
(\ref{eq-inyectiva_no}) no se satisface. Para ello, basta con encontrar
un par de valores \(x_1\) y \(x_2\) en el dominio tales que
\(x_1\neq x_2\) pero que \(f(x_1)=f(x_2)\). Si consideramos \(x_1=-2\) y
\(x_2=2\) claramente es \(x_1\neq x_2\). Sin embargo

\[
f(x_1)=f(-2)=(-2)^2=4=2^2=f(2)=f(x_2)
\]

con lo que \(f(x_1)=f(x_2)\). De esta forma, \(f\) no puede ser una
función uno a uno.

\end{tcolorbox}

\begin{tcolorbox}[enhanced jigsaw, opacityback=0, left=2mm, opacitybacktitle=0.6, toprule=.15mm, bottomrule=.15mm, colbacktitle=quarto-callout-note-color!10!white, titlerule=0mm, bottomtitle=1mm, title=\textcolor{quarto-callout-note-color}{\faInfo}\hspace{0.5em}{Definición (Función inversa)}, toptitle=1mm, colframe=quarto-callout-note-color-frame, colback=white, breakable, rightrule=.15mm, arc=.35mm, leftrule=.75mm, coltitle=black]

Sea \(f\) una función uno a uno con dominio \(A\) y rango \(B\).
Entonces su \textbf{función inversa} tiene dominio \(B\) y rango \(A\),
y se define de manera que

\begin{equation}\phantomsection\label{eq-funcion_inversa}{
f^{-1}(y)=x\quad \Longleftrightarrow\quad f(x)=y,
}\end{equation}

para cualquier \(y\) en \(B\).

\end{tcolorbox}

Notemos que es importante que \(f\) sea uno a uno: de lo contrario, al
intentar definir \(f^{-1}\) podría ocurrir que esta ley no corresponda a
una función (\(f^{-1}\) no puede tomar más de un valor en cada \(y\) en
\(B\)).

Al comparar \(f\) con \(f^{-1}\) vemos que cada una invierte el efecto
de la otra. Notar que

\[
\operatorname{Dom}(f^{-1})=\operatorname{Im}(f)
\]

y también

\[
\operatorname{Im}(f^{-1})=\operatorname{Dom}(f).
\]

Por ejemplo, si \(f(x)=x^3\) entonces \(f^{-1}(y)=y^{1/3}\). Pues, si
\(y=x^3\) entonces

\[
f^{-1}(y)=f^{-1}(x^3)=(x^3)^{1/3}=x.
\]

Recíprocamente, si \(x=y^{1/3}\) entonces

\[
f(x)=f(y^{1/3})=(y^{1/3})^3=y.
\]

El siguiente diagrama muestra que \(f^{-1}\) invierte el efecto de \(f\)

\begin{center}
\includegraphics[width=0.6\linewidth,height=\textheight,keepaspectratio]{index_files/mediabag/imagenes_cap1/1.6-funcion-inversa.pdf}
\end{center}

Si sustituimos \(x\) e \(y\) en las correspondientes igualdades de
(\ref{eq-funcion_inversa}) obtenemos

\begin{equation}\phantomsection\label{eq-cancelacion-inversafunc}{
f^{-1}(f(x))=x, \quad \text{ para } \quad  x\in A
}\end{equation}

\begin{equation}\phantomsection\label{eq-cancelacion-funcinversa}{
f(f^{-1}(y))=y, \quad \text{ para } \quad y\in B
}\end{equation}

llamadas \textbf{ecuaciones de cancelación}.

\begin{tcolorbox}[enhanced jigsaw, opacityback=0, left=2mm, opacitybacktitle=0.6, toprule=.15mm, bottomrule=.15mm, colbacktitle=quarto-callout-caution-color!10!white, titlerule=0mm, bottomtitle=1mm, title=\textcolor{quarto-callout-caution-color}{\faFire}\hspace{0.5em}{Importante}, toptitle=1mm, colframe=quarto-callout-caution-color-frame, colback=white, breakable, rightrule=.15mm, arc=.35mm, leftrule=.75mm, coltitle=black]

No confundir el \(-1\) de \(f^{-1}\) con un exponente. Ésta es
simplemente una notación para referirse a la función inversa de \(f\), y
no significa \(1/f(x)\).

\end{tcolorbox}

\subsection{¿Cómo encontrar la función inversa de una función uno a
uno?}\label{cuxf3mo-encontrar-la-funciuxf3n-inversa-de-una-funciuxf3n-uno-a-uno}

Para ello podemos seguir los siguientes pasos:

\begin{enumerate}
\def\labelenumi{\arabic{enumi}.}
\tightlist
\item
  Escribir \(y=f(x)\).
\item
  Resolver la ecuación para \(x\).
\item
  Intercambiar \(x\) e \(y\).
\end{enumerate}

Ejemplo

Encontrar la función inversa de \(f(x)=x^3+2\).

\begin{tcolorbox}[enhanced jigsaw, opacityback=0, left=2mm, opacitybacktitle=0.6, toprule=.15mm, bottomrule=.15mm, colbacktitle=quarto-callout-tip-color!10!white, titlerule=0mm, bottomtitle=1mm, title=\textcolor{quarto-callout-tip-color}{\faLightbulb}\hspace{0.5em}{Solución}, toptitle=1mm, colframe=quarto-callout-tip-color-frame, colback=white, breakable, rightrule=.15mm, arc=.35mm, leftrule=.75mm, coltitle=black]

Primero verifiquemos que \(f\) es una función uno a uno. Sean \(x_1\) y
\(x_2\) dos números en el dominio de \(f\) (es decir, en \(\mathbb{R}\))
tales que \(f(x_1)=f(x_2)\). Entonces tenemos

\[
x_1^3+2=x_2^3+2.
\]

Restando 2 a ambos miembros y luego aplicando raíz cúbica, concluimos
que \(x_1=x_2\), y entonces \(f\colon \mathbb{R}\to \mathbb{R}\) es un
función uno a uno, en virtud de (\ref{eq-inyectiva_si}).

Ahora procedemos a calcular su función inversa realizando los pasos
antes descriptos. Si intercambiamos \(y\) con \(x\) resulta

\[
x=y^3+2
\]

y a continuación, despejando \(y\) obtenemos

\[
y=\sqrt[3]{x-2}.
\]

De esta manera, \(f^{-1}\colon \mathbb{R}\to \mathbb{R}\) está dada por

\[
f^{-1}(x)=\sqrt[3]{x-2}.
\]

\end{tcolorbox}

Al intercambiar \(y\) con \(x\) podemos hallar el gráfico de \(f^{-1}\)
a partir del de \(f\). Por definición, sabemos que \(f(a)=b\) si y sólo
si \(f^{-1}(b)=a\). Entonces

\[
(a,b)\in \text{Gráf}\,(f)\quad \text{ si y sólo si }\quad (b,a)\in \text{Gráf}\,(f^{-1}).
\]

Como los puntos \((a,b)\) y \((b,a)\) son simétricos respecto de la
recta \(y=x\), el gráfico de \(f^{-1}\) puede obtenerse reflejando el de
\(f\) sobre la recta \(y=x\).

\begin{center}
\includegraphics[width=0.4\linewidth,height=\textheight,keepaspectratio]{index_files/mediabag/imagenes_cap1/1.6-graficas-funcion-e-inversa.pdf}
\end{center}

Ejemplo

Trazar el gráfico de \(f(x)=\sqrt{-1-x}\) y el de su función inversa en
un mismo sistema de ejes.

\begin{tcolorbox}[enhanced jigsaw, opacityback=0, left=2mm, opacitybacktitle=0.6, toprule=.15mm, bottomrule=.15mm, colbacktitle=quarto-callout-tip-color!10!white, titlerule=0mm, bottomtitle=1mm, title=\textcolor{quarto-callout-tip-color}{\faLightbulb}\hspace{0.5em}{Solución}, toptitle=1mm, colframe=quarto-callout-tip-color-frame, colback=white, breakable, rightrule=.15mm, arc=.35mm, leftrule=.75mm, coltitle=black]

Notemos que \(D_f=(-\infty,-1]\) e \(\operatorname{Im}(f)=[0,+\infty)\).
Veamos que \(f\) es efectivamente una función uno a uno. Usaremos
nuevamente la definición (\ref{eq-inyectiva_si}). Si \(x_1\) y \(x_2\)
son dos números en el dominio tales que \(f(x_1)=f(x_2)\), entonces
tenemos que

\[
\sqrt{-1-x_1}=\sqrt{-1-x_2}.
\]

Elevando al cuadrado y sumando uno en ambos miembros resulta

\[
-x_1=-x_2,
\]

y al multiplicar ambos lados de la igualdad por \(-1\) concluimos que
\(x_1=x_2\).

Ahora calculamos \(f^{-1}\) procediendo como en el ejemplo anterior.
Intercambiando \(y\) con \(x\) obtenemos

\[
x=\sqrt{-1-y}
\]

con lo cual

\[
f^{-1}(x)=-x^2-1.
\]

La función obtenida luce, a primera vista, como una parábola. Sin
embargo, recordar que \(D_{f^{-1}}=[0,+\infty)\) por lo tanto la gráfica
de la función inversa corresponde sólamente a la ``mitad derecha'' de la
parábola.

\includegraphics[width=0.4\linewidth,height=\textheight,keepaspectratio]{index_files/mediabag/imagenes_cap1/1.6-ejemplo-inversa-raiz.pdf}

\end{tcolorbox}

\subsection{Funciones logarítmicas}\label{funciones-logaruxedtmicas-1}

Si \(a>0\) y \(a\neq 1\) la función exponencial \(y=a^x\) es uno a uno.
Entonces tiene función inversa, llamada \textbf{función logarítmica} de
base \(a\) y denotada como

\[
y=\log_a x.
\]

Utilizando la \hyperref[eq-funcion_inversa]{definición de función
inversa} tenemos que

\[
\log_a x=y\quad \Longleftrightarrow \quad a^y=x.
\]

En definitiva, si \(x>0\), \(\log_a x\) es el \textbf{exponente} al que
debe elevarse a \(a\) para obtener \(x\). Por ejemplo, \(\log_2 8=3\) ya
que \(2^3=8\).

Las ecuaciones de cancelación en este caso son

\[
\log_a (a^x)=x, \quad \text{ para } \quad  x\in \mathbb{R}
\]

y también

\[
a^{\log_a x}=x, \quad \text{ para }  \quad x>0.
\]

La función \(y=\log_a x\) tiene dominio \((0,\infty)\) y rango
\(\mathbb{R}\). Su gráfico puede obtenerse al reflejar el de \(y=a^x\)
respecto de \(y=x\).

\begin{center}
\includegraphics[width=0.5\linewidth,height=\textheight,keepaspectratio]{index_files/mediabag/imagenes_cap1/1.6-exponencial-logaritmo.pdf}
\end{center}

Como la gráfica de \(y=a^x\) pasa por el punto \((0,1)\), la de su
inversa \(y=\log_a x\) pasa por \((1,0)\).

\phantomsection\label{propo-leyes-logaritmos}
\textbf{Proposición (Leyes de los logaritmos)}

Sean \(x\) e \(y\) números reales positivos y \(a>0, a\neq 1\). Entonces

\textbf{(1)} \(\log_a(xy)=\log_a x+\log_a y\).

\textbf{(2)} \(\log_a\!\left(\frac{x}{y}\right)=\log_a x-\log_a y\).

\textbf{(3)} \(\log_a(x^r)=r\log_a x,\quad r\in\mathbb{R}\).

Ejemplo

Evalúe \(\log_2 80-\log_2 5\) utilizando las leyes de los logaritmos.

\begin{tcolorbox}[enhanced jigsaw, opacityback=0, left=2mm, opacitybacktitle=0.6, toprule=.15mm, bottomrule=.15mm, colbacktitle=quarto-callout-tip-color!10!white, titlerule=0mm, bottomtitle=1mm, title=\textcolor{quarto-callout-tip-color}{\faLightbulb}\hspace{0.5em}{Solución}, toptitle=1mm, colframe=quarto-callout-tip-color-frame, colback=white, breakable, rightrule=.15mm, arc=.35mm, leftrule=.75mm, coltitle=black]

Si escribimos \(80=16\cdot 5=2^4\cdot 5\), entonces aplicando las
\hyperref[propo-leyes-logaritmos]{leyes de los logaritmos} (1) y (3)
resulta

\[
\log_2 80-\log_2 5=\log_2(2^4\cdot 5)-\log_2 5=\log_2(2^4)+\log_2 5-\log_2 5=4\log_2 2=4\cdot 1=4.
\]

\end{tcolorbox}

\subsection{Logaritmos naturales}\label{logaritmos-naturales}

Cuando la base \(a\) es el número \(e\), el logaritmo recibe el nombre
de \textbf{logaritmo natural} o \textbf{neperiano} y se denota por
\(\log_e x=\ln x\). En este caso, usando la
\hyperref[eq-funcion_inversa]{definición de inversa} vemos que

\begin{equation}\phantomsection\label{eq-relacioninv_explog}{
\ln x=y \quad\Longleftrightarrow\quad e^y=x.
}\end{equation}

Las ecuaciones de cancelación son

\begin{equation}\phantomsection\label{eq-cancelacion_lnexp}{
\ln (e^x)=x, \quad \text{ para } \quad  x\in \mathbb{R}
}\end{equation}

y también

\begin{equation}\phantomsection\label{eq-cancelacion_expln}{
e^{\ln x}=x, \quad \text{ para } \quad x>0.
}\end{equation}

Ejemplo

Encuentre \(x\) si \(\ln x=5\).

\begin{tcolorbox}[enhanced jigsaw, opacityback=0, left=2mm, opacitybacktitle=0.6, toprule=.15mm, bottomrule=.15mm, colbacktitle=quarto-callout-tip-color!10!white, titlerule=0mm, bottomtitle=1mm, title=\textcolor{quarto-callout-tip-color}{\faLightbulb}\hspace{0.5em}{Solución}, toptitle=1mm, colframe=quarto-callout-tip-color-frame, colback=white, breakable, rightrule=.15mm, arc=.35mm, leftrule=.75mm, coltitle=black]

Utilizando la equivalencia (\ref{eq-relacioninv_explog}) obtenemos que

\[
\ln x=5 \quad \Longleftrightarrow \quad x=e^5.
\]

\end{tcolorbox}

Ejemplo

Resuelva la ecuación \(e^{5-3x}=10\).

\begin{tcolorbox}[enhanced jigsaw, opacityback=0, left=2mm, opacitybacktitle=0.6, toprule=.15mm, bottomrule=.15mm, colbacktitle=quarto-callout-tip-color!10!white, titlerule=0mm, bottomtitle=1mm, title=\textcolor{quarto-callout-tip-color}{\faLightbulb}\hspace{0.5em}{Solución}, toptitle=1mm, colframe=quarto-callout-tip-color-frame, colback=white, breakable, rightrule=.15mm, arc=.35mm, leftrule=.75mm, coltitle=black]

Aplicamos logaritmo natural a ambos miembros de la igualdad

\[
\ln\left(e^{5-3x}\right)=\ln 10
\]

y luego utilizamos la \hyperref[eq-cancelacion_lnexp]{ecuación de
cancelación}, obteniendo

\[
5-3x=\ln 10, \quad \text{ lo que implica que }\quad x=\frac{5-\ln 10}{3}.
\]

\end{tcolorbox}

Ejemplo

Expresar a \(\ln a+\tfrac{1}{2}\ln b\) como un solo logaritmo.

\begin{tcolorbox}[enhanced jigsaw, opacityback=0, left=2mm, opacitybacktitle=0.6, toprule=.15mm, bottomrule=.15mm, colbacktitle=quarto-callout-tip-color!10!white, titlerule=0mm, bottomtitle=1mm, title=\textcolor{quarto-callout-tip-color}{\faLightbulb}\hspace{0.5em}{Solución}, toptitle=1mm, colframe=quarto-callout-tip-color-frame, colback=white, breakable, rightrule=.15mm, arc=.35mm, leftrule=.75mm, coltitle=black]

Utilizamos las \hyperref[propo-leyes-logaritmos]{leyes de los
logaritmos} (1) y (3) podemos escribir

\[
\ln a+\frac{1}{2}\ln b=\ln a+\ln b^{1/2}=\ln\left(ab^{1/2}\right)=\ln\left(a\sqrt{b}\right).
\]

\end{tcolorbox}

Si \(a>0\), \(a\neq 1\) entonces podemos escribir

\[
\log_a x=\frac{\ln x}{\ln a},
\]

conocida como \textbf{fórmula de cambio de base}. Más generalmente, si
\(a\) y \(b\) son positivos y distintos de \(1\) entonces

\[
\log_b x=\frac{\log_a x}{\log_a b}\quad \text{(conversión de base $b$ a base $a$)}
\]

Ejemplo

Trazar la gráfica de \(y=\ln(x-2)-1\).

\begin{tcolorbox}[enhanced jigsaw, opacityback=0, left=2mm, opacitybacktitle=0.6, toprule=.15mm, bottomrule=.15mm, colbacktitle=quarto-callout-tip-color!10!white, titlerule=0mm, bottomtitle=1mm, title=\textcolor{quarto-callout-tip-color}{\faLightbulb}\hspace{0.5em}{Solución}, toptitle=1mm, colframe=quarto-callout-tip-color-frame, colback=white, breakable, rightrule=.15mm, arc=.35mm, leftrule=.75mm, coltitle=black]

Partiendo del gráfico de \(y=\ln x\), podemos obtener la curva deseada
desplazando la anterior dos unidades para la derecha y luego una unidad
hacia abajo.

\includegraphics[width=0.6\linewidth,height=\textheight,keepaspectratio]{index_files/mediabag/imagenes_cap1/1.6-ejemplo-logaritmo.pdf}

\end{tcolorbox}

\subsection{Funciones trigonométricas
inversas}\label{funciones-trigonomuxe9tricas-inversas}

Si intentamos calcular las funciones inversas de las trigonométricas nos
encontramos con el problema de que no son uno a uno. Entonces debemos
restringir sus dominios a un intervalo adecuado donde tomen todos los
valores en los respectivos rangos y sean uno a uno.

Por ejemplo, si restringimos la función \(y=f(x)=\operatorname{sen} x\)
al intervalo \(I=[-\tfrac{\pi}{2},\tfrac{\pi}{2}]\), entonces \(f\) toma
todos los valores del intervalo \([-1,1]\) (su rango) y resulta uno a
uno en \(I\).

\begin{center}
\includegraphics[width=0.65\linewidth,height=\textheight,keepaspectratio]{index_files/mediabag/imagenes_cap1/1.6-arcoseno.pdf}
\end{center}

La función inversa del seno restringido a este intervalo existe y se
denota mediante \(\operatorname{sen}^{-1}\) o \(\operatorname{arcsen}\),
y se conoce como \textbf{función inversa del seno} o \textbf{función
arco seno}.

Por la \hyperref[eq-funcion_inversa]{definición de inversa} tenemos que

\[
\operatorname{sen}^{-1} x=y\quad \Longleftrightarrow\quad \operatorname{sen} y=x \qquad -\frac{\pi}{2}\leq y\leq \frac{\pi}{2}.
\]

Si \(-1\leq x\leq 1\), \(\operatorname{sen}^{-1} x\) es un número en el
intervalo \([-\tfrac{\pi}{2},\tfrac{\pi}{2}]\) cuyo seno es \(x\).

En este caso, las ecuaciones de cancelación resultan

\[
\operatorname{sen}^{-1}(\operatorname{sen} x)=x, \quad para \quad  -\tfrac{\pi}{2}\leq x\leq \tfrac{\pi}{2}
\]

y también

\[
\operatorname{sen}(\operatorname{sen}^{-1} x)=x, \quad para \quad -1\leq x\leq 1.
\]

Para el caso de la función coseno se procede de forma similar. La
función \(y=\cos x\) es uno a uno y toma todos los valores del rango si
nos restringimos al intervalo \([0,\pi]\). Entonces existe su función
inversa, denotada por \(\cos^{-1}\) o \(\operatorname{arccos}\) y
llamada \textbf{función inversa del coseno} o \textbf{función arco
coseno}.

\begin{center}
\includegraphics[width=0.65\linewidth,height=\textheight,keepaspectratio]{index_files/mediabag/imagenes_cap1/1.6-arcocoseno.pdf}
\end{center}

En virtud de la \hyperref[eq-funcion_inversa]{definición de inversa}
tenemos que

\[
\cos^{-1}x=y\quad\Longleftrightarrow\quad \cos y=x\qquad 0\leq y\leq \pi.
\]

El dominio de \(y=\operatorname{arccos} x\) es el intervalo {[}-1,1{]} y
el rango el intervalo \([0,\pi]\).

Las ecuaciones de cancelación son

\[
\cos^{-1}(\cos x)=x, \quad \text{ para } \quad  0\leq x\leq \pi
\]

como también

\[
\cos(\cos^{-1} x)=x, \quad \text{ para } \quad -1\leq x\leq 1.
\]

Para el caso de la función tangente, resulta uno a uno si nos
restringimos al intervalo
\(\displaystyle\left(-\frac{\pi}{2},\frac{\pi}{2}\right)\).

\begin{center}
\includegraphics[width=0.65\linewidth,height=\textheight,keepaspectratio]{index_files/mediabag/imagenes_cap1/1.6-arcotangente.pdf}
\end{center}

En este intervalo \(y=\tan x\) posee función inversa, llamada
\textbf{función tangente inversa} o \textbf{arcotangente}, denotada
\(y=\tan^{-1} x\) o \(y=\operatorname{arctan} x\).

Usando la \hyperref[eq-funcion_inversa]{definición de inversa} obtenemos
la relación

\[
\tan^{-1} x=y\quad\Longleftrightarrow\quad x=\tan y\qquad -\frac{\pi}{2}<y<\frac{\pi}{2}.
\]

Este gráfico presenta dos asíntotas horizontales, cuyas ecuaciones son
\(y=-\tfrac{\pi}{2}\) e \(y=\tfrac{\pi}{2}\). Esto se corresponde con el
hecho de que la función tangente tiene las asíntotas verticales
\(x=-\tfrac{\pi}{2}\) y \(x=-\tfrac{\pi}{2}\).

Las ecuaciones de cancelación son

\[
\tan^{-1}(\tan x)=x, \quad \text{ para }  \quad   -\tfrac{\pi}{2}< x<\tfrac{\pi}{2}
\]

y

\[
\tan(\tan^{-1} x)=x, \quad \text{ para } \quad x\in \mathbb{R}.
\]

Ejemplo

Simplificar la expresión \(\cos(\operatorname{arctan} x)\).

\begin{tcolorbox}[enhanced jigsaw, opacityback=0, left=2mm, opacitybacktitle=0.6, toprule=.15mm, bottomrule=.15mm, colbacktitle=quarto-callout-tip-color!10!white, titlerule=0mm, bottomtitle=1mm, title=\textcolor{quarto-callout-tip-color}{\faLightbulb}\hspace{0.5em}{Solución}, toptitle=1mm, colframe=quarto-callout-tip-color-frame, colback=white, breakable, rightrule=.15mm, arc=.35mm, leftrule=.75mm, coltitle=black]

Sea \(y=\operatorname{arctan} x\). Entonces \(\tan y =x\). Dado que
queremos calcular \(\cos y\), debemos expresar el coseno del número
\(y\) en términos de la tangente. Es sencillo verificar que

\[
1+\tan^2 y = \frac{1}{\cos^2 y}
\]

y utilizando esa relación, obtenemos

\[
\cos^2 y =\frac{1}{1+\tan^2 y}.
\]

Dado que \(y\in \left(-\frac{\pi}{2},\frac{\pi}{2}\right)\),
\(\cos y>0\) y despejando deducimos que

\[
\cos(\operatorname{arctan} x) = \cos y =\sqrt{\frac{1}{1+\tan^2 y}}=\frac{1}{\sqrt{1+x^2}}.
\]

\end{tcolorbox}

Las inversas de las funciones trigonométricas restantes no aparecen con
frecuencia. Las enunciamos a continuación.

\begin{itemize}
\tightlist
\item
  \textbf{Función cosecante inversa} o \textbf{arcocosecante}.
\end{itemize}

\[
y=\operatorname{csc}^{-1} x \,\, \quad (|x|\geq 1)\quad \Longleftrightarrow \quad\csc y=x, \,\,\,\, y\in (0,\tfrac{\pi}{2}]\cup (\pi,\tfrac{3}{2}\pi].
\]

\begin{itemize}
\tightlist
\item
  \textbf{Función secante inversa} o \textbf{arcosecante}.
\end{itemize}

\[
y=\operatorname{sec}^{-1} x \,\, \quad (|x|\geq 1)\quad \Longleftrightarrow \quad\sec y=x, \,\,\,\, y\in [0,\tfrac{\pi}{2})\cup [\pi,\tfrac{3}{2}\pi).
\]

\begin{itemize}
\tightlist
\item
  \textbf{Función cotangente inversa} o \textbf{arcocotangente}.
\end{itemize}

\[
y=\operatorname{cot}^{-1} x \,\, \quad (x\in\mathbb{R})\quad \Longleftrightarrow \quad\cot y=x, \,\,\,\, y\in (0,\pi).
\]

\hyperref[seccion_1.6]{↑ Volver al inicio de la sección}

\bookmarksetup{startatroot}

\chapter{2. Límites, continuidad y
derivadas}\label{luxedmites-continuidad-y-derivadas}

\section{2.1. La tangente y los problemas de la
velocidad}\label{seccion_2.1}

\subsection{El problema de la
tangente}\label{el-problema-de-la-tangente}

La palabra \textbf{tangente} deriva del latín \emph{tangens} que
significa ``tocar''. Así, podramos decir que una recta tangente a una
curva es una recta que toca a la curva.

\begin{center}
\includegraphics[width=0.6\linewidth,height=\textheight,keepaspectratio]{index_files/mediabag/imagenes_cap2/2.1-tangentes.pdf}
\end{center}

Cuando la curva es una circunferencia, la recta tangente es aquella que
la interseca en un solo punto. Sin embargo, para curvas más generales
esta definición no parece adecuada: la recta \(\ell\) corta a la curva
una sola vez, pero no se parece a lo que consideramos una tangente. Por
otro lado, la recta \(t\) se parece a una tangente, pero la interseca
más de una vez.

Ejemplo

Encontrar la ecuación de la recta tangente a la parbola \(y=x^2\) en el
punto \(P=(1,1)\).

\begin{tcolorbox}[enhanced jigsaw, opacityback=0, left=2mm, opacitybacktitle=0.6, toprule=.15mm, bottomrule=.15mm, colbacktitle=quarto-callout-tip-color!10!white, titlerule=0mm, bottomtitle=1mm, title=\textcolor{quarto-callout-tip-color}{\faLightbulb}\hspace{0.5em}{Solución}, toptitle=1mm, colframe=quarto-callout-tip-color-frame, colback=white, breakable, rightrule=.15mm, arc=.35mm, leftrule=.75mm, coltitle=black]

Para dar la ecuación de la recta necesitamos conocer su pendiente. Como
sólo se cuenta con un punto, tomamos \(Q=(x,x^2)\) con \(x\neq 1\) y
estimamos la pendiente \(m\) con la de la recta secante \(PQ\), dada por

\[
m_{PQ}=\frac{x^2-1}{x-1}.
\]

En la siguiente tabla observamos los valores de \(m_{PQ}\) para
distintos valores de \(x\).

\begin{longtable}[]{@{}ll@{}}
\toprule\noalign{}
\(x\) & \(m_{PQ}\) \\
\midrule\noalign{}
\endhead
\bottomrule\noalign{}
\endlastfoot
0 & 1 \\
0.5 & 1.5 \\
0.9 & 1.9 \\
0.99 & 1.99 \\
0.999 & 1.999 \\
1.001 & 2.001 \\
1.01 & 2.01 \\
1.1 & 2.1 \\
1.5 & 2.5 \\
2 & 3 \\
\end{longtable}

Podemos ver que, conforme \(x\) está más cerca de \(1\), \(m_{PQ}\)
estará más cerca de 2. Conjeturamos entonces que la pendiente buscada es
\(2\). Decimos que la pendiente \(m\) es el \textbf{límite} de las
pendientes de las rectas secantes \(m_{PQ}\), y escribimos

\[
m=\lim_{Q\to P} m_{PQ} \quad \text{ y también }\quad    m=\lim_{x\to 1} \frac{x^2-1}{x-1}=2.
\]

Utilizando este valor de \(m\) podemos ver que la recta buscada tiene
ecuación

\[
y-1=2(x-1) \quad \text{ o equivalentemente }\quad  y=2x-1.
\]

\includegraphics[width=0.4\linewidth,height=\textheight,keepaspectratio]{index_files/mediabag/imagenes_cap2/2.1-parabola.pdf}

\end{tcolorbox}

\subsection{El problema de la
velocidad}\label{el-problema-de-la-velocidad}

En física resulta sencillo calcular una velocidad como el cociente entre
el espacio recorrido y el tiempo que se tarda en recorrerlo. Sin
embargo, en este caso se supone que la velocidad es constante durante el
experimento. ¿Qué sucede entonces cuando la velocidad varía? ¿Cómo puede
determinarse en cada instante de tiempo?

Ejemplo

Se deja caer una pelota desde la plataforma de observación de la Torre
CN en Toronto, a 450 metros por encima del nivel del suelo. Encontrar la
velocidad de la pelota una vez que transcurren cinco segundos.

\begin{tcolorbox}[enhanced jigsaw, opacityback=0, left=2mm, opacitybacktitle=0.6, toprule=.15mm, bottomrule=.15mm, colbacktitle=quarto-callout-tip-color!10!white, titlerule=0mm, bottomtitle=1mm, title=\textcolor{quarto-callout-tip-color}{\faLightbulb}\hspace{0.5em}{Solución}, toptitle=1mm, colframe=quarto-callout-tip-color-frame, colback=white, breakable, rightrule=.15mm, arc=.35mm, leftrule=.75mm, coltitle=black]

Si \(s(t)\) representa la distancia en metros recorrida por la pelota, a
tiempo \(t\) medido en segundos, entonces la \textbf{ley de Galileo}
establece que

\[
s(t)=4.9\, t^2
\]

Podemos comenzar calculando la velocidad promedio en distintos
intervalos de tiempo. Si \(h>0\), la velocidad promedio de la pelota en
el intervalo \([t, t+h]\) viene dada por

\[
\frac{s(t+h)-s(t)}{h}.
\]

En la siguiente tabla tenemos el cálculo de velocidades promedio para
intervalos de diferentes amplitudes.

\begin{longtable}[]{@{}ll@{}}
\toprule\noalign{}
Intervalo de tiempo & velocidad promedio \((m/s)\) \\
\midrule\noalign{}
\endhead
\bottomrule\noalign{}
\endlastfoot
{[}5,6{]} & 53.9 \\
{[}5,5.1{]} & 49.49 \\
{[}5,5.05{]} & 49.245 \\
{[}5,5.01{]} & 49.049 \\
{[}5,5.001{]} & 49.0049 \\
\end{longtable}

Conforme se achica el intervalo las velocidades promedio se aproximan a
49 \(m/s\). La \textbf{velocidad instantánea} se define como el límite
de las velocidades promedio cuando la amplitud de los intervalos se
achica. Esta idea se precisará mejor más adelante. Con esto, podemos
conjeturar que la velocidad instantánea de la pelota después de cinco
segundos es de 49 \(m/s\).

\end{tcolorbox}

Existe una relación entre los dos problemas presentados. Si dibujamos la
función del ejemplo anterior

\[
s(t)=4.9\, t^2
\]

y consideramos los puntos \(P=(a, 4.9\,a^2)\) y
\(Q=(a+h, 4.9\,(a+h)^2)\) entonces la pendiente de la secante \(PQ\)
está dada por

\[
m_{PQ}=\frac{4.9\,(a+h)^2- 4.9\,a^2}{(a+h)-a}=\frac{4.9\,(a+h)^2- 4.9\,a^2}{h}.
\]

\begin{center}
\includegraphics[width=0.7\linewidth,height=\textheight,keepaspectratio]{index_files/mediabag/imagenes_cap2/2.1-velocidades.pdf}
\end{center}

Luego, la velocidad instantánea representa la pendiente de la recta
tangente en el punto \(P\), que es el límite de las velocidades promedio
dadas por las pendientes de las rectas secantes \(m_{PQ}\) conforme
\(Q\) se aproxima a \(P\).

\hyperref[seccion_2.1]{↑ Volver al inicio de la sección}

\section{2.2. Límite de una función}\label{seccion_2.2}

En esta sección definiremos el concepto de \textbf{límite} de una
función, surgido al intentar calcular la recta tangente a una curva
dada. Consideremos \(f(x)=x^2-x+2\) y observemos los valores de \(f\)
cuando \(x\) está cerca de 2.

\begin{longtable}[]{@{}llll@{}}
\toprule\noalign{}
\textbf{\(x\)} & \textbf{\(f(x)\)} & \textbf{\(x\)} &
\textbf{\(f(x)\)} \\
\midrule\noalign{}
\endhead
\bottomrule\noalign{}
\endlastfoot
1.0 & 2.000000 & 3.0 & 8.000000 \\
1.5 & 2.750000 & 2.5 & 5.750000 \\
1.8 & 3.440000 & 2.2 & 4.640000 \\
1.9 & 3.710000 & 2.1 & 4.310000 \\
1.95 & 3.852500 & 2.05 & 4.152500 \\
1.99 & 3.970100 & 2.01 & 4.030100 \\
1.995 & 3.985025 & 2.005 & 4.015025 \\
1.999 & 3.997001 & 2.001 & 4.003001 \\
& & & \\
\end{longtable}

De la tabla anterior podemos deducir que cuando \(x\) está cerca de 2
(por la izquierda o la derecha), \(f(x)\) está cerca de 4.

\begin{center}
\includegraphics[width=0.3\linewidth,height=\textheight,keepaspectratio]{index_files/mediabag/imagenes_cap2/2.2-limite-parabola.pdf}
\end{center}

De hecho, se pueden acercar los valores de \(f(x)\) tanto a 4 como se
desee, conforme se elija \(x\) suficientemente cercano a 2.

Esto se expresa diciendo que: ``el límite de la función
\(f(x)=x^2-x+2\), cuando \(x\) tiende a 2, es igual a 4'' y se escribe

\[
\lim_{x\to 2} \,(x^2-x+2) =4.
\]

\begin{tcolorbox}[enhanced jigsaw, opacityback=0, left=2mm, opacitybacktitle=0.6, toprule=.15mm, bottomrule=.15mm, colbacktitle=quarto-callout-note-color!10!white, titlerule=0mm, bottomtitle=1mm, title=\textcolor{quarto-callout-note-color}{\faInfo}\hspace{0.5em}{Definición (Límite de una función)}, toptitle=1mm, colframe=quarto-callout-note-color-frame, colback=white, breakable, rightrule=.15mm, arc=.35mm, leftrule=.75mm, coltitle=black]

Decimos que el \textbf{límite} cuando \(x\) tiende a \(a\) de \(f(x)\)
es \(L\), y lo escribimos

\[
\lim_{x\to a} f(x)=L
\]

si los valores de \(f\) pueden acercarse arbitrariamente a \(L\)
conforme \(x\) se elija suficientemente cerca de \(a\), pero
\textbf{\(x\neq a\)}.

También podemos decir esto escribiendo \(f(x)\to L\) cuando \(x\to a\).

\end{tcolorbox}

\begin{tcolorbox}[enhanced jigsaw, opacityback=0, left=2mm, opacitybacktitle=0.6, toprule=.15mm, bottomrule=.15mm, colbacktitle=quarto-callout-caution-color!10!white, titlerule=0mm, bottomtitle=1mm, title=\textcolor{quarto-callout-caution-color}{\faFire}\hspace{0.5em}{Importante}, toptitle=1mm, colframe=quarto-callout-caution-color-frame, colback=white, breakable, rightrule=.15mm, arc=.35mm, leftrule=.75mm, coltitle=black]

Notar que en la definición consideramos \(x\neq a\). Esto significa que
\(f\) no necesita estar definida en \(x=a\) para calcular su límite para
\(x\to a\). Lo importante es que \(f\) esté definida \emph{cerca} de
\(a\).

\end{tcolorbox}

En los tres gráficos de abajo se cumple que \(\lim_{x\to a}f(x)=L\).

\begin{center}
\includegraphics[width=0.6\linewidth,height=\textheight,keepaspectratio]{index_files/mediabag/imagenes_cap2/2.2-limites.pdf}
\end{center}

Ejemplo

Conjeturar el valor del límite
\(\displaystyle \lim_{x\to 1} \frac{x-1}{x^2-1}\).

\begin{tcolorbox}[enhanced jigsaw, opacityback=0, left=2mm, opacitybacktitle=0.6, toprule=.15mm, bottomrule=.15mm, colbacktitle=quarto-callout-tip-color!10!white, titlerule=0mm, bottomtitle=1mm, title=\textcolor{quarto-callout-tip-color}{\faLightbulb}\hspace{0.5em}{Solución}, toptitle=1mm, colframe=quarto-callout-tip-color-frame, colback=white, breakable, rightrule=.15mm, arc=.35mm, leftrule=.75mm, coltitle=black]

Notemos que \(f\) no está definida en \(x=1\). Sin embargo, esto no
importa pues a los efectos de calcular el límite debemos considerar
\(x\) cerca de \(1\), pero \(x\neq 1\).

La siguiente tabla muestra los valores de \(f(x)\) cuando \(x\) está
cerca de \(1\), por izquierda y por derecha.

\begin{longtable}[]{@{}ll@{}}
\toprule\noalign{}
\textbf{\(x<1\)} & \textbf{\(f(x)\)} \\
\midrule\noalign{}
\endhead
\bottomrule\noalign{}
\endlastfoot
0.5 & 0.666667 \\
0.9 & 0.526316 \\
0.99 & 0.502513 \\
0.999 & 0.500250 \\
0.9999 & 0.500025 \\
& \\
\end{longtable}

\begin{longtable}[]{@{}ll@{}}
\toprule\noalign{}
\textbf{\(x>1\)} & \textbf{\(f(x)\)} \\
\midrule\noalign{}
\endhead
\bottomrule\noalign{}
\endlastfoot
1.5 & 0.400000 \\
1.1 & 0.476190 \\
1.019 & 0.497512 \\
1.001 & 0.499750 \\
1.0001 & 0.499975 \\
& \\
\end{longtable}

Esto sugiere que \(\operatorname{lím}_{x\to 1}f(x)=\frac{1}{2}\), como
comprobaremos luego.

Aquí podemos ver el gráfico de \(f\) junto al de la función \[
g(x)=
\begin{cases}
f(x) & \text{ si } x\neq 1,\\
2 & \text{ si } x=1.
\end{cases}
\]

\includegraphics[width=0.6\linewidth,height=\textheight,keepaspectratio]{index_files/mediabag/imagenes_cap2/2.2-limite-comparacion.pdf}

Ambas tiene el mismo límite para \(x\to 1\) ya que sólo difieren en
\(x=1\).

\end{tcolorbox}

Ejemplo

Estimar el valor de \(\displaystyle \lim_{t\to 0} f(t)\), siendo
\(\displaystyle f(t)=\frac{\sqrt{t^2+9}-3}{t^2}\).

\begin{tcolorbox}[enhanced jigsaw, opacityback=0, left=2mm, opacitybacktitle=0.6, toprule=.15mm, bottomrule=.15mm, colbacktitle=quarto-callout-tip-color!10!white, titlerule=0mm, bottomtitle=1mm, title=\textcolor{quarto-callout-tip-color}{\faLightbulb}\hspace{0.5em}{Solución}, toptitle=1mm, colframe=quarto-callout-tip-color-frame, colback=white, breakable, rightrule=.15mm, arc=.35mm, leftrule=.75mm, coltitle=black]

Nuevamente realizamos una tabla evaluando la función en valores de \(t\)
cercanos a \(0\), \(t\neq 0\).

\begin{longtable}[]{@{}ll@{}}
\toprule\noalign{}
\(t\) & \(f(t)\) \\
\midrule\noalign{}
\endhead
\bottomrule\noalign{}
\endlastfoot
\(\pm\,\) 1.0 & 0.16228 \\
\(\pm\,\) 0.5 & 0.16553 \\
\(\pm\,\) 0.1 & 0.16662 \\
\(\pm\,\) 0.05 & 0.16666 \\
\(\pm\,\) 0.01 & 0.16667 \\
\end{longtable}

Conforme \(t\) se acerca a \(0\), los valores de la función parecen
acercarse a \(1/6\), con lo que conjeturamos que

\[
 \lim_{t\to 0}\frac{\sqrt{t^2+9}-3}{t^2}=\frac{1}{6}.
\]

\end{tcolorbox}

¿Qué sucedería si eligiésemos valores de \(t\) más cercanos a \(0\)?

\begin{longtable}[]{@{}ll@{}}
\toprule\noalign{}
\(t\) & \(f(t)\) \\
\midrule\noalign{}
\endhead
\bottomrule\noalign{}
\endlastfoot
\(\pm\,\) 0.0005 & 0.16800 \\
\(\pm\,\) 0.0001 & 0.20000 \\
\(\pm\,\) 0.00005 & 0.00000 \\
\(\pm\,\) 0.00001 & 0.00000 \\
\end{longtable}

En general, si elegimos valores de \(t\) suficientemente chicos, habrá
un momento en que la calculadora nos arroje como resultado \(0\). Esto
no significa que el límite es cero, sino que la diferencia del numerador
es tan pequeña que hace que esta expresión parezca \(0\) para la
precisión que la calculadora maneja.

\begin{tcolorbox}[enhanced jigsaw, opacityback=0, left=2mm, opacitybacktitle=0.6, toprule=.15mm, bottomrule=.15mm, colbacktitle=quarto-callout-caution-color!10!white, titlerule=0mm, bottomtitle=1mm, title=\textcolor{quarto-callout-caution-color}{\faFire}\hspace{0.5em}{Importante}, toptitle=1mm, colframe=quarto-callout-caution-color-frame, colback=white, breakable, rightrule=.15mm, arc=.35mm, leftrule=.75mm, coltitle=black]

Si bien una tabla de valores muchas veces ayuda a entender el
comportamiento de una función al analizar límites, \textbf{no debe
usarse como herramienta formal}.

\end{tcolorbox}

Ejemplo

Investigar el valor de
\(\displaystyle \lim_{x\to 0}\operatorname{sen}\left(\frac{\pi}{x}\right)\).

\begin{tcolorbox}[enhanced jigsaw, opacityback=0, left=2mm, opacitybacktitle=0.6, toprule=.15mm, bottomrule=.15mm, colbacktitle=quarto-callout-tip-color!10!white, titlerule=0mm, bottomtitle=1mm, title=\textcolor{quarto-callout-tip-color}{\faLightbulb}\hspace{0.5em}{Solución}, toptitle=1mm, colframe=quarto-callout-tip-color-frame, colback=white, breakable, rightrule=.15mm, arc=.35mm, leftrule=.75mm, coltitle=black]

Nuevamente \(f(x)=\operatorname{sen}(\pi/x)\) no está definida en
\(x=0\). Si \(n\in\mathbb{N}\), observemos que

\[
f\left(\frac{1}{n}\right)=\operatorname{sen}(n\pi)=0.
\]

Como esto vale para cualquier \(n\in\mathbb{N}\) (más aún, es cierto
para todo \(n\) entero no nulo), podríamos estar tentados a suponer que

\[
\lim_{x\to 0} f(x)=0.
\]

Sin embargo, esto es erróneo pues podemos comprobar que \[
f\left(\frac{2}{4n+1}\right)=1 \quad \text { y }\quad f\left(\frac{2}{4n-1}\right)=-1 
\]

para cualquier \(n\in\mathbb{N}\).

Esto muestra que los valores de \(f(x)\) oscilan constantemente cuando
se eligen distintos valores de \(x\) cercanos a cero. En este caso, el
límite \textbf{no existe} por oscilación, como lo muestra la gráfica.

\includegraphics[width=0.6\linewidth,height=\textheight,keepaspectratio]{index_files/mediabag/imagenes_cap2/2.2-seno-oscilacion.pdf}

\end{tcolorbox}

Ejemplo

Encontrar el valor de
\(\displaystyle \lim_{x\to 0} \left(x^3+\frac{\cos(5x)}{10000}\right)\).

\begin{tcolorbox}[enhanced jigsaw, opacityback=0, left=2mm, opacitybacktitle=0.6, toprule=.15mm, bottomrule=.15mm, colbacktitle=quarto-callout-tip-color!10!white, titlerule=0mm, bottomtitle=1mm, title=\textcolor{quarto-callout-tip-color}{\faLightbulb}\hspace{0.5em}{Solución}, toptitle=1mm, colframe=quarto-callout-tip-color-frame, colback=white, breakable, rightrule=.15mm, arc=.35mm, leftrule=.75mm, coltitle=black]

Realizando una tabla de valores, tenemos lo siguiente

\begin{longtable}[]{@{}ll@{}}
\toprule\noalign{}
\(x\) & \(f(x)\) \\
\midrule\noalign{}
\endhead
\bottomrule\noalign{}
\endlastfoot
1 & 1.000028 \\
0.5 & 0.124920 \\
0.1 & 0.001088 \\
0.05 & 0.000222 \\
0.01 & 0.000101 \\
\end{longtable}

Esto parece indicar que el lmite es cero. Sin embargo, si consideramos
valores más pequeños de \(x\) obtenemos

\begin{longtable}[]{@{}ll@{}}
\toprule\noalign{}
\(x\) & \(f(x)\) \\
\midrule\noalign{}
\endhead
\bottomrule\noalign{}
\endlastfoot
0.005 & 0.00010009 \\
0.001 & 0.00010000 \\
\end{longtable}

En realidad, veremos más adelante que \(\lim_{x\to 0}\cos(5x)=1\), con
lo cual

\[
\lim_{x\to 0} \left(x^3+\frac{\cos(5x)}{10000}\right)=\frac{1}{10000}=0.0001.
\]

\end{tcolorbox}

La \textbf{función de Heaviside} \(H\) se define mediante la expresión

\[
H(t)=
\begin{cases}
0 & \text{ si } t<0,\\
1 & \text{ si } t\geq 1.
\end{cases}
\]

\begin{center}
\includegraphics[width=0.4\linewidth,height=\textheight,keepaspectratio]{index_files/mediabag/imagenes_cap2/2.2-heaviside.pdf}
\end{center}

Notemos que \(H(t)\) se aproxima a 1 cuando \(t\) se aproxima a cero por
la derecha, y se aproxima a 0 cuando \(t\) se acerca a cero por la
izquierda. Esto podemos escribirlo como

\[
\lim_{t\to 0^+} H(t)=1 \quad \text{ y } \quad \lim_{t\to 0^-} H(t)=0.
\]

No hay un único valor \(L\) al que \(H(t)\) se aproxime cuando \(t\)
tiende a cero, por lo que el límite \(\lim_{t\to 0} H(t)\) no existe.

\begin{tcolorbox}[enhanced jigsaw, opacityback=0, left=2mm, opacitybacktitle=0.6, toprule=.15mm, bottomrule=.15mm, colbacktitle=quarto-callout-note-color!10!white, titlerule=0mm, bottomtitle=1mm, title=\textcolor{quarto-callout-note-color}{\faInfo}\hspace{0.5em}{Definición (Límites laterales)}, toptitle=1mm, colframe=quarto-callout-note-color-frame, colback=white, breakable, rightrule=.15mm, arc=.35mm, leftrule=.75mm, coltitle=black]

\begin{itemize}
\tightlist
\item
  Decimos que el límite de \(f(x)\) cuando \(x\) se aproxima a \(a\) por
  la \textbf{izquierda} es \(L\) y lo denotamos
\end{itemize}

\[
\lim_{x\to a^-} f(x)=L
\]

si los valores \(f(x)\) pueden aproximarse al número \(L\) tanto como se
desee, siempre que se elija \(x\) lo suficientemente cerca de \(a\),
pero \textbf{\(x<a\)}.

\begin{itemize}
\tightlist
\item
  Análogamente, decimos que el límite de \(f(x)\) cuando \(x\) tiende a
  \(a\) por la \textbf{derecha} es \(L\) y lo expresamos
\end{itemize}

\[
\lim_{x\to a^+} f(x)=L
\]

si \(f(x)\) puede aproximarse al número \(L\) tanto como se quiera,
siempre que elijamos \(x\) lo suficientemente cercano a \(a\), pero
\textbf{\(x>a\)}.

\end{tcolorbox}

Al comparar la definición de límite con la de límites laterales, tenemos
que

\[
\lim_{x\to a} f(x)=L \quad \Longleftrightarrow \quad \lim_{x\to a^-} f(x)=L \quad \text{ y } \quad \lim_{x\to a^+} f(x)=L.
\]

Ejemplo

Dar los valores de los siguientes límites (si existen) observando la
gráfica.

\begin{itemize}
\tightlist
\item
  \(\lim_{x\to 2^-}\,g(x)\)
\item
  \(\lim_{x\to 2^+}\,g(x)\)
\item
  \(\lim_{x\to 2}\,g(x)\)
\end{itemize}

\begin{itemize}
\tightlist
\item
  \(\lim_{x\to 5^-}\, g(x)\)
\item
  \(\lim_{x\to 5^+}\, g(x)\)
\item
  \(\lim_{x\to 5}\, g(x)\)
\end{itemize}

\begin{center}
\includegraphics[width=0.5\linewidth,height=\textheight,keepaspectratio]{index_files/mediabag/imagenes_cap2/2.2-ejemplo-limites.pdf}
\end{center}

\begin{tcolorbox}[enhanced jigsaw, opacityback=0, left=2mm, opacitybacktitle=0.6, toprule=.15mm, bottomrule=.15mm, colbacktitle=quarto-callout-tip-color!10!white, titlerule=0mm, bottomtitle=1mm, title=\textcolor{quarto-callout-tip-color}{\faLightbulb}\hspace{0.5em}{Solución}, toptitle=1mm, colframe=quarto-callout-tip-color-frame, colback=white, breakable, rightrule=.15mm, arc=.35mm, leftrule=.75mm, coltitle=black]

Observando el gráfico vemos que \(g\) se aproxima a 3 cuando \(x\) se
acerca a 2 por la izquierda, y a 1 cuando \(x\) se aproxima a 2 por la
derecha. En consecuencia

\[
\lim_{x\to 2^-}g(x)=3\quad \text{ y } \quad \lim_{x\to 2^+} g(x)=1.
\]

Como los límites laterales son distintos, concluimos que el límite
\(\lim_{x\to 2} g(x)\) no existe.

Por otra parte, también podemos ver que \(g\) se aproxima a 2 cuando
\(x\) se acerca a 5, tanto por la derecha como por la izquierda, con lo
cual

\[
\lim_{x\to 5^-} g(x)=2=\lim_{x\to 5^+} g(x),
\]

y en conclusión \(\lim_{x\to 5} g(x)=2\). Notar que \(g(5)=1\neq 2\),
sin embargo esto no infliuye en el cálculo del límite pues al hacerlo
consideramos \(x\neq 5\).

\end{tcolorbox}

\subsection{Límites infinitos}\label{luxedmites-infinitos}

Vamos a investigar a continuación el comportamiento de la función

\[
\displaystyle f(x)=\frac{1}{x^2}
\]

cuando \(x\) se aproxima a cero. Cuando \(x\) tiende a cero, tanto por
la derecha como por la izquierda, \(1/x^2\) se hace cada vez más grande.
Más aún, podemos hacer que \(1/x^2\) sea tan grande como deseemos,
siempre que elijamos \(x\) lo suficientemente cerca de cero. En la
siguiente tabla se muestran algunos valores de \(f\) cuando \(x\) toma
valores cercanos a cero.

\begin{longtable}[]{@{}ll@{}}
\toprule\noalign{}
\(x\) & \(f(x)\) \\
\midrule\noalign{}
\endhead
\bottomrule\noalign{}
\endlastfoot
\(\pm \,\) 1 & 1 \\
\(\pm \,\) 0.5 & 4 \\
\(\pm \,\) 0.2 & 25 \\
\(\pm \,\) 0.1 & 100 \\
\(\pm \,\) 0.05 & 400 \\
\(\pm \,\) 0.01 & 10000 \\
\(\pm \,\) 0.001 & 1000000 \\
\end{longtable}

Este comportamiento se puede observar también con el gráfico de la
función.

\begin{center}
\includegraphics[width=0.4\linewidth,height=\textheight,keepaspectratio]{index_files/mediabag/imagenes_cap2/2.2-limite-infinito.pdf}
\end{center}

En este caso, decimos que \(\lim_{x\to 0} f(x)\) no existe \textbf{por
infinitud}. Para indicar este tipo de comportamiento, escribiremos

\[
\lim_{x\to 0} \frac{1}{x^2}=\infty.
\]

\begin{tcolorbox}[enhanced jigsaw, opacityback=0, left=2mm, opacitybacktitle=0.6, toprule=.15mm, bottomrule=.15mm, colbacktitle=quarto-callout-important-color!10!white, titlerule=0mm, bottomtitle=1mm, title=\textcolor{quarto-callout-important-color}{\faExclamation}\hspace{0.5em}{Importante}, toptitle=1mm, colframe=quarto-callout-important-color-frame, colback=white, breakable, rightrule=.15mm, arc=.35mm, leftrule=.75mm, coltitle=black]

Esto no significa que consideremos a \(\infty\) como un número ni que el
límite dado exista. La notación sólamente establece de qué manera el
límite no existe.

\end{tcolorbox}

\begin{tcolorbox}[enhanced jigsaw, opacityback=0, left=2mm, opacitybacktitle=0.6, toprule=.15mm, bottomrule=.15mm, colbacktitle=quarto-callout-note-color!10!white, titlerule=0mm, bottomtitle=1mm, title=\textcolor{quarto-callout-note-color}{\faInfo}\hspace{0.5em}{Definición (Límites infinitos)}, toptitle=1mm, colframe=quarto-callout-note-color-frame, colback=white, breakable, rightrule=.15mm, arc=.35mm, leftrule=.75mm, coltitle=black]

Sea \(f\) una función definida en ambos lados de \(a\), excepto quizás
en \(a\).

\begin{itemize}
\tightlist
\item
  La expresión
\end{itemize}

\[
\lim_{x\to a} f(x)=\infty
\]

significa que \(f(x)\) puede hacerse arbitrariamente grande conforme
\(x\) se elija lo suficientemente cerca de \(a\), pero \(x\neq a\).

\begin{itemize}
\tightlist
\item
  Por otro lado
\end{itemize}

\[
\lim_{x\to a} f(x)=-\infty
\]

quiere decir que \(f(x)\) puede hacerse arbitrariamente grande y
negativa siempre que \(x\) se elija lo suficientemente cerca de \(a\),
pero \(x\neq a\).

\end{tcolorbox}

Otra notación para \(\lim_{x\to  a} f(x)=\infty\) es: \(f(x)\to \infty\)
cuando \(x\to a\), y se lee

\begin{itemize}
\tightlist
\item
  ``el límite cuando \(x\to a\) de \(f(x)\) es infinito'', o
\item
  ``\(f(x)\) se vuelve infinita cuando \(x\) se aproxima a \(a\)'', o
  bien
\item
  ``\(f(x)\) se incrementa sin límite cuando \(x\) tiende a \(a\)''.
\end{itemize}

\begin{center}
\includegraphics[width=0.4\linewidth,height=\textheight,keepaspectratio]{index_files/mediabag/imagenes_cap2/2.2-limite-infinito2.pdf}
\end{center}

Para \(\lim_{x\to a} f(x)=-\infty\) podemos escribir:
\(f(x)\to -\infty\) cuando \(x\to a\), y leerlo

\begin{itemize}
\tightlist
\item
  ``el límite cuando \(x\to a\) de \(f(x)\) es el infinito negativo'', o
\item
  ``\(f(x)\) decrece sin cota cuando \(x\) se aproxima a \(a\)''.
\end{itemize}

\begin{center}
\includegraphics[width=0.4\linewidth,height=\textheight,keepaspectratio]{index_files/mediabag/imagenes_cap2/2.2-limite-infinito3.pdf}
\end{center}

Definiciones similares pueden darse para

\[
\lim_{x\to a^-} f(x)=\infty \quad \text{ y }\quad \lim_{x\to a^+} f(x)=\infty.
\]

\begin{center}
\includegraphics[width=0.7\linewidth,height=\textheight,keepaspectratio]{index_files/mediabag/imagenes_cap2/2.2-limite-infinito4.pdf}
\end{center}

Y también para

\[
\lim_{x\to a^-} f(x)=-\infty \quad \text{ y }\quad \lim_{x\to a^+} f(x)=-\infty.
\]

\begin{center}
\includegraphics[width=0.7\linewidth,height=\textheight,keepaspectratio]{index_files/mediabag/imagenes_cap2/2.2-limite-infinito5.pdf}
\end{center}

\begin{tcolorbox}[enhanced jigsaw, opacityback=0, left=2mm, opacitybacktitle=0.6, toprule=.15mm, bottomrule=.15mm, colbacktitle=quarto-callout-note-color!10!white, titlerule=0mm, bottomtitle=1mm, title=\textcolor{quarto-callout-note-color}{\faInfo}\hspace{0.5em}{Definición (Asíntota vertical)}, toptitle=1mm, colframe=quarto-callout-note-color-frame, colback=white, breakable, rightrule=.15mm, arc=.35mm, leftrule=.75mm, coltitle=black]

La recta \(x=a\) se llama \textbf{asíntota vertical} de la curva
\(y=f(x)\) si al menos uno de los siguientes enunciados es verdadero

\begin{itemize}
\tightlist
\item
  \(\lim_{x\to a}\,f(x)=\infty\)
\item
  \(\lim_{x\to a^-}\,f(x)=\infty\)
\item
  \(\lim_{x\to a^+}\,f(x)=\infty\)
\end{itemize}

\begin{itemize}
\tightlist
\item
  \(\lim_{x\to a}\,f(x)=-\infty\)
\item
  \(\lim_{x\to a^-}\,f(x)=-\infty\)
\item
  \(\lim_{x\to a^+}\,f(x)=-\infty\)
\end{itemize}

\end{tcolorbox}

Ejemplo

Determinar

\[
\lim_{x\to 3^+} \frac{2x}{x-3} \quad \text{ y } \quad  \lim_{x\to 3^-} \frac{2x}{x-3}.
\]

\begin{tcolorbox}[enhanced jigsaw, opacityback=0, left=2mm, opacitybacktitle=0.6, toprule=.15mm, bottomrule=.15mm, colbacktitle=quarto-callout-tip-color!10!white, titlerule=0mm, bottomtitle=1mm, title=\textcolor{quarto-callout-tip-color}{\faLightbulb}\hspace{0.5em}{Solución}, toptitle=1mm, colframe=quarto-callout-tip-color-frame, colback=white, breakable, rightrule=.15mm, arc=.35mm, leftrule=.75mm, coltitle=black]

Veamos primero \(\lim_{x\to 3^+} \frac{2x}{x-3}\). Notemos que el
denominador se aproxima a cero y el numerador a \(6\neq 0\), con lo que
el límite resultará infinito. Para saber si es \(\infty\) o \(-\infty\),
observemos que \(x\to 3^+\) implica que \(x>3\), lo que equivale a que
\(x-3>0\). Entonces al ser numerador y denominador positivos para todo
\(x>3\) resulta

\[
\lim_{x\to 3^+} \frac{2x}{x-3}=\infty.
\]

Para el otro caso, procediendo de forma similar y notando que
\(x\to 3^-\) implica que \(x-3<0\) para todo \(x<3\) obtenemos

\[
\lim_{x\to 3^-} \frac{2x}{x-3}=-\infty.
\]

\end{tcolorbox}

Ejemplo

Hallar las asíntotas verticales de \(f(x)=\tan x\).

\begin{tcolorbox}[enhanced jigsaw, opacityback=0, left=2mm, opacitybacktitle=0.6, toprule=.15mm, bottomrule=.15mm, colbacktitle=quarto-callout-tip-color!10!white, titlerule=0mm, bottomtitle=1mm, title=\textcolor{quarto-callout-tip-color}{\faLightbulb}\hspace{0.5em}{Solución}, toptitle=1mm, colframe=quarto-callout-tip-color-frame, colback=white, breakable, rightrule=.15mm, arc=.35mm, leftrule=.75mm, coltitle=black]

Recordemos que \(\tan x=\frac{\operatorname{sen} x}{\cos x}\), por lo
tanto tendremos posibles asíntotas verticales cuando \(\cos x=0\). Esto
ocurre, por ejemplo, cuando \(x=\pi/2\). En este caso, si \(x\) está
cerca de \(\pi/2\) y \(x>\pi/2\), \(\cos x\) se aproxima a cero con
valores negativos, mientras que \(\operatorname{sen} x\) se aproxima a
1. Con lo cual

\[
\lim_{x\to (\pi/2)^+} \tan x=-\infty.
\]

Por otra parte, si \(x\) está cerca de \(\pi/2\) y \(x<\pi/2\),
\(\cos x\) se aproxima a cero con valores negativos, y
\(\operatorname{sen} x\) sigue aproximándose a 1. Por lo tanto

\[
\lim_{x\to (\pi/2)^-} \tan x=\infty.
\]

Estos cálculos muestran que \(x=\pi/2\) es una asíntota vertical de
\(f\). De manera similar podemos ver que todas las asíntotas de \(f\)
son de la forma

\[
x=\left(2n+1\right)\frac{\pi}{2} \quad \text{ con }\quad n\in\mathbb{Z}.
\]

\end{tcolorbox}

\hyperref[seccion_2.2]{↑ Volver al inicio de la sección}

\section{2.3. Leyes de los límites}\label{seccion_2.3}

Hasta ahora hemos conjeturado el valor de ciertos límites a partir de
una tabla o de observar el comportamiento de sus gráficos.

En esta sección calcularemos los límites utilizando propiedades, lo que
nos conducirá a los valores exactos, cuando estos límites existen.

\begin{theorem}[]\protect\hypertarget{thm-leyes-limites}{}\label{thm-leyes-limites}

Propiedad (Leyes de los límites)

Sea \(c\) una constante y supongamos que los límites \[
\lim_{x\to a} f(x) \quad \text{ y } \quad \lim_{x\to a}g(x)
\]

existen. Entonces

\begin{enumerate}
\def\labelenumi{\arabic{enumi}.}
\item
  \(\lim_{x\to a}\,[f(x)+g(x)]=\lim_{x\to a}f(x)+\lim_{x\to a} g(x)\).
\item
  \(\lim_{x\to a}\, [f(x)-g(x)]=\lim_{x\to a} f(x)-\lim_{x\to a} g(x)\).
\item
  \(\lim_{x\to a} \,[cf(x)]=c\lim_{x\to a} f(x)\).
\item
  \(\lim_{x\to a} [f(x)g(x)]=\lim_{x\to a} f(x) \lim_{x\to a} g(x)\).
\item
  \(\displaystyle \lim_{x\to a} \frac{f(x)}{g(x)}= \frac{\lim_{x\to a} f(x)}{\lim_{x\to a} g(x)}\),
  siempre que \(\lim_{x\to a} g(x)\neq 0\).
\end{enumerate}

\end{theorem}

Aplicando la ley \textbf{(4)} con \(f=g\) repetidas veces, obtenemos

\begin{enumerate}
\def\labelenumi{\arabic{enumi}.}
\setcounter{enumi}{5}
\tightlist
\item
  \(\lim_{x\to a}\,(f(x))^n=\left(\lim_{x\to a}f(x)\right)^n\), con
  \(n\) entero positivo.
\end{enumerate}

Dos límites especiales que vale la pena mencionar son

\begin{enumerate}
\def\labelenumi{\arabic{enumi}.}
\setcounter{enumi}{6}
\item
  \(\lim_{x\to a}\, c=c\).
\item
  \(\lim_{x\to a}\, x=a\).
\end{enumerate}

Combinando las leyes \textbf{(6)} y \textbf{(8)} obtenemos

\begin{enumerate}
\def\labelenumi{\arabic{enumi}.}
\setcounter{enumi}{8}
\tightlist
\item
  \(\lim_{x\to a}\, x^n=a^n\).
\end{enumerate}

Para el caso de races, si \(n\in\mathbb{N}\) tenemos que

\begin{enumerate}
\def\labelenumi{\arabic{enumi}.}
\setcounter{enumi}{9}
\item
  \(\lim_{x\to a}\, \sqrt[n]{x}=\sqrt[n]{a}\quad\) (\(a>0\) si \(n\) es
  par).
\item
  \(\lim_{x\to a}\, \sqrt[n]{f(x)}=\sqrt[n]{\lim_{x\to a}f(x)}\quad\)
  (\(\lim_{x\to a} f(x)>0\) si \(n\) es par).
\end{enumerate}

Ejemplo

Evaluar los siguientes lmites y justificar cada paso.

\begin{enumerate}
\def\labelenumi{\arabic{enumi}.}
\item
  \(\lim_{x\to 5}\, (2x^2-3x+4)\).
\item
  \(\displaystyle \lim_{x\to -2}\, \frac{x^3+2x^2-1}{5-3x}\).
\end{enumerate}

\begin{tcolorbox}[enhanced jigsaw, opacityback=0, left=2mm, opacitybacktitle=0.6, toprule=.15mm, bottomrule=.15mm, colbacktitle=quarto-callout-tip-color!10!white, titlerule=0mm, bottomtitle=1mm, title=\textcolor{quarto-callout-tip-color}{\faLightbulb}\hspace{0.5em}{Solución}, toptitle=1mm, colframe=quarto-callout-tip-color-frame, colback=white, breakable, rightrule=.15mm, arc=.35mm, leftrule=.75mm, coltitle=black]

Comencemos con el primer límite. Utilizando la ley \textbf{(6)} con
\(n=2\) junto a la ley \textbf{(8)} obtenemos que

\[
\lim_{x\to 5} x^2 = 5^2 =25,
\]

y por la ley \textbf{(3)}

\[
\lim_{x\to 5} 2x^2 = 2\left(\lim_{x\to 5} x^2\right) =2\cdot 25=50.
\]

Por otra parte, combinando las leyes \textbf{(3)} y \textbf{(8)}

\[
\lim_{x\to 5} 3x = 3\left(\lim_{x\to 5} x\right) =3\cdot 5=15.
\]

Finalmente, por las leyes \textbf{(1)}, \textbf{(2)} y \textbf{(7)}

\[
\lim_{x\to 5}\, (2x^2-3x+4)=\lim_{x\to 5}\, 2x^2 - \lim_{x\to 5}\, 3x+ \lim_{x\to 5} 4 = 50-15+4=39.
\]

Para calcular el segundo límite, notemos que se trata de una función
racional, es decir, cociente de polinomios. Procediendo como en el caso
anterior podemos ver que

\[
\lim_{x\to -2}\, (x^3+2x^2-1) = (-2)^3+2(-2)^2-1=-1
\]

y también

\[
\lim_{x\to -2}\, (5-3x) = 5-3(-2)=11.
\]

Entonces, utilizando la ley \textbf{(5)} concluimos que

\[
\lim_{x\to -2}\, \frac{x^3+2x^2-1}{5-3x}=\frac{\lim_{x\to -2}\, (x^3+2x^2-1)}{\lim_{x\to -2}\, (5-3x)}=-\frac{1}{11}.
\]

\end{tcolorbox}

El procedimiento utilizado en el ejemplo anterior puede repetirse para
cualquier función polinómica o racional, lo que nos conduce a la
siguiente propiedad.

\phantomsection\label{teo-sust-directa}
Propiedad de sustitución directa

Si \(f\) es un \textbf{polinomio} o una \textbf{función racional} y
\(a\) está en el dominio de \(f\), entonces

\[
\lim_{x\to a}\, f(x)=f(a).
\]

Las funciones que cumplen la igualdad de arriba se dicen
\textbf{continuas} en \(a\). Este concepto lo trataremos más adelante.

\phantomsection\label{ejemplo-cuadratica-lineal}
Ejemplo

Hallar \(\displaystyle \lim_{x\to 1} \frac{x^2-1}{x-1}\).

\begin{tcolorbox}[enhanced jigsaw, opacityback=0, left=2mm, opacitybacktitle=0.6, toprule=.15mm, bottomrule=.15mm, colbacktitle=quarto-callout-tip-color!10!white, titlerule=0mm, bottomtitle=1mm, title=\textcolor{quarto-callout-tip-color}{\faLightbulb}\hspace{0.5em}{Solución}, toptitle=1mm, colframe=quarto-callout-tip-color-frame, colback=white, breakable, rightrule=.15mm, arc=.35mm, leftrule=.75mm, coltitle=black]

Recordemos que a los efectos de operar con límites importan los valores
de \(f(x)\) para \(x\) cerca de \(1\), pero \(x\neq 1\). Si tomamos
\(x\neq 1\) entonces podemos escribir

\[
\frac{x^2-1}{x-1}=\frac{(x-1)(x+1)}{x-1}=x+1.
\]

La igualdad de arriba es cierta siempre que \(x\neq 1\), por lo tanto

\[
\lim_{x\to 1} \frac{x^2-1}{x-1}=\lim_{x\to 1}\, (x+1)
\]

pero como \(x+1\) es un polinomio, la
\hyperref[teo-sust-directa]{propiedad de sustitución directa} nos da

\[
\lim_{x\to 1}\, (x+1) =2,
\]

con lo cual

\[
\lim_{x\to 1} \frac{x^2-1}{x-1}=2.
\]

\end{tcolorbox}

Ejemplo

Encontrar \(\displaystyle \lim_{x\to 1} g(x)\), siendo

\[
g(x)=
\begin{cases}
x+1 & \text{ si } x\neq 1,\\
\pi & \text{ si } x=1.
\end{cases}
\]

\begin{tcolorbox}[enhanced jigsaw, opacityback=0, left=2mm, opacitybacktitle=0.6, toprule=.15mm, bottomrule=.15mm, colbacktitle=quarto-callout-tip-color!10!white, titlerule=0mm, bottomtitle=1mm, title=\textcolor{quarto-callout-tip-color}{\faLightbulb}\hspace{0.5em}{Solución}, toptitle=1mm, colframe=quarto-callout-tip-color-frame, colback=white, breakable, rightrule=.15mm, arc=.35mm, leftrule=.75mm, coltitle=black]

En este caso sabemos que \(g(x)=x+1\) para todo \(x\neq 1\). Por lo
tanto

\[
\lim_{x\to 1} g(x)=\lim_{x\to 1}\, (x+1)=2,
\]

por la propíedad de sustitución directa.

\end{tcolorbox}

Ejemplo

Calcular \(\displaystyle \lim_{h\to 0} \frac{(3+h)^2-9}{h}\).

\begin{tcolorbox}[enhanced jigsaw, opacityback=0, left=2mm, opacitybacktitle=0.6, toprule=.15mm, bottomrule=.15mm, colbacktitle=quarto-callout-tip-color!10!white, titlerule=0mm, bottomtitle=1mm, title=\textcolor{quarto-callout-tip-color}{\faLightbulb}\hspace{0.5em}{Solución}, toptitle=1mm, colframe=quarto-callout-tip-color-frame, colback=white, breakable, rightrule=.15mm, arc=.35mm, leftrule=.75mm, coltitle=black]

Procedemos como antes. Primero consideramos \(h\neq 0\) y escribimos

\[
\frac{(3+h)^2-9}{h}=\frac{9+6h+h^2-9}{h}=\frac{h(6+h)}{h}=6+h.
\]

Entonces

\[
\lim_{h\to 0} \frac{(3+h)^2-9}{h} =\lim_{h\to 0}\, (6+h) = 6.
\]

\end{tcolorbox}

Ejemplo

Encontrar \(\displaystyle \lim_{t\to 0} \frac{\sqrt{t^2+9}-3}{t^2}\).

\begin{tcolorbox}[enhanced jigsaw, opacityback=0, left=2mm, opacitybacktitle=0.6, toprule=.15mm, bottomrule=.15mm, colbacktitle=quarto-callout-tip-color!10!white, titlerule=0mm, bottomtitle=1mm, title=\textcolor{quarto-callout-tip-color}{\faLightbulb}\hspace{0.5em}{Solución}, toptitle=1mm, colframe=quarto-callout-tip-color-frame, colback=white, breakable, rightrule=.15mm, arc=.35mm, leftrule=.75mm, coltitle=black]

Consideremos \(t\neq 0\). Multiplicando y dividiendo por
\(\sqrt{t^2+9}+3\) obtenemos

\[
\frac{\sqrt{t^2+9}-3}{t^2}=\frac{\sqrt{t^2+9}-3}{t^2}\frac{\sqrt{t^2+9}+3}{\sqrt{t^2+9}+3}=\frac{\left(\sqrt{t^2+9}\right)^2-3^2}{t^2(\sqrt{t^2+9}+3)}=\frac{t^2}{t^2(\sqrt{t^2+9}+3)}=\frac{1}{\sqrt{t^2+9}+3}.
\]

Con esto

\[
\lim_{t\to 0} \frac{\sqrt{t^2+9}-3}{t^2} =\lim_{t\to 0} \frac{1}{\sqrt{t^2+9}+3},
\]

por lo que es suficiente con calcular este último límite. Para ello, por
la \hyperref[teo-sust-directa]{propiedad de sustitución directa}
observemos primero que

\[
\lim_{t\to 0} (t^2+9)=9>0
\]

y utilizando la ley \textbf{(11)}

\[
\lim_{t\to 0} \sqrt{t^2+9}=\sqrt{9}=3.
\]

Esto nos indica que

\[
\lim_{t\to 0} \left(\sqrt{t^2+9}+3\right)=3+3=6.
\]

Finalmente, utilizando la ley \textbf{(5)}

\[
\lim_{t\to 0} \frac{1}{\sqrt{t^2+9}+3}=\frac{\lim_{t\to 0} 1}{\lim_{t\to 0}\left(\sqrt{t^2+9}+3\right)}=\frac{1}{6}.
\]

\end{tcolorbox}

A veces resulta útil estudiar un límite a partir de los correspondientes
límites laterales. Recordemos el siguiente resultado de la
\hyperref[seccion_2.2]{Sección 2.2}.

Relación entre límite y límites laterales

\[
\lim_{x\to a} f(x) = L \quad \Longleftrightarrow \quad \lim_{x\to a^-} f(x) = L\quad \text{ y } \quad  \lim_{x\to a^+} f(x) = L.
\]

Las leyes de los límites vistas anteriormente \textbf{también valen}
cambiando \(x\to a\) por \(x\to a^+\) o \(x\to a^-\).

Ejemplo

Demostrar que \(\displaystyle \lim_{x\to 0} |x|=0\).

\begin{tcolorbox}[enhanced jigsaw, opacityback=0, left=2mm, opacitybacktitle=0.6, toprule=.15mm, bottomrule=.15mm, colbacktitle=quarto-callout-tip-color!10!white, titlerule=0mm, bottomtitle=1mm, title=\textcolor{quarto-callout-tip-color}{\faLightbulb}\hspace{0.5em}{Solución}, toptitle=1mm, colframe=quarto-callout-tip-color-frame, colback=white, breakable, rightrule=.15mm, arc=.35mm, leftrule=.75mm, coltitle=black]

Vale la pena recordar la definición de valor absoluto vista en la
Sección 1.1

\[
|x|=
\begin{cases}
x & \text{ si } x\geq 0,\\
-x & \text{ si } x<0.
\end{cases}
\]

Como la definición de \(|x|\) depende de si nos encontramos a la derecha
o a la izquierda de cero, debemos analizar los límites laterales.

\[
\lim_{x\to 0^+} |x|=\lim_{x\to 0^+} x = 0
\]

y también

\[
\lim_{x\to 0^-} |x|=\lim_{x\to 0^-} (-x) = 0.
\]

Como los límites laterales existen y coinciden, concluimos que

\[
\lim_{x\to 0} |x|=0.
\]

\end{tcolorbox}

Ejemplo

Comprobar que el límite \(\displaystyle \lim_{x\to 0} \frac{|x|}{x}\) no
existe.

\begin{tcolorbox}[enhanced jigsaw, opacityback=0, left=2mm, opacitybacktitle=0.6, toprule=.15mm, bottomrule=.15mm, colbacktitle=quarto-callout-tip-color!10!white, titlerule=0mm, bottomtitle=1mm, title=\textcolor{quarto-callout-tip-color}{\faLightbulb}\hspace{0.5em}{Solución}, toptitle=1mm, colframe=quarto-callout-tip-color-frame, colback=white, breakable, rightrule=.15mm, arc=.35mm, leftrule=.75mm, coltitle=black]

Nuevamente estudiamos los límites laterales, ya que aparece \(|x|\) en
la expresión de la función.

\[
\lim_{x\to 0^+} \frac{|x|}{x}=\lim_{x\to 0^+} \frac{x}{x} = \lim_{x\to 0^+} 1=1.
\]

Por otro lado,

\[
\lim_{x\to 0^-} \frac{|x|}{x}=\lim_{x\to 0^-} \frac{-x}{x} = \lim_{x\to 0^-} -1=-1.
\]

Es decir, los límites laterales existen pero son distintos. Entonces
concluimos que el límite dado no existe.

\end{tcolorbox}

Ejemplo

Sea \(f\) la función definida por \[
f(x)=
\begin{cases}
\sqrt{x-4} & \text{ si } x>4,\\
 8-2x & \text{ si } x<4.
\end{cases}
\]

Determinar, si existe, \(\lim_{x\to 4} f(x)\).

\begin{tcolorbox}[enhanced jigsaw, opacityback=0, left=2mm, opacitybacktitle=0.6, toprule=.15mm, bottomrule=.15mm, colbacktitle=quarto-callout-tip-color!10!white, titlerule=0mm, bottomtitle=1mm, title=\textcolor{quarto-callout-tip-color}{\faLightbulb}\hspace{0.5em}{Solución}, toptitle=1mm, colframe=quarto-callout-tip-color-frame, colback=white, breakable, rightrule=.15mm, arc=.35mm, leftrule=.75mm, coltitle=black]

La función \(f\) está definida de formas distintas a la derecha e
izquierda de \(x=4\), con lo cual analizamos los límites laterales.

\[
\lim_{x\to 4^-} f(x)=\lim_{x\to 4^-} (8-2x) = 8-2\cdot 4=0
\]

y, por la ley \textbf{(11)}

\[
\lim_{x\to 4^+} f(x)=\lim_{x\to 4^+} \sqrt{x-4}= 8-2\cdot 4=0.
\]

Como los límites laterales existen y son iguales, concluímos que
\(\displaystyle \lim_{x\to 4} f(x)=0\).

\end{tcolorbox}

Teorema

Sean \(f\) y \(g\) dos funciones tales que \(f(x)\leq g(x)\) cuando
\(x\) está cerca de \(a\) (excepto quizás en \(a\)) y los límites \[
\lim_{x\to a} f(x) \quad \text{ y } \quad \lim_{x\to a} g(x)
\]

existen, entonces se cumple que

\[
\lim_{x\to a} f(x)\leq \lim_{x\to a} g(x).
\]

\phantomsection\label{teo-compresion}
Teorema de la compresión

Sean \(f\), \(g\) y \(h\) tres funciones tales que \[
f(x)\leq  g(x)\leq h(x)
\]

para \(x\) cerca de \(a\) (excepto, quizás en \(x=a\)) y además

\[
\lim_{x\to a} f(x)=\lim_{x\to a} h(x)=L.
\] Entonces \[
\lim_{x\to a} g(x)=L.
\]

\begin{center}
\includegraphics[width=0.5\linewidth,height=\textheight,keepaspectratio]{index_files/mediabag/imagenes_cap2/2.3-teo-sandwich.pdf}
\end{center}

Ejemplo

Demostrar que
\(\displaystyle \lim_{x\to 0} x^2\operatorname{sen}\left(\frac{1}{x}\right)=0\).

\begin{tcolorbox}[enhanced jigsaw, opacityback=0, left=2mm, opacitybacktitle=0.6, toprule=.15mm, bottomrule=.15mm, colbacktitle=quarto-callout-tip-color!10!white, titlerule=0mm, bottomtitle=1mm, title=\textcolor{quarto-callout-tip-color}{\faLightbulb}\hspace{0.5em}{Solución}, toptitle=1mm, colframe=quarto-callout-tip-color-frame, colback=white, breakable, rightrule=.15mm, arc=.35mm, leftrule=.75mm, coltitle=black]

Sea \(x\neq 0\). Observemos primero que

\[
-1\leq \operatorname{sen}\left(\frac{1}{x}\right)\leq 1
\]

dado que la función seno toma valores entre \(-1\) y \(1\).
Multiplicando cada miembro por \(x^2\geq 0\) resulta

\[
-x^2\leq x^2\operatorname{sen}\left(\frac{1}{x}\right)\leq x^2
\]

para cualquier \(x\neq 0\). Como

\[
\lim_{x\to 0} x^2=0 \quad \text{ y }\quad \lim_{x\to 0} (-x^2)=0
\]

utilizando el \hyperref[teo-compresion]{teorema de la compresión}
podemos concluir que

\[
\lim_{x\to 0} x^2\operatorname{sen}\left(\frac{1}{x}\right)=0.
\]

\end{tcolorbox}

\hyperref[seccion_2.3]{↑ Volver al inicio de la sección}

\section{2.5. Continuidad}\label{seccion_2.5}

\begin{tcolorbox}[enhanced jigsaw, opacityback=0, left=2mm, opacitybacktitle=0.6, toprule=.15mm, bottomrule=.15mm, colbacktitle=quarto-callout-note-color!10!white, titlerule=0mm, bottomtitle=1mm, title=\textcolor{quarto-callout-note-color}{\faInfo}\hspace{0.5em}{Definición (continuidad)}, toptitle=1mm, colframe=quarto-callout-note-color-frame, colback=white, breakable, rightrule=.15mm, arc=.35mm, leftrule=.75mm, coltitle=black]

Una función \(f\) es \textbf{continua} en \(a\) si se cumple que

\[
\lim_{x\to a} f(x)=f(a).
\]

Esta definición involucra implícitamente tres cosas:

\begin{itemize}
\tightlist
\item
  \(f(a)\) existe (es decir, \(a\) está en el dominio de \(f\)).
\item
  \(\lim_{x\to a} f(x)\) existe.
\item
  \(\lim_{x\to a} f(x)=f(a)\).
\end{itemize}

Esta condición implica que \(f(x)\) está cerca de \(f(a)\) cuando \(x\)
está cerca de \(a\). Es decir, un cambio pequeño en \(x\) sólo produce
una alteración pequeña en \(f(x)\).

Si \(f\) está definida cerca de \(a\) (excepto quizás en \(a\)), decimos
que \(f\) es \textbf{discontinua} en \(a\) si \(f\) no es continua en
\(a\).

\end{tcolorbox}

El gráfico de una función continua no presenta saltos ni interrupciones
abruptas.

\begin{center}
\includegraphics[width=0.5\linewidth,height=\textheight,keepaspectratio]{index_files/mediabag/imagenes_cap2/2.5-continuidad.pdf}
\end{center}

Ejemplo

En la figura se muestra el gráfico de \(y=f(x)\). ¿En qué puntos es
\(f\) discontinua? ¿Por qué?

\begin{center}
\includegraphics[width=0.5\linewidth,height=\textheight,keepaspectratio]{index_files/mediabag/imagenes_cap2/2.5-ejemplo-continuidad.pdf}
\end{center}

\begin{tcolorbox}[enhanced jigsaw, opacityback=0, left=2mm, opacitybacktitle=0.6, toprule=.15mm, bottomrule=.15mm, colbacktitle=quarto-callout-tip-color!10!white, titlerule=0mm, bottomtitle=1mm, title=\textcolor{quarto-callout-tip-color}{\faLightbulb}\hspace{0.5em}{Solución}, toptitle=1mm, colframe=quarto-callout-tip-color-frame, colback=white, breakable, rightrule=.15mm, arc=.35mm, leftrule=.75mm, coltitle=black]

Observando el gráfico podemos decir que \(f\) es discontinua en \(x=1\)
dado que \(f(1)\) no existe.

Por otro lado, notemos que \(f(3)\) existe pero
\(\displaystyle \lim_{x\to 3} f(x)\) no, con lo cual \(f\) también es
discontinua en \(x=3\).

Por último \(f\) resulta discontinua en \(x=5\) ya que, si bien existen
\(f(5)\) y \(\lim_{x\to 5} f(x)\), no coinciden. En los demás puntos
\(f\) resulta continua.

\end{tcolorbox}

Ejemplo

¿En qué puntos son continuas cada una de las siguientes funciones?

\begin{enumerate}
\def\labelenumi{\arabic{enumi}.}
\item
  \(\displaystyle f_1(x)=\frac{x^2-x-2}{x-2}\).
\item
  \(\displaystyle f_2(x)=\begin{cases}
  \frac{1}{x^2} & \text{ si } x\neq 0,\\
  1 & \text{ si } x=0.
  \end{cases}\)
\end{enumerate}

\begin{enumerate}
\def\labelenumi{\arabic{enumi}.}
\setcounter{enumi}{2}
\item
  \(\displaystyle f_3(x)=\begin{cases}
  \frac{x^2-x-2}{x-2} & \text{ si } x\neq 2,\\
  1 & \text{ si } x=2.
  \end{cases}\)
\item
  \(\displaystyle f_4(x)=\frac{|x|}{x}\).
\end{enumerate}

\begin{tcolorbox}[enhanced jigsaw, opacityback=0, left=2mm, opacitybacktitle=0.6, toprule=.15mm, bottomrule=.15mm, colbacktitle=quarto-callout-tip-color!10!white, titlerule=0mm, bottomtitle=1mm, title=\textcolor{quarto-callout-tip-color}{\faLightbulb}\hspace{0.5em}{Solución}, toptitle=1mm, colframe=quarto-callout-tip-color-frame, colback=white, breakable, rightrule=.15mm, arc=.35mm, leftrule=.75mm, coltitle=black]

Comencemos con \(f_1\). Si factorizamos el polinomio del numerador
obtenemos

\[
f_1(x)=\frac{(x-2)(x+1)}{x-2}=x+1
\]

para todo \(x\neq 2\). Como \(f(2)\) no existe, \(f\) es discontinua en
\(x=2\). Si \(a\neq 2\) entonces por la
\hyperref[teo-sust-directa]{propiedad de sustitución directa}

\[
\lim_{x\to a} f_1(x)=\lim_{x\to a} (x+1)=a+1=f_1(a), 
\]

con lo que \(f_1\) resulta continua en \(a\). De hecho, notar que la
función \(f_1\) resulta ser la recta \(y=x+1\), con la excepción de que
en \(x=2\) tenemos un punto faltante, ya que \(f_1(2)\) no existe.

\includegraphics[width=0.4\linewidth,height=\textheight,keepaspectratio]{index_files/mediabag/imagenes_cap2/2.5-ejemplo-f1.pdf}

Ahora analicemos \(f_2\). Si bien \(f_2(0)=1\), tenemos que
\(\lim_{x\to 0} f(x)\) no existe por infinitud, con lo cual \(f_2\) es
discontinua en \(x=0\). Si \(a\neq 0\), entonces nuevamente por la
\hyperref[teo-sust-directa]{propiedad de sustitución directa}

\[
\lim_{x\to a} f_2(x)=\lim_{x\to a} \frac{1}{x^2}=\frac{1}{a^2}=f_2(a),
\]

con lo cual \(f_2\) es continua en \(a\).

\includegraphics[width=0.5\linewidth,height=\textheight,keepaspectratio]{index_files/mediabag/imagenes_cap2/2.5-ejemplo-f2.pdf}

Para el caso de \(f_3\), notar que se trata de la misma función que
\(f_1\), excepto que ahora sí tenemos definido su valor en \(x=2\). Como
vimos que \(f_1\) era continua en todo \(x\neq 2\), lo mismo ocurrirá
con \(f_3\). Veamos si ésta es continua en \(x=2\). Factorizando como
antes obtenemos

\[
\lim_{x\to 2} f_3(x)=\lim_{x\to 2}\frac{(x-2)(x+1)}{x-2}=\lim_{x\to 2} (x+1)=3, 
\]

sin embargo \(f_3(2)=1\neq 3\), con lo que \(f_3\) es discontinua en
\(x=2\).

\includegraphics[width=0.35\linewidth,height=\textheight,keepaspectratio]{index_files/mediabag/imagenes_cap2/2.5-ejemplo-f3.pdf}

Por último, para \(f_4\) observemos que esta función puede ser reescrita
como

\[
f_4(x)=\begin{cases} 
1 & \text{ si } x> 0,\\
-1 & \text{ si } x<0.
\end{cases}
\]

Como \(y=x\) e \(y=-x\) son funciones continuas, \(f_4\) resulta
continua en todo \(x\neq 0\). Como \(f_4(0)\) no existe, \(f_4\) es
discontinua en \(x=0\).

\includegraphics[width=0.35\linewidth,height=\textheight,keepaspectratio]{index_files/mediabag/imagenes_cap2/2.5-ejemplo-f4.pdf}

\end{tcolorbox}

El tipo de discontinuidad que aparece en \(f_1\) y \(f_3\) del ejemplo
anterior se conoce como discontinuidad \textbf{removible} o
\textbf{evitable}, pues podría eliminarse al cambiar (redefinir) el
valor de la función en ese punto por su límite en ese punto.

La función \(f_2\) presenta una discontinuidad \textbf{de salto
infinito} y \(f_4\) una discontinuidad de \textbf{salto finito}.

\begin{tcolorbox}[enhanced jigsaw, opacityback=0, left=2mm, opacitybacktitle=0.6, toprule=.15mm, bottomrule=.15mm, colbacktitle=quarto-callout-note-color!10!white, titlerule=0mm, bottomtitle=1mm, title=\textcolor{quarto-callout-note-color}{\faInfo}\hspace{0.5em}{Definición (continuidad lateral)}, toptitle=1mm, colframe=quarto-callout-note-color-frame, colback=white, breakable, rightrule=.15mm, arc=.35mm, leftrule=.75mm, coltitle=black]

Una función \(f\) es \textbf{continua por derecha} en \(a\) si se cumple
que

\[
\lim_{x\to a^+}f(x)=f(a).
\]

De manera similar, \(f\) es \textbf{continua por izquierda} en \(a\) si
se verifica que

\[
\lim_{x\to a^-}f(x)=f(a).
\]

\end{tcolorbox}

\begin{tcolorbox}[enhanced jigsaw, opacityback=0, left=2mm, opacitybacktitle=0.6, toprule=.15mm, bottomrule=.15mm, colbacktitle=quarto-callout-note-color!10!white, titlerule=0mm, bottomtitle=1mm, title=\textcolor{quarto-callout-note-color}{\faInfo}\hspace{0.5em}{Definición (continuidad en un intervalo)}, toptitle=1mm, colframe=quarto-callout-note-color-frame, colback=white, breakable, rightrule=.15mm, arc=.35mm, leftrule=.75mm, coltitle=black]

Una función \(f\) es \textbf{continua sobre el intervalo \((a,b)\)} si
es continua en \(x\), para todo \(a<x<b\).

Decimos que \(f\) es \textbf{continua sobre \([a,b]\)} si es continua en
\((a,b)\), y además resulta continua por la derecha en \(a\) y por la
izquierda en \(b\).

\end{tcolorbox}

Una definición análoga puede darse para los intervalos \([a,b)\) y
\((a,b]\).

Ejemplo

Mostrar que \(f(x)=1-\sqrt{1-x^2}\) es continua sobre el intervalo
\([-1,1]\).

\begin{tcolorbox}[enhanced jigsaw, opacityback=0, left=2mm, opacitybacktitle=0.6, toprule=.15mm, bottomrule=.15mm, colbacktitle=quarto-callout-tip-color!10!white, titlerule=0mm, bottomtitle=1mm, title=\textcolor{quarto-callout-tip-color}{\faLightbulb}\hspace{0.5em}{Solución}, toptitle=1mm, colframe=quarto-callout-tip-color-frame, colback=white, breakable, rightrule=.15mm, arc=.35mm, leftrule=.75mm, coltitle=black]

Consideremos primero \(-1<a<1\). Dado que \(1-a^2>0\), por la ley
\textbf{(11)} de límites obtenemos que

\[
\lim_{x\to a} \sqrt{1-x^2}=\sqrt{1-a^2},
\]

y de esta manera, por la ley \textbf{(2)}

\[
\lim_{x\to a} f(x)=\lim_{x\to a}\,1 -\lim_{x\to a} \sqrt{1-x^2}=1-\sqrt{1-a^2}=f(a),
\]

lo que nos indica que \(f\) es continua en \(a\). Para terminar, veamos
que \(f\) es continua por derecha en \(x=-1\) y por izquierda en
\(x=1\).

Si \(-1<x<0\), entonces \(x^2<1\) y \(1-x^2>0\). Utilizando las leyes
\textbf{(11)} y \textbf{(2)} resulta

\[
\lim_{x\to -1^+} f(x)=\lim_{x\to -1^+} \,\left(1-\sqrt{1-x^2}\right)=1-0=1=f(-1).
\]

Si \(0<x<1\), entonces también es cierto que \(x^2<1\), o
equivalentemente \(1-x^2>0\). La ley \textbf{(11)} combinada con la
\textbf{(2)} en este caso nos da

\[
\lim_{x\to 1^-} f(x)=\lim_{x\to 1^-} \,\left(1-\sqrt{1-x^2}\right)=1-0=1=f(1).
\]

Por lo tanto, \(f\) es continua en todo el intervalo \([-1,1]\).

\end{tcolorbox}

\phantomsection\label{teo-continuas}
Teorema

Sean \(f\) y \(g\) funciones continuas en \(a\) y \(c\) una constante
real. Entonces las siguientes funciones también son continuas en \(a\)

\begin{itemize}
\item
  \(f+g\).
\item
  \(f-g\).
\item
  \(cf\).
\end{itemize}

\begin{itemize}
\item
  \(fg\).
\item
  \(\displaystyle \frac{1}{g}\), si \(g(a)\neq  0\).
\end{itemize}

Teorema

\begin{itemize}
\item
  Cualquier polinomio es continuo en su dominio, es decir, es continuo
  sobre \(\mathbb{R}\).
\item
  Cualquier función racional es continua siempre que esté definida, es
  decir, es continua en su dominio.
\end{itemize}

Ejemplo

Encontrar \(\lim_{x\to -2} f(x)\), siendo

\[
f(x)=\frac{x^3+2x^2-1}{5-3x}.
\]

\begin{tcolorbox}[enhanced jigsaw, opacityback=0, left=2mm, opacitybacktitle=0.6, toprule=.15mm, bottomrule=.15mm, colbacktitle=quarto-callout-tip-color!10!white, titlerule=0mm, bottomtitle=1mm, title=\textcolor{quarto-callout-tip-color}{\faLightbulb}\hspace{0.5em}{Solución}, toptitle=1mm, colframe=quarto-callout-tip-color-frame, colback=white, breakable, rightrule=.15mm, arc=.35mm, leftrule=.75mm, coltitle=black]

Notemos que \(f\) es racional y \(-2\in D_f\), por lo que usando la
\hyperref[teo-sust-directa]{propiedad de sustitución directa} obtenemos

\[
\lim_{x\to -2} f(x)=f(-2)=\frac{(-2)^3+2(-2)^2-1}{5-3(-2)}=-\frac{1}{11}.
\]

\end{tcolorbox}

La mayor parte de las funciones estudiadas hasta aquí son continuas,
como establece el siguiente teorema.

Teorema

Los siguientes tipos de funciones son continuas en sus respectivos
dominios.

\begin{itemize}
\item
  Polinomios.
\item
  Funciones racionales.
\item
  Funciones raíz.
\item
  Trigonométricas.
\end{itemize}

\begin{itemize}
\item
  Trigonométricas inversas.
\item
  Exponenciales.
\item
  Logarítmicas.
\end{itemize}

Ejemplo

¿Dónde es continua la función
\(\displaystyle f(x)=\frac{\ln x+\arctan x}{x^2-1}\)?

\begin{tcolorbox}[enhanced jigsaw, opacityback=0, left=2mm, opacitybacktitle=0.6, toprule=.15mm, bottomrule=.15mm, colbacktitle=quarto-callout-tip-color!10!white, titlerule=0mm, bottomtitle=1mm, title=\textcolor{quarto-callout-tip-color}{\faLightbulb}\hspace{0.5em}{Solución}, toptitle=1mm, colframe=quarto-callout-tip-color-frame, colback=white, breakable, rightrule=.15mm, arc=.35mm, leftrule=.75mm, coltitle=black]

Notemos que \(y=\ln x\) e \(y=\arctan x\) son continuas en sus
respectivos dominios, es decir, en \((0,\infty)\) y \(\mathbb{R}\). Por
lo tanto el numerador de \(f\) es continuo en \((0,\infty)\).

Por otro lado, \(y=x^2-1\) es continua en \(\mathbb{R}\) por ser un
polinomio, y el cociente será continuo en todos los puntos de
\((0,\infty)\) tales que \(x^2-1\neq 0\). Por lo tanto, \(f\) es
continua en \((0,1)\cup (1,\infty)\).

\end{tcolorbox}

Ejemplo

Hallar
\(\displaystyle \lim_{x\to \pi} \frac{\operatorname{sen} x}{2+\cos x}\).

\begin{tcolorbox}[enhanced jigsaw, opacityback=0, left=2mm, opacitybacktitle=0.6, toprule=.15mm, bottomrule=.15mm, colbacktitle=quarto-callout-tip-color!10!white, titlerule=0mm, bottomtitle=1mm, title=\textcolor{quarto-callout-tip-color}{\faLightbulb}\hspace{0.5em}{Solución}, toptitle=1mm, colframe=quarto-callout-tip-color-frame, colback=white, breakable, rightrule=.15mm, arc=.35mm, leftrule=.75mm, coltitle=black]

Como \(y=\operatorname{sen} x\) e \(y=2+\cos x\) son funciones continuas
en \(x=\pi\) y \(2+\cos \pi =1\neq 0\), el cociente es continuo y por lo
tanto

\[
\lim_{x\to \pi} \frac{\operatorname{sen} x}{2+\cos x}=\frac{\operatorname{sen} \pi}{2+\cos \pi}=\frac{1}{2-1}=1.
\]

\end{tcolorbox}

\phantomsection\label{teo-limite-composicion}
Teorema (límite de la composición de funciones)

Si \(f\) es continua en \(b\) y \(\lim_{x\to a} g(x)=b\) entonces

\[
\lim_{x\to a}f(g(x))=f(b).
\]

Esto también puede reformularse como sigue

\[
\lim_{x\to a}f(g(x))=f\left(\lim_{x\to a}g(x)\right).
\]

Ejemplo

Calcular
\(\displaystyle \lim_{x\to 1} \operatorname{arcsen}\left(\frac{1-\sqrt{x}}{1-x}\right)\).

\begin{tcolorbox}[enhanced jigsaw, opacityback=0, left=2mm, opacitybacktitle=0.6, toprule=.15mm, bottomrule=.15mm, colbacktitle=quarto-callout-tip-color!10!white, titlerule=0mm, bottomtitle=1mm, title=\textcolor{quarto-callout-tip-color}{\faLightbulb}\hspace{0.5em}{Solución}, toptitle=1mm, colframe=quarto-callout-tip-color-frame, colback=white, breakable, rightrule=.15mm, arc=.35mm, leftrule=.75mm, coltitle=black]

Sean \(f(x)=\operatorname{arcsen} x\) y
\(g(x)=\displaystyle \frac{1-\sqrt{x}}{1-x}\). Calculamos primero
\(\lim_{x\to 1} g(x)\). En efecto, si \(x\neq 1\) escribimos

\[
\frac{1-\sqrt{x}}{1-x}=\frac{1-\sqrt{x}}{1-x}\,\frac{1+\sqrt{x}}{1+\sqrt{x}}=\frac{1^2-(\sqrt{x})^2}{(1-x)(1+\sqrt{x})}=\frac{1-x}{(1-x)(1+\sqrt{x})}=\frac{1}{1+\sqrt{x}}.
\]

Luego

\[
\lim_{x\to 1} \frac{1-\sqrt{x}}{1-x} =\lim_{x\to 1} \frac{1}{1+\sqrt{x}}=\frac{1}{1+\sqrt{1}}=\frac{1}{2},
\] dado que \(y_1=1\) e \(y_2=1+\sqrt{x}\) son continuas en \(x=1\) e
\(y_2(1)\neq 0\).

Ahora bien, como \(y=\operatorname{arcsen} x\) es continua en \(x=1/2\),
por el \hyperref[teo-limite-composicion]{teorema del límite de la
composición de funciones} obtenemos

\[
\lim_{x\to 1} \operatorname{arcsen}\left(\frac{1-\sqrt{x}}{1-x}\right)= \operatorname{arcsen}\left(\lim_{x\to 1} \frac{1-\sqrt{x}}{1-x}\right)=\operatorname{arcsen}\left(\frac{1}{2}\right)=\frac{\pi}{6}.
\]

\end{tcolorbox}

\phantomsection\label{teo-continuidad-composicion}
Teorema (continuidad de la composición de funciones)

Si \(g\) es continua en \(a\) y \(f\) es continua en \(g(a)\), entonces
la composición \((f\circ g)(x)=f(g(x))\) es continua en \(a\).

Ejemplo

¿Dónde son continuas las siguientes funciones?

\begin{itemize}
\tightlist
\item
  \(h(x)=\operatorname{sen}(x^2)\).
\item
  \(F(x)=\ln(1+\cos x)\).
\end{itemize}

\begin{tcolorbox}[enhanced jigsaw, opacityback=0, left=2mm, opacitybacktitle=0.6, toprule=.15mm, bottomrule=.15mm, colbacktitle=quarto-callout-tip-color!10!white, titlerule=0mm, bottomtitle=1mm, title=\textcolor{quarto-callout-tip-color}{\faLightbulb}\hspace{0.5em}{Solución}, toptitle=1mm, colframe=quarto-callout-tip-color-frame, colback=white, breakable, rightrule=.15mm, arc=.35mm, leftrule=.75mm, coltitle=black]

Podemos expresar \(h(x)=f(g(x))\), donde \(f(x)=\operatorname{sen} x\) y
\(g(x)=x^2\). Si \(x\in \mathbb{R}\), entonces \(g\) es continua en
\(x\) por ser un polinomio y además \(g(x)\in [0,+\infty)\). Como \(f\)
es continua en \(\mathbb{R}\), en particular es continua en
\([0,+\infty)\). Por lo tanto, el
\hyperref[teo-continuidad-composicion]{teorema de la continuidad de la
composición} establece que \(h\) es continua en \(\mathbb{R}\).

Para el segundo caso, escribimos \(F(x)=f(g(x))\) siendo \(f(x)=\ln x\)
y \(g(x)=1+\cos x\). Notemos que \(g\) es continua en \(\mathbb{R}\) y
\(f\) es continua en \((0,+\infty)\), por lo tanto la composición \(F\)
será continua en todos los puntos \(x\) tales que \(1+\cos x>0\). Dado
que \(\cos x\geq -1\), la desigualdad anterior se cumple siempre y
cuando \(1+\cos x\neq 0\), es decir, cuando \(\cos x \neq -1\).

Observar que

\[
\cos x=-1 \quad \Longleftrightarrow \quad x=(2n+1)\pi, n\in\mathbb{Z}.
\]

Entonces \(f\) es continua en
\((0,+\infty)\setminus \{x: x=(2n+1)\pi, n\in\mathbb{Z}\}\).

\end{tcolorbox}

\phantomsection\label{teo-valor-intermedio}
Teorema (teorema del valor intermedio)

Sea \(f\) una función continua en el intervalo cerrado \([a,b]\) y sea
\(N\) un número entre \(f(a)\) y \(f(b)\) (suponemos que
\(f(a)\neq f(b)\)). Entonces existe un número \(c \in(a,b)\) tal que
\(f(c)=N\).

Este teorema nos dice que una función continua en \([a,b]\) toma todos
los valores intermedios comprendidos entre \(f(a)\) y \(f(b)\).

\begin{center}
\includegraphics[width=0.7\linewidth,height=\textheight,keepaspectratio]{index_files/mediabag/imagenes_cap2/2.5-continuidad-intermedio.pdf}
\end{center}

Ejemplo

Demostrar que existe una raíz \(r\) de la ecuación \(4x^3-6x^2+3x-2=0\)
que cumple \(1<r<2\).

\begin{tcolorbox}[enhanced jigsaw, opacityback=0, left=2mm, opacitybacktitle=0.6, toprule=.15mm, bottomrule=.15mm, colbacktitle=quarto-callout-tip-color!10!white, titlerule=0mm, bottomtitle=1mm, title=\textcolor{quarto-callout-tip-color}{\faLightbulb}\hspace{0.5em}{Solución}, toptitle=1mm, colframe=quarto-callout-tip-color-frame, colback=white, breakable, rightrule=.15mm, arc=.35mm, leftrule=.75mm, coltitle=black]

Definimos \(f(x)=4x^3-6x^2+3x-2\). Entonces \(f\) es continua por ser un
polinomio. En particular, \(f\) es continua en el intervalo \([1,2]\).
Además \(f(1)=-1<0\) y \(f(2)=12>0\), es decir, \(f(1)<0<f(2)\). Por el
\hyperref[teo-valor-intermedio]{teorema del valor intermedio} existe un
número \(c\in (1,2)\) tal que \(f(c)=0\) y esto nos dice que \(c\) es
una raíz de la ecuación dada.

\end{tcolorbox}

\hyperref[seccion_2.5]{↑ Volver al inicio de la sección}

\section{2.6. Límites al infinito, asíntotas
horizontales}\label{seccion_2.6}

Consideremos la función

\[
f(x)= \frac{x^2-1}{x^2+1}
\]

cuya gráfica se muestra a continuación.

\begin{center}
\includegraphics[width=0.45\linewidth,height=\textheight,keepaspectratio]{index_files/mediabag/imagenes_cap2/2.6-racional.pdf}
\end{center}

Si le damos distintos valores a \(x\) podemos ver que los valores de la
función se estacionan conforme \(x\) se hace grande en magnitud, como se
observa en la siguiente tabla.

\begin{longtable}[]{@{}ll@{}}
\toprule\noalign{}
\(x\) & \(f(x)\) \\
\midrule\noalign{}
\endhead
\bottomrule\noalign{}
\endlastfoot
0 & 1 \\
\(\pm\,\) 1 & 0 \\
\(\pm\,\) 2 & 0.600000 \\
\(\pm\,\) 3 & 0.800000 \\
\(\pm\,\) 4 & 0.882353 \\
\(\pm\,\) 5 & 0.923077 \\
\(\pm\,\) 10 & 0.980198 \\
\(\pm\,\) 50 & 0.999200 \\
\(\pm\,\) 100 & 0.999800 \\
\(\pm\,\) 1000 & 0.999998 \\
\end{longtable}

Parece que podemos acercar \(f(x)\) tanto a \(1\) como deseemos, siempre
que elijamos \(x\) lo suficientemente grande en magnitud. Esto lo
expresamos escribiendo

\[
\lim_{x\to\infty} \frac{x^2-1}{x^2+1}=1.
\]

\begin{tcolorbox}[enhanced jigsaw, opacityback=0, left=2mm, opacitybacktitle=0.6, toprule=.15mm, bottomrule=.15mm, colbacktitle=quarto-callout-note-color!10!white, titlerule=0mm, bottomtitle=1mm, title=\textcolor{quarto-callout-note-color}{\faInfo}\hspace{0.5em}{Definición (límite en el infinito)}, toptitle=1mm, colframe=quarto-callout-note-color-frame, colback=white, breakable, rightrule=.15mm, arc=.35mm, leftrule=.75mm, coltitle=black]

Sea \(f\) una función definida en algún intervalo de la forma
\((a,\infty)\). Entonces \[
\lim_{x\to \infty} f(x)=L 
\]

significa que \(f(x)\) puede aproximarse a \(L\) tanto como se quiera,
siempre que se elija \(x\) suficientemente grande.

También escribiremos: \(f(x)\to L\) cuando \$x\to\infty \$. La expresión
de arriba se lee

\begin{itemize}
\tightlist
\item
  ``el límite de \(f(x)\) cuando \(x\) tiende a infinito es \(L\)'', o
\item
  ``el límite de \(f(x)\) cuando \(x\) se hace infinito es \(L\)''\,``,
  o
\item
  ``\(f(x)\) tiende a \(L\) cuando \(x\) crece sin cota''.
\end{itemize}

\end{tcolorbox}

Vemos que \(f\) puede aproximarse a \(L\) de diferentes maneras, como
muestra la siguiente figura.

\begin{center}
\includegraphics[width=0.7\linewidth,height=\textheight,keepaspectratio]{index_files/mediabag/imagenes_cap2/2.6-asintotas-horizontales.pdf}
\end{center}

\begin{tcolorbox}[enhanced jigsaw, opacityback=0, left=2mm, opacitybacktitle=0.6, toprule=.15mm, bottomrule=.15mm, colbacktitle=quarto-callout-note-color!10!white, titlerule=0mm, bottomtitle=1mm, title=\textcolor{quarto-callout-note-color}{\faInfo}\hspace{0.5em}{Definición (límite en el infinito negativo)}, toptitle=1mm, colframe=quarto-callout-note-color-frame, colback=white, breakable, rightrule=.15mm, arc=.35mm, leftrule=.75mm, coltitle=black]

Sea \(f\) una función definida en algún intervalo de la forma
\((-\infty,a)\). Entonces

\[
\lim_{x\to -\infty} f(x)=L 
\]

significa que \(f(x)\) puede aproximarse a \(L\) tanto como se quiera,
siempre que se elija \(x\) negativo y suficientemente grande en
magnitud.

También escribiremos: \(f(x)\to L\) cuando \(x \to -\infty\) y se lee
``el límite de \(f(x)\) cuando \(x\) tiende a infinito negativo es
\(L\)''.

\end{tcolorbox}

\begin{tcolorbox}[enhanced jigsaw, opacityback=0, left=2mm, opacitybacktitle=0.6, toprule=.15mm, bottomrule=.15mm, colbacktitle=quarto-callout-note-color!10!white, titlerule=0mm, bottomtitle=1mm, title=\textcolor{quarto-callout-note-color}{\faInfo}\hspace{0.5em}{Definición (asíntota horizontal)}, toptitle=1mm, colframe=quarto-callout-note-color-frame, colback=white, breakable, rightrule=.15mm, arc=.35mm, leftrule=.75mm, coltitle=black]

La recta \(y=L\) se llama \textbf{asíntota horizontal} de la curva
\(y=f(x)\) si \[
\lim_{x\to\infty} f(x)=L \quad \text{ o bien si } \quad \lim_{x\to -\infty} f(x)=L. 
\]

\end{tcolorbox}

Un ejemplo de función con dos asíntotas horizontales es
\(y= \arctan x\).

\begin{center}
\includegraphics[width=0.5\linewidth,height=\textheight,keepaspectratio]{index_files/mediabag/imagenes_cap2/2.6-arcotangente.pdf}
\end{center}

En este caso tenemos que

\[
\lim_{x\to \infty } \arctan x=\frac{\pi}{2} \quad \text{ y }\quad  \lim_{x\to -\infty } \arctan x=-\frac{\pi}{2}.
\]

La mayor parte de las \hyperref[thm-leyes-limites]{leyes de los límites}
que estudiamos en la \hyperref[seccion_2.3]{Sección 2.3} también se
cumple para límites en el infinito. Las leyes vistas (con excepción de
la \textbf{(9)} y la \textbf{(10)}) siguen siendo válidas reemplazando
\(x\to a\) por \(x \to \infty\) o \(x\to -\infty\).

Teorema

Sea \(r>0\) un número racional. Entonces

\begin{itemize}
\item
  \(\displaystyle\lim_{x\to \infty} \frac{1}{x^r}=0\).
\item
  Si \(x^r\) está definido para todo \(x\), entonces
  \(\displaystyle\lim_{x\to -\infty} \frac{1}{x^r}=0\).
\end{itemize}

Ejemplo

Calcular \[
\lim_{x\to \infty}\frac{3x^2-x-2}{5x^2+4x+1}
\]

e indicar las propiedades que se utilizan en cada paso.

\begin{tcolorbox}[enhanced jigsaw, opacityback=0, left=2mm, opacitybacktitle=0.6, toprule=.15mm, bottomrule=.15mm, colbacktitle=quarto-callout-tip-color!10!white, titlerule=0mm, bottomtitle=1mm, title=\textcolor{quarto-callout-tip-color}{\faLightbulb}\hspace{0.5em}{Solución}, toptitle=1mm, colframe=quarto-callout-tip-color-frame, colback=white, breakable, rightrule=.15mm, arc=.35mm, leftrule=.75mm, coltitle=black]

Para \(x>0\) podemos escribir

\[
\frac{3x^2-x-2}{5x^2+4x+1}=\frac{x^2\left(3-\frac{x}{x^2}-\frac{2}{x^2}\right)}{x^2\left(5+\frac{4x}{x^2}+\frac{1}{x^2}\right)}=\frac{3-\frac{1}{x}-\frac{2}{x^2}}{5+\frac{4}{x}+\frac{1}{x^2}}.
\]

Por el teorema anterior tenemos que

\[
\lim_{x\to \infty} \frac{1}{x}=0 \quad \text{ y }\quad \lim_{x\to \infty} \frac{1}{x^2}=0.
\]

Por las leyes \textbf{(2)} y \textbf{(3)}

\[
\lim_{x\to \infty} \left(3-\frac{1}{x}-\frac{2}{x^2}\right)=\lim_{x\to \infty} 3-\lim_{x\to \infty} \frac{1}{x}-2\lim_{x\to \infty} \frac{1}{x^2}=3-0-0=3.
\]

De la misma forma

\[
\lim_{x\to \infty} \left(5+\frac{4}{x}+\frac{1}{x^2}\right)=\lim_{x\to \infty} 5+4\lim_{x\to \infty} \frac{1}{x}+\lim_{x\to \infty} \frac{1}{x^2}=5-0-0=5.
\]

Finalmente, usando la ley \textbf{(5)} del cociente obtenemos

\[
\lim_{x\to \infty} \frac{3-\frac{1}{x}-\frac{2}{x^2}}{5+\frac{4}{x}+\frac{1}{x^2}}=\frac{\lim_{x\to\infty}\left(3-\frac{1}{x}-\frac{2}{x^2}\right)}{\lim_{x\to\infty} \left(5+\frac{4}{x}+\frac{1}{x^2}\right)}=\frac{3}{5}.
\]

\end{tcolorbox}

Ejemplo

Determinar las asíntotas horizontales y verticales de la gráfica de la
función \[
f(x)=\frac{\sqrt{2x^2+1}}{3x-5}.
\]

\begin{tcolorbox}[enhanced jigsaw, opacityback=0, left=2mm, opacitybacktitle=0.6, toprule=.15mm, bottomrule=.15mm, colbacktitle=quarto-callout-tip-color!10!white, titlerule=0mm, bottomtitle=1mm, title=\textcolor{quarto-callout-tip-color}{\faLightbulb}\hspace{0.5em}{Solución}, toptitle=1mm, colframe=quarto-callout-tip-color-frame, colback=white, breakable, rightrule=.15mm, arc=.35mm, leftrule=.75mm, coltitle=black]

Observemos que el denominador se anula cuando \(x=5/3\). Si \(x\) se
aproxima a \(5/3\) por la derecha, entonces el numerador se aproxima a
\(\sqrt{59}/3>0\) y \(3x-5\) se aproxima a cero con valores positivos,
con lo cual

\[
\lim_{x\to (5/3)^+} f(x)=+\infty,
\]

y un análisis similar muestra que

\[
\lim_{x\to (5/3)^-} f(x)=-\infty.
\]

Por lo tanto \(\displaystyle x=\frac{5}{3}\) es una asíntota vertical de
\(f\). Dado que \(f\) es continua en todo \(x\neq 5/3\), ésta es la
única asíntota vertical de la función.

Para calcular las asíntotas horizontales, debemos hacer
\(x\to \pm \infty\). Primero notemos que, si \(x>0\)

\[
\frac{\sqrt{2x^2+1}}{3x-5}=\frac{\sqrt{x^2\left(2+\frac{1}{x^2}\right)}}{x\left(3-\frac{5}{x}\right)}=\frac{|x|\sqrt{2+\frac{1}{x^2}}}{x\left(3-\frac{5}{x}\right)}=\frac{\sqrt{2+\frac{1}{x^2}}}{3-\frac{5}{x}}.
\]

Por las leyes \textbf{(1)} y \textbf{(11)} tenemos que

\[
\lim_{x\to \infty} \sqrt{2+\frac{1}{x^2}}=\sqrt{\lim_{x\to \infty} \left(2+\frac{1}{x^2}\right)}=\sqrt{2}.
\]

Además, por la ley \textbf{(2)}

\[
\lim_{x\to \infty} \left(3-\frac{5}{x}\right)=\lim_{x\to \infty} 3- \lim_{x\to \infty} \frac{5}{x}=3-0=3.
\]

Finalmente, por la ley \textbf{(5)} del cociente

\[
\lim_{x\to \infty} \frac{\sqrt{2x^2+1}}{3x-5}=\lim_{x\to\infty}\frac{\sqrt{2+\frac{1}{x^2}}}{3-\frac{5}{x}}=\frac{\lim_{x\to\infty} \sqrt{2+\frac{1}{x^2}}}{\lim_{x\to\infty}\left(3-\frac{5}{x}\right)}=\frac{\sqrt{2}}{3}.
\]

Con esto, \(y=\sqrt{2}/3\) es asíntota horizontal de \(f\). Para
\(x\to-\infty\) consideramos \(x<0\) y escribimos

\[
\frac{\sqrt{2x^2+1}}{3x-5}=\frac{\sqrt{x^2\left(2+\frac{1}{x^2}\right)}}{x\left(3-\frac{5}{x}\right)}=\frac{|x|\sqrt{2+\frac{1}{x^2}}}{x\left(3-\frac{5}{x}\right)}=-\frac{\sqrt{2+\frac{1}{x^2}}}{3-\frac{5}{x}}.
\]

Repitiendo el proceso anterior podemos concluir que

\[
\lim_{x\to -\infty} \frac{\sqrt{2x^2+1}}{3x-5}=\lim_{x\to-\infty}-\frac{\sqrt{2+\frac{1}{x^2}}}{3-\frac{5}{x}}=-\frac{\sqrt{2}}{3},
\]

por lo que \(y=-\sqrt{2}/3\) es otra asíntota horizontal de \(f\).

\includegraphics[width=0.5\linewidth,height=\textheight,keepaspectratio]{index_files/mediabag/imagenes_cap2/2.6-ejemplo-asintotas.pdf}

\end{tcolorbox}

Ejemplo

Calcular
\(\displaystyle \lim_{x\to \infty} \left(\sqrt{x^2+1}-x\right)\).

\begin{tcolorbox}[enhanced jigsaw, opacityback=0, left=2mm, opacitybacktitle=0.6, toprule=.15mm, bottomrule=.15mm, colbacktitle=quarto-callout-tip-color!10!white, titlerule=0mm, bottomtitle=1mm, title=\textcolor{quarto-callout-tip-color}{\faLightbulb}\hspace{0.5em}{Solución}, toptitle=1mm, colframe=quarto-callout-tip-color-frame, colback=white, breakable, rightrule=.15mm, arc=.35mm, leftrule=.75mm, coltitle=black]

Considerando \(x>0\), multiplicamos y dividimos por \(\sqrt{x^2+1}+x\)

\[
\sqrt{x^2+1}-x=\left(\sqrt{x^2+1}-x\right)\frac{\sqrt{x^2+1}+x}{\sqrt{x^2+1}+x}=\frac{\left(\sqrt{x^2+1}\right)^2-x^2}{\sqrt{x^2+1}+x}=\frac{1}{\sqrt{x^2+1}+x}.
\]

Factorizando el denominador de la expresión resultante, tenemos que

\[
\frac{1}{\sqrt{x^2+1}+x}=\frac{1}{\sqrt{x^2\left(1+\frac{1}{x^2}\right)}+x}=\frac{1}{\sqrt{x^2}\sqrt{1+\frac{1}{x^2}}+x}=\frac{1}{x}\,\frac{1}{\sqrt{1+\frac{1}{x^2}}}.
\]

Ahora, sabiendo que

\[
\lim_{x\to \infty} \frac{1}{x}=0 \quad \text{ y } \quad \lim_{x\to\infty} \frac{1}{\sqrt{1+\frac{1}{x^2}}}=1
\]

por la ley \textbf{(3)} del producto concluímos que

\[
\lim_{x\to \infty} \left(\sqrt{x^2+1}-x\right) = \left(\lim_{x\to \infty} \frac{1}{x}\right) \left(\lim_{x\to\infty} \frac{1}{\sqrt{1+\frac{1}{x^2}}}\right)=0\cdot 1=0.
\]

\end{tcolorbox}

En base al gráfico de \(y=e^x\), podemos ver que la recta \(y=0\) es
asíntota horizontal de su gráfico.

\begin{center}
\includegraphics[width=0.4\linewidth,height=\textheight,keepaspectratio]{index_files/mediabag/imagenes_cap2/2.6-exponencial.pdf}
\end{center}

Esto nos dice que

\phantomsection\label{form-limite-exp}
\[
\lim_{x\to -\infty} e^{x}=0. 
\]

Ejemplo

Evaluar \(\lim_{x\to 0-} e^{1/x}\).

\begin{tcolorbox}[enhanced jigsaw, opacityback=0, left=2mm, opacitybacktitle=0.6, toprule=.15mm, bottomrule=.15mm, colbacktitle=quarto-callout-tip-color!10!white, titlerule=0mm, bottomtitle=1mm, title=\textcolor{quarto-callout-tip-color}{\faLightbulb}\hspace{0.5em}{Solución}, toptitle=1mm, colframe=quarto-callout-tip-color-frame, colback=white, breakable, rightrule=.15mm, arc=.35mm, leftrule=.75mm, coltitle=black]

Sea \(t=1/x\). Dado que

\[
\lim_{x\to 0^-} \frac{1}{x} = -\infty,
\]

tenemos que \(t\to -\infty\) cuando \(x\to 0^-\). Expresando el límite
en términos de la variable \(t\) resulta

\[
\lim_{x\to 0-} e^{1/x}=\lim_{t\to -\infty} e^{t}=0,
\]

de acuerdo a la \hyperref[form-limite-exp]{fórmula} anterior.

\end{tcolorbox}

\subsection{Límites infinitos en el
infinito}\label{luxedmites-infinitos-en-el-infinito}

\begin{tcolorbox}[enhanced jigsaw, opacityback=0, left=2mm, opacitybacktitle=0.6, toprule=.15mm, bottomrule=.15mm, colbacktitle=quarto-callout-note-color!10!white, titlerule=0mm, bottomtitle=1mm, title=\textcolor{quarto-callout-note-color}{\faInfo}\hspace{0.5em}{Definición (límite infinito en infinito)}, toptitle=1mm, colframe=quarto-callout-note-color-frame, colback=white, breakable, rightrule=.15mm, arc=.35mm, leftrule=.75mm, coltitle=black]

Al escribir

\[
\lim_{x\to \infty} f(x)=\infty
\]

indicamos que \(f(x)\) crece sin cota cuando \(x\) se hace grande.

De manera similar se definen las expresiones

\[
\lim_{x\to -\infty}f(x)=\infty, \quad \lim_{x\to\infty }f(x)=-\infty \quad  \text{ y } \quad \lim_{x\to -\infty}f(x)=-\infty.
\]

\end{tcolorbox}

Ejemplo

Determinar \(\lim_{x\to \infty } x^3\) y \(\lim_{x\to -\infty} x^3\).

\begin{tcolorbox}[enhanced jigsaw, opacityback=0, left=2mm, opacitybacktitle=0.6, toprule=.15mm, bottomrule=.15mm, colbacktitle=quarto-callout-tip-color!10!white, titlerule=0mm, bottomtitle=1mm, title=\textcolor{quarto-callout-tip-color}{\faLightbulb}\hspace{0.5em}{Solución}, toptitle=1mm, colframe=quarto-callout-tip-color-frame, colback=white, breakable, rightrule=.15mm, arc=.35mm, leftrule=.75mm, coltitle=black]

Si \(x\) crece sin cota, entonces su cubo también (de hecho, \(x^3>x\)
para \(x>1\)), por lo tanto \(\lim_{x\to \infty } x^3=\infty\).

Cuando \(x\) se hace grande y negativo, su cubo también es negativo y
grande en magnitud, por lo tanto \(\lim_{x\to -\infty } x^3=-\infty\).

\end{tcolorbox}

Si observamos nuevamente el gráfico de \(y=e^x\), también podemos
concluir que

\[
\lim_{x\to \infty} e^{x}=\infty. 
\]

Ejemplo

Calcular \(\lim_{x\to \infty} \left(x^2-x\right)\).

\begin{tcolorbox}[enhanced jigsaw, opacityback=0, left=2mm, opacitybacktitle=0.6, toprule=.15mm, bottomrule=.15mm, colbacktitle=quarto-callout-tip-color!10!white, titlerule=0mm, bottomtitle=1mm, title=\textcolor{quarto-callout-tip-color}{\faLightbulb}\hspace{0.5em}{Solución}, toptitle=1mm, colframe=quarto-callout-tip-color-frame, colback=white, breakable, rightrule=.15mm, arc=.35mm, leftrule=.75mm, coltitle=black]

En este caso conviene factorizar de la siguiente manera

\[
x^2-x=x(x-1).
\]

Cuando \(x\) se hace grande, \(x-1\) también. Por lo tanto el producto
crecerá sin cota. Así

\[
\lim_{x\to \infty} \left(x^2-x\right)=\lim_{x\to \infty} x(x-1)=\infty.
\]

\end{tcolorbox}

Ejemplo

Calcular \(\lim_{x\to \infty} \displaystyle \frac{x^2+x}{3-x}\).

\begin{tcolorbox}[enhanced jigsaw, opacityback=0, left=2mm, opacitybacktitle=0.6, toprule=.15mm, bottomrule=.15mm, colbacktitle=quarto-callout-tip-color!10!white, titlerule=0mm, bottomtitle=1mm, title=\textcolor{quarto-callout-tip-color}{\faLightbulb}\hspace{0.5em}{Solución}, toptitle=1mm, colframe=quarto-callout-tip-color-frame, colback=white, breakable, rightrule=.15mm, arc=.35mm, leftrule=.75mm, coltitle=black]

Extrayendo la potencia de mayor orden de factor común en el numerador y
el denominador, para \(x>0\) resulta

\[
\frac{x^2\left(1+\frac{1}{x}\right)}{x\left(\frac{3}{x}-1\right)}=x\,\frac{1+\frac{1}{x}}{\frac{3}{x}-1}.
\]

Ahora bien, \[
\lim_{x\to \infty} \frac{1+\frac{1}{x}}{\frac{3}{x}-1}=-1
\]

y \(\lim_{x\to \infty} x=\infty\), por lo tanto el límite del producto
es \(-\infty\).

\end{tcolorbox}

\hyperref[seccion_2.6]{↑ Volver al inicio de la sección}

\section{2.7. Derivadas y razones de cambio}\label{seccion_2.7}

\subsection{Tangentes}\label{tangentes}

Consideremos el problema de encontrar la recta tangente a una curva
\(C\) que es el gráfico de la función \(y=f(x)\) en el punto
\(P=(a,f(a))\).

Para ello, consideramos un punto cercano \(Q=(x,f(x))\) con \(x\neq a\)
y calculamos la pendiente de la recta secante \(PQ\) \[
m_{PQ}=\frac{f(x)-f(a)}{x-a}.
\]

Si dejamos que \(x\) se aproxime a \(a\), la recta secante se aproxima a
la recta tangente buscada. Si \(m_{PQ}\) se aproxima a un número \(m\),
definimos la \textbf{recta tangente} como la recta que pasa por \(P\) y
tiene pendiente \(m\).

\begin{tcolorbox}[enhanced jigsaw, opacityback=0, left=2mm, opacitybacktitle=0.6, toprule=.15mm, bottomrule=.15mm, colbacktitle=quarto-callout-note-color!10!white, titlerule=0mm, bottomtitle=1mm, title=\textcolor{quarto-callout-note-color}{\faInfo}\hspace{0.5em}{Definición (recta tangente)}, toptitle=1mm, colframe=quarto-callout-note-color-frame, colback=white, breakable, rightrule=.15mm, arc=.35mm, leftrule=.75mm, coltitle=black]

La \textbf{recta tangente} a la curva \(y=f(x)\) en el punto
\(P=(a,f(a))\) es la recta que pasa por \(P\) con pendiente

\[
m=\lim_{x\to a}\frac{f(x)-f(a)}{x-a}
\] siempre y cuando este límite exista.

\includegraphics[width=0.5\linewidth,height=\textheight,keepaspectratio]{index_files/mediabag/imagenes_cap2/2.7-derivada.pdf}

\end{tcolorbox}

Ejemplo

Encontrar una ecuación de la recta tangente a la parábola \(y=f(x)=x^2\)
en el punto \(P=(1,1)\).

\begin{tcolorbox}[enhanced jigsaw, opacityback=0, left=2mm, opacitybacktitle=0.6, toprule=.15mm, bottomrule=.15mm, colbacktitle=quarto-callout-tip-color!10!white, titlerule=0mm, bottomtitle=1mm, title=\textcolor{quarto-callout-tip-color}{\faLightbulb}\hspace{0.5em}{Solución}, toptitle=1mm, colframe=quarto-callout-tip-color-frame, colback=white, breakable, rightrule=.15mm, arc=.35mm, leftrule=.75mm, coltitle=black]

La recta buscada tiene pendiente

\[
m=\lim_{x\to 1} \frac{f(x)-f(1)}{x-1}=\lim_{x\to 1}\frac{x^2-1}{x-1}=2,
\]

según vimos en el \hyperref[ejemplo-cuadratica-lineal]{Ejemplo 11}. Como
además pasa por el punto \((1,1)\) utilizando la
\hyperref[recta-punto-pendiente]{ecuación punto-pendiente de la recta}
que pasa por \((x_1,y_1)\) con pendiente \(m\) obtenemos

\phantomsection\label{recta-punto-pendiente}
\begin{quote}
\textbf{📐 Ecuación punto - pendiente de la recta}

\[
y-y_1=m(x-x_1)
\]
\end{quote}

\[
y-1=2(x-1) \quad \text{ o de forma equivalente }\quad y=2x-1.
\]

\includegraphics[width=0.4\linewidth,height=\textheight,keepaspectratio]{index_files/mediabag/imagenes_cap2/2.1-parabola.pdf}

\end{tcolorbox}

A menudo se hace referencia a la pendiente de la recta tangente a una
curva en un punto como la \textbf{pendiente de la curva} en ese punto.
La idea es que, acercándonos lo suficiente al punto, la curva se parece
mucho a su recta tangente.

Podemos escribir la pendiente de la recta tangente de una forma
equivalente: si \(x\to a\), entonces \(h=x-a\to 0\). Con este cambio,
\(x=a+h\) y cuando \(x\) se aproxima a \(a\), \(h\) se aproxima a 0. De
esta manera tenemos

\[
m_{PQ}=\frac{f(a+h)-f(a)}{h} \quad \text{ y }\quad m=\lim_{h\to 0} \frac{f(a+h)-f(a)}{h}.
\]

Ejemplo

Encontrar una ecuación de la recta tangente a
\(\displaystyle y=f(x)=\frac{3}{x}\) en el punto \(P=(3,1)\).

\begin{tcolorbox}[enhanced jigsaw, opacityback=0, left=2mm, opacitybacktitle=0.6, toprule=.15mm, bottomrule=.15mm, colbacktitle=quarto-callout-tip-color!10!white, titlerule=0mm, bottomtitle=1mm, title=\textcolor{quarto-callout-tip-color}{\faLightbulb}\hspace{0.5em}{Solución}, toptitle=1mm, colframe=quarto-callout-tip-color-frame, colback=white, breakable, rightrule=.15mm, arc=.35mm, leftrule=.75mm, coltitle=black]

La recta buscada tiene pendiente

\[
m=\lim_{h\to 0} \frac{f(3+h)-f(3)}{h}=\lim_{h\to 0} \frac{\frac{3}{(3+h)}-1}{h}=\lim_{h\to 0} \frac{\frac{3-(3+h)}{(3+h)}}{h}=\lim_{h\to 0} -\frac{h}{h(3+h)}=\lim_{h\to 0} -\frac{1}{3+h}=-\frac{1}{3}.
\]

Como además la recta pasa por \((3,1)\), usando nuevamente la
\hyperref[recta-punto-pendiente]{ecuación punto-pendiente de la recta}
resulta

\[
y-1=-\frac{1}{3}(x-3) \quad \text{ o bien } \quad y=-\frac{1}{3}x+2.
\]

\includegraphics[width=0.4\linewidth,height=\textheight,keepaspectratio]{index_files/mediabag/imagenes_cap2/2.7-ejemplo-reciprocaa.pdf}

\end{tcolorbox}

\subsection{Velocidades}\label{velocidades}

Supongamos que un objeto se mueve a lo largo de una línea recta y su
\textbf{función de posición} es \(s=f(t)\), donde \(s\) es el
desplazamiento respecto del origen en el instante \(t\). El cambio de
posición en el intervalo de tiempo desde \(t=a\) hasta \(t=a+h\) es
\(f(a+h)-f(a)\). La velocidad promedio es

\[
\text{vel. promedio} =\frac{\text{desplazamiento}}{\text{tiempo}}=\frac{f(a+h)-f(a)}{(a+h)-a}=\frac{f(a+h)-f(a)}{h},
\]

que podemos interpretar como la pendiente de la recta secante que une
los puntos \(P=(a,f(a))\) y \(Q=(a+h,f(a+h))\). Al hacer la amplitud de
los intervalos más y más chica, se puede definir la \textbf{velocidad} o
\textbf{velocidad instantánea} como el límite de estas velocidades
promedio

\[
v(a)=\lim_{h\to 0} \frac{f(a+h)-f(a)}{h}.
\]

La velocidad instantánea puede interpretarse como la pendiente de la
recta tangente en el punto \(P\).

\begin{center}
\includegraphics[width=0.5\linewidth,height=\textheight,keepaspectratio]{index_files/mediabag/imagenes_cap2/2.7-velocidades.pdf}
\end{center}

Ejemplo

Se deja caer una pelota desde la plataforma superior de observación de
la Torre CN, a 450 \(m\) sobre el nivel del suelo.

\begin{itemize}
\tightlist
\item
  ¿Cuál es la velocidad de la pelota después de 5 segundos?
\item
  ¿Con qué velocidad viaja cuando choca contra el suelo?
\end{itemize}

\begin{tcolorbox}[enhanced jigsaw, opacityback=0, left=2mm, opacitybacktitle=0.6, toprule=.15mm, bottomrule=.15mm, colbacktitle=quarto-callout-tip-color!10!white, titlerule=0mm, bottomtitle=1mm, title=\textcolor{quarto-callout-tip-color}{\faLightbulb}\hspace{0.5em}{Solución}, toptitle=1mm, colframe=quarto-callout-tip-color-frame, colback=white, breakable, rightrule=.15mm, arc=.35mm, leftrule=.75mm, coltitle=black]

Podemos modelar el movimiento de la pelota con la función de posición
\(f(t)=4.9\,t^2\). Calculemos \(v(a)\), la velocidad en un instante de
tiempo genérico \(a\). Es decir, escribimos

\[
v(a)=\lim_{h\to 0} \frac{f(a+h)-f(a)}{h}=\lim_{h\to 0} \frac{4.9\,(a+h)^2-4.9\,a^2}{h}
\]

Si \(h\neq 0\) entonces

\[
\frac{4.9\,(a+h)^2-4.9\,a^2}{h}=\frac{4.9\, a^2+9.8\,ah+4.9\,h^2-4.9\,a^2}{h}=\frac{h(9.8\,a+4.9\,h)}{h}=9.8\,a+4.9\,h.
\]

Con esto, obtenemos que

\[
\lim_{h\to 0} \frac{4.9\,(a+h)^2-4.9\,a^2}{h}=\lim_{h\to 0} (9.8\,a+4.9\,h)=9.8\,a.
\]

Es decir, \(v(a)=9.8\, a\). En particular, si \(a=5\), la velocidad
después de 5 segundos es \(v(5)=9.8\cdot 5 =49\).

Para la segunda pregunta, primero debemos determinar \(t_0\) el instante
de tiempo en que la pelota toca el suelo. Para ello, debe ocurrir que el
espacio recorrido coincida con la altura de la torre, es decir
\(f(t_0)=450\), con lo cual

\[
450=4.9\, t_0^2,
\]

y despejando obtenemos \(t_0\approx 9.58\), con lo cual la velocidad de
la pelota en ese instante es \(v(t_0)\approx v(9.58)\approx 46.9\)
m/seg.

\end{tcolorbox}

\subsection{Derivadas}\label{derivadas}

Con las motivaciones vistas hasta ahora, estamos en condiciones de
definir el concepto de \textbf{derivada}.

\begin{tcolorbox}[enhanced jigsaw, opacityback=0, left=2mm, opacitybacktitle=0.6, toprule=.15mm, bottomrule=.15mm, colbacktitle=quarto-callout-note-color!10!white, titlerule=0mm, bottomtitle=1mm, title=\textcolor{quarto-callout-note-color}{\faInfo}\hspace{0.5em}{Definición (derivada)}, toptitle=1mm, colframe=quarto-callout-note-color-frame, colback=white, breakable, rightrule=.15mm, arc=.35mm, leftrule=.75mm, coltitle=black]

La \textbf{derivada de una función \(f\) en un número \(a\)} se denota
\(f'(a)\) y se define por \[
f'(a)=\lim_{h\to 0}\frac{f(a+h)-f(a)}{h},
\]

siempre que este límite exista. De manera equivalente, podemos escribir

\[
f'(a)=\lim_{x\to a}\frac{f(x)-f(a)}{x-a}.
\]

\end{tcolorbox}

Ejemplo

Hallar \(f'(a)\) si \(f(x)=x^2-8x+9\).

\begin{tcolorbox}[enhanced jigsaw, opacityback=0, left=2mm, opacitybacktitle=0.6, toprule=.15mm, bottomrule=.15mm, colbacktitle=quarto-callout-tip-color!10!white, titlerule=0mm, bottomtitle=1mm, title=\textcolor{quarto-callout-tip-color}{\faLightbulb}\hspace{0.5em}{Solución}, toptitle=1mm, colframe=quarto-callout-tip-color-frame, colback=white, breakable, rightrule=.15mm, arc=.35mm, leftrule=.75mm, coltitle=black]

Por definición es

\[
f'(a)=\lim_{h\to 0} \frac{f(a+h)-f(a)}{h}.
\]

Si \(h\neq 0\), entonces

\[
\frac{f(a+h)-f(a)}{h}=\frac{(a+h)^2-8(a+h)+9-[a^2-8a+9]}{h}=\frac{2ah+h^2-8h}{h}=2a+h-8.
\]

Entonces

\[
\lim_{h\to 0} \frac{f(a+h)-f(a)}{h}=\lim_{h\to 0}\,(2a+h-8)=2a-8.
\]

\end{tcolorbox}

La \textbf{recta tangente} a la curva \(y=f(x)\) en el punto
\(P=(a,f(a))\) es la recta que pasa por \(P\) y cuya pendiente es
\(f'(a)\), es decir, la derivada de \(f\) en \(a\).

Usando la \hyperref[recta-punto-pendiente]{ecuación punto-pendiente}
podemos encontrar la ecuación de esta recta.

\phantomsection\label{form-recta-tangente}
Ecuación de la recta tangente al gráfico de \(y=f(x)\) en el punto
\((a,f(a))\)

\[
 y=f(a)+f'(a)(x-a).
\]

\begin{center}
\includegraphics[width=0.5\linewidth,height=\textheight,keepaspectratio]{index_files/mediabag/imagenes_cap2/2.7-recta-tangente.pdf}
\end{center}

Ejemplo

Hallar la ecuación de la recta tangente al gráfico de \(y=x^2-8x+9\) en
\((3,-6)\).

\begin{tcolorbox}[enhanced jigsaw, opacityback=0, left=2mm, opacitybacktitle=0.6, toprule=.15mm, bottomrule=.15mm, colbacktitle=quarto-callout-tip-color!10!white, titlerule=0mm, bottomtitle=1mm, title=\textcolor{quarto-callout-tip-color}{\faLightbulb}\hspace{0.5em}{Solución}, toptitle=1mm, colframe=quarto-callout-tip-color-frame, colback=white, breakable, rightrule=.15mm, arc=.35mm, leftrule=.75mm, coltitle=black]

En el ejemplo anterior vimos que \(f'(a)=2a-8\). Usando la
\hyperref[form-recta-tangente]{ecuación de la recta tangente} con
\(a=3\) obtenemos

\[
y=f(3)+f'(3)(x-3)
\]

y como \(f'(3)=2\cdot 3-8=-2\), resulta \[
y=-6-2(x-3) \quad \text{o bien}\quad y=-2x.
\]

\end{tcolorbox}

\subsection{Relaciones de cambio}\label{relaciones-de-cambio}

Supongamos que \(y\) depende de \(x\) mediante la relación \(y=f(x)\).
Si \(x\) cambia de \(x_1\) a \(x_2\), el \textbf{incremento en \(x\)}
está dado por

\[
\Delta x=x_2-x_1
\]

mientras que el correspondiente \textbf{incremento en \(y\)} es

\[
\Delta y=f(x_2)-f(x_1).
\]

El cociente de diferencias

\[
\frac{\Delta y}{\Delta x}=\frac{f(x_2)-f(x_1)}{x_2-x_1}
\]

se llama \textbf{razón de cambio promedio} de \(y\) con respecto a \(x\)
en el intervalo \([x_1,x_2]\). Ésta puede interpretarse como la
pendiente de la recta secante \(PQ\) del siguiente gráfico

\begin{center}
\includegraphics[width=0.5\linewidth,height=\textheight,keepaspectratio]{index_files/mediabag/imagenes_cap2/2.7-relaciones-cambio.pdf}
\end{center}

Si consideramos intervalos \([x_1,x_2]\) cada vez más pequeños (esto es,
hacemos \(x_2\) acercarse a \(x_1\)) el límite de estas razones de
cambio promedio (si existe) se llama \textbf{razón de cambio
instantánea} de \(y\) con respecto a \(x\) en \(x=x_1\), y puede
interpretarse como la pendiente de la recta tangente a la curva en el
punto \(P\).

\[
\text{Razón de cambio instantánea} =\lim_{\Delta x\to  0}\frac{\Delta y}{\Delta x}=\lim_{x_2\to x_1}\frac{f(x_2)-f(x_1)}{x_2-x_1}.
\]

El límite de arriba resulta ser \(f'(x_1)\). Esto nos da otra
interpretación para la derivada:

La derivada \(f'(a)\) es la razón de cambio instantánea de \(y=f(x)\)
con respecto a \(x\), cuando \(x=a\).

Si dibujamos la curva \(y=f(x)\), la razón de cambio instantánea es la
pendiente de la tangente a esta curva en el punto \((a,f(a))\). Cuando
la derivada es grande, los valores de \(y\) cambian rápidamente respecto
de los de \(x\) y en consecuencia la curva es escarpada, como sucede en
el punto \(P\). Cuando la derivada es pequeña, los valores de \(y\)
cambian lentamente en relación a \(x\) y la curva es relativamente
plana, como en el punto \(Q\).

\begin{center}
\includegraphics[width=0.4\linewidth,height=\textheight,keepaspectratio]{index_files/mediabag/imagenes_cap2/2.7-relaciones-cambio2.pdf}
\end{center}

Cuando \(s=f(t)\) es la función posición de una partícula a tiempo \(t\)
que se mueve a lo largo de una línea recta, \(f'(a)\) es la
\textbf{velocidad} de la partícula en \(t=a\).

La \textbf{rapidez} de la partícula es el valor absoluto de la
velocidad, es decir \(|f'(a)|\).

Esto muestra que la velocidad de una partícula puede ser positiva,
negativa o cero, mientras que la rapidez siempre es no negativa.

\hyperref[seccion_2.7]{↑ Volver al inicio de la sección}

\section{2.8. La derivada como una función}\label{seccion_2.8}

Hemos definido la derivada de \(f\) en \(a\) como el límite

\[
f'(a)=\lim_{h\to 0}\frac{f(a+h)-f(a)}{h}
\]

siempre que este límite exista. Si dejamos que el número \(a\) varíe,
podemos escribir

\[
f'(x)=\lim_{h\to 0}\frac{f(x+h)-f(x)}{h}.
\]

De modo que \(f'(x)\) está bien definida para todo número \(x\) tal que
el límite anterior existe. Entonces podemos considerar a \(f'\) como una
nueva función, definida como el límite de arriba, y llamada
\textbf{derivada} de \(f\).

El valor que toma \(f'\) en \(x\), denotado \(f'(x)\) representa la
pendiente de la recta tangente al gráfico de \(y=f(x)\) en el punto
\((x,f(x))\).

Observemos que

\[
D_{f'}=\{x\in D_f: f'(x) \text{ existe}\}.
\]

En general, \(D_{f'}\subseteq D_{f}\), pudiendo esta contención ser
estricta.

Ejemplo

En la siguiente figura se muestra el gráfico de una función \(f\).
Usarla para dibujar el gráfico de \(f'\).

\begin{center}
\includegraphics[width=0.5\linewidth,height=\textheight,keepaspectratio]{index_files/mediabag/imagenes_cap2/2.8-ejemplo-derivada.pdf}
\end{center}

\begin{tcolorbox}[enhanced jigsaw, opacityback=0, left=2mm, opacitybacktitle=0.6, toprule=.15mm, bottomrule=.15mm, colbacktitle=quarto-callout-tip-color!10!white, titlerule=0mm, bottomtitle=1mm, title=\textcolor{quarto-callout-tip-color}{\faLightbulb}\hspace{0.5em}{Solución}, toptitle=1mm, colframe=quarto-callout-tip-color-frame, colback=white, breakable, rightrule=.15mm, arc=.35mm, leftrule=.75mm, coltitle=black]

Para realizar el gráfico de \(f'\) podemos seleccionar distintos puntos
\((x,f(x))\) en el gráfico de \(f\) e intentar estimar la pendiente de
la recta tangente allí. Este método, si bien nos ayuda a visualizar el
gráfico de \(f'\) de manera aproximada, no es exacto y la técnica es un
poco engorrosa.

Más adelante veremos herramientas que nos permitirán dar un esbozo más
preciso del gráfico de una función.

\includegraphics[width=0.6\linewidth,height=\textheight,keepaspectratio]{index_files/mediabag/imagenes_cap2/2.8-ejemplo-derivada2.pdf}

\end{tcolorbox}

\phantomsection\label{ejemplo-derivada-cubica}
Ejemplo

Si \(f(x)=x^3-x\), encontrar una fórmula para \(f'(x)\) e ilustrarla
comparando las gráficas de \(f\) y \(f'\).

\begin{tcolorbox}[enhanced jigsaw, opacityback=0, left=2mm, opacitybacktitle=0.6, toprule=.15mm, bottomrule=.15mm, colbacktitle=quarto-callout-tip-color!10!white, titlerule=0mm, bottomtitle=1mm, title=\textcolor{quarto-callout-tip-color}{\faLightbulb}\hspace{0.5em}{Solución}, toptitle=1mm, colframe=quarto-callout-tip-color-frame, colback=white, breakable, rightrule=.15mm, arc=.35mm, leftrule=.75mm, coltitle=black]

Por definición, debemos calcular

\[
\lim_{h\to 0} \frac{f(x+h)-f(x)}{h}
\]

para cualquier \(x\in\mathbb{R}\).

Fijado \(x\), si \(h\neq 0\), entonces tenemos que

\[ \small
\frac{f(x+h)-f(x)}{h}=\frac{(x+h)^3-(x+h)-[x^3-x]}{h}=\frac{3x^2h+3xh^2+h^3-h}{h}=3x^2+3xh+h^2-1.
\]

Entonces

\[
\lim_{h\to 0} \frac{f(x+h)-f(x)}{h}=\lim_{h\to 0} (3x^2+3xh+h^2-1)=3x^2-1.
\]

Por lo tanto, \(f'(x)=3x^2-1\).

\includegraphics[width=0.5\linewidth,height=\textheight,keepaspectratio]{index_files/mediabag/imagenes_cap2/2.8-ejemplo-funcyder.pdf}

\end{tcolorbox}

Ejemplo

Si \(f(x)=\sqrt{x}\), calcular \(f'(x)\) y dar su dominio.

\begin{tcolorbox}[enhanced jigsaw, opacityback=0, left=2mm, opacitybacktitle=0.6, toprule=.15mm, bottomrule=.15mm, colbacktitle=quarto-callout-tip-color!10!white, titlerule=0mm, bottomtitle=1mm, title=\textcolor{quarto-callout-tip-color}{\faLightbulb}\hspace{0.5em}{Solución}, toptitle=1mm, colframe=quarto-callout-tip-color-frame, colback=white, breakable, rightrule=.15mm, arc=.35mm, leftrule=.75mm, coltitle=black]

Vamos a calcular

\[
\lim_{h\to 0} \frac{f(x+h)-f(x)}{h}
\]

para \(x\geq 0\).

Fijado \(x\), tomando \(h\neq 0\) escribimos

\[
\frac{f(x+h)-f(x)}{h}=\frac{\sqrt{x+h}-\sqrt{x}}{h}=\frac{\sqrt{x+h}-\sqrt{x}}{h}\,\frac{\sqrt{x+h}+\sqrt{x}}{\sqrt{x+h}+\sqrt{x}}=\frac{1}{\sqrt{x+h}+\sqrt{x}}.
\]

Utilizando las leyes \textbf{(11)} y \textbf{(1)} obtenemos

\[
\lim_{h\to 0} \frac{f(x+h)-f(x)}{h}=\lim_{h\to 0} \frac{1}{\sqrt{x+h}+\sqrt{x}}=\frac{1}{2\sqrt{x}},
\]

por lo tanto \(f'(x)=\displaystyle \frac{1}{2\sqrt{x}}\), siempre que
\(x>0\).

¿Es \(f\) derivable en \(x=0\)? Para responder esta pregunta, analizamos

\[
\lim_{h\to 0^+} \frac{f(0+h)-f(0)}{h}=\lim_{h\to 0^+} \frac{\sqrt{h}-0}{h}=\lim_{h\to 0^+} \frac{1}{\sqrt{h}}=\infty,
\]

con lo cual el límite no existe por infinitud y \(f\) no es derivable en
\(x=0\). Entonces \(D_{f'}=(0,\infty)\).

\end{tcolorbox}

Ejemplo

Encontrar \(f'(x)\) si \(\displaystyle f(x)=\frac{1-x}{2+x}\).

\begin{tcolorbox}[enhanced jigsaw, opacityback=0, left=2mm, opacitybacktitle=0.6, toprule=.15mm, bottomrule=.15mm, colbacktitle=quarto-callout-tip-color!10!white, titlerule=0mm, bottomtitle=1mm, title=\textcolor{quarto-callout-tip-color}{\faLightbulb}\hspace{0.5em}{Solución}, toptitle=1mm, colframe=quarto-callout-tip-color-frame, colback=white, breakable, rightrule=.15mm, arc=.35mm, leftrule=.75mm, coltitle=black]

Notemos que \(D_f=\mathbb{R}\backslash {-2}\). Para \(x\neq -2\),
calculamos

\[
\lim_{h\to 0} \frac{f(x+h)-f(x)}{h}.
\]

Fijemos \(x\) y \(h\neq 0\). Entonces

\[ \small
\frac{f(x+h)-f(x)}{h}=\frac{\frac{1-(x+h)}{2+(x+h)}-\frac{1-x}{2+x}}{h}=\frac{1}{h}\, \frac{(2+x)(1-x-h)-(2+x+h)(1-x)}{(2+x+h)(2+x)}=-\frac{3}{(2+x+h)(2+x)}.
\]

Por la \hyperref[teo-sust-directa]{propiedad de sustitución directa} y
las leyes \textbf{(4)} y \textbf{(5)} del producto y del cociente
resulta

\[
\lim_{h\to 0} \frac{f(x+h)-f(x)}{h}=\lim_{h\to 0} -\frac{3}{(2+x+h)(2+x)}=-\frac{3}{(2+x)^2}.
\]

\end{tcolorbox}

\subsection{Otras notaciones}\label{otras-notaciones}

Si escribimos \(y=f(x)\), donde \(x\) es la variable independiente e
\(y\) la dependiente, otras formas de representar a la derivada de \(f\)
en \(x\) son

\[
f'(x)=y'={dx}={dx}={dx}f(x)=Df(x)=D_xf(x).
\]

Los símbolos \(D\) y \(d/dx\) se llaman \textbf{operadores de
derivación}. El símbolo \(dy/dx\) introducido por Leibniz es otra forma
de escribir \(f'(x)\). En término de incrementos, tenemos que

\[
\frac{dy}{dx}=\lim_{\Delta x\to 0} \frac{\Delta y}{\Delta x}.
\]

\begin{tcolorbox}[enhanced jigsaw, opacityback=0, left=2mm, opacitybacktitle=0.6, toprule=.15mm, bottomrule=.15mm, colbacktitle=quarto-callout-note-color!10!white, titlerule=0mm, bottomtitle=1mm, title=\textcolor{quarto-callout-note-color}{\faInfo}\hspace{0.5em}{Definición (Función derivable)}, toptitle=1mm, colframe=quarto-callout-note-color-frame, colback=white, breakable, rightrule=.15mm, arc=.35mm, leftrule=.75mm, coltitle=black]

Una función \(f\) es \textbf{derivable en \(a\)} si \(f'(a)\) existe.

Si \$-\infty\leq a\textless b\leq \infty \$, decimos que \(f\) es
\textbf{derivable en \((a,b)\)} si \(f\) es derivable en todo punto del
intervalo.

\end{tcolorbox}

Ejemplo

¿Dónde es derivable la función \(f(x)=|x|\)?

\begin{tcolorbox}[enhanced jigsaw, opacityback=0, left=2mm, opacitybacktitle=0.6, toprule=.15mm, bottomrule=.15mm, colbacktitle=quarto-callout-tip-color!10!white, titlerule=0mm, bottomtitle=1mm, title=\textcolor{quarto-callout-tip-color}{\faLightbulb}\hspace{0.5em}{Solución}, toptitle=1mm, colframe=quarto-callout-tip-color-frame, colback=white, breakable, rightrule=.15mm, arc=.35mm, leftrule=.75mm, coltitle=black]

Consideremos diferentes casos. Supongamos primero que \(x>0\) y
\(h\neq 0\) es suficientemente pequeño como para que \(x+h>0\), esto es,
elegimos \(h>-x\). Entonces

\[
\frac{f(x+h)-f(x)}{h}=\frac{|x+h|-|x|}{h}=\frac{x+h-x}{h}=1,
\] dado que tanto \(x\) como \(x+h\) son positivos.

Luego

\[
\lim_{h\to 0} \frac{f(x+h)-f(x)}{h} = \lim_{h\to 0} 1=1.
\]

Con esto concluimos que \(f'(x)=1\) para cualquier \(x>0\).

Ahora consideremos \(x<0\) y \(h\neq 0\) suficientemente pequeño en
magnitud como para que \(x+h<0\). Entonces

\[
\frac{f(x+h)-f(x)}{h}=\frac{|x+h|-|x|}{h}=\frac{-x-h+x}{h}=-1,
\] pues ambos \(x\) y \(x+h\) son negativos. Así,

\[
\lim_{h\to 0} \frac{f(x+h)-f(x)}{h} = \lim_{h\to 0} -1=-1.
\]

Por lo tanto, \(f'(x)=-1\) para todo \(x<0\).

Finalmente, estudiamos el caso \(x=0\). Como \(f\) cambia su definición
en \(x=0\), estudiamos los límites laterales

\[
\lim_{h\to 0^+} \frac{f(0+h)-f(0)}{h}=\lim_{h\to 0^+} \frac{h-0}{h}= 1
\]

y

\[
\lim_{h\to 0^-} \frac{f(0+h)-f(0)}{h}=\lim_{h\to 0^+} \frac{-h-0}{h}= -1.
\]

Como los límites laterales no coinciden, el límite

\[
\lim_{h\to 0} \frac{f(0+h)-f(0)}{h}
\]

no existe y por lo tanto \(f\) no es derivable en \(x=0\). Es decir,
\(f(x)=|x|\) es derivable en todo \(x\neq 0\) y

\[
f'(x)=
\begin{cases}
1 & \text{ si } x>0,\\
-1 & \text{ si } x<0.
\end{cases}
\]

\end{tcolorbox}

\phantomsection\label{teo-derivable-continua}
Teorema

Si \(f\) es derivable en \(a\), entonces \(f\) es continua en \(a\).

\begin{tcolorbox}[enhanced jigsaw, opacityback=0, left=2mm, opacitybacktitle=0.6, toprule=.15mm, bottomrule=.15mm, colbacktitle=quarto-callout-caution-color!10!white, titlerule=0mm, bottomtitle=1mm, title=\textcolor{quarto-callout-caution-color}{\faFire}\hspace{0.5em}{Importante}, toptitle=1mm, colframe=quarto-callout-caution-color-frame, colback=white, breakable, rightrule=.15mm, arc=.35mm, leftrule=.75mm, coltitle=black]

El recíproco del \hyperref[teo-derivable-continua]{teorema} anterior es
\textbf{falso}. En efecto, la función \(f(x)=|x|\) es continua en
\(x=0\) pero no es derivable en este punto.

\end{tcolorbox}

\subsection{¿Cómo deja de ser derivable una
función?}\label{cuxf3mo-deja-de-ser-derivable-una-funciuxf3n}

Si la gráfica de una función presenta ``esquinas'' o ``rizos'' en
ciertos puntos, \(f\) no tiene recta tangente allí y resulta no
derivable. Esto se debe a que al intentar calcular \(f'(a)\), los
límites laterales son distintos.

\begin{center}
\includegraphics[width=0.4\linewidth,height=\textheight,keepaspectratio]{index_files/mediabag/imagenes_cap2/2.8-esquina.pdf}
\end{center}

El \hyperref[teo-derivable-continua]{teorema} anterior da otra forma en
que una función no es derivable en un punto: si \(f\) no es continua en
\(a\), entonces no puede ser derivable en este punto. Si \(f\) tiene una
discontinuidad en \(a\) (como por ejemplo un salto) entonces \(f'(a)\)
no existe.

\begin{center}
\includegraphics[width=0.4\linewidth,height=\textheight,keepaspectratio]{index_files/mediabag/imagenes_cap2/2.8-discontinuidad.pdf}
\end{center}

Una tercera posibilidad es que \(f\) presente una \textbf{recta tangente
vertical} en \(x=a\), es decir, \(f\) es continua en \(a\) y se cumple
que

\[
\lim_{x\to a}|f'(x)|=\infty.
\]

Esto significa que las tangentes se vuelven más empinadas conforme
\(x\to a\).

\begin{center}
\includegraphics[width=0.4\linewidth,height=\textheight,keepaspectratio]{index_files/mediabag/imagenes_cap2/2.8-pendienteinf.pdf}
\end{center}

\subsection{Derivadas superiores}\label{derivadas-superiores}

Si \(f\) es una función derivable, \(f'\) también es una función y
podemos derivarla. La derivada de la derivada de \(f\) (\((f')'=f''\))
se denomina \textbf{derivada segunda} de \(f\). En la notación de
Leibniz, si \(y=f(x)\) podemos escribir

\[
f''(x)=\frac{d}{dx}\left(\frac{df}{dx}\right)=\frac{d^2f}{dx^2}.
\]

\phantomsection\label{ejemplo-derivada-segunda}
Ejemplo

Si \(f(x)=x^3-x\), hallar e interpretar \(f''(x)\).

\begin{tcolorbox}[enhanced jigsaw, opacityback=0, left=2mm, opacitybacktitle=0.6, toprule=.15mm, bottomrule=.15mm, colbacktitle=quarto-callout-tip-color!10!white, titlerule=0mm, bottomtitle=1mm, title=\textcolor{quarto-callout-tip-color}{\faLightbulb}\hspace{0.5em}{Solución}, toptitle=1mm, colframe=quarto-callout-tip-color-frame, colback=white, breakable, rightrule=.15mm, arc=.35mm, leftrule=.75mm, coltitle=black]

En el \hyperref[ejemplo-derivada-cubica]{Ejemplo 41} vimos que
\(f'(x)=3x^2-1\). Usando la \hyperref[definicion-derivada]{definición de
derivada}, tenemos que

\[
f''(x)=\lim_{h\to 0} \frac{f'(x+h)-f'(x)}{h}.
\]

Para \(x\) fijo y \(h\neq 0\), calculamos

\[
\frac{f'(x+h)-f'(x)}{h}=\frac{3(x+h)^2-1-(3x^2-1)}{h}=\frac{6xh+3h^2}{h}=\frac{h(6x+3h)}{h}=6x+3h.
\]

Luego

\[
\lim_{h\to 0} \frac{f'(x+h)-f'(x)}{h}=\lim_{h\to 0} (6x+3h)=6x.
\]

Por lo tanto, \(f''(x)=6x\). Como \(f''\) resulta ser la derivada de
\(f'\), \(f''(x)\) nos da la pendiente de la recta tangente al gráfico
de \(f'\) en el punto \((x,f'(x))\).

\includegraphics[width=0.4\linewidth,height=\textheight,keepaspectratio]{index_files/mediabag/imagenes_cap2/2.8-ejemplo-derivadas.pdf}

\end{tcolorbox}

En general podemos interpretar a la segunda derivada como una razón de
cambio de una razón de cambio. Si \(s=s(t)\) es la función de posición
de un objeto que se mueve en línea recta, entonces

\[
v(t)=s'(t)=\frac{ds}{dt}
\]

es la velocidad del objeto. La razón de cambio de la velocidad con
respecto al tiempo es la \textbf{aceleración} del objeto. Ésta es la
derivada de la velocidad y, en consecuencia, la derivada segunda de la
posición

\[
a(t)=\frac{d}{dt}\left(\frac{ds}{dt}\right)=\frac{d^2s}{dt^2}
\] en la notación de Leibniz.

La \textbf{tercera derivada} \(f'''\) es la derivada de la derivada
segunda de \(f\), es decir

\[
y'''=f'''=\frac{d}{dx}\left(\frac{d^2y}{dx^2}\right)=\frac{d^3y}{dx^3}.
\]

Este proceso puede continuar siempre y cuando el límite del cociente de
diferencias exista. En general, la \textbf{derivada \(n-\)ésima} se
denota por \(f^{(n)}\), o bien

\[
y^{(n)}=f^{(n)}(x)=\frac{d^nf}{dx^n}.
\]

Ejemplo

Si \(f(x)=x^3-x\), hallar \(f'''(x)\) e interpretar \(f^{(4)}(x)\).

\begin{tcolorbox}[enhanced jigsaw, opacityback=0, left=2mm, opacitybacktitle=0.6, toprule=.15mm, bottomrule=.15mm, colbacktitle=quarto-callout-tip-color!10!white, titlerule=0mm, bottomtitle=1mm, title=\textcolor{quarto-callout-tip-color}{\faLightbulb}\hspace{0.5em}{Solución}, toptitle=1mm, colframe=quarto-callout-tip-color-frame, colback=white, breakable, rightrule=.15mm, arc=.35mm, leftrule=.75mm, coltitle=black]

En el \hyperref[ejemplo-derivada-segunda]{Ejemplo 45} obtuvimos que
\(f''(x)=6x\). Usando la \hyperref[definicion-derivada]{definición de
derivada} con \(f''\), obtenemos

\[
f'''(x)=\lim_{h\to 0} \frac{f''(x+h)-f''(x)}{h}.
\]

Si \(x\in\mathbb{R}\) está fijo y \(h\neq 0\) tenemos que

\[
\frac{f''(x+h)-f''(x)}{h}=\frac{6(x+h)-6x}{h}=\frac{6h}{h}=6,
\]

con lo cual

\[
\lim_{h\to 0} \frac{f''(x+h)-f''(x)}{h}=\lim_{h\to 0} 6= 6
\]

y en consecuencia \(f'''(x)=6\), es decir, \(f'''\) es una función
constante. Si calculamos la derivada de ésta última función, obtenemos

\[
f^{(4)}(x)=\lim_{h\to 0} \frac{f'''(x+h)-f'''(x)}{h}=\lim_{h\to 0} \frac{6-6}{h}=0.
\]

Por lo tanto, \(f^{(4)}(x)=0\). En el siguiente gráfico se muestra
\(f''\) junto a \(f'''\) y \(f^{(4)}\).

\includegraphics[width=0.4\linewidth,height=\textheight,keepaspectratio]{index_files/mediabag/imagenes_cap2/2.8-ejemplo-derivadas2.pdf}

\end{tcolorbox}

Para el caso de la función posición de un objeto, la tercera derivada
\(s'''\) se denomina \textbf{jerk} (impulso) y es la derivada de la
aceleración. Tenemos que

\[
j=\frac{da}{dt}=\frac{d^3s}{dt^3}.
\]

\hyperref[seccion_2.8]{↑ Volver al inicio de la sección}

\bookmarksetup{startatroot}

\chapter{3. Derivadas}\label{derivadas-1}

\section{3.1. Derivadas de polinomios y funciones
exponenciales}\label{seccion_3.1}

Comenzamos considerando la función constante \(f(x)=c\), donde \(c\) es
un número real.

\begin{center}
\includegraphics[width=0.25\linewidth,height=\textheight,keepaspectratio]{index_files/mediabag/imagenes_cap3/3.1-constante.pdf}
\end{center}

El gráfico de \(f\) es simplemente una recta horizontal, cuya pendiente
es cero. Entonces debería ocurrir que \(f'(x)=0\). En efecto,

\[
f'(x)=\lim_{h\to 0} \frac{f(x+h)-f(x)}{h}=\lim_{h\to 0} \frac{c-c}{h}=\lim_{h\to 0}0=0.
\]

\[
\frac{d}{dx}(c)=0.
\]

\subsection{Funciones potencia}\label{funciones-potencia}

Si consideramos \(f(x)=x\), el gráfico es la recta identidad, que tiene
pendiente uno.

\begin{center}
\includegraphics[width=0.3\linewidth,height=\textheight,keepaspectratio]{index_files/mediabag/imagenes_cap3/3.1-lineal.pdf}
\end{center}

Entonces

\[
f'(x)=\lim_{h\to 0} \frac{f(x+h)-f(x)}{h}=\lim_{h\to 0} \frac{x+h-x}{h}=\lim_{h\to 0}\frac{h}{h}=1.
\]

Por lo tanto

\[
\frac{d}{dx}(x)=1.
\]

Se puede comprobar fácilmente que

\[
\frac{d}{dx}\left(x^2\right)=2x \quad \text { y } \quad \frac{d}{dx}\left(x^3\right)=3x^2. 
\]

Para el caso \(f(x)=x^4\) podemos comprobar que

\[
f'(x)=\lim_{h\to 0} \frac{f(x+h)-f(x)}{h}=\lim_{h\to 0} \frac{(x+h)^4-x^4}{h}=\lim_{h\to 0}\left(4x^3+6x^2h+4xh^2+h^3\right)=4x^3.
\]

Por lo tanto

\[
\frac{d}{dx}\left(x^4\right)=4x^3.
\]

En general, tenemos la siguiente \textbf{regla de la potencia}: si \(n\)
es un entero positivo, entonces

\begin{equation}\phantomsection\label{eq-regla-potencia}{
\frac{d}{dx}\left(x^n\right)=nx^{n-1}.
}\end{equation}

Ejemplo

Utilizar la regla de la potencia para hallar la derivada en cada caso.

\begin{itemize}
\item
  \(f(x)=x^6\).
\item
  \(y=x^{1000}\).
\end{itemize}

\begin{itemize}
\item
  \(y=t^4\).
\item
  \(g(r)=r^3\).
\end{itemize}

\begin{tcolorbox}[enhanced jigsaw, opacityback=0, left=2mm, opacitybacktitle=0.6, toprule=.15mm, bottomrule=.15mm, colbacktitle=quarto-callout-tip-color!10!white, titlerule=0mm, bottomtitle=1mm, title=\textcolor{quarto-callout-tip-color}{\faLightbulb}\hspace{0.5em}{Solución}, toptitle=1mm, colframe=quarto-callout-tip-color-frame, colback=white, breakable, rightrule=.15mm, arc=.35mm, leftrule=.75mm, coltitle=black]

Aplicando la \hyperref[eq-regla-potencia]{regla de la potencia} en cada
caso, obtenemos

\begin{itemize}
\item
  \(\,\,\displaystyle \frac{df}{dx}(x)=\frac{d}{dx}(x^6)=6x^{6-1}=6x^5\).
\item
  \(\,\,\displaystyle \frac{dy}{dx}(x)=\frac{d}{dx}(x^1000)=1000x^{1000-1}=1000x^999\).
\item
  \(\,\,\displaystyle \frac{dy}{dt}(t)=\frac{d}{dt}(t^4)=4t^{4-1}=4t^3\).
\item
  \(\,\,\displaystyle \frac{dg}{dr}(r)=\frac{d}{dr}(r^3)=3r^{3-1}=3r^2\).
\end{itemize}

\end{tcolorbox}

\subsection{\texorpdfstring{¿Qué ocurre cuando \(n\) es entero
negativo?}{¿Qué ocurre cuando n es entero negativo?}}\label{quuxe9-ocurre-cuando-n-es-entero-negativo}

\begin{itemize}
\tightlist
\item
  Consideremos primero \(f(x)=1/x\). Podemos ver que
\end{itemize}

\[
\frac{f(x+h)-f(x)}{h}=\frac{\frac{1}{x+h}-\frac{1}{x}}{h}=-\frac{1}{x(x+h)},
\]

para todo \(h\neq 0\). Entonces

\[
\lim_{h\to 0}\frac{f(x+h)-f(x)}{h}=\lim_{h\to 0}\left(-\frac{1}{x(x+h)}\right)=-\frac{1}{x^2}=(-1)x^{-1-1}.
\]

\begin{itemize}
\tightlist
\item
  En la Sección 2.8 vimos que
\end{itemize}

\[
\frac{d}{dx}\left(\sqrt{x}\right)=\frac{1}{2\sqrt{x}}=\frac{1}{2}x^{1/2-1}.
\]

Esto parece indicar que la regla de la potencia vale para cualquier
\(n\in\mathbb{R}\). Es decir, si \(f(x)=x^n\) con \(n\in\mathbb{R}\),
entonces

\begin{equation}\phantomsection\label{eq-potencia-general}{
\frac{d}{dx}\left(x^n\right)=nx^{n-1}.
}\end{equation}

Esto lo comprobaremos más adelante.

Ejemplo

Derivar las funciones

\begin{itemize}
\tightlist
\item
  \(f(x)=1/x^2\)
\item
  \(y=\sqrt[3]{x^2}\)
\end{itemize}

\begin{tcolorbox}[enhanced jigsaw, opacityback=0, left=2mm, opacitybacktitle=0.6, toprule=.15mm, bottomrule=.15mm, colbacktitle=quarto-callout-tip-color!10!white, titlerule=0mm, bottomtitle=1mm, title=\textcolor{quarto-callout-tip-color}{\faLightbulb}\hspace{0.5em}{Solución}, toptitle=1mm, colframe=quarto-callout-tip-color-frame, colback=white, breakable, rightrule=.15mm, arc=.35mm, leftrule=.75mm, coltitle=black]

Utilizando la \hyperref[eq-potencia-general]{regla general de la
potencia}, obtenemos

\[
\frac{df}{dx}(x)=\frac{d}{dx}\left(x^{-2}\right)=(-2)x^{-2-1}=-2x^{-3}=-\frac{2}{x^3}
\]

y también

\[
\frac{dy}{dx}(x)=\frac{d}{dx}\left(\sqrt[3]{x^2}\right)=\frac{d}{dx}\left(x^{2/3}\right)=\frac{2}{3}x^{2/3-1}=\frac{2}{3}x^{-1/3}=\frac{2}{3 \sqrt[3]{x}}.
\]

\end{tcolorbox}

\subsection{Rectas tangente y normal}\label{rectas-tangente-y-normal}

A partir de la derivada de una función podemos encontrar también rectas
normales a una curva en un punto dado. La \textbf{recta normal} a una
curva \(C\) en un punto \(P\) es aquella que pasa por \(P\) y es
ortogonal a la tangente por \(C\) en ese punto.

\begin{center}
\includegraphics[width=0.4\linewidth,height=\textheight,keepaspectratio]{index_files/mediabag/imagenes_cap3/3.1-tangente-normal.pdf}
\end{center}

Supongamos que la curva \(C\) está dada por el gráfico de \(y=f(x)\) y
\(f\) es derivable en \(a\). Recordemos que la ecuación de la recta
tangente a \(C\) en el punto \((a,f(a))\) es

\[
y=f(a)+f'(a)(x-a).
\]

La recta normal es aquella perpendicular a la tangente y que pasa por
\((a,f(a))\). Su ecuación es

\[
y=f(a)-\frac{1}{f'(a)}(x-a)
\]

siempre que \(f'(a)\neq 0\). Cuando \(f'(a)=0\), la recta normal será
simplemente \(x=a\).

Ejemplo

Hallar la ecuación de la recta tangente y de la recta normal a la curva
\(y=x\sqrt{x}\) en el punto \((1,1)\). Dibujar ambas rectas junto al
gráfico de la función.

\begin{tcolorbox}[enhanced jigsaw, opacityback=0, left=2mm, opacitybacktitle=0.6, toprule=.15mm, bottomrule=.15mm, colbacktitle=quarto-callout-tip-color!10!white, titlerule=0mm, bottomtitle=1mm, title=\textcolor{quarto-callout-tip-color}{\faLightbulb}\hspace{0.5em}{Solución}, toptitle=1mm, colframe=quarto-callout-tip-color-frame, colback=white, breakable, rightrule=.15mm, arc=.35mm, leftrule=.75mm, coltitle=black]

Para encontrar las rectas deseadas necesitamos calcular \(y'(1)\).
Utilizando la \hyperref[eq-regla-potencia]{regla general de la potencia}
resulta

\[
y'(x)=(x\sqrt{x})'=\left(x^{3/2}\right)'=\frac{3}{2} x^{1/2}.
\]

La recta tangente tiene pendiente \(y'(1)=\frac{3}{2}\) y pasa por
\((1,1)\). Su ecuación es

\[
y-1=\frac{3}{2}(x-1) \quad \text{ o equivalentemente } \quad y=\frac{3}{2}x-\frac{1}{2}.
\]

La recta normal en \(P\) tiene pendiente
\(m=-\frac{1}{y'(1)}=-\frac{2}{3}\) y pasa también por \((1,1)\). Su
ecuación es

\[
y-1=-\frac{2}{3}(x-1) \quad \text{ o también } \quad y=-\frac{2}{3}x+\frac{5}{3}.
\]

\includegraphics[width=0.4\linewidth,height=\textheight,keepaspectratio]{index_files/mediabag/imagenes_cap3/3.1-tangente-normal-ejemplo.pdf}

\end{tcolorbox}

\subsection{Nuevas derivadas a partir de otras
conocidas}\label{nuevas-derivadas-a-partir-de-otras-conocidas}

\phantomsection\label{teo-propiedades-derivada}
Propiedad

Sean \(f\) y \(g\) dos funciones derivables y \(c\) un número real.
Entonces valen las siguientes propiedades

\begin{itemize}
\item
  \(\displaystyle \frac{d}{dx}(cf(x))=c\frac{df}{dx}(x)\quad\)
  \textbf{(regla del múltiplo constante)}
\item
  \(\displaystyle \frac{d}{dx}(f(x)+g(x))=\frac{df}{dx}(x)+\frac{dg}{dx}(x)\quad\)\textbf{(regla
  de la suma)}
\item
  \(\displaystyle \frac{d}{dx}(f(x)-g(x))=\frac{df}{dx}(x)-\frac{dg}{dx}(x)\quad\)\textbf{(regla
  de la resta)}
\end{itemize}

Ejemplo

Calcular \(f'\) si \(f(x)=x^8+12x^5-4x^4+10x^3-8x+5\).

\begin{tcolorbox}[enhanced jigsaw, opacityback=0, left=2mm, opacitybacktitle=0.6, toprule=.15mm, bottomrule=.15mm, colbacktitle=quarto-callout-tip-color!10!white, titlerule=0mm, bottomtitle=1mm, title=\textcolor{quarto-callout-tip-color}{\faLightbulb}\hspace{0.5em}{Solución}, toptitle=1mm, colframe=quarto-callout-tip-color-frame, colback=white, breakable, rightrule=.15mm, arc=.35mm, leftrule=.75mm, coltitle=black]

Utilizamos las reglas de la
\hyperref[teo-propiedades-derivada]{propiedad} anterior para escribir

\[
f'(x)=(x^8+12x^5-4x^4+10x^3-8x+5)'=\left(x^8\right)'+12\left(x^5\right)'-4\left(x^4\right)'+10\left(x^3\right)'-8x'+5'.
\]

Ahora, por la \hyperref[eq-regla-potencia]{regla de la potencia}
obtenemos

\[
f'(x)=8x^7+60x^4-16x^3+30x^2-8.
\]

\end{tcolorbox}

Ejemplo

Encontrar los puntos de la curva \(y=x^4-6x^2+4\) donde la recta
tangente es horizontal.

\begin{tcolorbox}[enhanced jigsaw, opacityback=0, left=2mm, opacitybacktitle=0.6, toprule=.15mm, bottomrule=.15mm, colbacktitle=quarto-callout-tip-color!10!white, titlerule=0mm, bottomtitle=1mm, title=\textcolor{quarto-callout-tip-color}{\faLightbulb}\hspace{0.5em}{Solución}, toptitle=1mm, colframe=quarto-callout-tip-color-frame, colback=white, breakable, rightrule=.15mm, arc=.35mm, leftrule=.75mm, coltitle=black]

Una recta tangente horizontal tiene pendiente cero, es decir, queremos
encontrar los puntos \(x_0\) tales que \(y'(x_0)=0\). Combinando las
\hyperref[teo-propiedades-derivada]{reglas de derivación} para la suma,
resta y múltiplo constante con la \hyperref[eq-regla-potencia]{regla de
la potencia} obtenemos

\[
y'(x)=4x^3-12x.
\]

Entonces \(y'(x)=0\) si y sólo si \(4x(x^2-3)=0\), de donde deducimos
que \(x=-\sqrt{3}\), \(x=0\) y \(x=\sqrt{3}\). Luego, los puntos
buscados son \((-\sqrt{3},y(-\sqrt{3}))=(-\sqrt{3},-5)\),
\((0,y(0))=(0,4)\) y \((\sqrt{3},y(\sqrt{3}))=(\sqrt{3},-5)\).

\end{tcolorbox}

Ejemplo

La ecuación del movimiento de una partícula es \(s=2t^3-5t^2+3t+4\),
donde \(s\) se mide en centímetros y \(t\) en segundos. Hallar la
aceleración como una función del tiempo. ¿Cuál es la aceleración después
de dos segundos?

\begin{tcolorbox}[enhanced jigsaw, opacityback=0, left=2mm, opacitybacktitle=0.6, toprule=.15mm, bottomrule=.15mm, colbacktitle=quarto-callout-tip-color!10!white, titlerule=0mm, bottomtitle=1mm, title=\textcolor{quarto-callout-tip-color}{\faLightbulb}\hspace{0.5em}{Solución}, toptitle=1mm, colframe=quarto-callout-tip-color-frame, colback=white, breakable, rightrule=.15mm, arc=.35mm, leftrule=.75mm, coltitle=black]

Recordemos que la aceleración es la derivada de la velocidad, es decir,
la derivada segunda de la posición. Entonces

\[
a(t)=\frac{d^2 s}{dt^2}(t)=\frac{d}{dt}\left(\frac{ds}{dt}(t)\right)=\frac{d}{dt}\left(6t^2-10t+3\right)=12t-10.
\]

La aceleración después de dos segundos es
\(a(2)=12\cdot 2-10=14\, cm/s^2\).

\end{tcolorbox}

\subsection{Función exponencial}\label{funciuxf3n-exponencial}

Si \(f(x)=a^x\), al calcular \(f'\) por definición tenemos que

\[
f'(x)=\lim_{h\to 0} \frac{f(x+h)-f(x)}{h}=\lim_{h\to 0} \frac{a^{x+h}-a^x}{h}=\lim_{h\to 0} a^x\left(\frac{a^h-1}{h}\right).
\]

Notemos que si el límite

\[
\lim_{h\to 0} \frac{a^h-1}{h}
\]

existe, entonces es \(f'(0)\) y en consecuencia

\[
f'(x)=f'(0) \,a^x=f'(0)\,f(x).
\]

Esto nos dice que si una función exponencial es derivable en cero,
entonces es derivable en todo \(x\) y la derivada es proporcional a la
misma función.

Recordando el desarrollo de la Sección 1.5, hemos definido al número
\(e\) como aquella base que hace que el gráfico de \(y=e^x\) en el punto
\((0,1)\) tenga recta tangente con pendiente \(1\). En otras palabras,
el número \(e\) verifica que

\[
\lim_{h\to 0}\frac{e^h-1}{h}=1.
\]

Usando este resultado, podemos concluir que

\[
\frac{d}{dx}\left(e^x\right)=e^x.
\]

Ejemplo

Si \(f(x)=e^x-x\), encontrar \(f'\) y \(f''\). Comparar el gráfico de
\(f\) con el de \(f'\).

\begin{tcolorbox}[enhanced jigsaw, opacityback=0, left=2mm, opacitybacktitle=0.6, toprule=.15mm, bottomrule=.15mm, colbacktitle=quarto-callout-tip-color!10!white, titlerule=0mm, bottomtitle=1mm, title=\textcolor{quarto-callout-tip-color}{\faLightbulb}\hspace{0.5em}{Solución}, toptitle=1mm, colframe=quarto-callout-tip-color-frame, colback=white, breakable, rightrule=.15mm, arc=.35mm, leftrule=.75mm, coltitle=black]

Utilizando la derivada de \(y=e^x\) y de \(y=x\) escribimos

\[
f'(x)=\left(e^x-x\right)'=\left(e^x\right)'-x'=e^x-1
\]

y también

\[
f''(x)=\left(f'(x)\right)'=\left(e^x-1\right)'=\left(e^x\right)'-1'=e^x.
\]

\includegraphics[width=0.35\linewidth,height=\textheight,keepaspectratio]{index_files/mediabag/imagenes_cap3/3.1-exp-y-derivada.pdf}

\end{tcolorbox}

Ejemplo

¿En qué punto de la curva \(y=e^x\) la recta tangente es paralela a
\(y=2x\)?

\begin{tcolorbox}[enhanced jigsaw, opacityback=0, left=2mm, opacitybacktitle=0.6, toprule=.15mm, bottomrule=.15mm, colbacktitle=quarto-callout-tip-color!10!white, titlerule=0mm, bottomtitle=1mm, title=\textcolor{quarto-callout-tip-color}{\faLightbulb}\hspace{0.5em}{Solución}, toptitle=1mm, colframe=quarto-callout-tip-color-frame, colback=white, breakable, rightrule=.15mm, arc=.35mm, leftrule=.75mm, coltitle=black]

La recta tangente en \((x_0, y(x_0))\) será paralela a \(y=2x\) si tiene
la misma pendiente que esta recta. Como la pendiente es la derivada,
estamos buscando los puntos \(x_0\) tales que \(y'(x_0)=2\). Como
\(y'(x)=e^x\), entonces buscamos \(x_0\) tal que

\[
e^{x_0}=2,
\]

lo que nos da \(x_0=\ln 2\). El punto de la curva buscado es
\(P=(\ln 2, y(\ln 2))=(\ln 2,2)\).

\end{tcolorbox}

\hyperref[seccion_3.1]{↑ Volver al inicio de la sección}

\section{3.2. Regla del producto y del cociente}\label{seccion_3.2}

\subsection{Derivada del producto}\label{derivada-del-producto}

Las propiedades vistas hasta ahora podrían sugerir que la derivada de un
producto de funciones es el producto de las derivadas de cada una. Sin
embargo, esto es \textbf{erróneo}: si \(f(x)=x\) y \(g(x)=x^2\),
entonces \((fg)(x)=x^3\) y

\[
\frac{d}{dx}(fg)=\frac{d}{dx}(x^3)=3x^2
\]

mientras que

\[
\left(\frac{df}{dx}(x)\right)\left(\frac{dg}{dx}(x)\right)=1\cdot 2x=2x.
\]

Consideremos dos funciones \(f\) y \(g\), derivables en \(x\). Entonces

\begin{equation}\phantomsection\label{eq-derivada-producto}{
\begin{aligned}
(fg)'(x)&=\lim_{h\to 0}\frac{(fg)(x+h)-(fg)(x)}{h} \\
\\
&=\lim_{h\to 0}
\left(\frac{f(x+h)g(x+h)-\color{red}{f(x+h)g(x)}}{h}
+\frac{{\color{red}{f(x+h)g(x)}}-f(x)g(x)}{h}\right)\\
\\
&=\lim_{h\to 0}
\left(f(x+h)\,\frac{g(x+h)-g(x)}{h}
+g(x)\,\frac{f(x+h)-f(x)}{h}\right).
\end{aligned}
}\end{equation}

Como tanto \(f\) como \(g\) son derivables en \(x\), tenemos que

\[
\lim_{h\to 0}\frac{f(x+h)-f(x)}{h}=f'(x)\quad\text{ y }\quad \lim_{h\to 0}\frac{g(x+h)-g(x)}{h}=g'(x).
\]

Por otra parte, como \(f\) es derivable en \(x\) también es continua en
\(x\). Por lo tanto

\[
\lim_{h\to 0} f(x+h)=f\left(\lim_{h\to 0}[x+h]\right)=f(x).
\]

Combinando estos resultados con las leyes de los límites para el
producto y la suma, volviendo a (\ref{eq-derivada-producto}) resulta

\[
(fg)'(x)=f(x)g'(x)+g(x)f'(x).
\]

En conclusión, la \textbf{regla del producto} establece que si \(f\) y
\(g\) son derivables en \(x\), entonces

\begin{equation}\phantomsection\label{eq-regla-producto}{
\frac{d}{dx}(fg)(x)=f(x)\frac{dg}{dx}(x)+g(x)\frac{df}{dx}(x).
}\end{equation}

Ejemplo

Si \(f(x)=x\,e^x\) hallar \(f'(x)\). Encontrar también la \(n\)-ésima
derivada de \(f\), \(f^{(n)}(x)\).

\begin{tcolorbox}[enhanced jigsaw, opacityback=0, left=2mm, opacitybacktitle=0.6, toprule=.15mm, bottomrule=.15mm, colbacktitle=quarto-callout-tip-color!10!white, titlerule=0mm, bottomtitle=1mm, title=\textcolor{quarto-callout-tip-color}{\faLightbulb}\hspace{0.5em}{Solución}, toptitle=1mm, colframe=quarto-callout-tip-color-frame, colback=white, breakable, rightrule=.15mm, arc=.35mm, leftrule=.75mm, coltitle=black]

Por la \hyperref[eq-regla-producto]{regla del producto} obtenemos

\[
f'(x)=x'e^x+x\left(e^x\right)'=e^x+xe^x=(1+x)e^x.
\]

Aplicando nuevamente esta regla

\[
f''(x)=(1+x)'e^x+(1+x)\left(e^x\right)'=e^x+(1+x)e^x=(2+x)e^x.
\]

Repitiendo este procedimiento para más derivadas, en general tendremos
que

\[
f^{(n)}(x)=(n+x)e^x.
\]

\end{tcolorbox}

Ejemplo

Derivar la función \(f(t)=\sqrt{t}\,(a+tb)\).

\begin{tcolorbox}[enhanced jigsaw, opacityback=0, left=2mm, opacitybacktitle=0.6, toprule=.15mm, bottomrule=.15mm, colbacktitle=quarto-callout-tip-color!10!white, titlerule=0mm, bottomtitle=1mm, title=\textcolor{quarto-callout-tip-color}{\faLightbulb}\hspace{0.5em}{Solución}, toptitle=1mm, colframe=quarto-callout-tip-color-frame, colback=white, breakable, rightrule=.15mm, arc=.35mm, leftrule=.75mm, coltitle=black]

Utilizando la \hyperref[eq-regla-producto]{regla del producto}
escribimos

\[
\frac{df}{dt}(t)=\frac{d}{dt}\left(\sqrt{t}\right)(a+tb)+\sqrt{t}\frac{d}{dt}(a+tb)=\frac{1}{2}t^{-1/2}(a+tb)+t^{1/2}b=\frac{1}{2}t^{-1/2}a+\frac{3}{2}t^{1/2}b.
\]

\end{tcolorbox}

Ejemplo

Si \(f(x)=\sqrt{x}\,g(x)\), donde \(g(4)=2\) y \(g'(4)=3\), hallar
\(f'(4)\).

\begin{tcolorbox}[enhanced jigsaw, opacityback=0, left=2mm, opacitybacktitle=0.6, toprule=.15mm, bottomrule=.15mm, colbacktitle=quarto-callout-tip-color!10!white, titlerule=0mm, bottomtitle=1mm, title=\textcolor{quarto-callout-tip-color}{\faLightbulb}\hspace{0.5em}{Solución}, toptitle=1mm, colframe=quarto-callout-tip-color-frame, colback=white, breakable, rightrule=.15mm, arc=.35mm, leftrule=.75mm, coltitle=black]

Primero calculamos \(f'(x)\) usando la
\hyperref[eq-regla-producto]{regla del producto}. Tenemos que

\[
f'(x)=\left(\sqrt{x}\right)'g(x)+\sqrt{x}g'(x)=\frac{1}{2}x^{-1/2}\,g(x)+\sqrt{x}\,g'(x).
\]

Ahora evaluamos en \(x=4\) utilizando la información dada

\[
f'(4)=\frac{1}{2} 4^{-1/2}\,g(4)+\sqrt{4}\,g'(4)=\frac{1}{4}\cdot 2+2\cdot 3=\frac{13}{2}.
\]

\end{tcolorbox}

\subsection{Derivada del cociente}\label{derivada-del-cociente}

De manera similar a lo que ocurre con el producto, podemos verificar que
\((f/g)'\neq f'/g'\), es decir, la derivada de un cociente de funciones
\textbf{no es} el cociente de las derivadas de cada una. Para encontrar
una fórmula procederemos nuevamente por definición.

Si \(f\) y \(g\) son funciones derivables en \(x\), y además
\(g(x)\neq 0\), entonces

\begin{equation}\phantomsection\label{eq-derivada-cociente}{
\begin{aligned}
\left(\frac{f}{g}\right)'(x)&=\lim_{h\to 0}\frac{\left(\frac{f}{g}\right)(x+h)-\left(\frac{f}{g}\right)(x)}{h}\\
\\
&=\lim_{h\to 0} \frac{1}{h}\left(\frac{f(x+h)}{g(x+h)}{\color{red}{-\frac{f(x)}{g(x+h)}}+\frac{f(x)}{g(x+h)}}-\frac{f(x)}{g(x)}\right)\\
\\
&=\lim_{h\to 0} \left[\frac{1}{g(x+h)}\left(\frac{f(x+h)-f(x)}{h}\right)-\frac{f(x)}{g(x+h)g(x)}\frac{g(x+h)-g(x)}{h}\right].
\end{aligned}
}\end{equation}

Como \(f\) y \(g\) son derivables en \(x\), por definición de derivada
sabemos que

\[
\lim_{h\to 0}\frac{f(x+h)-f(x)}{h}=f'(x)\quad\text{ y }\quad \lim_{h\to 0}\frac{g(x+h)-g(x)}{h}=g'(x).
\]

Además, \(g\) es derivable en \(x\) y por lo tanto también es continua
en \(x\). De esta manera

\[
\lim_{h\to 0}\frac{1}{g(x+h)}=\frac{1}{g(x)},
\]

dado que el denominador tiene como límite \(g(x)\neq 0\) por hipótesis.
Volviendo a la expresión calculada en (\ref{eq-derivada-cociente})
resulta

\[
\left(\frac{f}{g}\right)'(x)=\frac{f'(x)}{g(x)}-\frac{f(x)g'(x)}{(g(x))^2}=\frac{f'(x)g(x)-f(x)g'(x)}{(g(x))^2}.
\]

En conclusión, la \textbf{regla del cociente} establece que si \(f\) y
\(g\) son derivables en \(x\) y \(g(x)\neq 0\), entonces

\begin{equation}\phantomsection\label{eq-regla-cociente}{
\frac{d}{dx}\left(\frac{f}{g}\right)(x)=\frac{\frac{df}{dx}(x)g(x)-f(x)\frac{dg}{dx}(x)}{g^2(x)}.
}\end{equation}

Ejemplo

Sea \(\displaystyle y=\frac{x^2+x-2}{x^3+6}\). Calcular \(y'\).

\begin{tcolorbox}[enhanced jigsaw, opacityback=0, left=2mm, opacitybacktitle=0.6, toprule=.15mm, bottomrule=.15mm, colbacktitle=quarto-callout-tip-color!10!white, titlerule=0mm, bottomtitle=1mm, title=\textcolor{quarto-callout-tip-color}{\faLightbulb}\hspace{0.5em}{Solución}, toptitle=1mm, colframe=quarto-callout-tip-color-frame, colback=white, breakable, rightrule=.15mm, arc=.35mm, leftrule=.75mm, coltitle=black]

Aplicando la \hyperref[eq-regla-cociente]{regla del cociente} podemos
escribir

\[
\begin{aligned}
y'=\frac{(x^2+x-2)'(x^3+6)-(x^2+x-2)(x^3+6)'}{(x^3+6)^2}&=\frac{(2x+1)(x^3+6)-(x^2+x-2)(3x^2)}{(x^3+6)^2}\\
\\
&=\frac{-x^4-2x^3+6x^2+12x+6}{(x^3+6)^2}.
\end{aligned}
\]

\end{tcolorbox}

Ejemplo

Encontrar una ecuación de la recta tangente a

\[
\displaystyle y=\frac{e^x}{1+x^2}
\]

en el punto \((1,e/2)\).

\begin{tcolorbox}[enhanced jigsaw, opacityback=0, left=2mm, opacitybacktitle=0.6, toprule=.15mm, bottomrule=.15mm, colbacktitle=quarto-callout-tip-color!10!white, titlerule=0mm, bottomtitle=1mm, title=\textcolor{quarto-callout-tip-color}{\faLightbulb}\hspace{0.5em}{Solución}, toptitle=1mm, colframe=quarto-callout-tip-color-frame, colback=white, breakable, rightrule=.15mm, arc=.35mm, leftrule=.75mm, coltitle=black]

Primero calculamos \(y'(x)\) usando la
\hyperref[eq-regla-cociente]{regla del cociente}

\[
y'(x)=\frac{\left(e^x\right)'(1+x^2)-e^x(1+x^2)'}{(1+x^2)^2}=\frac{e^x(1+x^2)-2x\,e^x}{(1+x^2)^2}=\frac{e^x(x-1)^2}{(1+x^2)^2}.
\]

La recta tangente buscada pasa por \((1,e/2)\) y tiene pendiente
\(y'(1)=0\). Por lo tanto, su ecuación es

\[
y-\frac{e}{2}=0(x-1) \quad \text{ o equivalentemente } \quad y=\frac{e}{2}.
\]

\end{tcolorbox}

\hyperref[seccion_3.2]{↑ Volver al inicio de la sección}

\section{3.3. Derivadas de las funciones
trigonométricas}\label{seccion_3.3}

\subsection{Derivada de la función
seno}\label{derivada-de-la-funciuxf3n-seno}

Comencemos considerando la función \(f(x)=\operatorname{sen} x\). Esta
función nos da el seno de un ángulo cuya medida en radianes es \(x\).
Queremos calcular

\[
f'(x)=\lim_{h\to 0}\frac{f(x+h)-f(x)}{h}.
\]

Considerando \(h\neq 0\) y aplicando la identidad

\[
\operatorname{sen}(\alpha+\beta)=\operatorname{sen} \alpha\cos \beta+\cos\alpha \operatorname{sen} \beta 
\]

escribimos

\begin{equation}\phantomsection\label{eq-derivada-seno}{
\begin{aligned}
\frac{f(x+h)-f(x)}{h}=\frac{\operatorname{sen}(x+h)-\operatorname{sen} x}{h}&=
\frac{\operatorname{sen} x\cos h+\cos x\operatorname{sen} h-\operatorname{sen} x}{h}\\
\\
&=\operatorname{sen} x\left(\frac{\cos h-1}{h}\right)+\cos x\left(\frac{\operatorname{sen} h}{h}\right).
\end{aligned}
}\end{equation}

Demostraremos que

\begin{equation}\phantomsection\label{eq-limite-seno}{
\lim_{h\to 0}\frac{\operatorname{sen} h}{h}=1.
}\end{equation}

y también que

\begin{equation}\phantomsection\label{eq-limite-coseno}{
\lim_{h\to 0}\frac{\cos h-1}{h}=0.
}\end{equation}

Utilizando estos límites en (\ref{eq-derivada-seno}) resulta que

\[
\begin{aligned}
\lim_{h\to 0}   \frac{f(x+h)-f(x)}{h}&=\lim_{h\to 0} \left[\operatorname{sen} x\left(\frac{\cos h-1}{h}\right)+\cos x\left(\frac{\operatorname{sen} h}{h}\right)\right]\\
\\
&=\operatorname{sen} x\left(\lim_{h\to 0}\frac{\cos h-1}{h}\right)+\cos x\left(\lim_{h\to 0}\frac{\operatorname{sen} h }{h}\right)\\
\\
&=\operatorname{sen} x \cdot 0+\cos x\cdot 1\\
\\
&=\cos x. 
\end{aligned}
\]

Concluimos que

\phantomsection\label{form-derivada-seno}
\[
\frac{d}{dx}(\operatorname{sen} x)=\cos x.
\]

Ahora procedemos a demostrar la validez del límite
(\ref{eq-limite-seno}). Consideremos \(0<\theta<\pi/2\) y un sector
circular con centro \(O\), radio 1 y ángulo central \(\theta\).

\begin{center}
\includegraphics[width=0.45\linewidth,height=\textheight,keepaspectratio]{index_files/mediabag/imagenes_cap3/3.3-limite-seno.pdf}
\end{center}

\begin{itemize}
\item
  Si prolongamos el segmento que va desde el origen al punto \(B\), éste
  corta a la tangente a la circunferencia por \(C\) en un punto \(D\).
\item
  Trazamos la tangente a la circunferencia por el punto \(B\), que corta
  al segmento \(\overline{CD}\) en un punto \(E\).
\item
  Observemos también que \(|\overline{AB}|=\operatorname{sen} \theta\) y
  \(|\overline{CD}|=\tan \theta\).
\end{itemize}

Entonces resulta que

\[
\operatorname{sen}\theta =|\overline{AB}|\leq \text{longitud arco}(BC)=\theta
\]

y adem'as

\[
\theta = \text{longitud arco}(BC)\leq |\overline{CE}|+|\overline{EB}|\leq |\overline{CE}|+|\overline{ED}|=|\overline{CD}|=\tan \theta.
\]

Es decir, hemos probado que

\[
\operatorname{sen}{\theta} \leq \theta \leq \tan\theta, 
\]

de donde podemos deducir que

\[
\cos \theta \leq \frac{\operatorname{sen} \theta}{\theta}\leq 1
\]

para \(0<\theta<\pi/2\).

Si \(-\pi/2<\theta<0\), entonces \(0<-\theta<\pi/2\). Aplicando la
desigualdad de arriba para \(-\theta\) resulta

\[
\cos (-\theta) \leq \frac{\operatorname{sen} (-\theta)}{-\theta}\leq 1.
\]

Sabiendo que el seno es una función impar y el coseno es par, lo
reescribimos como

\[
\cos \theta \leq \frac{\operatorname{sen} (\theta)}{\theta}\leq 1.
\]

Por lo tanto, la desigualdad de arriba es válida para todo
\(\theta\in (-\pi/2,\pi/2)\), \(\theta\neq 0\).

Dado que

\[
\lim_{\theta\to 0} \cos \theta=1 \quad\text{ y }\quad\lim_{\theta \to 0}1=1
\]

el \hyperref[teo-compresion]{teorema de la compresión} nos permite
deducir que

\[
\lim_{\theta\to 0} \frac{\operatorname{sen}\theta}{\theta}=1,
\]

lo que nos da la validez de la fórmula (\ref{eq-limite-seno}). Ahora, la
fórmula (\ref{eq-limite-coseno}) es consecuencia directa de ésta última.
En efecto, para \(h\neq 0\) tenemos que

\[
\frac{\cos h-1}{h}=\frac{\cos h-1}{h}\,\,\frac{\cos h+1}{\cos h+1}=-\frac{\operatorname{sen}^2 h}{h(\cos h+1)}=-\frac{\operatorname{sen} h}{h}\frac{\operatorname{sen} h}{\cos h +1}.
\]

Y usando el límite anterior junto a las leyes de los límites

\[
\lim_{h\to 0} \frac{\cos h-1}{h}=\left(\lim_{h\to 0}-\frac{\operatorname{sen} h}{h}\right)\left(\lim_{h\to 0}\frac{\operatorname{sen} h}{\cos h+1}\right)=(-1)\cdot 0=0,
\]

lo que muestra la validez de la fórmula (\ref{eq-limite-coseno}).

Ejemplo

Derivar \(y=x^2\operatorname{sen} x\).

\begin{tcolorbox}[enhanced jigsaw, opacityback=0, left=2mm, opacitybacktitle=0.6, toprule=.15mm, bottomrule=.15mm, colbacktitle=quarto-callout-tip-color!10!white, titlerule=0mm, bottomtitle=1mm, title=\textcolor{quarto-callout-tip-color}{\faLightbulb}\hspace{0.5em}{Solución}, toptitle=1mm, colframe=quarto-callout-tip-color-frame, colback=white, breakable, rightrule=.15mm, arc=.35mm, leftrule=.75mm, coltitle=black]

Calculamos \(\displaystyle\frac{dy}{dx}\) combinando la
\hyperref[eq-regla-producto]{regla del producto} con la
\hyperref[form-derivada-seno]{fórmula para la derivada del seno}

\[
\frac{dy}{dx}(x)=\frac{d}{dx}\left(x^2\right) \operatorname{sen} x + x^2 \frac{d}{dx}\left(\operatorname{sen} x\right)= 2x\operatorname{sen} x+x^2\cos x=x(2\operatorname{sen} x+x\cos x).
\]

\end{tcolorbox}

\subsection{Derivada de las funciones coseno y
tangente}\label{derivada-de-las-funciones-coseno-y-tangente}

Consideremos ahora \(f(x)=\cos x\). Si \(h\neq 0\), aplicamos la
identidad

\[
\cos(\alpha+\beta)=\cos \alpha \cos \beta -\operatorname{sen} \alpha \operatorname{sen} \beta
\]

para escribir

\[
\frac{\cos(x+h)-\cos(x)}{h}=\frac{\cos x\cos h-\operatorname{sen} x\operatorname{sen} h-\cos x}{h}
    =\cos x\left(\frac{\cos h-1}{h}\right)-\operatorname{sen} x\frac{\operatorname{sen} h}{h}.
\]

Utilizando los límites dados por las fórmulas (\ref{eq-limite-seno}) y
(\ref{eq-limite-coseno}) tenemos que \[
\begin{aligned}
\lim_{h\to 0}\frac{\cos(x+h)-\cos(x)}{h}
    &=\lim_{h\to 0}\left[\cos x\left(\frac{\cos h-1}{h}\right)-\operatorname{sen} x\frac{\operatorname{sen} h}{h}\right]\\
    \\
    &=\cos x\lim_{h\to 0}\frac{\cos h-1}{h}-\operatorname{sen} x\lim_{h\to 0}\frac{\operatorname{sen} h}{h}\\
    \\
    &=\cos x\cdot 0-\operatorname{sen} x\cdot 1\\
    \\
    &=-\operatorname{sen} x.
\end{aligned}
\]

Como conclusión, obtenemos que

\[
\frac{d}{dx}(\cos x)=-\operatorname{sen} x.
\]

Para \(f(x)=\tan x\), como \(\tan x=\operatorname{sen} x/\cos x\),
podemos utilizar la \hyperref[eq-regla-cociente]{regla del cociente}.
Así

\[
f'(x)=\frac{(\operatorname{sen} x)'\cos x-\operatorname{sen} x(\cos x)'}{\cos^2 x}=\frac{\cos^2 x+\operatorname{sen}^2 x}{\cos^2 x}=\frac{1}{\cos^2 x}=\sec^2 x.
\]

En definitiva

\[
\frac{d}{dx}(\tan x)=\sec^2 x.
\]

Como aplicación de la regla del cociente podemos deducir también las
siguientes fórmulas

\[
\frac{d}{dx}(\csc x)=-\csc x\cot x.
\]

\[
\frac{d}{dx}(\sec x)=\sec x\tan x.
\]

\[
\frac{d}{dx}(\cot x)=-\csc^2 x.
\]

Ejemplo

Derivar la función \(\displaystyle f(x)=\frac{\sec x}{1+\tan x}\). ¿Para
qué valores de \(x\) la gráfica de \(f\) tiene una tangente horizontal?

\begin{tcolorbox}[enhanced jigsaw, opacityback=0, left=2mm, opacitybacktitle=0.6, toprule=.15mm, bottomrule=.15mm, colbacktitle=quarto-callout-tip-color!10!white, titlerule=0mm, bottomtitle=1mm, title=\textcolor{quarto-callout-tip-color}{\faLightbulb}\hspace{0.5em}{Solución}, toptitle=1mm, colframe=quarto-callout-tip-color-frame, colback=white, breakable, rightrule=.15mm, arc=.35mm, leftrule=.75mm, coltitle=black]

Utilizamos las reglas de \(\sec x\) y \(\tan x\) junto a la del
cociente, es decir

\[
f'(x)=\frac{(\sec x)'(1+\tan x)-\sec x (1+\tan x)'}{(1+\tan x)^2}=\frac{\sec x\tan x (1+\tan x)-\sec x \sec^2 x}{(1+\tan x)^2}=\frac{\operatorname{sen} x-\cos x}{1+2\operatorname{sen}\cos x}.
\]

Ahora, \(f'(x)=0\) si y sólo si \(\operatorname{sen} x=\cos x\), es
decir, cuando \(\tan x=1\). Los valores de \(x\) buscados son de la
forma

\[
x=\frac{\pi}{4}+k\pi, \quad \text{ para } k \text{ entero}.
\]

\end{tcolorbox}

Ejemplo

Un objeto que se encuentra en el extremo de un resorte vertical se
desplaza hacia abajo 4 \(cm\) más allá de su posición de reposo, para
estirar el resorte, y se libera en el instante \(t=0\). Su posición
\(s\) en el instante \(t\) es \(s=f(t)=4\cos t\). Hallar la velocidad y
la aceleración en el instante \(t\) y usarlas para analizar el
movimiento del objeto.

\begin{center}
\includegraphics[width=0.5\linewidth,height=\textheight,keepaspectratio]{index_files/mediabag/imagenes_cap3/3.3-resorte.pdf}
\end{center}

\begin{tcolorbox}[enhanced jigsaw, opacityback=0, left=2mm, opacitybacktitle=0.6, toprule=.15mm, bottomrule=.15mm, colbacktitle=quarto-callout-tip-color!10!white, titlerule=0mm, bottomtitle=1mm, title=\textcolor{quarto-callout-tip-color}{\faLightbulb}\hspace{0.5em}{Solución}, toptitle=1mm, colframe=quarto-callout-tip-color-frame, colback=white, breakable, rightrule=.15mm, arc=.35mm, leftrule=.75mm, coltitle=black]

Recordemos que la velocidad y la aceleración son las derivadas primera y
segunda de la posición, respectivamente. Entonces resulta

\[
v(t)=\frac{ds}{dt}(t)=-4\operatorname{sen} t
\]

y

\[
a(t)=\frac{d^2 s}{dt^2}(t)=\frac{dv}{dt}(t)=-4\cos t.
\]

En la siguiente figura podemos observar los gráficos de la posición, la
velocidad y la aceleración. Notemos que cuando el resorte alcanza las
posiciones extremas (4 \(cm\) arriba o abajo de la posición de
equilibrio) la velocidad es cero, ya que en ese instante la masa está
cambiando el sentido de recorrido. Por otra parte, al cruzar la posición
de equilibrio, correspondiente a \(s=0\), la \emph{rapidez} \(|v(t)|\)
resulta ser máxima.

\includegraphics[width=0.45\linewidth,height=\textheight,keepaspectratio]{index_files/mediabag/imagenes_cap3/3.3-curvas-resorte.pdf}

\end{tcolorbox}

\hyperref[seccion_3.2]{↑ Volver al inicio de la sección}

\section{3.4. La regla de la cadena}\label{seccion_3.4}

En esta sección veremos cómo derivar funciones compuestas.

\phantomsection\label{teo-regla-cadena}
Regla de la cadena

Si \(g\) es derivable en \(x\) y \(f\) es derivable en \(g(x)\),
entonces la función compuesta \(F=f\circ g\) dada por \(F(x)=f(g(x))\)
es derivable en \(x\) y además \[
F'(x)=f'(g(x))g'(x).
\]

En la notación de Leibniz, si \(y=f(u)\) y \(u=g(x)\) son
diferenciables, entonces

\[
\frac{dy}{dx}=\frac{dy}{du}\frac{du}{dx}.
\]

Ejemplo

Encontrar \(F'(x)\) si \(F(x)=\sqrt{x^2+1}\).

\begin{tcolorbox}[enhanced jigsaw, opacityback=0, left=2mm, opacitybacktitle=0.6, toprule=.15mm, bottomrule=.15mm, colbacktitle=quarto-callout-tip-color!10!white, titlerule=0mm, bottomtitle=1mm, title=\textcolor{quarto-callout-tip-color}{\faLightbulb}\hspace{0.5em}{Solución}, toptitle=1mm, colframe=quarto-callout-tip-color-frame, colback=white, breakable, rightrule=.15mm, arc=.35mm, leftrule=.75mm, coltitle=black]

Observemos que podemos escribir \(F(x)=f(g(x))\), donde
\(f(x)=\sqrt{x}\) y \(g(x)=x^2+1\). En este caso
\(f'(x)=\frac{1}{2\sqrt{x}}\) y \(g'(x)=2x\). Luego, utilizando la
\hyperref[teo-regla-cadena]{regla de la cadena} obtenemos

\[
F'(x)=f'(g(x))g'(x)=\frac{1}{2\sqrt{x^2+1}} 2x =\frac{x}{\sqrt{x^2+1}}.
\]

\end{tcolorbox}

La regla de la cadena establece que la derivada de \(f(g(x))\) puede
calcularse derivando la función externa \(f\) (evaluada en \(g(x)\)) y
multiplicándola por la derivada de \(g\), evaluada en \(x\).

Ejemplo

Derivar \(y=\operatorname{sen}(x^2)\) e \(y=\operatorname{sen}^2x\).

\begin{tcolorbox}[enhanced jigsaw, opacityback=0, left=2mm, opacitybacktitle=0.6, toprule=.15mm, bottomrule=.15mm, colbacktitle=quarto-callout-tip-color!10!white, titlerule=0mm, bottomtitle=1mm, title=\textcolor{quarto-callout-tip-color}{\faLightbulb}\hspace{0.5em}{Solución}, toptitle=1mm, colframe=quarto-callout-tip-color-frame, colback=white, breakable, rightrule=.15mm, arc=.35mm, leftrule=.75mm, coltitle=black]

Para la primer función tenemos que \(y=f(g(x))\), donde
\(f(x)=\operatorname{sen} x\) y \(g(x)=x^2\). Entonces

\[
y'=f'(g(x))g'(x)=\cos(x^2)2x=2x\cos(x^2).
\]

La segunda función puede expresarse como \(y=g(f(x))\). En este caso
resulta

\[
y'=g'(f(x))f'(x)=2\operatorname{sen} x \cos x.
\]

\end{tcolorbox}

Veamos un caso especial de aplicación de la regla de la cadena.

\phantomsection\label{form-potencia-general}
Regla de la potencia combinada con regla de la cadena

Si \(n\in\mathbb{R}\) y \(u=g(x)\) es derivable, entonces \[
\frac{d}{dx}\left(u^n\right)=nu^{n-1}\frac{du}{dx}\quad \text{ o bien} \quad \frac{d}{dx}\left(g(x)\right)^n=ng(x)^{n-1}g'(x).
\]

Ejemplo

Derivar \(y=(x^3-1)^{100}\).

\begin{tcolorbox}[enhanced jigsaw, opacityback=0, left=2mm, opacitybacktitle=0.6, toprule=.15mm, bottomrule=.15mm, colbacktitle=quarto-callout-tip-color!10!white, titlerule=0mm, bottomtitle=1mm, title=\textcolor{quarto-callout-tip-color}{\faLightbulb}\hspace{0.5em}{Solución}, toptitle=1mm, colframe=quarto-callout-tip-color-frame, colback=white, breakable, rightrule=.15mm, arc=.35mm, leftrule=.75mm, coltitle=black]

Aplicando la \hyperref[form-potencia-general]{formula anterior} resulta

\[
y'=100(x^3-1)^{99}(x^3-1)'=300(x^3-1)^{99}x^2.
\]

\end{tcolorbox}

Ejemplo

Encontrar \(f'(x)\) si
\(\displaystyle f(x)=\frac{1}{\sqrt[3]{x^2+x+1}}\).

\begin{tcolorbox}[enhanced jigsaw, opacityback=0, left=2mm, opacitybacktitle=0.6, toprule=.15mm, bottomrule=.15mm, colbacktitle=quarto-callout-tip-color!10!white, titlerule=0mm, bottomtitle=1mm, title=\textcolor{quarto-callout-tip-color}{\faLightbulb}\hspace{0.5em}{Solución}, toptitle=1mm, colframe=quarto-callout-tip-color-frame, colback=white, breakable, rightrule=.15mm, arc=.35mm, leftrule=.75mm, coltitle=black]

Escribimos \(f(x)=(x^2+x+1)^{-1/3}\), con lo cual

\[
\frac{df}{dx}(x)=-\frac{1}{3}(x^2+x+1)^{-4/3}\frac{d}{dx}(x^2+x+1)=-\frac{1}{3}(2x+1)(x^2+x+1)^{-4/3}.
\]

\end{tcolorbox}

Ejemplo

Calcular la derivada de
\(\displaystyle g(t)=\left(\frac{t-2}{2t+1}\right)^9\).

\begin{tcolorbox}[enhanced jigsaw, opacityback=0, left=2mm, opacitybacktitle=0.6, toprule=.15mm, bottomrule=.15mm, colbacktitle=quarto-callout-tip-color!10!white, titlerule=0mm, bottomtitle=1mm, title=\textcolor{quarto-callout-tip-color}{\faLightbulb}\hspace{0.5em}{Solución}, toptitle=1mm, colframe=quarto-callout-tip-color-frame, colback=white, breakable, rightrule=.15mm, arc=.35mm, leftrule=.75mm, coltitle=black]

Aplicando nuevamente la \hyperref[form-potencia-general]{formula de la
potencia} resulta

\[
g'(t)=9\left(\frac{t-2}{2t+1}\right)^8\left(\frac{t-2}{2t+1}\right)'=9\left(\frac{t-2}{2t+1}\right)^8\left(\frac{2t+1-(t-2)2}{(2t+1)^2}\right)=45\frac{(t-2)^8}{(2t+1)^{10}}.
\]

\end{tcolorbox}

Ejemplo

Derivar \(y=(2x+1)^5(x^3-x+1)^4\).

\begin{tcolorbox}[enhanced jigsaw, opacityback=0, left=2mm, opacitybacktitle=0.6, toprule=.15mm, bottomrule=.15mm, colbacktitle=quarto-callout-tip-color!10!white, titlerule=0mm, bottomtitle=1mm, title=\textcolor{quarto-callout-tip-color}{\faLightbulb}\hspace{0.5em}{Solución}, toptitle=1mm, colframe=quarto-callout-tip-color-frame, colback=white, breakable, rightrule=.15mm, arc=.35mm, leftrule=.75mm, coltitle=black]

Combinando la \hyperref[regla-producto]{regla del producto} con la
\hyperref[form-potencia-general]{formula de la potencia} resulta

\[
\begin{aligned}
y'&=5(2x+1)^4\, 2(x^3-x+1)^4+(2x+1)^5\, 4(x^3-x+1)^3(3x^2-1)\\
\\
&=2(2x+1)^4(x^3-x+1)^3(5(x^3-x+1)+2(2x+1)(3x^2-1))\\
\\
&=2(2x+1)^4(x^3-x+1)^3(17x^3+6x^2-9x+3).
\end{aligned}
\]

\end{tcolorbox}

Ejemplo

Derivar \(y=e^{\operatorname{sen} x}\).

\begin{tcolorbox}[enhanced jigsaw, opacityback=0, left=2mm, opacitybacktitle=0.6, toprule=.15mm, bottomrule=.15mm, colbacktitle=quarto-callout-tip-color!10!white, titlerule=0mm, bottomtitle=1mm, title=\textcolor{quarto-callout-tip-color}{\faLightbulb}\hspace{0.5em}{Solución}, toptitle=1mm, colframe=quarto-callout-tip-color-frame, colback=white, breakable, rightrule=.15mm, arc=.35mm, leftrule=.75mm, coltitle=black]

En este caso tenemos que

\[
\frac{dy}{dx}(x)=e^{\operatorname{sen}x}\cos x.
\]

\end{tcolorbox}

\subsection{Derivadas de funciones exponenciales
generales}\label{derivadas-de-funciones-exponenciales-generales}

La regla de la cadena nos permite ahora encontrar la derivada de
\(y=a^x\), para cualquier \(a>0\), \(a\neq 1\). Recordando que
\(a=e^{\ln a}\), tenemos que

\[
a^x=\left(e^{\ln a}\right)^x=e^{(\ln a) x }.
\]

Aplicando la regla de la cadena y la derivada de la función \(f(x)=e^x\)
obtenemos

\[
\frac{d}{dx}(a^x)=\frac{d}{dx}\left(e^{(\ln a) x }\right)=e^{(\ln a) x }\frac{d}{dx}\left((\ln a)x\right)=a^x\ln a.
\]

En conclusión

\[
\frac{d}{dx}(a^x)=a^x\ln a.
\]

La regla de la cadena puede extenderse a composiciones con más
funciones: si \(y=f(u)\), \(u=g(x)\) y \(x=h(t)\) son funciones tales
que \(f\) es derivable en \(g(h(t))\), \(g\) es derivable en \(h(t)\) y
\(h\) es derivable en \(t\), entonces

\[
\frac{dy}{dt}=\frac{dy}{du}\frac{du}{dx}\frac{dx}{dt}=f'(g(h(t)))\,g'(h(t))\,h'(t).
\]

Ejemplo

Si \(f(x)=\operatorname{sen}(\cos(\tan x))\), calcular \(f'(x)\).

\begin{tcolorbox}[enhanced jigsaw, opacityback=0, left=2mm, opacitybacktitle=0.6, toprule=.15mm, bottomrule=.15mm, colbacktitle=quarto-callout-tip-color!10!white, titlerule=0mm, bottomtitle=1mm, title=\textcolor{quarto-callout-tip-color}{\faLightbulb}\hspace{0.5em}{Solución}, toptitle=1mm, colframe=quarto-callout-tip-color-frame, colback=white, breakable, rightrule=.15mm, arc=.35mm, leftrule=.75mm, coltitle=black]

Aplicamos la regla de la cadena dos veces

\[
\begin{aligned}
f'(x)=\cos(\cos(\tan x))(\cos(\tan x))'&=\cos(\cos(\tan x))(-\operatorname{sen}(\tan x))(\tan x)'\\
\\
&=-\cos(\cos(\tan x))\operatorname{sen}(\tan x)\sec^2 x.
\end{aligned}
\]

\end{tcolorbox}

Ejemplo

Derivar \(y=e^{\sec(3\theta)}\).

\begin{tcolorbox}[enhanced jigsaw, opacityback=0, left=2mm, opacitybacktitle=0.6, toprule=.15mm, bottomrule=.15mm, colbacktitle=quarto-callout-tip-color!10!white, titlerule=0mm, bottomtitle=1mm, title=\textcolor{quarto-callout-tip-color}{\faLightbulb}\hspace{0.5em}{Solución}, toptitle=1mm, colframe=quarto-callout-tip-color-frame, colback=white, breakable, rightrule=.15mm, arc=.35mm, leftrule=.75mm, coltitle=black]

Aplicando dos veces la regla de la cadena resulta

\[
\frac{dy}{d\theta}(\theta)=e^{\sec(3\theta)}\frac{d}{d\theta}(\sec(3\theta))=e^{\sec(3\theta)}\sec(3\theta)\tan(3\theta)(3\theta)'=3\sec(3\theta)\tan(3\theta)e^{\sec(3\theta)}.
\]

\end{tcolorbox}

\hyperref[seccion_3.4]{↑ Volver al inicio de la sección}

\section{3.5. Derivación implícita}\label{seccion_3.5}

Hasta ahora hemos visto funciones de la forma \(y=f(x)\), donde \(y\) se
expresa explícitamente en términos de \(x\). Sin embargo, muchas veces
\(x\) e \(y\) se relacionan de forma implícita, tal como en las
expresiones

\[
x^2+y^2=25\quad \text{o también }\quad x^3+y^3=6xy.
\]

En algunos casos es posible despejar \(y\) en términos de \(x\),
obteniendo una (o varias) funciones de \(x\). Por ejemplo, la ecuación

\[
x^2+y^2=25
\]

describe una circunferencia de radio 5 y centrada en el origen de
coordenadas.

Despejando \(y\) obtenemos \(y^2=25-x^2\), con lo cual

\[
|y|=\sqrt{25-x^2} \quad \text{ o bien }\quad y=\pm\sqrt{25-x^2},
\]

que se corresponden con la mitad superior e inferior de la
circunferencia.

\begin{center}
\includegraphics[width=0.6\linewidth,height=\textheight,keepaspectratio]{index_files/mediabag/imagenes_cap3/3.5-circunferencia.pdf}
\end{center}

En general, no es necesario despejar \(y\) en términos de \(x\) para
hallar la derivada de \(y\). En su lugar, podemos aplicar el
\textbf{método de derivación implícita}, que consiste en derivar ambos
miembros de una ecuación con respecto a \(x\) y luego resolver para
\(y'\).

Ejemplo

Considerar la ecuación \(x^2+y^2=25\).

\begin{itemize}
\tightlist
\item
  Calcular \(\displaystyle\frac{dy}{dx}\).
\item
  Encontrar la ecuación de la recta tangente a la circunferencia en el
  punto \((3,4)\).
\end{itemize}

\begin{tcolorbox}[enhanced jigsaw, opacityback=0, left=2mm, opacitybacktitle=0.6, toprule=.15mm, bottomrule=.15mm, colbacktitle=quarto-callout-tip-color!10!white, titlerule=0mm, bottomtitle=1mm, title=\textcolor{quarto-callout-tip-color}{\faLightbulb}\hspace{0.5em}{Solución}, toptitle=1mm, colframe=quarto-callout-tip-color-frame, colback=white, breakable, rightrule=.15mm, arc=.35mm, leftrule=.75mm, coltitle=black]

Si derivamos miembro a miembro la ecuación \(x^2+y^2=25\), recordando
que \(y\) es implícitamente una función de \(x\), por la regla de la
cadena obtenemos

\[
2x+2y\frac{dy}{dx}=0,
\]

y despejando la derivada de \(y\)

\[
\frac{dy}{dx}=-\frac{x}{y},
\]

para todo \((x,y)\) tal que \(y\neq 0\). En particular, en el punto
\((3,4)\) obtenemos

\[
\frac{dy}{dx}=-\frac{3}{4}.
\]

La recta tangente tiene pendiente \(m=-3/4\) y pasa por \((3,4)\), con
lo que su ecuación es

\[
y-4=-\frac{3}{4}(x-3) \quad \text{ o, equivalentemente, } \quad y=-\frac{3}{4}x+\frac{25}{4}.
\]

\end{tcolorbox}

Ejemplo

Considerar la curva \(x^3+y^3=6xy\), conocida como \textbf{folio de
Descartes}.

\begin{itemize}
\tightlist
\item
  Encontrar \(y'\).
\item
  Hallar la recta tangente al folio de Descartes en \((3,3)\).
\end{itemize}

\begin{center}
\includegraphics[width=0.35\linewidth,height=\textheight,keepaspectratio]{index_files/mediabag/imagenes_cap3/3.5-folio-Descartes.pdf}
\end{center}

\begin{tcolorbox}[enhanced jigsaw, opacityback=0, left=2mm, opacitybacktitle=0.6, toprule=.15mm, bottomrule=.15mm, colbacktitle=quarto-callout-tip-color!10!white, titlerule=0mm, bottomtitle=1mm, title=\textcolor{quarto-callout-tip-color}{\faLightbulb}\hspace{0.5em}{Solución}, toptitle=1mm, colframe=quarto-callout-tip-color-frame, colback=white, breakable, rightrule=.15mm, arc=.35mm, leftrule=.75mm, coltitle=black]

Primero calculamos \(y'\) derivando implícitamente

\[
3x^2+3y^2y'=6y+6xy'
\]

con lo cual

\[
y'=\frac{2y-x^2}{y^2-2x}.
\]

Al evaluar en el punto \((3,3)\), la pendiente de la recta tangente es

\[
m=\frac{2\cdot 3-3^2}{3^2-2\cdot 3}=-1
\]

y la ecuación de la recta es

\[
y-3=-1(x-3) \quad \text{ o también }\quad y=-x+6.
\]

\end{tcolorbox}

Ejemplo

Encontrar \(y'\) si \(\operatorname{sen}(x+y)=y^2\cos x\).

\begin{tcolorbox}[enhanced jigsaw, opacityback=0, left=2mm, opacitybacktitle=0.6, toprule=.15mm, bottomrule=.15mm, colbacktitle=quarto-callout-tip-color!10!white, titlerule=0mm, bottomtitle=1mm, title=\textcolor{quarto-callout-tip-color}{\faLightbulb}\hspace{0.5em}{Solución}, toptitle=1mm, colframe=quarto-callout-tip-color-frame, colback=white, breakable, rightrule=.15mm, arc=.35mm, leftrule=.75mm, coltitle=black]

Derivando implícitamente resulta

\[
\cos(x+y)(1+y')=2yy'\cos x-y^2\operatorname{sen}x.
\]

Reacomodando términos y despejando \(y'\) obtenemos

\[
y'=\frac{y^2\operatorname{sen} x+\cos(x+y)}{2y\cos x-\cos(x+y)}.
\]

\end{tcolorbox}

Ejemplo

Hallar \(y''\) si \(x^4+y^4=16\).

\begin{tcolorbox}[enhanced jigsaw, opacityback=0, left=2mm, opacitybacktitle=0.6, toprule=.15mm, bottomrule=.15mm, colbacktitle=quarto-callout-tip-color!10!white, titlerule=0mm, bottomtitle=1mm, title=\textcolor{quarto-callout-tip-color}{\faLightbulb}\hspace{0.5em}{Solución}, toptitle=1mm, colframe=quarto-callout-tip-color-frame, colback=white, breakable, rightrule=.15mm, arc=.35mm, leftrule=.75mm, coltitle=black]

Primero calculamos \(y'\)

\[
4x^3+4y^3y'=0 \quad \text{ o también } \quad y'=-\left(\frac{x}{y}\right)^3.
\]

Luego

\[
y''=-3\left(\frac{x}{y}\right)^{2}\left(\frac{y-xy'}{y^2}\right)=-3\frac{x^2}{y^4}\left(y+\frac{x^4}{y^3}\right)=-3\frac{x^2}{y^4}\frac{16}{y^3}=-48 \frac{x^2}{y^7}.
\]

\end{tcolorbox}

\subsection{Derivadas de funciones trigonométricas
inversas}\label{derivadas-de-funciones-trigonomuxe9tricas-inversas}

Investigaremos las derivadas de las funciones trigonométricas inversas.
Es razonable suponer que éstas son derivables, pues sus funciones
inversas son derivables (de hecho, si \(f\) es una función derivable y
uno a uno, puede mostrarse que su inversa \(f^{-1}\) es derivable en
todo punto excepto donde sus tangentes son verticales).

Consideremos primero la función arcoseno. Tenemos que

\[
y=\text{arcsen } x \quad\Longleftrightarrow\quad \operatorname{sen} y=x,\quad -\frac{\pi}{2}\leq y\leq \frac{\pi}{2}.
\]

Derivando implícitamente, obtenemos

\[
(\cos y)\, y'=1\quad \text{ o bien }\quad y'=\frac{1}{\cos y}.
\]

Como \(\cos^2y+\operatorname{sen}^2y=1\), obtenemos
\(|\cos y|=\sqrt{1-\operatorname{sen}^2 y}\). Luego, como
\(y\in[-\pi/2,\pi/2]\), tenemos que \(\cos y\geq 0\) y entonces

\[
y'=\frac{1}{\cos y}=\frac{1}{\sqrt{1-\operatorname{sen}^2 y}}=\frac{1}{\sqrt{1-x^2 }}.
\]

Como conclusión obtenemos la fórmula para la derivada del arcoseno

\phantomsection\label{form-derivada-arcoseno}
\[
\frac{d}{dx}(\text{arcsen }x)=\frac{1}{\sqrt{1-x^2}}.
\]

Para las demás funciones trigonométricas inversas podemos proceder de
manera similar. En el caso de la función arcocoseno tenemos

\[
y=\arccos x \quad\Longleftrightarrow\quad \cos y=x,\quad 0\leq y\leq \pi.
\]

Derivando implícitamente obtenemos

\[
(-\operatorname{sen} y)y'=1.
\]

En consecuencia, como \(\operatorname{sen} y\geq 0\) en el intervalo
\([0,\pi]\)

\[
y'=-\frac{1}{\operatorname{sen} y}=-\frac{1}{\sqrt{1-\cos^2 y}}=-\frac{1}{\sqrt{1-x^2}}.
\]

Obtenemos así la fórmula para la derivada del arcocoseno

\phantomsection\label{form-derivada-arcocoseno}
\[
\frac{d}{dx}(\arccos x)=-\frac{1}{\sqrt{1-x^2}}.
\]

Para el caso de la función arcotangente, recordemos que

\[
y=\arctan x \quad \Longleftrightarrow\quad \tan y=x,\quad -\frac{\pi}{2}<y<\frac{\pi}{2}.
\]

Derivando implícitamente resulta \((\sec^2 y)y'=1\), o bien

\[
y'=-\frac{1}{\sec^2 y}=\frac{1}{1+\tan^2y}=\frac{1}{1+x^2}.
\]

En conclusión

\begin{quote}
{[}Gráfico tikz omitido{]}
\end{quote}

\phantomsection\label{form-derivada-arcotangente}
\[
\frac{d}{dx}(\arctan x)=\frac{1}{1+x^2}.
\]

Ejemplo

Derivar las funciones:

\begin{itemize}
\tightlist
\item
  \(y=\displaystyle (\operatorname{arcsen} x)^{-1}\)
\item
  \(f(x)=x\,\arctan(\sqrt{x})\)
\end{itemize}

\begin{tcolorbox}[enhanced jigsaw, opacityback=0, left=2mm, opacitybacktitle=0.6, toprule=.15mm, bottomrule=.15mm, colbacktitle=quarto-callout-tip-color!10!white, titlerule=0mm, bottomtitle=1mm, title=\textcolor{quarto-callout-tip-color}{\faLightbulb}\hspace{0.5em}{Solución}, toptitle=1mm, colframe=quarto-callout-tip-color-frame, colback=white, breakable, rightrule=.15mm, arc=.35mm, leftrule=.75mm, coltitle=black]

Para la primera función utilizamos la
\hyperref[form-derivada-arcoseno]{derivada del arcoseno}

\[
y'=(-1)(\operatorname{arcsen} x)^{-2} (\operatorname{arcsen} x)'=-\frac{1}{(\operatorname{arcsen} x)^{2}}\frac{1}{\sqrt{1-x^2}}.
\]

En el segundo caso aplicamos la
\hyperref[form-derivada-arcotangente]{derivada del arcotangente}

\[
f'(x)=\arctan(\sqrt{x})+x\frac{1}{1+(\sqrt{x})^2}\frac{1}{2\sqrt{x}}=\arctan(\sqrt{x})+\frac{\sqrt{x}}{2(1+x)}.
\]

\end{tcolorbox}

Las derivadas de las funciones trigonométricas inversas restantes pueden
deducirse de forma similar.

\[
\frac{d}{dx}(\text{arccsc } x)=-\frac{1}{x\sqrt{x^2-1}}.            
\]

\[
\frac{d}{dx}(\text{arcsec } x)=\frac{1}{x\sqrt{x^2-1}}.
\]

\[
\frac{d}{dx}(\text{arccot } x)=-\frac{1}{1+x^2}.
\]

\hyperref[seccion_3.5]{↑ Volver al inicio de la sección}

\section{3.6. Derivadas de funciones logarítmicas}\label{seccion_3.6}

En esta sección utilizaremos el método de derivación implícita para
deducir la derivada de funciones logarítmicas. El hecho de que estas
funciones resultan derivables es razonable por la forma que presentan
sus gráficas.

Dado \(a>0\), \(a\neq 1\) consideramos \(y=\log_a x\). Por definición de
función inversa tenemos que

\[
y=\log_a x\quad\Longleftrightarrow\quad a^y=x.
\]

Derivando implícitamente, obtenemos \(a^y\,(\ln a)\, y'=1\), de donde

\[
y'=\frac{1}{\ln a}\frac{1}{a^y}=\frac{1}{(\ln a)x}.
\]

En conclusión

\phantomsection\label{form-derivada-log}
\[
\frac{d}{dx}(\log_a x)=\frac{1}{x\,\ln a}.
\]

En particular, si \(a=e\) tenemos que \(\ln e=1\) y \(\log_a x=\ln x\),
con lo cual

\phantomsection\label{form-derivada-ln}
\[
\frac{d}{dx}(\ln x)=\frac{1}{x}.
\]

Ejemplo

Derivar la función \(y=\ln(x^3+1)\).

\begin{tcolorbox}[enhanced jigsaw, opacityback=0, left=2mm, opacitybacktitle=0.6, toprule=.15mm, bottomrule=.15mm, colbacktitle=quarto-callout-tip-color!10!white, titlerule=0mm, bottomtitle=1mm, title=\textcolor{quarto-callout-tip-color}{\faLightbulb}\hspace{0.5em}{Solución}, toptitle=1mm, colframe=quarto-callout-tip-color-frame, colback=white, breakable, rightrule=.15mm, arc=.35mm, leftrule=.75mm, coltitle=black]

Aplicamos la \hyperref[form-derivada-ln]{fórmula de la derivada del
logaritmo} junto a la regla de la cadena para obtener

\[
y'=\frac{1}{x^3+1}(x^3+1)'=\frac{3x^2}{x^3+1}.
\]

\end{tcolorbox}

Ejemplo

Encontrar
\(\displaystyle\frac{d}{dx}\big(\ln(\operatorname{sen} x)\big)\).

\begin{tcolorbox}[enhanced jigsaw, opacityback=0, left=2mm, opacitybacktitle=0.6, toprule=.15mm, bottomrule=.15mm, colbacktitle=quarto-callout-tip-color!10!white, titlerule=0mm, bottomtitle=1mm, title=\textcolor{quarto-callout-tip-color}{\faLightbulb}\hspace{0.5em}{Solución}, toptitle=1mm, colframe=quarto-callout-tip-color-frame, colback=white, breakable, rightrule=.15mm, arc=.35mm, leftrule=.75mm, coltitle=black]

Aplicamos nuevamente la \hyperref[form-derivada-ln]{fórmula de la
derivada del logaritmo} combinada con la regla de la cadena

\[
\frac{d}{dx}\big(\ln(\operatorname{sen} x)\big)=\frac{1}{\operatorname{sen} x}\frac{d}{dx}(\operatorname{sen} x)=\frac{\cos x}{\operatorname{sen} x}=\cot x.
\]

\end{tcolorbox}

Ejemplo

Derivar \(f(x)=\sqrt{\ln x}\).

\begin{tcolorbox}[enhanced jigsaw, opacityback=0, left=2mm, opacitybacktitle=0.6, toprule=.15mm, bottomrule=.15mm, colbacktitle=quarto-callout-tip-color!10!white, titlerule=0mm, bottomtitle=1mm, title=\textcolor{quarto-callout-tip-color}{\faLightbulb}\hspace{0.5em}{Solución}, toptitle=1mm, colframe=quarto-callout-tip-color-frame, colback=white, breakable, rightrule=.15mm, arc=.35mm, leftrule=.75mm, coltitle=black]

\[
f'(x)=\frac{1}{2\sqrt{\ln x}}\frac{1}{x}=\frac{1}{2x\sqrt{\ln x}}.
\]

\end{tcolorbox}

Ejemplo

Calcular la derivada de \(f(x)=\log_{10} (2+\operatorname{sen} x)\).

\begin{tcolorbox}[enhanced jigsaw, opacityback=0, left=2mm, opacitybacktitle=0.6, toprule=.15mm, bottomrule=.15mm, colbacktitle=quarto-callout-tip-color!10!white, titlerule=0mm, bottomtitle=1mm, title=\textcolor{quarto-callout-tip-color}{\faLightbulb}\hspace{0.5em}{Solución}, toptitle=1mm, colframe=quarto-callout-tip-color-frame, colback=white, breakable, rightrule=.15mm, arc=.35mm, leftrule=.75mm, coltitle=black]

Utilizando la \hyperref[form-derivada-log]{fórmula de la derivada del
logaritmo} junto a la regla de la cadena resulta

\[
\frac{df}{dx}(x)=\frac{1}{(2+\operatorname{sen} x)\ln 10}\frac{d}{dx}(2+\operatorname{sen}x)=\frac{\cos x}{(2+\operatorname{sen} x)\ln 10}.
\]

\end{tcolorbox}

Ejemplo

Hallar
\(\displaystyle \frac{d}{dx}\ln\left(\frac{x+1}{\sqrt{x-2}}\right)\).

\begin{tcolorbox}[enhanced jigsaw, opacityback=0, left=2mm, opacitybacktitle=0.6, toprule=.15mm, bottomrule=.15mm, colbacktitle=quarto-callout-tip-color!10!white, titlerule=0mm, bottomtitle=1mm, title=\textcolor{quarto-callout-tip-color}{\faLightbulb}\hspace{0.5em}{Solución}, toptitle=1mm, colframe=quarto-callout-tip-color-frame, colback=white, breakable, rightrule=.15mm, arc=.35mm, leftrule=.75mm, coltitle=black]

En este caso resulta de utilidad aplicar primero las propiedades del
logaritmo para reescribir a la función como sigue

\[
f(x)=\ln\left(\frac{x+1}{\sqrt{x-2}}\right)=\ln(x+1)-\ln(\sqrt{x-2})=\ln(x+1)-\frac{1}{2}\ln(x-2).
\]

Ahora derivamos utilizando la \hyperref[form-derivada-ln]{fórmula de la
derivada del logaritmo}

\[
f'(x)=\frac{1}{x+1}(x+1)'-\frac{1}{2}\frac{1}{(x-2)}(x-2)'=\frac{1}{x+1}-\frac{1}{2(x-2)}=\frac{x-5}{2(x+1)(x-2)}.
\]

\end{tcolorbox}

Ejemplo

Encontrar \(f'(x)\) si \(f(x)=\ln|x|\).

\begin{tcolorbox}[enhanced jigsaw, opacityback=0, left=2mm, opacitybacktitle=0.6, toprule=.15mm, bottomrule=.15mm, colbacktitle=quarto-callout-tip-color!10!white, titlerule=0mm, bottomtitle=1mm, title=\textcolor{quarto-callout-tip-color}{\faLightbulb}\hspace{0.5em}{Solución}, toptitle=1mm, colframe=quarto-callout-tip-color-frame, colback=white, breakable, rightrule=.15mm, arc=.35mm, leftrule=.75mm, coltitle=black]

Primero observemos que \(f\) no está definida en \(x=0\), con lo cual
\(f'(0)\) no existe. Si \(x>0\), entonces \(f(x)=\ln x\) y en
consecuencia \(\displaystyle f'(x)=\frac{1}{x}\).

Si \(x<0\), \(f(x)=\ln(-x)\) y por la regla de la cadena obtenemos
\(\displaystyle f'(x)=\frac{1}{-x}(-x)'=\frac{1}{x}\).

Por lo tanto

\[
\frac{d}{dx}\left(\ln |x|\right)=\frac{1}{x}, \quad \text{ para todo }\quad x\neq 0. 
\]

\end{tcolorbox}

\hyperref[seccion_3.6]{↑ Volver al inicio de la sección}

\section{3.7. Razones de cambio en las ciencias}\label{seccion_3.7}

Si \(y\) depende de \(x\) mediante la relación \(y=f(x)\), entonces la
derivada \(\displaystyle \frac{dy}{dx}\) representa la razón de cambio
instantánea de \(y\) con respecto a \(x\).

Recordando lo estudiado en la \hyperref[seccion_2.7]{Sección 2.7}, si
\(x\) cambia de \(x_1\) a \(x_2\), el incremento en \(x\) es

\[
\Delta x=x_2-x_1,
\]

mientras que el correspondiente incremento en \(y\) es

\[
\Delta y=f(x_2)-f(x_1).
\]

El cociente de diferencias

\[
\frac{\Delta y}{\Delta x}=\frac{f(x_2)-f(x_1)}{x_2-x_1}
\]

es la razón de cambio promedio de \(y\) con respecto a \(x\) en el
intervalo \([x_1,x_2]\). Esta razón de cambio promedio puede
interpretarse como la pendiente de la recta secante \(PQ\) del siguiente
gráfico

\begin{center}
\includegraphics[width=0.5\linewidth,height=\textheight,keepaspectratio]{index_files/mediabag/imagenes_cap3/3.7-tangente-y-secante.pdf}
\end{center}

El límite de estas razones de cambio promedio (si existe) es la razón de
cambio instantánea de \(y\) con respecto a \(x\) en \(x=x_1\), y puede
interpretarse como la pendiente de la recta tangente a la curva en el
punto \(P\).

Utilizando la notación de Leibniz, podemos escribir

\[
\frac{dy}{dx}= \lim_{\Delta x\to 0} \frac{\Delta y}{\Delta x}.
\]

Si \(y=f(x)\) tiene alguna interpretación específica modelando una
situación particular de alguna ciencia, su derivada tendrá la
correspondiente interpretación como razón de cambio. Veamos algunos
ejemplos concretos.

\subsection{Razones de cambio en
física}\label{razones-de-cambio-en-fuxedsica}

Ejemplo

Una partícula se desplaza en línea recta. Su posición en función del
tiempo transcurrido está dada por la expresión

\[
s=f(t)=t^3-6t^2+9t,
\]

donde \(t\) se mide en segundos y \(s\) en metros.

\begin{enumerate}
\def\labelenumi{\arabic{enumi}.}
\item
  Encontrar la velocidad de la partícula en el instante \(t\).
\item
  ¿Cuál es la velocidad luego de 2 y de 4 segundos?
\item
  ¿Cuándo está en reposo la partícula?
\item
  ¿Cuándo se mueve hacia adelante, es decir, en la dirección positiva?
\item
  Dibujar un diagrama que represente el movimiento de la partícula.
\item
  Encontrar la distancia total recorrida por la partícula durante los
  primero cinco segundos.
\item
  Hallar la aceleración a tiempo \(t\) y luego de cuatro segundos.
\item
  Graficar las funciones posición, velocidad y aceleración para
  \(0\leq t\leq 5\).
\item
  ¿Cuándo incrementa su rapidez la partícula? ¿Cuándo la disminuye?
\end{enumerate}

\begin{tcolorbox}[enhanced jigsaw, opacityback=0, left=2mm, opacitybacktitle=0.6, toprule=.15mm, bottomrule=.15mm, colbacktitle=quarto-callout-tip-color!10!white, titlerule=0mm, bottomtitle=1mm, title=\textcolor{quarto-callout-tip-color}{\faLightbulb}\hspace{0.5em}{Solución}, toptitle=1mm, colframe=quarto-callout-tip-color-frame, colback=white, breakable, rightrule=.15mm, arc=.35mm, leftrule=.75mm, coltitle=black]

\begin{enumerate}
\def\labelenumi{\arabic{enumi}.}
\tightlist
\item
  La velocidad de la partícula es la derivada de la posición. Entonces
\end{enumerate}

\[
v(t)=\frac{ds}{dt}(t)=3t^2-12t+9,
\]

para cada instante de tiempo \(t\).

\begin{enumerate}
\def\labelenumi{\arabic{enumi}.}
\setcounter{enumi}{1}
\tightlist
\item
  Simplemente evaluamos \(v\) para \(t=2\) y \(t=4\). Es decir
\end{enumerate}

\[
v(2)=3\cdot 2^2-12\cdot 2+9=-3 \quad \text{ y }\quad v(4)=3\cdot 4^2-12\cdot 4+9=9.
\]

La velocidad a los dos segundos es de \(-3\) m/seg y a los cuatro
segundos es de \(9\) m/seg.

\begin{enumerate}
\def\labelenumi{\arabic{enumi}.}
\setcounter{enumi}{2}
\tightlist
\item
  La partícula estará en reposo en los instantes \(t\) tales que
  \(v(t)=0\). Dado que
\end{enumerate}

\[
3t^2-12t+9=0 \quad \text{ si y sólo si }\quad t=1 \text{ o }  t=3 
\]

la partícula estará en reposo al transcurrir uno o tres segundos del
inicio de su movimiento.

\begin{enumerate}
\def\labelenumi{\arabic{enumi}.}
\setcounter{enumi}{3}
\tightlist
\item
  Se mueve hacia adelante cuando el cambio en la posición respecto del
  tiempo es positivo, esto es, cuando \(f'(t)>0\). Notar que
\end{enumerate}

\[
3t^2-12t+9=3(t-1)(t-3) > 0 \quad \text{ si y sólo si }\quad t\in [0,1)\cup(3,+\infty).
\]

Es decir, la partícula se mueve hacia adelante durante el primer
segundo, y luego a partir de los tres segundos de recorrido.

\begin{enumerate}
\def\labelenumi{\arabic{enumi}.}
\setcounter{enumi}{4}
\tightlist
\item
  En la siguiente figura representamos el movimiento de la partícula. El
  movimiento hacia adelante o hacia atrás se simboliza con la dirección
  de las flechas.
\end{enumerate}

\includegraphics[width=0.4\linewidth,height=\textheight,keepaspectratio]{index_files/mediabag/imagenes_cap3/3.7-ejemplo-fisica.pdf}

\begin{enumerate}
\def\labelenumi{\arabic{enumi}.}
\setcounter{enumi}{5}
\tightlist
\item
  Para calcular la distancia total es preciso tener en cuenta si el
  movimiento es hacia adelante (derecha) o hacia atrás (izquierda). Como
  los cambios de sentido de recorrido se dan cuando \(t=1\) y \(t=3\),
  podemos calcular la distancia total de la siguiente manera
\end{enumerate}

\[
d=|s(1)-s(0)|+|s(3)-s(1)|+|s(5)-s(3)|=|4-0|+|0-4|+|20-0|=4+4+20=28.
\]

Durante los primeros cinco segundos la partícula recorre 28 m.

\begin{enumerate}
\def\labelenumi{\arabic{enumi}.}
\setcounter{enumi}{6}
\tightlist
\item
  La aceleración es la derivada de la velocidad, con lo cual
\end{enumerate}

\[
a(t)=\frac{dv}{dt}(t)=\frac{d^2s}{dt^2}(t)=6t-12.
\]

Después de cuatro segundos, la aceleración es de \(a(4)=6\cdot 4-12=12\)
m/seg2.

\begin{enumerate}
\def\labelenumi{\arabic{enumi}.}
\setcounter{enumi}{7}
\tightlist
\item
  Las gráficas de la posición, velocidad y aceleración se muestran a
  continuación
\end{enumerate}

\includegraphics[width=0.35\linewidth,height=\textheight,keepaspectratio]{index_files/mediabag/imagenes_cap3/3.7-ejemplo-funciones.pdf}

\begin{enumerate}
\def\labelenumi{\arabic{enumi}.}
\setcounter{enumi}{8}
\tightlist
\item
  Recordemos que la rapidez es el módulo de la velocidad, esto es
  \(|v(t)|\). Si partimos de la igualdad trivial
\end{enumerate}

\[
|v(t)|^2=v^2(t)
\]

y derivamos implícitamente, resulta

\[
2|v(t)|\frac{d}{dt}\left(|v(t)|\right)=2v(t)v'(t)
\]

de donde

\[
\frac{d}{dt}\left(|v(t)|\right) = \frac{v(t)v'(t)}{|v(t)|}=\frac{v(t)\,a(t)}{|v(t)|},
\]

siempre que \(|v(t)|\neq 0\).

Ahora bien, la rapidez se incrementa cuando su derivada es positiva, es
decir, cuando \(v\) y \(a\) tienen el mismo signo. Esto nos dice también
que la rapidez disminuye cuando \(v\) y \(a\) tienen signos opuestos.

\end{tcolorbox}

\subsection{Razones de cambio en
química}\label{razones-de-cambio-en-quuxedmica}

Consideremos una reacción química donde se forman una o más sustancias
(llamadas \textbf{productos}) a partir de otras sustancias (llamadas
\textbf{reactivos}). Esta reacción puede representarse con la ecuación
química

\[
A+B \rightarrow C
\]

donde \(A\) y \(B\) son los reactivos y \(C\) es el producto.

La \textbf{concentración} de una sustancia \(S\) es la cantidad de moles
(donde 1 mol equivale aproximadamente a \(6\text{.}022\times 10^{23}\)
moléculas) por cada litro y la denotaremos como \([S]\).

Como la concentración de los elementos varía durante la reacción,
tenemos que \([A], [B]\) y \([C]\) son funciones del tiempo de reacción.

La \textbf{velocidad de reacción promedio} del producto \(C\) en un
intervalo \([t_1,t_2]\) está dada por

\[
\frac{\Delta[C]}{\Delta t}=\frac{[C](t_2)-[C](t_1)}{t_2-t_1}.
\]

Es de mayor interés conocer la \textbf{velocidad de reacción
instantánea}, que se obtiene como el límite de las velocidades de
reacción promedio sobre intervalos cuya amplitud \(\Delta t\) tiende a
cero, es decir

\[
\text{velocidad de reacción}=\lim_{\Delta t\to 0}\frac{\Delta [C]}{\Delta t}=\frac{d[C]}{dt}.
\]

En una reacción, la concentración del producto aumenta conforme la de
los reactivos disminuye. Entonces tenemos que \(d[C]/dt\) es positiva
mientras que \(d[A]/dt\) y \(d[B]/dt\) son negativas. Además, \(A\) y
\(B\) disminuyen su concentración a la misma velocidad que \(C\) aumenta
la suya, con lo cual

\[
\frac{d[C]}{dt}=-\frac{d[A]}{dt}=-\frac{d[B]}{dt}.
\]

Ejemplo

En este modelo, supongamos que la concentración del producto \(C\) viene
dada por la fórmula

\[
[C](t)=\frac{a^2kt}{akt+1},
\]

donde \(k\) es una constante positiva, \(a\) es la concentración inicial
de los reactivos \(A\) y \(B\), y \(t\) es el tiempo de reacción medido
en segundos.

\begin{enumerate}
\def\labelenumi{\arabic{enumi}.}
\item
  Hallar la velocidad de reacción en cada instante \(t\).
\item
  Mostrar que \(\displaystyle \frac{d[C]}{dt}=k(a-[C])^2\).
\item
  ¿Qué sucede con la concentración cuando el tiempo tiende a infinito?
\item
  ¿Qué ocurre con la velocidad de reacción cuando el tiempo tiende a
  infinito?
\end{enumerate}

\begin{tcolorbox}[enhanced jigsaw, opacityback=0, left=2mm, opacitybacktitle=0.6, toprule=.15mm, bottomrule=.15mm, colbacktitle=quarto-callout-tip-color!10!white, titlerule=0mm, bottomtitle=1mm, title=\textcolor{quarto-callout-tip-color}{\faLightbulb}\hspace{0.5em}{Solución}, toptitle=1mm, colframe=quarto-callout-tip-color-frame, colback=white, breakable, rightrule=.15mm, arc=.35mm, leftrule=.75mm, coltitle=black]

\begin{enumerate}
\def\labelenumi{\arabic{enumi}.}
\tightlist
\item
  Calculamos la derivada de \([C]\) utilizando la regla del cociente
\end{enumerate}

\[
\frac{d[C]}{dt}(t)=\frac{a^2k(akt+1)-a^2kt(ak)}{(akt+1)^2}=\frac{a^2k}{(akt+1)^2}.
\]

\begin{enumerate}
\def\labelenumi{\arabic{enumi}.}
\setcounter{enumi}{1}
\tightlist
\item
  Notemos que
\end{enumerate}

\[
k(a-[C](t))^2=k\left(a-\frac{a^2kt}{akt+1}\right)^2=k\left(\frac{a^2kt+a-a^2kt}{akt+1}\right)=\frac{a^2k}{(akt+1)^2}=\frac{d[C]}{dt}(t).
\]

\begin{enumerate}
\def\labelenumi{\arabic{enumi}.}
\setcounter{enumi}{2}
\tightlist
\item
  Calculamos
\end{enumerate}

\[
\lim_{t\to \infty} [C](t)=\lim_{t\to \infty} \frac{a^2k}{ak+\frac{1}{t}}=\frac{a^2k}{ak}=a.
\]

Es decir, la concentración del producto tiende a ser la concentración
inicial de los reactivos.

\begin{enumerate}
\def\labelenumi{\arabic{enumi}.}
\setcounter{enumi}{3}
\tightlist
\item
  En este caso debemos hallar el límite de \(d[C]/dt\) para
  \(t\to\infty\). En efecto,
\end{enumerate}

\[
\lim_{t\to \infty} \frac{d[C]}{dt}(t)=\lim_{t\to \infty} \frac{a^2k}{(akt+1)^2}=\lim_{t\to \infty} \frac{a^2k}{t^2(ak+\frac{1}{t})^2}=\lim_{t\to \infty} \frac{1}{t^2}\,\frac{a^2k}{(ak+\frac{1}{t})^2}=0\cdot a=0.
\]

\end{tcolorbox}

\hyperref[seccion_3.7]{↑ Volver al inicio de la sección}

\section{3.9. Relaciones afines}\label{seccion_3.9}

En esta sección estudiaremos problemas en los que tenemos ciertas
cantidades relacionadas entre sí y que varían con el tiempo. La idea es
determinar la rapidez de cambio (o razón de cambio) de alguna de estas
cantidades conociendo la de otras que se relacionan con ésta. Veremos
diferentes ejemplos para ilustrar este tipo de situaciones.

Ejemplo

Se infla un globo esférico y su volumen se incrementa a razón de 100
cm3/seg. ¿Qué tan rápido aumenta el radio del globo cuando el diámetro
es de 50 cm?

\begin{tcolorbox}[enhanced jigsaw, opacityback=0, left=2mm, opacitybacktitle=0.6, toprule=.15mm, bottomrule=.15mm, colbacktitle=quarto-callout-tip-color!10!white, titlerule=0mm, bottomtitle=1mm, title=\textcolor{quarto-callout-tip-color}{\faLightbulb}\hspace{0.5em}{Solución}, toptitle=1mm, colframe=quarto-callout-tip-color-frame, colback=white, breakable, rightrule=.15mm, arc=.35mm, leftrule=.75mm, coltitle=black]

El globo tiene forma esférica, con lo cual el volumen \(V\) a tiempo
\(t\) viene dado por la fórmula

\[
V(t)=\frac{4}{3}\pi (r(t))^3,
\]

siendo \(r(t)\) el radio del globo a tiempo \(t\). Si derivamos a ambos
miembros con respecto a \(t\) y utilizamos la regla de la cadena
(recordar que tanto \(V\) como \(r\) son funciones de \(t\)) obtenemos

\[
\frac{dV}{dt}(t)=\frac{4}{3}\pi\,\, 3(r(t))^2\,\frac{dr}{dt}(t)=4\pi\,\, (r(t))^2\,\frac{dr}{dt}(t).
\]

Queremos conocer \(\displaystyle \frac{dr}{dt}(t_0)\), siendo \(t_0\) un
instante de tiempo para el cual se cumple que

\[
\frac{dV}{dt}(t_0)=100\, \mathrm{cm^3/seg} \quad \text{ y } \quad r(t_0)=25\, \mathrm{cm}.
\]

Reemplazando estos datos en la fórmula anterior resulta

\[
100\, \mathrm{cm^3/seg}=\frac{4}{3}\pi\,\, 3(25 \mathrm{cm})^2\,\frac{dr}{dt}(t_0),
\]

de donde despejando obtenemos

\[
\frac{dr}{dt}(t_0)=\frac{1}{25 \pi}\, \mathrm{cm/seg} \approx 0.0127\, \mathrm{cm/seg}.
\]

\end{tcolorbox}

Ejemplo

Una escalera de 10 pies de largo está apoyada contra un muro vertical.
Si la parte inferior de la escalera se desliza alejándose de la pared en
una proporción de 1 pie/seg, ¿qué tan rápido la parte superior de la
escalera resbala hacia abajo por la pared cuando la parte inferior de la
escalera está a 6 pies del muro?

\begin{center}
\includegraphics[width=0.35\linewidth,height=\textheight,keepaspectratio]{index_files/mediabag/imagenes_cap3/3.9-escalera.pdf}
\end{center}

\begin{tcolorbox}[enhanced jigsaw, opacityback=0, left=2mm, opacitybacktitle=0.6, toprule=.15mm, bottomrule=.15mm, colbacktitle=quarto-callout-tip-color!10!white, titlerule=0mm, bottomtitle=1mm, title=\textcolor{quarto-callout-tip-color}{\faLightbulb}\hspace{0.5em}{Solución}, toptitle=1mm, colframe=quarto-callout-tip-color-frame, colback=white, breakable, rightrule=.15mm, arc=.35mm, leftrule=.75mm, coltitle=black]

Sean \(y(t)\) y \(x(t)\) la altura del extremo de la escalera apoyado
sobre el muro y la distancia desde la pared al pie de la escalera
expresada en pies, respectivamente. Como la escalera mide 10 pies, para
todo instante de tiempo \(t\) se verifica que

\begin{equation}\phantomsection\label{eq-ejemplo-escalera1}{
(x(t))^2+(y(t))^2=100
}\end{equation}

y derivando implícitamente resulta

\begin{equation}\phantomsection\label{eq-ejemplo-escalera2}{
2 x(t) \frac{dx}{dt}(t) + 2y(t)\frac{dy}{dt}(t)=0 \quad \text{o también} \quad   x(t) \frac{dx}{dt}(t) + y(t)\frac{dy}{dt}(t)=0.
}\end{equation}

Ahora nos interesa encontrar \(\displaystyle \frac{dy}{dt}(t_0)\),
siendo \(t_0\) un instante de tiempo en el cual

\[
x(t_0)= 6\, \mathrm{pies} \quad \text{ y } \quad \frac{dx}{dt}(t_0)= 1 \, \mathrm{pie/seg}.
\]

Para calcular \(y(t_0)\) utilizamos la fórmula
(\ref{eq-ejemplo-escalera1}), obteniendo

\[
y(t_0)=\sqrt{100-(x(t_0))^2}=8\, \mathrm{pies}. 
\]

Ahora despejamos \(\displaystyle \frac{dy}{dt}\) de
(\ref{eq-ejemplo-escalera2}), con lo cual

\[
\frac{dy}{dt}(t_0)=-\frac{x(t_0)}{y(t_0)}\frac{dx}{dt}(t_0)=-\frac{6\,\mathrm{pies}}{8\,\mathrm{pies}} \cdot 1 \,\mathrm{pies/seg} = -\frac{3}{4}\,\mathrm{pies/seg}.
\]

La parte superior de la escalera está descendiendo a \(0.75\) pies/seg.

\end{tcolorbox}

Ejemplo

Un depósito para agua tiene la forma de un cono circular invertido. El
radio de la base es de 2 m y la altura es de 4 m. Si el agua se bombea
hacia el depósito a una razón de 2 m3/min, determine la rapidez a la
cual el nivel del agua sube cuando el agua tiene 3 m de profundidad.

\begin{center}
\includegraphics[width=0.35\linewidth,height=\textheight,keepaspectratio]{index_files/mediabag/imagenes_cap3/3.9-cono.pdf}
\end{center}

\begin{tcolorbox}[enhanced jigsaw, opacityback=0, left=2mm, opacitybacktitle=0.6, toprule=.15mm, bottomrule=.15mm, colbacktitle=quarto-callout-tip-color!10!white, titlerule=0mm, bottomtitle=1mm, title=\textcolor{quarto-callout-tip-color}{\faLightbulb}\hspace{0.5em}{Solución}, toptitle=1mm, colframe=quarto-callout-tip-color-frame, colback=white, breakable, rightrule=.15mm, arc=.35mm, leftrule=.75mm, coltitle=black]

Si denotamos con \(r(t)\) y \(h(t)\) al radio de la base y la altura de
la columna de líquido en el tanque a tiempo \(t\) y expresados en
metros, respectivamente, entonces el volumen de agua dentro del
recipiente a tiempo \(t\) es

\begin{equation}\phantomsection\label{eq-ejemplo-tanque1}{
V(t)= \frac{\pi}{3}(r(t))^2h(t).
}\end{equation}

Por semejanza de triángulos, podemos observar también que

\[
\frac{r(t)}{h(t)}=\frac{2}{4}=\frac{1}{2},
\]

de donde resulta que \(\displaystyle r(t)=\frac{h(t)}{2}\). Sustituyendo
en (\ref{eq-ejemplo-tanque1}) obtenemos

\[
V(t)=\frac{\pi}{12}(h(t))^3.
\]

Derivando con respecto a \(t\)

\[
\frac{dV}{dt}(t)=\frac{\pi}{4}(h(t))^2 \frac{dh}{dt}(t).
\]

Ahora determinamos \(\displaystyle \frac{dh}{dt}(t_0)\), sabiendo que
\(h(t_0)=3\, \mathrm{m}\) y
\(\displaystyle \frac{dV}{dt}(t_0)=2\, \mathrm {m^3/min}\). En efecto,

\[
\frac{dh}{dt}(t_0)=\frac{4}{\pi} \frac{2\, \mathrm {m^3/min}}{(3\, \mathrm{m})^2}=\frac{8}{9\pi}\, \mathrm{m/min} \approx 0.283\, \mathrm{m/min}.
\]

\end{tcolorbox}

Ejemplo

Dos automóviles se dirigen de forma perpendicular a la intersección de
dos caminos. El automóvil \(A\) se dirige hacia el oeste a 50 millas/h y
el vehículo \(B\) viaja hacia el norte a 60 millas/h. ¿Con qué rapidez
se aproximan los vehículos entre sí cuando el automóvil \(A\) está a 0.3
millas y el vehículo \(B\) está a 0.4 millas de la intersección?

\begin{center}
\includegraphics[width=0.35\linewidth,height=\textheight,keepaspectratio]{index_files/mediabag/imagenes_cap3/3.9-vehiculos.pdf}
\end{center}

\begin{tcolorbox}[enhanced jigsaw, opacityback=0, left=2mm, opacitybacktitle=0.6, toprule=.15mm, bottomrule=.15mm, colbacktitle=quarto-callout-tip-color!10!white, titlerule=0mm, bottomtitle=1mm, title=\textcolor{quarto-callout-tip-color}{\faLightbulb}\hspace{0.5em}{Solución}, toptitle=1mm, colframe=quarto-callout-tip-color-frame, colback=white, breakable, rightrule=.15mm, arc=.35mm, leftrule=.75mm, coltitle=black]

Sean \(x(t)\) e \(y(t)\) las distancias de los vehículos \(A\) y \(B\)
al punto de intersección de los caminos expresadas en millas,
respectivamente. Sea también \(z(t)\) la distancia en millas entre ambos
vehículos. En todo instante de tiempo \(t\) tenemos que

\begin{equation}\phantomsection\label{eq-ejemplo-automoviles1}{
x^2(t)+y^2(t)=z^2(t),
}\end{equation}

y al derivar de forma implícita

\begin{equation}\phantomsection\label{eq-ejemplo-automoviles2}{
2x(t)x'(t)+2y(t)y'(t)=2z(t)z'(t) \quad \text{ o bien } \quad x(t)x'(t)+y(t)y'(t)=z(t)z'(t).
}\end{equation}

Queremos encontrar \(z'(t_0)\), sabiendo que \(x(t_0)=0.3\),
\(y(t_0)=0.4\), \(x'(t_0)=-50\) e \(y'(t_0)=-60\). Para determinar
\(z(t_0)\) usamos la ecuación (\ref{eq-ejemplo-automoviles1})

\[
z(t_0)=\sqrt{x^2(t_0)+y^2(t_0)}=\sqrt{(0.3\, \mathrm{millas})^2+(0.4\, \mathrm{millas})^2}=0.5\, \mathrm{millas}.
\]

Despejando de (\ref{eq-ejemplo-automoviles2}) obtenemos

\[
\begin{aligned}
z'(t_0)=\frac{x(t_0)x'(t_0)+y(t_0)y'(t_0)}{z(t_0)}&=\frac{0.3\, \mathrm{millas} \cdot (-50 \, \mathrm{millas/h})+0.4\, \mathrm{millas} \cdot (-60 \, \mathrm{millas/h})}{0.5\, \mathrm{millas}}\\
\\
&=-78\, \mathrm{millas/h}.
\end{aligned}
\]

Es decir, los vehículos se están acercando con una rapidez de 78
millas/h.

\end{tcolorbox}

Ejemplo

Un hombre camina a lo largo de una trayectoria recta a una rapidez de 4
pies/seg. Un faro está situado sobre el nivel de la tierra a 20 pies de
la trayectoria y se mantiene enfocado hacia el hombre. ¿Con qué rapidez
el faro gira cuando el hombre está a 15 pies del punto sobre la
trayectoria más cercana a la fuente de luz?

\begin{center}
\includegraphics[width=0.35\linewidth,height=\textheight,keepaspectratio]{index_files/mediabag/imagenes_cap3/3.9-faro.pdf}
\end{center}

\begin{tcolorbox}[enhanced jigsaw, opacityback=0, left=2mm, opacitybacktitle=0.6, toprule=.15mm, bottomrule=.15mm, colbacktitle=quarto-callout-tip-color!10!white, titlerule=0mm, bottomtitle=1mm, title=\textcolor{quarto-callout-tip-color}{\faLightbulb}\hspace{0.5em}{Solución}, toptitle=1mm, colframe=quarto-callout-tip-color-frame, colback=white, breakable, rightrule=.15mm, arc=.35mm, leftrule=.75mm, coltitle=black]

Sea \(x(t)\) la distancia en pies desde el hombre al punto de la
trayectoria más cercano a la fuente de luz a tiempo \(t\), expresado en
segundos. Sea también \(\theta(t)\) el ángulo en radianes que se forma
entre la perpendicular a la trayectoria desde el faro y el rayo que
apunta al hombre, a tiempo \(t\). Observando el triángulo rectángulo que
se forma, podemos inferir que

\[
\tan(\theta(t))=\frac{x(t)}{20}.
\]

Derivando respecto del tiempo obtenemos

\begin{equation}\phantomsection\label{eq-ejemplo-faro}{
\sec^2(\theta(t))\frac{d\theta}{dt}(t)=\frac{1}{20}\frac{dx}{dt}(t).
}\end{equation}

Si \(x(t_0)=15\), la distancia del faro al hombre es de 25 pies, con lo
cual

\[
\cos(\theta(t_0))=\frac{4}{5}
\]

y despejando de (\ref{eq-ejemplo-faro}) resulta

\[
\frac{d\theta}{dt}(t_0)=\frac{1}{20}\cos^2(\theta(t_0))\frac{dx}{dt}(t_0)=\frac{1}{20}\left(\frac{4}{5}\right)^2\cdot 4\, \mathrm{ rad/seg}=0.128\, \mathrm{ rad/seg}.
\]

\end{tcolorbox}

\hyperref[seccion_3.9]{↑ Volver al inicio de la sección}

\bookmarksetup{startatroot}

\chapter{4. Aplicaciones de la
derivación}\label{aplicaciones-de-la-derivaciuxf3n}

\section{4.1 Valores máximos y mínimos}\label{seccion_4.1}

\begin{tcolorbox}[enhanced jigsaw, opacityback=0, left=2mm, opacitybacktitle=0.6, toprule=.15mm, bottomrule=.15mm, colbacktitle=quarto-callout-note-color!10!white, titlerule=0mm, bottomtitle=1mm, title=\textcolor{quarto-callout-note-color}{\faInfo}\hspace{0.5em}{Definición (Valores máximos y mínimos)}, toptitle=1mm, colframe=quarto-callout-note-color-frame, colback=white, breakable, rightrule=.15mm, arc=.35mm, leftrule=.75mm, coltitle=black]

Una función \(f\) tiene un \textbf{máximo absoluto} (o \textbf{máximo
global}) en \(c\) si

\[
f(x)\leq f(c) \quad \text{ para todo } \quad x\in D,
\]

siendo \(D\) el dominio de \(f\). Al número \(f(c)\) se lo llama
\textbf{valor máximo} de \(f\) en \(D\).

Análogamente, \(f\) tiene un \textbf{mínimo absoluto} (o \textbf{mínimo
global}) en \(c\) si

\[
f(x)\geq f(c) \quad \text{ para todo } \quad x\in D.
\]

En este caso, a \(f(c)\) se lo llama \textbf{valor mínimo} de \(f\) en
\(D\).

Los valores máximos y mínimos de \(f\) se conocen como \textbf{valores
extremos} de \(f\).

\end{tcolorbox}

En la siguiente gráfica vemos una función con máximo absoluto en \(d\) y
mínimo absoluto en \(a\). Esto significa que el punto \((d,f(d))\) es el
más alto de la gráfica y \((a,f(a))\) el más bajo.

\begin{center}
\includegraphics[width=0.4\linewidth,height=\textheight,keepaspectratio]{index_files/mediabag/imagenes_cap4/4.1-minymax.pdf}
\end{center}

Si solamente restringimos la atención al intervalo \((a,c)\), entonces
\(f(b)\) es el valor más grande de \(f\) en esa región. En este caso,
decimos que \(f(b)\) es un valor máximo local de \(f\). De manera
similar, \(f(c)\) y \(f(e)\) son valores mínimos locales de \(f\).

\begin{tcolorbox}[enhanced jigsaw, opacityback=0, left=2mm, opacitybacktitle=0.6, toprule=.15mm, bottomrule=.15mm, colbacktitle=quarto-callout-note-color!10!white, titlerule=0mm, bottomtitle=1mm, title=\textcolor{quarto-callout-note-color}{\faInfo}\hspace{0.5em}{Definición (Valores máximos y mínimos)}, toptitle=1mm, colframe=quarto-callout-note-color-frame, colback=white, breakable, rightrule=.15mm, arc=.35mm, leftrule=.75mm, coltitle=black]

Una función \(f\) tiene un \textbf{máximo local} (o \textbf{máximo
relativo}) en \(c\) si \(f(x)\leq f(c)\) para \(x\) \emph{cercano} a
\(c\) (es decir, si existe \(r>0\) tal que la desigualdad anterior se
cumple para todo \(x\in (c-r,c+r)\) que esté en el dominio de \(f\)).

De manera similar, \(f\) tiene un \textbf{mínimo local} en \(c\) si
\(f(x)\geq  f(c)\) para todo \(x\) cercano a \(c\).

\end{tcolorbox}

Ejemplo

Encontrar extremos locales y absolutos de la función \(f(x)=\cos x\).

\begin{tcolorbox}[enhanced jigsaw, opacityback=0, left=2mm, opacitybacktitle=0.6, toprule=.15mm, bottomrule=.15mm, colbacktitle=quarto-callout-tip-color!10!white, titlerule=0mm, bottomtitle=1mm, title=\textcolor{quarto-callout-tip-color}{\faLightbulb}\hspace{0.5em}{Solución}, toptitle=1mm, colframe=quarto-callout-tip-color-frame, colback=white, breakable, rightrule=.15mm, arc=.35mm, leftrule=.75mm, coltitle=black]

Recordemos que \(-1\leq \cos x\leq 1\) para todo \(x\in\mathbb{R}\). Por
lo tanto, \(1\) es el valor máximo absoluto de \(f\) (y, en
consecuencia, también local) y \(f\) toma este valor una cantidad
infinita de veces, en todo \(x=2n\pi\), con \(n\in \mathbb{Z}\).

Por otro lado, \(-1\) es el valor mínimo absoluto (y local) de \(f\),
que se alcanza siempre que \(x=(2n+1)\pi\), con \(n\in \mathbb{Z}\).

\includegraphics[width=0.5\linewidth,height=\textheight,keepaspectratio]{index_files/mediabag/imagenes_cap4/4.1-coseno.pdf}

\end{tcolorbox}

Ejemplo

Analizar extremos locales y absolutos de \(f(x)=x^2\).

\begin{tcolorbox}[enhanced jigsaw, opacityback=0, left=2mm, opacitybacktitle=0.6, toprule=.15mm, bottomrule=.15mm, colbacktitle=quarto-callout-tip-color!10!white, titlerule=0mm, bottomtitle=1mm, title=\textcolor{quarto-callout-tip-color}{\faLightbulb}\hspace{0.5em}{Solución}, toptitle=1mm, colframe=quarto-callout-tip-color-frame, colback=white, breakable, rightrule=.15mm, arc=.35mm, leftrule=.75mm, coltitle=black]

Observemos primero que \(f(0)=0\leq x^2=f(x)\), es decir,
\(f(0)\leq f(x)\) para cualquier \(x\in\mathbb{R}\). Esto nos dice que
\(0\) es el valor mínimo absoluto (y local) de la función \(f\), y se
alcanza cuando \(x=0\).

Sin embargo, \(f\) no tiene máximo absoluto, ya que

\[
\lim_{x\to \infty} x^2 =\infty \quad \text{ y } \quad \lim_{x\to -\infty} x^2=\infty.
\]

Inspeccionando el gráfico de \(f\) podemos ver que tampoco hay máximos
locales.

\includegraphics[width=0.25\linewidth,height=\textheight,keepaspectratio]{index_files/mediabag/imagenes_cap4/4.1-cuadratica.pdf}

\end{tcolorbox}

Ejemplo

Analizar extremos para la función \(f(x)=x^3\).

\begin{tcolorbox}[enhanced jigsaw, opacityback=0, left=2mm, opacitybacktitle=0.6, toprule=.15mm, bottomrule=.15mm, colbacktitle=quarto-callout-tip-color!10!white, titlerule=0mm, bottomtitle=1mm, title=\textcolor{quarto-callout-tip-color}{\faLightbulb}\hspace{0.5em}{Solución}, toptitle=1mm, colframe=quarto-callout-tip-color-frame, colback=white, breakable, rightrule=.15mm, arc=.35mm, leftrule=.75mm, coltitle=black]

Notemos que \(f\) no tiene valor extremos absolutos, puesto que

\[
\lim_{x\to \infty} x^3 =\infty \quad \text{ y } \quad \lim_{x\to -\infty} x^3=-\infty.
\]

Observando el gráfico de \(f\) podemos comprobar que tampoco existen
extremos locales.

\includegraphics[width=0.25\linewidth,height=\textheight,keepaspectratio]{index_files/mediabag/imagenes_cap4/4.1-cubica.pdf}

\end{tcolorbox}

Ejemplo

Analizar los extremos de la función

\[
f(x)=3x^4-16x^3+18x^2, \quad\text{ para }\quad  -1\leq x\leq 4.
\]

\begin{center}
\includegraphics[width=0.4\linewidth,height=\textheight,keepaspectratio]{index_files/mediabag/imagenes_cap4/4.1-ejemplo-minmax.pdf}
\end{center}

\begin{tcolorbox}[enhanced jigsaw, opacityback=0, left=2mm, opacitybacktitle=0.6, toprule=.15mm, bottomrule=.15mm, colbacktitle=quarto-callout-tip-color!10!white, titlerule=0mm, bottomtitle=1mm, title=\textcolor{quarto-callout-tip-color}{\faLightbulb}\hspace{0.5em}{Solución}, toptitle=1mm, colframe=quarto-callout-tip-color-frame, colback=white, breakable, rightrule=.15mm, arc=.35mm, leftrule=.75mm, coltitle=black]

Observando el gráfico podemos ver que \(f(1)=5\) es un valor máximo
local y \(f(-1)=37\) es el valor máximo absoluto de \(f\) sobre el
intervalo (éste no se considera máximo local por encontrarse en un
extremo del intervalo dado).

Por otra parte, \(f(0)=0\) es un mínimo local y \(f(3)=-27\) es el valor
mínimo absoluto, y también resulta local.

\end{tcolorbox}

El siguiente teorema será una herramienta importante para encontrar
extremos de funciones continuas sobre intervalos cerrados.

\phantomsection\label{teo-valores-ext}
Teorema del valor extremo

Si \(f\) es continua sobre el intervalo cerrado \([a,b]\), entonces
\(f\) alcanza un valor máximo absoluto \(f(c)\) y un valor mínimo
absoluto \(f(d)\) en algunos números \(c,d\) en \([a,b]\).

Este teorema nos asegura existencia de los puntos \(c\) y \(d\), pero
como vemos en la siguiente figura éstos podrían no ser únicos.

\begin{center}
\includegraphics[width=0.7\linewidth,height=\textheight,keepaspectratio]{index_files/mediabag/imagenes_cap4/4.1-teovalorextremo.pdf}
\end{center}

En general, si \(f\) no es continua o el intervalo no es cerrado, el
resultado anterior podría no valer. Aquí vemos algunos ejemplos donde
esto ocurre.

\begin{center}
\includegraphics[width=0.6\linewidth,height=\textheight,keepaspectratio]{index_files/mediabag/imagenes_cap4/4.1-teovalorextremo2.pdf}
\end{center}

A la izquierda, vemos que \(f\) tiene un valor mínimo \(f(2)=0\), pero
no hay valor máximo. A la derecha no hay valor mínimo ni máximo.

El \hyperref[teo-valores-ext]{teorema anterior} nos da condiciones para
asegurar la existencia de extremos de una función en un intervalo, pero
no dice cómo encontrarlos. Para responder a esta pregunta, veamos el
siguiente resultado.

\phantomsection\label{teo-fermat}
Teorema de Fermat

Si \(f\) tiene un extremo local en \(c\) y \(f'(c)\) existe, entonces
\(f'(c)=0\).

\begin{center}
\includegraphics[width=0.35\linewidth,height=\textheight,keepaspectratio]{index_files/mediabag/imagenes_cap4/4.1-teofermat.pdf}
\end{center}

\begin{tcolorbox}[enhanced jigsaw, opacityback=0, left=2mm, opacitybacktitle=0.6, toprule=.15mm, bottomrule=.15mm, colbacktitle=quarto-callout-caution-color!10!white, titlerule=0mm, bottomtitle=1mm, title=\textcolor{quarto-callout-caution-color}{\faFire}\hspace{0.5em}{Importante}, toptitle=1mm, colframe=quarto-callout-caution-color-frame, colback=white, breakable, rightrule=.15mm, arc=.35mm, leftrule=.75mm, coltitle=black]

El recíproco del \hyperref[teo-fermat]{teorema de Fermat} es
\textbf{falso}, como veremos en el siguiente ejemplo.

\end{tcolorbox}

Ejemplo

Mostrar que \(f(x)=x^3\) satisface \(f'(0)=0\) pero que \(f\) no alcanza
extremos locales en \(x=0\).

\begin{tcolorbox}[enhanced jigsaw, opacityback=0, left=2mm, opacitybacktitle=0.6, toprule=.15mm, bottomrule=.15mm, colbacktitle=quarto-callout-tip-color!10!white, titlerule=0mm, bottomtitle=1mm, title=\textcolor{quarto-callout-tip-color}{\faLightbulb}\hspace{0.5em}{Solución}, toptitle=1mm, colframe=quarto-callout-tip-color-frame, colback=white, breakable, rightrule=.15mm, arc=.35mm, leftrule=.75mm, coltitle=black]

Dado que \(f'(x)=3x^2\), es inmediato comprobar que \(f'(0)=0\). Por
otro lado, si consideramos el intervalo \((0-r,0+r)=(-r,r)\) para
cualquier \(r>0\), tenemos que

\[
f(x)<f(0) \quad \text{ para } -r<x<0
\]

y

\[
f(x)>f(0) \quad \text{ para } 0<x<r.
\]

En consecuencia, en \(x=0\) no puede haber ni un mínimo ni un máximo
local.

\end{tcolorbox}

Ejemplo

Mostrar que \(f(x)=|x|\) presenta un extremo (local y global) en \(x=0\)
pero \(f'(0)\) no existe.

\begin{tcolorbox}[enhanced jigsaw, opacityback=0, left=2mm, opacitybacktitle=0.6, toprule=.15mm, bottomrule=.15mm, colbacktitle=quarto-callout-tip-color!10!white, titlerule=0mm, bottomtitle=1mm, title=\textcolor{quarto-callout-tip-color}{\faLightbulb}\hspace{0.5em}{Solución}, toptitle=1mm, colframe=quarto-callout-tip-color-frame, colback=white, breakable, rightrule=.15mm, arc=.35mm, leftrule=.75mm, coltitle=black]

Dado que \(f(0)=0\leq |x|=f(x)\) para todo \(x\), \(f\) presenta un
mínimo local y global en \(x=0\), de valor 0. Sin embargo, según vimos
en el Ejemplo 44 de la \hyperref[seccion_2.8]{Sección 2.8}, \(f\) no es
derivable en \(x=0\).

\end{tcolorbox}

El teorema de Fermat establece que si \(f\) presenta algún extremo
local, éste debe ocurrir en un punto \(c\) tal que \(f'(c)=0\) o bien
que cumpla que \(f'(c)\) no exista. Introduciremos una definición que
será de utilidad en adelante.

\begin{tcolorbox}[enhanced jigsaw, opacityback=0, left=2mm, opacitybacktitle=0.6, toprule=.15mm, bottomrule=.15mm, colbacktitle=quarto-callout-note-color!10!white, titlerule=0mm, bottomtitle=1mm, title=\textcolor{quarto-callout-note-color}{\faInfo}\hspace{0.5em}{Definición (Número crítico)}, toptitle=1mm, colframe=quarto-callout-note-color-frame, colback=white, breakable, rightrule=.15mm, arc=.35mm, leftrule=.75mm, coltitle=black]

Si \(c\) es un número en el dominio de \(f\) tal que \(f'(c)=0\) o bien
\(f'(c)\) no existe, entonces \(c\) se llama \textbf{número crítico} o
\textbf{número estacionario} de \(f\).

\end{tcolorbox}

Ejemplo

Encontrar los números críticos de \(f(x)=x^{3/5}(4-x)\).

\begin{tcolorbox}[enhanced jigsaw, opacityback=0, left=2mm, opacitybacktitle=0.6, toprule=.15mm, bottomrule=.15mm, colbacktitle=quarto-callout-tip-color!10!white, titlerule=0mm, bottomtitle=1mm, title=\textcolor{quarto-callout-tip-color}{\faLightbulb}\hspace{0.5em}{Solución}, toptitle=1mm, colframe=quarto-callout-tip-color-frame, colback=white, breakable, rightrule=.15mm, arc=.35mm, leftrule=.75mm, coltitle=black]

La derivada de \(f\) es

\[
f'(x)=\frac{3}{5}x^{3/5-1}(4-x)+x^{3/5}(-1)=\frac{12-3x-5x}{5x^{2/5}}=\frac{12-8x}{5x^{2/5}}.
\]

Luego, \(f'(x)=0\) si y sólo si \(12-8x=0\), o bien \(x=3/2\). Por otro
lado, \(f'(x)\) no existe si \(x=0\). Por lo tanto, los números críticos
de \(f\) son \(x=0\) y \(x=3/2\).

\end{tcolorbox}

Con la definición de número crítico, el teorema de Fermat puede
refrasearse como sigue:

Si \(f\) tiene un máximo o un mínimo local en \(c\), entonces \(c\) es
un número crítico de \(f\).

Es decir, todos los extremos locales tienen lugar en números críticos.
Sin embargo, no todo número crítico da lugar a un extremo, por lo que
será necesario clasificarlos de alguna manera, como veremos luego.

\subsection{\texorpdfstring{¿Cómo encontrar los extremos de una función
continua \(f\) en un intervalo cerrado
\([a,b]\)?}{¿Cómo encontrar los extremos de una función continua f en un intervalo cerrado {[}a,b{]}?}}\label{cuxf3mo-encontrar-los-extremos-de-una-funciuxf3n-continua-f-en-un-intervalo-cerrado-ab}

Para ello podemos seguir los siguientes pasos:

\begin{enumerate}
\def\labelenumi{\arabic{enumi}.}
\tightlist
\item
  Hallar los valores de \(f\) en cada número crítico de \(f\) en
  \((a,b)\).
\item
  Calcular \(f(a)\) y \(f(b)\).
\item
  Comparar todos estos valores. El más grande de todos es el valor
  máximo absoluto y el más chico el valor mínimo absoluto.
\end{enumerate}

Ejemplo

Calcular los valores máximos y mínimos absolutos de \(f(x)=x^3-3x^2+1\)
en el itervalo \([-1/2, 4]\).

\begin{tcolorbox}[enhanced jigsaw, opacityback=0, left=2mm, opacitybacktitle=0.6, toprule=.15mm, bottomrule=.15mm, colbacktitle=quarto-callout-tip-color!10!white, titlerule=0mm, bottomtitle=1mm, title=\textcolor{quarto-callout-tip-color}{\faLightbulb}\hspace{0.5em}{Solución}, toptitle=1mm, colframe=quarto-callout-tip-color-frame, colback=white, breakable, rightrule=.15mm, arc=.35mm, leftrule=.75mm, coltitle=black]

Antes de comenzar, notemos que \(f\) es continua en el intervalo cerrado
\([-1/2,4]\) por tratarse de un polinomio. El
\hyperref[teo-valores-ext]{teorema del valor extremo} nos asegura
entonces que existen el máximo y el mínimo absolutos de \(f\) en el
intervalo.

Primero buscamos los números críticos de \(f\). Tenemos que
\(f'(x)=3x^2-6x=3x(x-2)\), por lo que \(f'(x)=0\) si y sólo si \(x=0\) o
\(x=2\). Estos son los únicos puntos críticos de \(f\) ya que no hay
valores de \(x\) para los que \(f'(x)\) no exista.

Siguiendo los pasos de la lista anterior, hallamos

\[
f(0)=0^3-3\cdot 0^2+1=1 \quad \text{ y } \quad f(2)=2^3-3\cdot 2^2+1=-3.
\]

Ahora calculamos la imagen por \(f\) de los extremos del intervalo

\[
f\left(-\frac{1}{2}\right)=\left(-\frac{1}{2}\right)^3-3\cdot \left(-\frac{1}{2}\right)^2+1=\frac{1}{8} \quad \text{ y } \quad f(4)=4^3-3\cdot 4^2+1=17.
\]

El valor máximo absoluto de \(f\) es 17 y se alcanza en \(x=4\), y el
valor mínimo absoluto es \(-3\), cuando \(x=2\).

\includegraphics[width=0.4\linewidth,height=\textheight,keepaspectratio]{index_files/mediabag/imagenes_cap4/4.1-ejemplo-cubica.pdf}

\end{tcolorbox}

Ejemplo

Hallar los valores máximos y mínimos de \(f(x)=x-2\operatorname{sen} x\)
en \(I=[0, 2\pi]\).

\begin{tcolorbox}[enhanced jigsaw, opacityback=0, left=2mm, opacitybacktitle=0.6, toprule=.15mm, bottomrule=.15mm, colbacktitle=quarto-callout-tip-color!10!white, titlerule=0mm, bottomtitle=1mm, title=\textcolor{quarto-callout-tip-color}{\faLightbulb}\hspace{0.5em}{Solución}, toptitle=1mm, colframe=quarto-callout-tip-color-frame, colback=white, breakable, rightrule=.15mm, arc=.35mm, leftrule=.75mm, coltitle=black]

Primero observemos que \(y=x\) e \(y=2\operatorname{sen} x\) son
funciones continuas en \(I\), con lo cual \(f\) es continua en \(I\). El
\hyperref[teo-valores-ext]{teorema del valor extremo} asegura la
existencia de extremos absolutos.

Notar que \(f'(x)=1-2\cos x\), con lo cual \(f'(x)=0\) si y sólo si
\(\cos x=\frac{1}{2}\). Los puntos de \(I\) que satisfacen esta ecuación
son \(x=\pi/3\) y \(x=5\pi/3\).

Evaluamos \(f\) en estos puntos

\[
f\left(\frac{\pi}{3}\right)=\frac{\pi}{3}-2\operatorname{sen}\left(\frac{\pi}{3}\right)=\frac{\pi}{3}-\sqrt{3}\approx -0.68 \quad \text{ y }\quad f\left(\frac{5\pi}{3}\right)=\frac{5\pi}{3}-2\operatorname{sen}\left(\frac{5\pi}{3}\right)=\frac{5\pi}{3}+\sqrt{3}\approx 6.97 
\]

y luego en los extremos de \(I\)

\[
f(0)=0-2\operatorname{sen} 0=0 \quad \text{ y } \quad f(2\pi)=2\pi-2\operatorname{sen}(2\pi)=2\pi \approx 6.28,
\]

con lo que el valor máximo absoluto de \(f\) en \(I\) es
\(5\pi/3+\sqrt{3}\) cuando \(x=5\pi/3\) y el valor mínimo absoluto es
\(\pi/3-\sqrt{3}\) cuando \(x=\pi/3\).

\includegraphics[width=0.4\linewidth,height=\textheight,keepaspectratio]{index_files/mediabag/imagenes_cap4/4.1-ejemplo-seno.pdf}

\end{tcolorbox}

\hyperref[seccion_4.1]{↑ Volver al inicio de la sección}

\section{4.3. ¿Cómo afectan las derivadas la forma de una
gráfica?}\label{seccion_4.3}

\subsection{\texorpdfstring{¿Qué dice \(f'\) con respecto a
\(f\)?}{¿Qué dice f\textquotesingle{} con respecto a f?}}\label{quuxe9-dice-f-con-respecto-a-f}

En la siguiente figura podemos ver que entre los puntos \(A\) y \(B\)
las rectas tangentes tienen pendiente positiva, con lo cual \(f'(x)>0\).
Esto mismo ocurre entre \(C\) y \(D\). Por otro lado, entre \(B\) y
\(C\) las pendientes son negativas, con lo que \(f'(x)<0\).

\begin{center}
\includegraphics[width=0.5\linewidth,height=\textheight,keepaspectratio]{index_files/mediabag/imagenes_cap4/4.3-crec-decrec.pdf}
\end{center}

\phantomsection\label{teo-prueba-cd}
Prueba creciente/decreciente

\begin{itemize}
\item
  Si \(f'(x)>0\) sobre un intervalo, entonces \(f\) es creciente en ese
  intervalo.
\item
  Si \(f'(x)<0\) sobre un intervalo, \(f\) es decreciente en ese
  intervalo.
\end{itemize}

Ejemplo

Encontrar dónde cree la función \(f(x)=3x^4-4x^3-12x^2+5\) y dónde
decrece.

\begin{tcolorbox}[enhanced jigsaw, opacityback=0, left=2mm, opacitybacktitle=0.6, toprule=.15mm, bottomrule=.15mm, colbacktitle=quarto-callout-tip-color!10!white, titlerule=0mm, bottomtitle=1mm, title=\textcolor{quarto-callout-tip-color}{\faLightbulb}\hspace{0.5em}{Solución}, toptitle=1mm, colframe=quarto-callout-tip-color-frame, colback=white, breakable, rightrule=.15mm, arc=.35mm, leftrule=.75mm, coltitle=black]

Tenemos que \(f'(x)=12x^3-12x^2-24x=12x(x^2-x-2)=12x(x-2)(x+1)\). Con
esto, vemos que \(f'\) se anula en \(x=-1\), \(x=0\) y \(x=2\). Para
determinar el signo de \(f'\) en cada uno de los intervalos que resultan
de dividir la recta en los puntos críticos de \(f\) realizamos la
siguiente tabla

\begin{longtable}[]{@{}lllll@{}}
\toprule\noalign{}
Intervalo & \((x+1)\) & \(12x\) & \((x-2)\) & \(f'(x)\) \\
\midrule\noalign{}
\endhead
\bottomrule\noalign{}
\endlastfoot
\((-\infty,-1)\) & \(-\) & \(-\) & \(-\) & \(-\) \\
\((-1,0)\) & \(+\) & \(-\) & \(-\) & \(+\) \\
\((0,2)\) & \(+\) & \(+\) & \(-\) & \(-\) \\
\((2,+\infty)\) & \(+\) & \(+\) & \(+\) & \(+\) \\
\end{longtable}

Observando la última columna de la tabla y utilizando
\hyperref[teo-prueba-cd]{la prueba creciente/decreciente}, obtenemos que
\(f\) es creciente en \((-1,0)\) y en \((2,+\infty)\), y decrece en
\((-\infty,-1)\) y en \((0,2)\).

\end{tcolorbox}

\begin{tcolorbox}[enhanced jigsaw, opacityback=0, left=2mm, opacitybacktitle=0.6, toprule=.15mm, bottomrule=.15mm, colbacktitle=quarto-callout-caution-color!10!white, titlerule=0mm, bottomtitle=1mm, title=\textcolor{quarto-callout-caution-color}{\faFire}\hspace{0.5em}{Importante}, toptitle=1mm, colframe=quarto-callout-caution-color-frame, colback=white, breakable, rightrule=.15mm, arc=.35mm, leftrule=.75mm, coltitle=black]

Notar que en el ejemplo anterior hemos escrito: \(f\) es creciente en
\((-1,0)\) y en \((2,+\infty)\), y \textbf{no escribimos} \(f\) es
creciente en \((-1,0)\cup (2,+\infty)\). Ésta última expresión es
incorrecta, pues de acuerdo con la definición de función creciente,
debería ocurrir que \(f(x_1)<f(x_2)\) siempre que \(x_1<x_2\). Al
considerar el conjunto \((-1,0)\cup (2,+\infty)\) ésto no se cumpliría
pues, por ejemplo, \(x_1=0.5\) y \(x_2=2.5\) están en este conjunto, y
\(x_1<x_2\), sin embargo \(f(x_1)=1.6875\) y \(f(x_2)=-15.3125\).

Una observación similar puede hacerse para el caso decreciente.

\end{tcolorbox}

El siguiente resultado nos permitirá clasificar los números críticos de
una función para saber si corresponden a extremos locales.

\phantomsection\label{teo-prueba-primerader}
Prueba de la primera derivada

Sea \(f\) continua, \(c\) un número crítico de \(f\) y supongamos que
\(f'(x)\) existe cerca de \(c\) (excepto quizás en \(c\)). Luego:

\begin{itemize}
\tightlist
\item
  si \(f'(x)>0\) para \(x<c\) y \(f'(x)<0\) para \(x>c\), entonces \(f\)
  tiene un máximo local en \(x=c\).
\item
  si \(f'(x)<0\) para \(x<c\) y \(f'(x)>0\) para \(x>c\), entonces \(f\)
  tiene un mínimo local en \(x=c\).
\item
  si \(f'\) es positiva (o negativa) a ambos lados de \(c\), entonces
  \(f\) no tiene máximo ni mínimo local en \(x=c\).
\end{itemize}

\begin{center}
\includegraphics[width=0.6\linewidth,height=\textheight,keepaspectratio]{index_files/mediabag/imagenes_cap4/4.3-primerader1.pdf}
\end{center}

\begin{center}
\includegraphics[width=0.6\linewidth,height=\textheight,keepaspectratio]{index_files/mediabag/imagenes_cap4/4.3-primerader2.pdf}
\end{center}

Ejemplo

Encontrar los valores máximos y mínimos locales de la función \(f\) del
Ejemplo 10.

\begin{tcolorbox}[enhanced jigsaw, opacityback=0, left=2mm, opacitybacktitle=0.6, toprule=.15mm, bottomrule=.15mm, colbacktitle=quarto-callout-tip-color!10!white, titlerule=0mm, bottomtitle=1mm, title=\textcolor{quarto-callout-tip-color}{\faLightbulb}\hspace{0.5em}{Solución}, toptitle=1mm, colframe=quarto-callout-tip-color-frame, colback=white, breakable, rightrule=.15mm, arc=.35mm, leftrule=.75mm, coltitle=black]

Ya sabemos que los números críticos de \(f\) son \(-1\), \(0\) y \(2\).
En \(x=-1\) el signo de \(f'\) pasa de negativo a positivo, con lo que
la \hyperref[teo-prueba-primerader]{prueba de la primera derivada} nos
dice que \(f\) tiene un mínimo local en \(x=-1\), cuyo valor es
\(f(-1)=0\).

De manera similar, \(f'\) pasa de positiva a negativa en \(x=0\), con lo
que en este punto \(f\) presenta un máximo local de valor \(f(0)=5\).

Por último, \(f'\) cambia de negativa a positiva en \(x=2\), con lo que
aquí \(f\) tiene un mínimo local, cuyo valor es \(f(2)=-27\).

\end{tcolorbox}

La siguiente propiedad será de utilidad en diversas ocasiones.

\phantomsection\label{teo-propiedad-signoconst}
Propiedad del signo constante

Supongamos que \(f\) es continua en \((a,b)\) y además \(f(x)\neq 0\)
para todo \(x\in (a,b)\). Entonces \(f\) tiene signo constante sobre el
intervalo, es decir, o bien se cumple que

\[f(x)>0 \quad \text{ para todo }x\in (a,b)\]

o bien

\[f(x)<0 \quad \text{ para todo }x\in (a,b).\]

Ejemplo

Determinar los valores máximo y mínimo de la función
\(g(x)=x+2\operatorname{sen} x\) en el intervalo \([0,2\pi]\).

\begin{tcolorbox}[enhanced jigsaw, opacityback=0, left=2mm, opacitybacktitle=0.6, toprule=.15mm, bottomrule=.15mm, colbacktitle=quarto-callout-tip-color!10!white, titlerule=0mm, bottomtitle=1mm, title=\textcolor{quarto-callout-tip-color}{\faLightbulb}\hspace{0.5em}{Solución}, toptitle=1mm, colframe=quarto-callout-tip-color-frame, colback=white, breakable, rightrule=.15mm, arc=.35mm, leftrule=.75mm, coltitle=black]

El \hyperref[teo-valor-ext]{teorema del valor extremo} asegura la
existencia de valor máximo y mínimo absoluto, ya que \(g\) es continua
en el intervalo cerrado \([0, 2\pi]\).

La derivada de \(g\) es \(g'(x)=1+2\cos x\), y \(g'(x)=0\) si y sólo si
\(x=2\pi/3\) o \(x=4\pi/3\).

Ahora bien, \(g'\) es continua en \((0,2\pi/3)\), \((2\pi/3,4\pi/3)\) y
en \((4\pi/3,2\pi)\) y además es no nula en cada uno de estos
intervalos.

Por la \hyperref[teo-propiedad-signoconst]{propiedad del signo
constante}, \(g'\) tiene signo constante en cada uno de estos
intervalos. Por lo tanto, para analizar el signo de \(g'\) en cada
intervalo basta con evaluar el signo de \(g'\) en cualquier punto de los
mismos. La siguiente tabla muestra estos cálculos, en la cual el
crecimiento o decrecimiento en cada caso es consecuencia de la
\hyperref[teo-prueba-cd]{prueba creciente/decreciente}.

\begin{longtable}[]{@{}llll@{}}
\toprule\noalign{}
Intervalo & \(x_0\) & \(g'(x_0)\) & \(g\) \\
\midrule\noalign{}
\endhead
\bottomrule\noalign{}
\endlastfoot
\((0,2\pi/3)\) & 1 & \(+\) & crece \\
\((2\pi/3,4\pi/3)\) & 3 & \(-\) & decrece \\
\((4\pi/3,2\pi)\) & 5 & \(+\) & crece \\
\end{longtable}

Utilizando ahora la \hyperref[teo-prueba-primerader]{prueba de la
primera derivada} podemos concluir que \(g\) presenta un máximo local en
\(x=2\pi/3\), de valor

\[
g\left(\frac{2\pi}{3}\right)=\frac{2\pi}{3}+\sqrt{3}\approx 3.83
\]

y en \(x=4\pi/3\) presenta un mínimo local cuyo valor es

\[
g\left(\frac{4\pi}{3}\right)=\frac{4\pi}{3}-\sqrt{3}\approx 2.46.
\]

Ahora analizamos los extremos del intervalo

\[
g(0)=0+2\operatorname{sen} 0=0 \quad \text{ y } \quad g(2\pi)=2\pi+2\operatorname{sen}(2\pi)=2\pi \approx 6.28,
\]

con lo que \(g\) alcanza sus valores mínimo y máximo absoluto en los
extremos izquierdo y derecho del intervalo, respectivamente.

\end{tcolorbox}

\subsection{\texorpdfstring{¿Qué dice \(f''\) con respecto a
\(f\)?}{¿Qué dice f\textquotesingle\textquotesingle{} con respecto a f?}}\label{quuxe9-dice-f-con-respecto-a-f-1}

En la siguiente figura vemos el gráfico de dos funciones crecientes que
unen el punto \(A\) con el punto \(B\), pero se flexionan en direcciones
distintas.

\begin{center}
\includegraphics[width=0.6\linewidth,height=\textheight,keepaspectratio]{index_files/mediabag/imagenes_cap4/4.3-concavidad1.pdf}
\end{center}

Si dibujamos algunas rectas tangentes en distintos puntos, vemos que la
gráfica de la izquierda queda por arriba de estas rectas, y la de la
derecha por debajo.

\begin{center}
\includegraphics[width=0.6\linewidth,height=\textheight,keepaspectratio]{index_files/mediabag/imagenes_cap4/4.3-concavidad2.pdf}
\end{center}

\begin{tcolorbox}[enhanced jigsaw, opacityback=0, left=2mm, opacitybacktitle=0.6, toprule=.15mm, bottomrule=.15mm, colbacktitle=quarto-callout-note-color!10!white, titlerule=0mm, bottomtitle=1mm, title=\textcolor{quarto-callout-note-color}{\faInfo}\hspace{0.5em}{Definición (Concavidad)}, toptitle=1mm, colframe=quarto-callout-note-color-frame, colback=white, breakable, rightrule=.15mm, arc=.35mm, leftrule=.75mm, coltitle=black]

Si la gráfica de \(f\) queda por arriba de todas sus tangentes en un
intervalo \(I\), entonces \(f\) se dice \textbf{cóncava hacia arriba} o
\textbf{convexa} en \(I\).

Si en cambio la gráfica de \(f\) queda por debajo de sus tangentes en
\(I\), \(f\) se dice \textbf{cóncava hacia abajo} o \textbf{cóncava} en
\(I\).

\end{tcolorbox}

En la primera gráfica vemos que la pendiente se incrementa de izquierda
a derecha. Esto quiere decir que \(f'\) es creciente y, por lo tanto, su
derivada \(f''\) es positiva. De manera similar, en la segunda gráfica
la pendiente disminuye de izquierda a derecha, con lo que \(f'\) es
decreciente y su derivada \(f''\) resulta negativa. Esta idea nos ayuda
a establecer el siguiente criterio de concavidad.

\phantomsection\label{teo-prueba-concavidad}
Prueba de la concavidad

\begin{itemize}
\item
  Si \(f''(x)>0\) para todo \(x\) en \(I=(a,b)\), entonces la gráfica de
  \(f\) es cóncava hacia arriba en \(I\).
\item
  Si \(f''(x)<0\) para todo \(x\in I\), entonces la gráfica de \(f\) es
  cóncava hacia abajo sobre \(I\).
\end{itemize}

\begin{tcolorbox}[enhanced jigsaw, opacityback=0, left=2mm, opacitybacktitle=0.6, toprule=.15mm, bottomrule=.15mm, colbacktitle=quarto-callout-note-color!10!white, titlerule=0mm, bottomtitle=1mm, title=\textcolor{quarto-callout-note-color}{\faInfo}\hspace{0.5em}{Definición (Punto de inflexión)}, toptitle=1mm, colframe=quarto-callout-note-color-frame, colback=white, breakable, rightrule=.15mm, arc=.35mm, leftrule=.75mm, coltitle=black]

Un punto \(P\) en la curva \(y=f(x)\) recibe el nombre de \textbf{punto
de inflexión} si \(f\) es continua allí y además la curva cambia de
cóncava hacia abajo a cóncava hacia arriba, o bien de cóncava hacia
arriba a cóncava hacia abajo en \(P\).

\end{tcolorbox}

En virtud de la prueba de concavidad, habrá puntos de inflexión donde la
derivada segunda cambie de signo.

Ejemplo

Trazar una posible gráfica de una función \(f\) que cumpla las
siguientes condiciones:

\begin{itemize}
\item
  \(f'(x)>0\) en \((-\infty,1)\), \(f'(x)<0\) en \((1,\infty)\),
\item
  \(f''(x)>0\) en \((-\infty,-2) \cup (2,\infty)\), \(f''(x)<0\) en
  \((-2,2)\),
\item
  \(\displaystyle \lim_{x\to -\infty} f(x)=-2\),
  \(\displaystyle \lim_{x\to \infty} f(x)=0\).
\end{itemize}

\begin{tcolorbox}[enhanced jigsaw, opacityback=0, left=2mm, opacitybacktitle=0.6, toprule=.15mm, bottomrule=.15mm, colbacktitle=quarto-callout-tip-color!10!white, titlerule=0mm, bottomtitle=1mm, title=\textcolor{quarto-callout-tip-color}{\faLightbulb}\hspace{0.5em}{Solución}, toptitle=1mm, colframe=quarto-callout-tip-color-frame, colback=white, breakable, rightrule=.15mm, arc=.35mm, leftrule=.75mm, coltitle=black]

En virtud de la \hyperref[teo-prueba-cd]{prueba creciente/decreciente},
la primera condición establece que \(f\) es creciente en \((-\infty,1)\)
y decreciente en \((1,\infty)\).

De la segunda condición, aplicando la
\hyperref[teo-prueba-concavidad]{prueba de la concavidad} concluimos que
la gráfica de \(f\) debe ser cóncava hacia arriba en
\((-\infty,-2)\cup (2,\infty)\) y cóncava hacia abajo en \((-2,2)\).

Por último, la tercera condición establece que \(y=-2\) e \(y=0\) son
las ecuaciones de las asíntotas horizontales de \(f\).

Con esta información trazamos un posible gráfico aproximado.

\includegraphics[width=0.4\linewidth,height=\textheight,keepaspectratio]{index_files/mediabag/imagenes_cap4/4.3-ejemplo-funcion.pdf}

\end{tcolorbox}

Como otra aplicación de la derivada segunda, podemos dar una forma
alternativa de determinar la naturaleza de un punto crítico.

\phantomsection\label{teo-prueba-segundader}
Prueba de la segunda derivada

Supongamos que \(f''\) es continua cerca de \(c\).

\begin{itemize}
\item
  Si \(f'(c)=0\) y \(f''(c)>0\), entonces \(f\) tiene un mínimo relativo
  en \(x=c\).
\item
  Si \(f'(c)=0\) y \(f''(c)<0\), entonces \(f\) tiene un máximo relativo
  en \(x=c\).
\end{itemize}

Ejemplo

Analizar la curva \(y=x^4-4x^3\) con respecto a la concavidad, puntos de
inflexión y extremos locales. Dibujar la curva.

\begin{tcolorbox}[enhanced jigsaw, opacityback=0, left=2mm, opacitybacktitle=0.6, toprule=.15mm, bottomrule=.15mm, colbacktitle=quarto-callout-tip-color!10!white, titlerule=0mm, bottomtitle=1mm, title=\textcolor{quarto-callout-tip-color}{\faLightbulb}\hspace{0.5em}{Solución}, toptitle=1mm, colframe=quarto-callout-tip-color-frame, colback=white, breakable, rightrule=.15mm, arc=.35mm, leftrule=.75mm, coltitle=black]

Primero calculamos la derivada primera \(y'(x)=4x^3-12x^2=4x^2(x-3)\).
Entonces \(y'(x)=0\) si y sólo si \(x=0\) o \(x=3\). Éstos son los
números críticos de la función.

Ahora calculamos la derivada segunda \(y''(x)=12x^2-24x=12x(x-2)\).
Observemos ahora que

\[
f''(3)=36>0,
\]

con lo que \(y\) presenta un mínimo local en \(x=3\) cuyo valor es
\(y(3)=-27\), en virtud de la \hyperref[teo-prueba-segundader]{prueba de
la segunda derivada}. Sin embargo,

\[
y''(0)=0
\]

con lo cual no podemos utilizar este mismo argumento para conocer la
naturaleza del punto \(x=0\). Como \(y'\) no cambia de signo en \(x=0\),
en este punto no tenemos extremo.

Para determinar la concavidad y puntos de inflexión hacemos la siguiente
tabla (recordar la \hyperref[teo-propiedad-signoconst]{propiedad del
signo constante}, que aplicamos a \(y''\)) y luego la
\hyperref[teo-prueba-concavidad]{prueba de concavidad}.

\begin{longtable}[]{@{}llll@{}}
\toprule\noalign{}
Intervalo & \(x_0\) & \(y''(x_0)\) & \(y\) \\
\midrule\noalign{}
\endhead
\bottomrule\noalign{}
\endlastfoot
\((-\infty,0)\) & \(-1\) & \(+\) & convexa \\
\((0,2)\) & 1 & \(-\) & cóncava \\
\((2,+\infty)\) & 3 & \(+\) & convexa \\
\end{longtable}

Luego la curva \(y=x^4-4x^3\) es convexa en los intervalos
\((-\infty,0)\) y \((2,+\infty)\), y cóncava en \((0,2)\).

Como \(y\) cambia de convexidad en \(x=0\) y en \(x=2\), y además
resulta continua allí, los puntos \(P_1=(0,y(0))=(0,0)\) y
\(P_2=(2,y(2))=(2,-16)\) son puntos de inflexión de \(y\).

\includegraphics[width=0.4\linewidth,height=\textheight,keepaspectratio]{index_files/mediabag/imagenes_cap4/4.3-ejemplo-funcion2.pdf}

\end{tcolorbox}

Ejemplo

Trazar la gráfica de \(f(x)=x^{2/3}(6-x)^{1/3}\), sabiendo que

\[
f'(x)=\frac{4-x}{x^{1/3}(6-x)^{2/3}} \quad\text{ y }\quad f''(x)=-\frac{8}{x^{4/3}(6-x)^{5/3}}.
\]

\begin{tcolorbox}[enhanced jigsaw, opacityback=0, left=2mm, opacitybacktitle=0.6, toprule=.15mm, bottomrule=.15mm, colbacktitle=quarto-callout-tip-color!10!white, titlerule=0mm, bottomtitle=1mm, title=\textcolor{quarto-callout-tip-color}{\faLightbulb}\hspace{0.5em}{Solución}, toptitle=1mm, colframe=quarto-callout-tip-color-frame, colback=white, breakable, rightrule=.15mm, arc=.35mm, leftrule=.75mm, coltitle=black]

Podemos ver que \(f'(x)=0\) si y sólo si \(x=4\). Además, \(f'\) no
existe cuando \(x=0\) o cuando \(x=6\). Entonces \(0\), \(4\) y \(6\)
son los tres números críticos de \(f\). Utilizando la
\hyperref[teo-propiedad-signoconst]{propiedad del signo constante} y
luego la\\
\hyperref[teo-prueba-primerader]{prueba de la primera derivada} armamos
la siguiente tabla.

\begin{longtable}[]{@{}llll@{}}
\toprule\noalign{}
Intervalo & \(x_0\) & \(f'(x_0)\) & \(f\) \\
\midrule\noalign{}
\endhead
\bottomrule\noalign{}
\endlastfoot
\((-\infty,0)\) & \(-1\) & \(-\) & decrece \\
\((0,4)\) & 1 & \(+\) & crece \\
\((4,6)\) & 5 & \(-\) & decrece \\
\((6,+\infty)\) & 7 & \(-\) & decrece \\
\end{longtable}

Por otro lado, observando \(f''\) vemos que nunca se anula, y no existe
cuando \(x=0\) o \(x=6\). La siguiente tabla nos da el signo de \(f''\)
y la convexidad, utilizando la \hyperref[teo-prueba-concavidad]{prueba
de la concavidad}.

\begin{longtable}[]{@{}llll@{}}
\toprule\noalign{}
Intervalo & \(x_0\) & \(f''(x_0)\) & \(f\) \\
\midrule\noalign{}
\endhead
\bottomrule\noalign{}
\endlastfoot
\((-\infty,0)\) & \(-1\) & \(-\) & cóncava \\
\((0,6)\) & \(1\) & \(-\) & cóncava \\
\((6,+\infty)\) & \(7\) & \(7\) & convexa \\
\end{longtable}

La curva cambia de convexidad en \(x=6\) y resulta continua allí, con lo
cual \(P=(6,f(6))=(6,0)\) es el único punto de inflexión de la curva. El
gráfico se muestra a continuación.

\includegraphics[width=0.4\linewidth,height=\textheight,keepaspectratio]{index_files/mediabag/imagenes_cap4/4.3-ejemplo-funcion3.pdf}

\end{tcolorbox}

\hyperref[seccion_4.3]{↑ Volver al inicio de la sección}

\section{4.4 Formas indeterminadas y la regla de
L'Hospital}\label{seccion_4.4}

Si tenemos un límite de la forma \[
\lim_{x\to a} \frac{f(x)}{g(x)},
\]

donde \(f(x)\to 0\) y \(g(x)\to 0\) cuando \(x\to a\), éste puede
existir o no, y se conoce como \textbf{forma indeterminada de tipo
\(\displaystyle \frac{0}{0}\)}.

Si en cambio \(f(x)\to \pm\infty\) y \(g(x)\to \pm\infty\) el límite del
cociente puede existir o no, y se conoce como \textbf{forma
indeterminada de tipo \(\displaystyle \frac{\infty}{\infty}\)}.

En esta sección veremos una forma de calcular estos límites utilizando
derivadas.

\phantomsection\label{teo-regla-LH}
Regla de L'Hospital

Sean \(f\) y \(g\) dos funciones derivables, de manera que
\(g'(x)\neq 0\) en un intervalo abierto \(I\) que contiene a \(a\)
(excepto quizás en \(a\)). Supongamos además que

\[
\lim_{x\to a}f(x)=0\quad \text{ y }\quad \lim_{x\to a}g(x)=0,
\]

o bien que

\[
\lim_{x\to a}f(x)=\pm \infty\quad \text{ y }\quad \lim_{x\to a}g(x)=\pm\infty.
\]

En otras palabras, tenemos una forma indeterminada del tipo \(0/0\) o
\(\infty/\infty\). Entonces

\[
\lim_{x\to a}\frac{f(x)}{g(x)}=\lim_{x\to a}\frac{f'(x)}{g'(x)},
\]

si este último límite existe (o si es \(\infty\) o \(-\infty\)).

Esta regla también es válida para límites laterales y límites en el
infinito. Es decir, sigue siendo válida si reemplazamos \(x\to a\) por
\(x\to a^+\), \(x\to a^-\), \(x\to\infty\) o \(x\to -\infty\).

Ejemplo

Encontrar \(\displaystyle \lim_{x\to 1} \frac{\ln x}{x-1}\).

\begin{tcolorbox}[enhanced jigsaw, opacityback=0, left=2mm, opacitybacktitle=0.6, toprule=.15mm, bottomrule=.15mm, colbacktitle=quarto-callout-tip-color!10!white, titlerule=0mm, bottomtitle=1mm, title=\textcolor{quarto-callout-tip-color}{\faLightbulb}\hspace{0.5em}{Solución}, toptitle=1mm, colframe=quarto-callout-tip-color-frame, colback=white, breakable, rightrule=.15mm, arc=.35mm, leftrule=.75mm, coltitle=black]

Si escribimos \(f(x)=\ln x\) y \(g(x)=x-1\), entonces tenemos que

\[
\lim_{x\to 1} f(x)=0 \quad \text{ y } \quad \lim_{x\to 1} g(x)=1.
\]

Además, \(g'(x)=1\neq 0\) para todo \(x\). Estudiamos el límite

\[
\lim_{x\to 1} \frac{f'(x)}{g'(x)}=\lim_{x\to 1} \frac{\frac{1}{x}}{1}=\lim_{x\to 1} \frac{1}{x}=1.
\]

Como el límite de \(f'/g'\) existe, la \hyperref[teo-regla-LH]{regla de
L'Hospital} establece que

\[
\lim_{x\to 1} \frac{f(x)}{g(x)}=\lim_{x\to 1} \frac{f'(x)}{g'(x)}=1.
\]

\end{tcolorbox}

Ejemplo

Calcular \(\displaystyle \lim_{x\to \infty} \frac{e^x}{x^2}\).

\begin{tcolorbox}[enhanced jigsaw, opacityback=0, left=2mm, opacitybacktitle=0.6, toprule=.15mm, bottomrule=.15mm, colbacktitle=quarto-callout-tip-color!10!white, titlerule=0mm, bottomtitle=1mm, title=\textcolor{quarto-callout-tip-color}{\faLightbulb}\hspace{0.5em}{Solución}, toptitle=1mm, colframe=quarto-callout-tip-color-frame, colback=white, breakable, rightrule=.15mm, arc=.35mm, leftrule=.75mm, coltitle=black]

Escribiendo \(f(x)=e^x\) y \(g(x)=x^2\), es inmediato que

\[
\lim_{x\to \infty} f(x)=\infty \quad \text{ y } \quad \lim_{x\to \infty} g(x)=\infty,
\]

por lo que se trata de una indeterminación del tipo \(\infty/\infty\).
Calculamos

\[
\lim_{x\to \infty} \frac{f'(x)}{g'(x)}=\lim_{x\to \infty} \frac{e^x}{2x},
\]

que vuelve a resultar una indeterminación del tipo \(\infty/\infty\).
Repitiendo el proceso una vez más

\[
\lim_{x\to \infty} \frac{f''(x)}{g''(x)}=\lim_{x\to \infty} \frac{e^x}{2}=\infty.
\]

Utilizando la \hyperref[teo-regla-LH]{regla de L'Hospital} concluimos
que

\[
\lim_{x\to \infty} \frac{f(x)}{g(x)}=\lim_{x\to \infty} \frac{f'(x)}{g'(x)}=\lim_{x\to \infty} \frac{f''(x)}{g''(x)}=\infty.
\]

\end{tcolorbox}

Ejemplo

Calcular \(\displaystyle \lim_{x\to \infty} \frac{\ln x}{\sqrt[3]{x}}\).

\begin{tcolorbox}[enhanced jigsaw, opacityback=0, left=2mm, opacitybacktitle=0.6, toprule=.15mm, bottomrule=.15mm, colbacktitle=quarto-callout-tip-color!10!white, titlerule=0mm, bottomtitle=1mm, title=\textcolor{quarto-callout-tip-color}{\faLightbulb}\hspace{0.5em}{Solución}, toptitle=1mm, colframe=quarto-callout-tip-color-frame, colback=white, breakable, rightrule=.15mm, arc=.35mm, leftrule=.75mm, coltitle=black]

Sean \(f(x)=\ln x\) y \(g(x)=\sqrt[3]{x}\). Dado que

\[
\lim_{x\to \infty} f(x)=\infty \quad \text{ y } \quad \lim_{x\to \infty} g(x)=\infty,
\]

tenemos una indeterminación de tipo \(\infty/\infty\). Calculamos

\[
\lim_{x\to \infty} \frac{f'(x)}{g'(x)}=\lim_{x\to \infty} \frac{\frac{1}{x}}{\frac{1}{3x^{2/3}}}=\lim_{x\to \infty} \frac{3x^{2/3}}{x}=\lim_{x\to \infty} \frac{3}{x^{1/3}}=0.
\]

Como el límite de \(f'/g'\) existe, por la \hyperref[teo-regla-LH]{regla
de L'Hospital} concluimos que

\[
\lim_{x\to \infty} \frac{f(x)}{g(x)}=\lim_{x\to \infty} \frac{f'(x)}{g'(x)}=0.
\]

\end{tcolorbox}

Ejemplo

Encontrar \(\displaystyle \lim_{x\to 0} \frac{\tan x-x}{x^3}\).

\begin{tcolorbox}[enhanced jigsaw, opacityback=0, left=2mm, opacitybacktitle=0.6, toprule=.15mm, bottomrule=.15mm, colbacktitle=quarto-callout-tip-color!10!white, titlerule=0mm, bottomtitle=1mm, title=\textcolor{quarto-callout-tip-color}{\faLightbulb}\hspace{0.5em}{Solución}, toptitle=1mm, colframe=quarto-callout-tip-color-frame, colback=white, breakable, rightrule=.15mm, arc=.35mm, leftrule=.75mm, coltitle=black]

Escribimos \(f(x)=\tan x-x\) y \(g(x)=x^3\). Es fácil comprobar que

\[
\lim_{x\to 0} f(x)=0 \quad \text{ y que } \quad \lim_{x\to 0} g(x)=0,
\]

con lo cual tenemos una indeterminación del tipo \(0/0\). Además,
\(g'(x)=3x^2\) con lo que \(g'(x)\neq 0\) para todo \(x\neq 0\).
Calculamos el límite del cociente de derivadas

\[
\lim_{x\to 0} \frac{f'(x)}{g'(x)}=\lim_{x\to 0} \frac{\sec^2 x-1}{3x^2},
\]

que resulta nuevamente indeterminado del tipo \(0/0\). Si calculamos el
límite del cociente de las derivadas segundas

\[
\lim_{x\to 0} \frac{f''(x)}{g''(x)}=\lim_{x\to 0} \frac{2\sec x\,\tan x}{6x}=\lim_{x\to 0} \frac{1}{3}\frac{\operatorname{sen} x}{x}\frac{1}{\cos^2 x}=\frac{1}{3},
\] en virtud de la \href{-@eq-limite-seno}{fórmula} obtenida en la
\hyperref[seccion_3.3]{Sección 3.3}.

Finalmente, por la \hyperref[teo-regla-LH]{regla de L'Hospital} tenemos

\[
\lim_{x\to 0} \frac{f(x)}{g(x)}=\lim_{x\to 0} \frac{f'(x)}{g'(x)}=\lim_{x\to 0} \frac{f''(x)}{g''(x)}=\frac{1}{3}.
\]

\end{tcolorbox}

Ejemplo

Hallar
\(\displaystyle \lim_{x\to \pi^-} \frac{\operatorname{sen} x}{1-\cos x}\).

\begin{tcolorbox}[enhanced jigsaw, opacityback=0, left=2mm, opacitybacktitle=0.6, toprule=.15mm, bottomrule=.15mm, colbacktitle=quarto-callout-tip-color!10!white, titlerule=0mm, bottomtitle=1mm, title=\textcolor{quarto-callout-tip-color}{\faLightbulb}\hspace{0.5em}{Solución}, toptitle=1mm, colframe=quarto-callout-tip-color-frame, colback=white, breakable, rightrule=.15mm, arc=.35mm, leftrule=.75mm, coltitle=black]

Notar que la función dada es continua en \(x=\pi\), por lo que una
sustitución directa nos da

\[
 \lim_{x\to \pi^-} \frac{\operatorname{sen} x}{1-\cos x}=\frac{\operatorname{sen} \pi}{1-\cos \pi}=\frac{0}{2}=0.
\]

\end{tcolorbox}

\begin{tcolorbox}[enhanced jigsaw, opacityback=0, left=2mm, opacitybacktitle=0.6, toprule=.15mm, bottomrule=.15mm, colbacktitle=quarto-callout-caution-color!10!white, titlerule=0mm, bottomtitle=1mm, title=\textcolor{quarto-callout-caution-color}{\faFire}\hspace{0.5em}{Importante}, toptitle=1mm, colframe=quarto-callout-caution-color-frame, colback=white, breakable, rightrule=.15mm, arc=.35mm, leftrule=.75mm, coltitle=black]

Si intentamos aplicar la regla de L'Hospital sin advertir que no se
trata de una indeterminación del tipo 0/0, obtendríamos

\[
\lim_{x\to \pi^-} \frac{\operatorname{sen} x}{1-\cos x}=\lim_{x\to \pi^-}\frac{\cos x}{\operatorname{sen} x}=+\infty,
\]

lo que nos da una respuesta equivocada.

\end{tcolorbox}

\subsection{Productos indeterminados}\label{productos-indeterminados}

Si tenemos \(\displaystyle \lim_{x\to a}f(x)=0\) y
\(\lim_{x\to a}g(x)=\pm \infty\), entonces el límite

\[
\lim_{x\to a}f(x)g(x)
\]

se llama \textbf{forma indeterminada de tipo \(0\cdot \infty\)}. Para
resolverlo, se convierte el producto en un cociente de la siguiente
manera

\[f\cdot g =\frac{f}{1/g}\quad \text{ o bien} \quad f\cdot g=\frac{g}{1/f},\]

y de esta forma se lo lleva a una indeterminación del tipo \(0/0\) o
\(\infty/\infty\).

Ejemplo

Evaluar \(\lim_{x\to 0^+}x\ln x\).

\begin{tcolorbox}[enhanced jigsaw, opacityback=0, left=2mm, opacitybacktitle=0.6, toprule=.15mm, bottomrule=.15mm, colbacktitle=quarto-callout-tip-color!10!white, titlerule=0mm, bottomtitle=1mm, title=\textcolor{quarto-callout-tip-color}{\faLightbulb}\hspace{0.5em}{Solución}, toptitle=1mm, colframe=quarto-callout-tip-color-frame, colback=white, breakable, rightrule=.15mm, arc=.35mm, leftrule=.75mm, coltitle=black]

Si \(f(x)=x\) y \(g(x)=\ln x\), entonces es claro que

\[
\lim_{x\to 0^+} f(x)=0 \quad \text{ y } \quad \lim_{x\to 0^+} g(x)=-\infty, 
\]

por lo que se trata de un producto indeterminado. Si reescribimos el
producto como un cociente obtenemos

\[
\lim_{x\to 0^+} \frac{\ln x}{\frac{1}{x}},
\]

que es una indeterminación de tipo \(\infty/\infty\). Planteamos
entonces

\[
\lim_{x\to 0^+} \frac{\frac{d}{dx}(\ln x)}{\frac{d}{dx}\left(\frac{1}{x}\right)}=\lim_{x\to 0^+} \frac{\frac{1}{x}}{-\frac{1}{x^2}}=\lim_{x\to 0^+} -\frac{x^2}{x}=0.
\]

Como éste último límite existe, podemos concluir por la
\hyperref[teo-regla-LH]{regla de L'Hospital} que

\[
\lim_{x\to 0^+}x\ln x = 0.
\]

\end{tcolorbox}

\subsection{Diferencias
indeterminadas}\label{diferencias-indeterminadas}

Si \(\displaystyle \lim_{x\to a}f(x)=\infty\) y
\(\lim_{x\to a}g(x)= \infty\), entonces el límite

\[
\lim_{x\to a}\left(f(x)-g(x)\right)
\]

se llama \textbf{forma indeterminada de tipo \(\infty - \infty\)}.

En este caso se intenta convertir la diferencia en un cociente (por
ejemplo, buscando un denominador común o haciendo una racionalización)
para obtener una indeterminación del tipo \(0/0\) o \(\infty/\infty\).

Ejemplo

Calcular \(\lim_{x\to (\pi/2)^-}\left(\sec x-\tan x\right)\).

\begin{tcolorbox}[enhanced jigsaw, opacityback=0, left=2mm, opacitybacktitle=0.6, toprule=.15mm, bottomrule=.15mm, colbacktitle=quarto-callout-tip-color!10!white, titlerule=0mm, bottomtitle=1mm, title=\textcolor{quarto-callout-tip-color}{\faLightbulb}\hspace{0.5em}{Solución}, toptitle=1mm, colframe=quarto-callout-tip-color-frame, colback=white, breakable, rightrule=.15mm, arc=.35mm, leftrule=.75mm, coltitle=black]

Tanto el minuendo como el sustraendo tienen límite \(\infty\), con lo
cual reescribimos la expresión como

\[
\sec x-\tan x=\frac{1}{\cos x}-\frac{\operatorname{sen} x}{\cos x}=\frac{1-\operatorname{sen} x}{\cos x}.
\]

Llamando \(f(x)=1-\operatorname{sen} x\) y \(g(x)=\cos x\) tenemos que

\[
\lim_{x\to (\pi/2)^-} f(x)=0 \quad \text{ y } \quad  \lim_{x\to (\pi/2)^-}g(x)= 0,
\]

por lo que resulta una indeterminación del tipo \(0/0\). Notemos que
\(g'(x)=-\operatorname{sen} x\), y es no nula cerca de \(\pi/2\).
Calculamos entonces

\[
\lim_{x\to (\pi/2)^-}\frac{f'(x)}{g'(x)}=\lim_{x\to (\pi/2)^-}\frac{-\cos x}{-\operatorname{sen} x}=0,
\]

y por la \hyperref[teo-regla-LH]{regla de L'Hospital} concluimos que

\[
\lim_{x\to (\pi/2)^-}\frac{f(x)}{g(x)}=\lim_{x\to (\pi/2)^-}\frac{f'(x)}{g'(x)}=0.
\]

\end{tcolorbox}

\subsection{Potencias indeterminadas}\label{potencias-indeterminadas}

Consideremos el límite \(\displaystyle \lim_{x\to a} f(x)^{g(x)}\).

\begin{itemize}
\tightlist
\item
  Cuando
\end{itemize}

\[
\lim_{x\to a}f(x)=0\quad\text{ y }\quad \lim_{x\to a}g(x)=0
\]

se llama \textbf{forma indeterminada de tipo \(0^0\)}.

\begin{itemize}
\tightlist
\item
  Si ocurre que
\end{itemize}

\[
\lim_{x\to a}f(x)=\infty\quad\text{ y }\quad \lim_{x\to a}g(x)=0,
\]

se dice \textbf{forma indeterminada de tipo \(\infty^0\)}.

\begin{itemize}
\tightlist
\item
  Si
\end{itemize}

\[
\lim_{x\to a}f(x)=1\quad\text{ y }\quad \lim_{x\to a}g(x)=\pm \infty
\]

el límite se conoce como \textbf{forma indeterminada de tipo
\(1^\infty\)}.

Para resolver este tipo de indeterminaciones, podemos utilizar el
logaritmo natural. Si \[
y=f(x)^{g(x)}
\]

entonces

\[
\ln y=g(x)\ln(f(x)).
\]

O bien, si usamos la exponencial

\[
f(x)^{g(x)}=e^{g(x)\ln(f(x))}.
\]

En ambos casos se obtiene el producto indeterminado \(g(x)\ln(f(x))\),
que es de la forma \(0\cdot \infty\).

Ejemplo

Calcular \(\lim_{x\to 0^+}(1+\operatorname{sen}(4x))^{\cot x}\).

\begin{tcolorbox}[enhanced jigsaw, opacityback=0, left=2mm, opacitybacktitle=0.6, toprule=.15mm, bottomrule=.15mm, colbacktitle=quarto-callout-tip-color!10!white, titlerule=0mm, bottomtitle=1mm, title=\textcolor{quarto-callout-tip-color}{\faLightbulb}\hspace{0.5em}{Solución}, toptitle=1mm, colframe=quarto-callout-tip-color-frame, colback=white, breakable, rightrule=.15mm, arc=.35mm, leftrule=.75mm, coltitle=black]

Observemos que la función en la base tiene límite \(1\) y la del
exponente \(\infty\), por lo que tenemos una indeterminación del tipo
\(1^\infty\).

Si llamamos \(F(x)=(1+\operatorname{sen}(4x))^{\cot x}\), aplicando
logaritmo natural resulta

\[
\ln F(x) =\ln\left((1+\operatorname{sen}(4x))^{\cot x}\right)=\cot x \ln((1+\operatorname{sen}(4x)))=\frac{\ln((1+\operatorname{sen}(4x)))}{\tan x}=\frac{f(x)}{g(x)}.
\]

Dado que ahora tanto \(f\) como \(g\) tienen límite \(0\) para
\(x\to 0^+\), y como \((\tan x)'=\sec^2 x\) es no nula cerca de \(0\),
calculamos

\[
\lim_{x\to 0^+}\frac{f'(x)}{g'(x)}=\lim_{x\to 0^+}\frac{\frac{4\cos(4x)}{1+\operatorname{sen}(4x)}}{\sec^2x}=\lim_{x\to 0^+} \frac{4\cos(4x)\cos^2 x}{1+\operatorname{sen}(4x)}=4.
\]

Utilizando la \hyperref[teo-regla-LH]{regla de L'Hospital}, concluimos
que

\[
\lim_{x\to 0^+}\ln F(x)=\lim_{x\to 0^+}\frac{f(x)}{g(x)}=\lim_{x\to 0^+}\frac{f'(x)}{g'(x)}=4.
\]

Finalmente

\[
\lim_{x\to 0^+} F(x)=\lim_{x\to 0^+} e^{\ln F(x)}=e^{\lim_{x\to 0^+} \ln F(x)}=e^4,
\]

dado que \(y=e^x\) es una función continua.

\end{tcolorbox}

Ejemplo

Encontrar \(\lim_{x\to 0^+}x^x\).

\begin{tcolorbox}[enhanced jigsaw, opacityback=0, left=2mm, opacitybacktitle=0.6, toprule=.15mm, bottomrule=.15mm, colbacktitle=quarto-callout-tip-color!10!white, titlerule=0mm, bottomtitle=1mm, title=\textcolor{quarto-callout-tip-color}{\faLightbulb}\hspace{0.5em}{Solución}, toptitle=1mm, colframe=quarto-callout-tip-color-frame, colback=white, breakable, rightrule=.15mm, arc=.35mm, leftrule=.75mm, coltitle=black]

En este caso tenemos una forma indeterminada del tipo \(0^0\). Si
\(F(x)=x^x\), aplicamos logaritmo natural como antes

\[
\ln F(x)=\ln x^x=x\ln x.
\]

En consecuencia,

\[
\lim_{x\to 0^+} \ln F(x)=\lim_{x\to 0^+} x\ln x=0,
\]

según vimos en el Ejemplo 21. Por lo tanto

\[
\lim_{x\to 0^+} F(x)=\lim_{x\to 0^+} e^{\ln F(x)}=e^{\lim_{x\to 0^+} \ln F(x)}=1,
\]

ya que \(y=e^x\) es una función continua.

\end{tcolorbox}

\hyperref[seccion_4.4]{↑ Volver al inicio de la sección}

\section{4.5. Resumen de trazo de curvas}\label{seccion_4.5}

En esta sección reuniremos todo lo aprendido hasta el momento para
realizar el estudio completo de una función dada. Para ello seguiremos
una serie de pasos que nos ayudarán a visualizar el gráfico de la
función, además de obtener los elementos más distintivos de la misma.

Pasos a seguir en el estudio de una función

\textbf{(A) Dominio} Recordemos que el dominio \(D_f\) de \(f\) consiste
en todos los valores de \(x\) para los cuales \(f\) está definida.

\textbf{(B) Intersecciones} La gráfica de \(y=f(x)\) corta al eje \(y\)
en \(f(0)\). Por otra parte, los valores de \(x\) donde el gráfico corta
al eje de abscisas se encuentran resolviendo la ecuación \(f(x)=0\).
(Este paso puede omitirse si la ecuación es difícil de resolver).

\textbf{(C) Simetría y periodicidad} En este paso podemos investigar si
la función resulta par o impar.

Recordemos que \(f\) es par si se cumple que

\[
f(-x)=f(x) \quad \text{ para todo }\quad x\in D_f.
\]

Por otro lado, \(f\) es impar si se cumple que

\[
f(-x)=-f(x) \quad \text{ para todo }\quad x\in D_f.
\]

Por otro lado, para estudiar la periodicidad debemos analizar si existe
un número \(p>0\) tal que

\[
f(x+p)=f(x) \quad \text{ para todo }\quad x\in D_f.
\]

Si esto pasa, entonces \(f\) es periódica. El número \(p\) más pequeño
que satisface esta igualdad se llama \textbf{período} de \(f\).

\textbf{(D) Asíntotas} La recta \(x=a\) es una asíntota vertical de
\(f\) si se cumple al menos una de las siguientes condiciones

\begin{itemize}
\tightlist
\item
  \(\displaystyle \lim_{x\to a^+}f(x)=\infty\),
\item
  \(\displaystyle \lim_{x\to a^+}f(x)=-\infty\),
\end{itemize}

\begin{itemize}
\tightlist
\item
  \(\displaystyle \lim_{x\to a^-}f(x)=\infty\),
\item
  \(\displaystyle \lim_{x\to a^-}f(x)=-\infty\).
\end{itemize}

La recta \(y=L\) es una asíntota horizontal de la curva \(y=f(x)\) si se
cumple que

\[
\lim_{x\to \infty}f(x)=L\quad\text{ o bien }\quad \lim_{x\to -\infty}f(x)=L.
\]

\textbf{(E) Intervalos de crecimiento y decrecimiento} Calculamos
\(f'(x)\) y determinamos en qué intervalos es positiva y en cuáles
negativa. La \hyperref[teo-prueba-cd]{prueba creciente/decreciente} nos
dará los intervalos donde \(f\) crece y donde decrece.

\textbf{(F) Valores extremos locales} Calculamos los números críticos de
\(f\), es decir, aquellos valores \(c\) para los que \(f'(c)=0\) o
\(f'(c)\) no existe. Para clasificarlos podemos utilizar la
\hyperref[teo-prueba-primerader]{prueba de la primera derivada} o la de
la \hyperref[teo-prueba-segundader]{segunda derivada}.

\textbf{(G) Concavidad y puntos de inflexión} Calculamos \(f''(x)\) y
aplicamos la \hyperref[teo-prueba-concavidad]{prueba de concavidad}. La
función resultará cóncava hacia arriba donde \(f''(x)>0\) y cóncava
hacia abajo donde \(f''(x)<0\). Los puntos de inflexión serán aquellos
donde \(f\) sea continua y cambie la concavidad de la curva.

\textbf{(H) Trazar la gráfica} Con toda la información obtenida hasta
ahora podemos hacer un esbozo aproximado de la curva, indicando los
elementos más distintivos.

Ejemplo

Graficar la función \(\displaystyle f(x)=\frac{2x^2}{x^2-1}\).

\begin{tcolorbox}[enhanced jigsaw, opacityback=0, left=2mm, opacitybacktitle=0.6, toprule=.15mm, bottomrule=.15mm, colbacktitle=quarto-callout-tip-color!10!white, titlerule=0mm, bottomtitle=1mm, title=\textcolor{quarto-callout-tip-color}{\faLightbulb}\hspace{0.5em}{Solución}, toptitle=1mm, colframe=quarto-callout-tip-color-frame, colback=white, breakable, rightrule=.15mm, arc=.35mm, leftrule=.75mm, coltitle=black]

Seguiremos los pasos indicados en el cuadro anterior.

\textbf{(A) Dominio} Al ser una función racional, sabemos que el dominio
consiste en todos los valores de \(x\) tales que el denominador es no
nulo. Como \(x^2-1=0\) si y sólo si \(x=\pm 1\), tenemos que
\(D_f =\mathbb{R}\setminus \{-1,1\}\).

\textbf{(B) Intersecciones}

\begin{itemize}
\tightlist
\item
  Para saber dónde la gráfica corta al eje \(x\) planteamos
\end{itemize}

\[
f(x)=0, \quad \text{ es decir,} \quad \frac{2x^2}{x^2-1}=0,
\]

de donde resulta \(x=0\).

\begin{itemize}
\tightlist
\item
  Para conocer la intersección con el eje \(y\) evaluamos
\end{itemize}

\[
f(0)=\frac{2\cdot 0^2}{0^2-1}=0.
\]

Entonces la gráfica interseca al eje \(y\) en \(y=0\).

En conclusión, la única intersección de la curva con los ejes se da en
el origen de coordenadas.

\textbf{(C) Simetría y periodicidad}

Para estudiar la simetría calculamos \(f(-x)\)

\[
f(-x)=\frac{2(-x)^2}{(-x)^2-1}=\frac{2x^2}{x^2-1}=f(x).
\]

Como \(f(-x)=f(x)\) para todo \(x\) en el dominio, \(f\) es una función
par.

Además \(f\) no resulta periódica.

\textbf{(D) Asíntotas}

\begin{itemize}
\tightlist
\item
  Comencemos buscando las asíntotas verticales. Como \(f\) es racional,
  es continua en su dominio, por lo tanto los valores de \(x\)
  candidatos a dar asíntotas verticales son \(x=-1\) y \(x=1\). Vamos a
  comprobarlo
\end{itemize}

\[
\lim_{x\to -1^+} f(x)=\lim_{x\to -1^+} \frac{2x^2}{(x-1)(x+1)}=-\infty,
\]

dado que el denominador tiende a cero con valores negativos y el
numerador se aproxima a \(2>0\). Haciendo un análisis similar, también
podemos ver que

\[
\lim_{x\to -1^-} f(x)=\infty, \quad \lim_{x\to 1^-} f(x)=-\infty \quad \text{ y } \quad \lim_{x\to 1^+} f(x)=\infty.
\]

Con esto, comprobamos que \(x=-1\) y \(x=1\) son ambas asíntotas
verticales de \(f\).

\begin{itemize}
\tightlist
\item
  Para las asíntotas horizontales, calculamos
\end{itemize}

\[
\lim_{x\to\infty} f(x)=\lim_{x\to \infty} \frac{2x^2}{x^2-1}=\lim_{x\to \infty} \frac{2x^2}{x^2\left(1-\frac{1}{x^2}\right)}=\lim_{x\to \infty} \frac{2}{1-\frac{1}{x^2}}=2,
\]

y también

\[
\lim_{x\to-\infty} f(x)=\lim_{x\to -\infty} \frac{2x^2}{x^2-1}==\lim_{x\to -\infty} \frac{2}{1-\frac{1}{x^2}}=2.
\]

Esto nos dice que \(y=2\) es asíntota horizontal del gráfico de \(f\).

\textbf{(E) Intervalos de crecimiento y decrecimiento}

Calculamos \(f'\) con la regla del cociente

\[
f'(x)=\frac{4x(x^2-1)-2x^2(2x)}{(x^2-1)^2}=-\frac{4x}{(x^2-1)^2}.
\]

Observemos que \(f'(x)=0\) si y sólo si \(x=0\). También vemos que
\(f'\) no existe en \(x=\pm 1\), sin embargo éstos puntos no son
críticos pues no están en el dominio de \(f\). Por lo tanto \(f\) tiene
un único punto crítico, \(x=0\). Recordando la
\hyperref[teo-propiedad-signoconst]{propiedad del signo constante},
\(f'\) tendrá signo constante en los intervalos \((-\infty, -1)\),
\((-1,0)\), \((0,1)\) y \((1,+\infty)\). En la siguiente tabla
determinamos el signo de \(f'\) en cada uno de ellos.

\begin{longtable}[]{@{}llll@{}}
\toprule\noalign{}
Intervalo & \(x_0\) & \(f'(x_0)\) & \(f\) \\
\midrule\noalign{}
\endhead
\bottomrule\noalign{}
\endlastfoot
\((-\infty,-1)\) & \(-2\) & \(+\) & crece \\
\((-1,0)\) & -0.5 & \(+\) & crece \\
\((0,1)\) & 0.5 & \(-\) & decrece \\
\((1,+\infty)\) & 2 & \(-\) & decrece \\
\end{longtable}

Mediante la \hyperref[teo-prueba-cd]{prueba creciente/decreciente}
podemos concluir que \(f\) crece en \((-\infty,-1)\) y en \((-1,0)\) y
decrece en \((0,1)\) y en \((0,+\infty)\).

\textbf{(F) Valores extremos locales}

Observando la tabla anterior, \(f'\) cambia de positiva a negativa en
\(x=0\), por lo que \(f\) presenta un máximo local en este punto (en
virtud de la \hyperref[teo-prueba-primerader]{prueba de la primera
derivada}) de valor \(f(0)=0\).

\textbf{(G) Concavidad y puntos de inflexión}

Calculamos la derivada segunda de \(f\), nuevamente utilizando la regla
del cociente

\[
f''(x)=-\frac{4(x^2-1)^2-4x\cdot 2(x^2-1)2x}{(x^2-1)^4}=-\frac{4(x^2-1)-16x^2}{(x^2-1)^3}=\frac{12x^2+4}{(x^2-1)^3}.
\]

Observemos que \(f''\) nunca se anula, y es continua en
\((-\infty,-1)\), \((-1,1)\) y en \((1,+\infty)\). Entonces la
\hyperref[teo-propiedad-signoconst]{propiedad del signo constante} nos
dice que tendrá signo constante en cada uno de estos intervalos, los
cuales calculamos en la siguiente tabla.

\begin{longtable}[]{@{}llll@{}}
\toprule\noalign{}
Intervalo & \(x_0\) & \(f''(x_0)\) & \(f\) \\
\midrule\noalign{}
\endhead
\bottomrule\noalign{}
\endlastfoot
\((-\infty,-1)\) & \(-2\) & \(+\) & convexa \\
\((-1,1)\) & 0 & \(-\) & cóncava \\
\((1,+\infty)\) & 2 & \(+\) & convexa \\
\end{longtable}

Entonces \(f\) resulta convexa en los intervalos \((-\infty,-1)\) y en
\((1,+\infty)\), y cóncava en \((-1,1)\), en virtud de la
\hyperref[teo-prueba-concavidad]{prueba de la concavidad}. El signo de
\(f''\) cambia en \(x=-1\) y en \(x=1\), pero estos puntos no están en
el dominio de \(f\), por lo que la gráfica no presenta puntos de
inflexión.

\textbf{(H) Trazar la gráfica} Con todo lo calculado trazamos la gráfica
aproximada de \(f\).

\includegraphics[width=0.45\linewidth,height=\textheight,keepaspectratio]{index_files/mediabag/imagenes_cap4/4.5-funcion1.pdf}

\end{tcolorbox}

Ejemplo

Trazar la gráfica de \(\displaystyle f(x)=\frac{x^2}{\sqrt{x+1}}\).

\begin{tcolorbox}[enhanced jigsaw, opacityback=0, left=2mm, opacitybacktitle=0.6, toprule=.15mm, bottomrule=.15mm, colbacktitle=quarto-callout-tip-color!10!white, titlerule=0mm, bottomtitle=1mm, title=\textcolor{quarto-callout-tip-color}{\faLightbulb}\hspace{0.5em}{Solución}, toptitle=1mm, colframe=quarto-callout-tip-color-frame, colback=white, breakable, rightrule=.15mm, arc=.35mm, leftrule=.75mm, coltitle=black]

Procedemos siguiendo los mismos pasos que antes.

\textbf{(A) Dominio} Por un lado, \(\sqrt{x+1}\) está bien definido
siempre que \(x+1\geq 0\). Sin embargo, al aparecer en el denominador el
valor cero no está permitido. Luego \(x+1>0\), o bien \(x>-1\) con lo
que \(D_f =(-1,+\infty)\).

\textbf{(B) Intersecciones}

\begin{itemize}
\tightlist
\item
  Para encontrar intersecciones con el eje \(x\) planteamos
\end{itemize}

\[
f(x)=0, \quad \text{ es decir,} \quad \frac{x^2}{\sqrt{x+1}}=0,
\]

de donde resulta \(x=0\).

\begin{itemize}
\tightlist
\item
  Para conocer la intersección con el eje \(y\) evaluamos
\end{itemize}

\[
f(0)=\frac{0^2}{\sqrt{0+1}}=0.
\]

Entonces la gráfica corta al eje \(y\) en \(y=0\).

Al igual que en el ejemplo anterior, la única intersección de la curva
con los ejes ocurre en el origen.

\textbf{(C) Simetría y periodicidad}

Notemos que

\[
f(-x)=\frac{(-x)^2}{\sqrt{-x+1}}=\frac{x^2}{\sqrt{1-x}},
\]

que no coincide con \(f(x)\) ni con \(-f(x)\), por lo que \(f\) no es
par ni impar. También observamos que \(f\) no es periódica.

\textbf{(D) Asíntotas}

\begin{itemize}
\tightlist
\item
  Calculemos las asíntotas verticales. Como \(f\) es continua en todo su
  dominio por ser cociente de funciones continuas, el único punto
  candidato a originar una asíntota vertical es \(x=-1\). Como \(f\)
  sólamente está definida a la derecha de \(-1\), calculamos
\end{itemize}

\[
\lim_{x\to -1^+} f(x)=\lim_{x\to -1^+} \frac{x^2}{\sqrt{x+1}}=\infty,
\]

pues el denominador tiende a cero con valores positivos y el numerador
se aproxima a \(1>0\). En conclusión, \(x=-1\) es la única asíntota
vertical del gráfico de \(f\).

\begin{itemize}
\tightlist
\item
  Recordando nuevamente que \(f\) existe sólo para \(x>-1\), calculamos
  el límite
\end{itemize}

\[
\lim_{x\to\infty} f(x)=\lim_{x\to \infty} \frac{x^2}{\sqrt{x+1}}=\lim_{x\to \infty} \frac{x^2}{\sqrt{x}\sqrt{1+\frac{1}{x}}}=\lim_{x\to \infty} \frac{x^{3/2}}{\sqrt{1+\frac{1}{x}}}=\infty.
\]

En consecuencia, \(f\) no presenta asíntotas horizontales.

\textbf{(E) Intervalos de crecimiento y decrecimiento}

Calculamos \(f'\) con la regla del cociente

\[
f'(x)=\frac{2x\sqrt{x+1}-x^2\frac{1}{2\sqrt{x+1}}}{\left(\sqrt{x+1}\right)^2}=\frac{x(3x+4)}{2(x+1)^{3/2}}.
\]

Notemos que \(f'(x)=0\) cuando \(x=0\) y \(x-=4/3\). Como éste último
valor no está en el dominio de \(f\), el único punto crítico de \(f\) es
\(x=0\). Aplicando la \hyperref[teo-propiedad-signoconst]{propiedad del
signo constante}, \(f'\) tendrá signo constante en los intervalos
\((-1, 0)\) y \((0,+\infty)\). En la siguiente tabla determinamos el
signo de \(f'\) en cada uno de ellos.

\begin{longtable}[]{@{}llll@{}}
\toprule\noalign{}
Intervalo & \(x_0\) & \(f'(x_0)\) & \(f\) \\
\midrule\noalign{}
\endhead
\bottomrule\noalign{}
\endlastfoot
\((-1,0)\) & \(-0.5\) & \(-\) & decrece \\
\((0,+\infty)\) & 1 & \(+\) & crece \\
\end{longtable}

Mediante la \hyperref[teo-prueba-cd]{prueba creciente/decreciente}
podemos concluir que \(f\) crece en \((0,+\infty)\) y decrece en
\((-1,0)\).

\textbf{(F) Valores extremos locales}

En la tabla anterior podemos ver que \(f'\) cambia de negativa a
positiva en \(x=0\), por lo que \(f\) presenta un mínimo local en este
punto de valor \(f(0)=0\), en virtud de la
\hyperref[teo-prueba-primerader]{prueba de la primera derivada}.

\textbf{(G) Concavidad y puntos de inflexión}

Calculamos la derivada segunda de \(f\), nuevamente utilizando la regla
del cociente

\[
\begin{aligned}
f''(x)=-\frac{(6x+4)2(x+1)^{3/2}-(3x^2+4x)2\frac{3}{2}(x+1)^{1/2}}{4(x+1)^3}&=\frac{(12x+8)(x+1)-3(3x^2+4x)}{4(x+1)^{5/2}}\\
\\
&=\frac{3x^2+8x+8}{4(x+1)^{5/2}}.
\end{aligned}
\]

Dado que \(f''>0\) en \((-1,+\infty)\), \(f\) resulta convexa en todo su
dominio por la \hyperref[teo-prueba-concavidad]{prueba de la
concavidad}. De esta manera, el gráfico de \(f\) no presenta puntos de
inflexión.

\textbf{(H) Trazar la gráfica} Trazamos la gráfica aproximada de \(f\).

\includegraphics[width=0.45\linewidth,height=\textheight,keepaspectratio]{index_files/mediabag/imagenes_cap4/4.5-funcion2.pdf}

\end{tcolorbox}

Ejemplo

Graficar \(\displaystyle f(x)=xe^x\).

\begin{tcolorbox}[enhanced jigsaw, opacityback=0, left=2mm, opacitybacktitle=0.6, toprule=.15mm, bottomrule=.15mm, colbacktitle=quarto-callout-tip-color!10!white, titlerule=0mm, bottomtitle=1mm, title=\textcolor{quarto-callout-tip-color}{\faLightbulb}\hspace{0.5em}{Solución}, toptitle=1mm, colframe=quarto-callout-tip-color-frame, colback=white, breakable, rightrule=.15mm, arc=.35mm, leftrule=.75mm, coltitle=black]

Procedemos como en los ejemplos anteriores.

\textbf{(A) Dominio} La función es el producto de un polinomio por una
exponencial, con lo cual \(D_f=\mathbb{R}\).

\textbf{(B) Intersecciones}

\begin{itemize}
\tightlist
\item
  Para encontrar intersecciones con el eje \(x\) escribimos
\end{itemize}

\[
f(x)=0, \quad \text{ es decir,} \quad x\,e^x=0,
\]

y como \(e^x>0\) para todo \(x\), obtenemos que \(x=0\).

\begin{itemize}
\tightlist
\item
  Para calcular la intersección con el eje \(y\) evaluamos
\end{itemize}

\[
f(0)=0\,e^0=0.
\]

Entonces la gráfica interseca al eje \(y\) en \(y=0\).

En conclusión, el única corte con los ejes se da en el origen de
coordenadas.

\textbf{(C) Simetría y periodicidad}

Comenzamos calculando \(f(-x)\)

\[
f(-x)=(-x)e^{-x}=-xe^{-x},
\]

que no coincide con \(f(x)\) ni con \(f(-x)\), por lo que la función no
es par ni impar. Vemos que \(f\) tampoco resulta periódica.

\textbf{(D) Asíntotas}

\begin{itemize}
\item
  La gráfica no presenta asíntotas verticales, ya que se trata de una
  función continua en todo \(\mathbb{R}\).
\item
  Para las asíntotas horizontales, calculamos
\end{itemize}

\[
\lim_{x\to\infty} f(x)=\lim_{x\to \infty} xe^x=\infty,
\]

ya que se trata de un producto donde cada factor se hace arbitrariamente
grande. Por otro lado, llamando \(y=-x\)

\[
\lim_{x\to-\infty} f(x)=\lim_{x\to -\infty} xe^x=\lim_{y\to \infty} (-y)e^{-y}=\lim_{y\to \infty} -\frac{y}{e^y}=0,
\]

utilizando la \hyperref[teo-regla-LH]{regla de L'Hospital}. Esto nos
dice que \(y=0\) es asíntota horizontal del gráfico de \(f\).

\textbf{(E) Intervalos de crecimiento y decrecimiento}

Calculamos \(f'\) con la regla del producto

\[
f'(x)=e^x+xe^x=e^x(x+1),
\]

de donde podemos ver que \(f'(x)=0\) si y sólo si \(x=-1\). Por lo tanto
\(f\) tiene un único punto crítico, \(x=-1\). Utilizando la
\hyperref[teo-propiedad-signoconst]{propiedad del signo constante},
\(f'\) tendrá signo constante en los intervalos \((-\infty, -1)\) y en
\((-1,+\infty)\). En la siguiente tabla determinamos el signo de \(f'\)
en cada uno de ellos.

\begin{longtable}[]{@{}llll@{}}
\toprule\noalign{}
Intervalo & \(x_0\) & \(f'(x_0)\) & \(f\) \\
\midrule\noalign{}
\endhead
\bottomrule\noalign{}
\endlastfoot
\((-\infty,-1)\) & \(-2\) & \(-\) & decrece \\
\((-1,+\infty)\) & 0 & \(+\) & crece \\
\end{longtable}

Con la \hyperref[teo-prueba-cd]{prueba creciente/decreciente} podemos
concluir que \(f\) crece en \((-1,+\infty)\) y que decrece en
\((-\infty,-1)\).

\textbf{(F) Valores extremos locales}

En la tabla anterior vemos que \(f'\) cambia de negativa a positiva en
\(x=-1\), por lo que \(f\) presenta un mínimo local en este punto (por
la \hyperref[teo-prueba-primerader]{prueba de la primera derivada}) y su
valor es \(f(-1)=-1\,e^{-1}=-1/e\approx -0.37\).

\textbf{(G) Concavidad y puntos de inflexión}

Calculamos la derivada segunda de \(f\), nuevamente con la regla del
producto

\[
f''(x)=e^x(x+1)+e^x=e^x(x+2).
\]

Así, \(f''(x)=0\) si y sólo si \(x=-2\). Por la
\hyperref[teo-propiedad-signoconst]{propiedad del signo constante},
\(f''\) tendrá signo constante en \((-\infty,-2)\) y en
\((-2,+\infty)\). En la siguiente tabla calculamos estos signos.

\begin{longtable}[]{@{}llll@{}}
\toprule\noalign{}
Intervalo & \(x_0\) & \(f''(x_0)\) & \(f\) \\
\midrule\noalign{}
\endhead
\bottomrule\noalign{}
\endlastfoot
\((-\infty,-2)\) & \(-3\) & \(-\) & cóncava \\
\((-2,+\infty)\) & -1 & \(+\) & convexa \\
\end{longtable}

Entonces \(f\) resulta convexa en \((-2,+\infty)\) y cóncava en
\((-\infty,-2)\) como consecuencia de la
\hyperref[teo-prueba-concavidad]{prueba de la concavidad}. El signo de
\(f''\) cambia en \(x=-2\), que resulta un punto de continuidad de
\(f\), por lo que \(P=(-2,f(-2))=(-2,-2/e^2)\) es el único punto de
inflexión de la gráfica.

\textbf{(H) Trazar la gráfica} Con lo calculado anteriormente trazamos
la gráfica aproximada de \(f\).

\includegraphics[width=0.4\linewidth,height=\textheight,keepaspectratio]{index_files/mediabag/imagenes_cap4/4.5-funcion3.pdf}

\end{tcolorbox}

Ejemplo

Realizar el gráfico de la función
\(\displaystyle f(x)=\frac{\cos x}{2+\operatorname{sen} x}\).

\begin{tcolorbox}[enhanced jigsaw, opacityback=0, left=2mm, opacitybacktitle=0.6, toprule=.15mm, bottomrule=.15mm, colbacktitle=quarto-callout-tip-color!10!white, titlerule=0mm, bottomtitle=1mm, title=\textcolor{quarto-callout-tip-color}{\faLightbulb}\hspace{0.5em}{Solución}, toptitle=1mm, colframe=quarto-callout-tip-color-frame, colback=white, breakable, rightrule=.15mm, arc=.35mm, leftrule=.75mm, coltitle=black]

\textbf{(A) Dominio} Observemos que
\(-1\leq \operatorname{sen} x\leq 1\), con lo que
\(2+\operatorname{sen}x\geq 1\) para todo \(x\), es decir, el
denominador nunca se anula. En consecuencia, \(D_f=\mathbb{R}\).

\textbf{(B) Intersecciones}

\begin{itemize}
\tightlist
\item
  Para encontrar intersecciones con el eje \(x\) escribimos
\end{itemize}

\[
f(x)=0, \quad \text{ es decir,} \quad\frac{\cos x}{2+\operatorname{sen} x}=0,
\]

que se cumple cuando \(\cos x=0\), es decir, si \(x=(2n+1)\pi/2\),
\(n\in\mathbb{Z}\).

\begin{itemize}
\tightlist
\item
  Para calcular la intersección con el eje \(y\) evaluamos
\end{itemize}

\[
f(0)=\frac{\cos 0}{2+\operatorname{sen} 0}=\frac{1}{2}.
\]

Entonces la gráfica interseca al eje \(y\) en \(y=1/2\).

\textbf{(C) Simetría y periodicidad}

Comenzamos calculando \(f(-x)\)

\[
f(-x)=\frac{\cos(-x)}{2+\operatorname{sen}(-x)}=\frac{\cos x}{2-\operatorname{sen}x},
\]

ya que \(y=\cos x\) es una función par e \(y=\operatorname{sen} x\) es
impar. Esta expresión no coincide con \(f(x)\) ni con \(f(-x)\), por lo
que la función no es par ni impar.

Para la periodicidad, como las funciones trigonométricas son periódicas,
esperamos obtener una función periódica. Para determinar el período, si
\(p>0\) está fijo planteamos la condición \(f(x+p)=f(x)\) e intentaremos
determinar \(p\).

Esta condición resulta

\[
\frac{\cos(x+p)}{2+\operatorname{sen}(x+p)}=\frac{\cos x}{2+\operatorname{sen} x}.
\]

Utilizando las fórmulas para el seno y el coseno de una suma de ángulos

\[
\operatorname{sen}(\alpha+\beta)=\operatorname{sen} \alpha \cos \beta +\cos\alpha \operatorname{sen} \beta \quad \text{ y } \quad \cos(\alpha+\beta)=\cos\alpha\cos\beta-\operatorname{sen} \alpha\operatorname{sen} \beta
\]

obtenemos

\[
\frac{\cos x\cos p-\operatorname{sen} x\operatorname{sen} p}{2+\operatorname{sen} x\cos p+\cos x\operatorname{sen} p}=\frac{\cos x}{2+\operatorname{sen} x}.
\]

Esta igualdad implica que

\[
(\cos x\cos p-\operatorname{sen} x\operatorname{sen} p)(2+\operatorname{sen} x)=(2+\operatorname{sen} x\cos p+\cos x\operatorname{sen} p)\cos x
\]

o de manera equivalente

\[
2\cos x\cos p+\cos x\operatorname{sen} x\cos p-2\operatorname{sen} x\operatorname{sen} p-\operatorname{sen}^2 x\operatorname{sen} p=2\cos x+\cos x\operatorname{sen} x\cos p+\cos^2x\operatorname{sen} p,
\]

que puede ser reducida a

\begin{equation}\phantomsection\label{eq-periodicidad}{
2\cos x(\cos p-1)=\operatorname{sen} p(1+2\operatorname{sen} x).
}\end{equation}

Ahora bien, si \(\operatorname{sen} p\neq 0\), entonces obtenemos

\[
\frac{\cos p-1}{\operatorname{sen} p}=\frac{1+2\operatorname{sen} x}{2\cos x},
\]

pero como esta igualdad debe cumplirse para todo \(x\) tal que
\(\cos x\neq 0\), es un absurdo. Entonces debe ser
\(\operatorname{sen} p=0\), y sustituyendo esto en
(\ref{eq-periodicidad}) resulta

\[
2\cos x(\cos p-1)=0,
\]

de donde concluimos que \(\cos p=1\). El valor positivo \(p\) más
pequeño que verifica \(\operatorname{sen} p=0\) y \(\cos p=1\) es
\(p=2\pi\), con lo cual \(f\) es periódica de período \(2\pi\). Así,
basta analizarla ahora en \([0,2\pi]\) y luego extender su gráfico a
todo \(\mathbb{R}\) de forma periódica.

\textbf{(D) Asíntotas}

\begin{itemize}
\item
  La gráfica no presenta asíntotas verticales, ya que se trata de una
  función continua en todo \(\mathbb{R}\).
\item
  No presenta asíntotas horizontales, ya que se trata de una función
  periódica de período \(2\pi\).
\end{itemize}

\textbf{(E) Intervalos de crecimiento y decrecimiento}

Como \(f\) es periódica de período \(2\pi\), Calculamos \(f'\) usando la
regla del cociente

\[
f'(x)=\frac{-\operatorname{sen} x(2+\operatorname{sen} x)-\cos x\cos x}{(2+\operatorname{sen} x)^2}=-\frac{1+2\operatorname{sen} x}{(2+\operatorname{sen} x)^2}.
\]

Luego \(f'(x)=0\) si y sólo si \(1+2\operatorname{sen} x=0\), o de forma
equivalente \(\operatorname{sen} x=-1/2\). Esto significa que los puntos
críticos de \(f\) son \(x=7\pi/6\) y \(11\pi/6\). Por la
\hyperref[teo-propiedad-signoconst]{propiedad del signo constante},
\(f'\) tendrá signo constante en los intervalos \((0, 7\pi/6)\),
\((7\pi/6,11\pi/6)\) y \((11\pi/6,2\pi)\). En la siguiente tabla
determinamos el signo de \(f'\) en cada uno de ellos.

\begin{longtable}[]{@{}
  >{\raggedright\arraybackslash}p{(\linewidth - 6\tabcolsep) * \real{0.2500}}
  >{\raggedright\arraybackslash}p{(\linewidth - 6\tabcolsep) * \real{0.2500}}
  >{\raggedright\arraybackslash}p{(\linewidth - 6\tabcolsep) * \real{0.2500}}
  >{\raggedright\arraybackslash}p{(\linewidth - 6\tabcolsep) * \real{0.2500}}@{}}
\toprule\noalign{}
\begin{minipage}[b]{\linewidth}\raggedright
Intervalo
\end{minipage} & \begin{minipage}[b]{\linewidth}\raggedright
\(x_0\)
\end{minipage} & \begin{minipage}[b]{\linewidth}\raggedright
\(f'(x_0)\)
\end{minipage} & \begin{minipage}[b]{\linewidth}\raggedright
\(f\)
\end{minipage} \\
\midrule\noalign{}
\endhead
\bottomrule\noalign{}
\endlastfoot
\(\displaystyle \left(0, \frac{7\pi}{6}\right)\) & \(\pi\) & \(-\) &
decrece \\
\(\displaystyle \left(\frac{7\pi}{6}, \frac{11\pi}{6}\right)\) &
\(\displaystyle\frac{3\pi}{2}\) & \(+\) & crece \\
\(\displaystyle \left(\frac{11\pi}{6}, 2\pi\right)\) &
\(\displaystyle\frac{23\pi}{12}\) & \(-\) & decrece \\
\end{longtable}

Con la \hyperref[teo-prueba-cd]{prueba creciente/decreciente} podemos
concluir que \(f\) crece en \((7\pi/6, 11\pi/6)\) y que decrece en
\((0, 7\pi/6)\) y en \((11\pi/6, 2\pi)\).

\textbf{(F) Valores extremos locales}

En la tabla anterior vemos que \(f'\) cambia de negativa a positiva en
\(x=7\pi/6\), por lo que \(f\) presenta un mínimo local en este punto
(por la \hyperref[teo-prueba-primerader]{prueba de la primera derivada})
y su valor es \(f(7\pi/6)=-\sqrt{3}/2\approx -0.577\). También vemos que
\(f'\) cambia de positiva a negativa en \(x=11\pi/6\), con lo cual \(f\)
presenta un máximo local en este punto de valor
\(f(11\pi/6)=\sqrt{3}/2\approx 0.577\).

\textbf{(G) Concavidad y puntos de inflexión}

Calculamos la derivada segunda de \(f\), nuevamente con la regla del
cociente

\[
f''(x)=-\frac{2\cos x(2+\operatorname{sen} x)^2-(1+2\operatorname{sen} x)2(2+\operatorname{sen} x)\cos x}{(2+\operatorname{sen} x)^4}=-\frac{2\cos x(1-\operatorname{sen} x)}{(2+\operatorname{sen} x)^3}.
\]

Así, \(f''(x)=0\) cuando \(\cos x=0\) o cuando
\(\operatorname{sen} x=1\). Los puntos son \(x=\pi/2\) y \(x=3\pi/2\)
para que se cumpla la primera, y \(x=\pi/2\) para la segunda. Por la
\hyperref[teo-propiedad-signoconst]{propiedad del signo constante}
\(f''\) tendrá signo constante en \((0,\pi/2)\), en \((\pi/2,3\pi/2)\) y
en \((3\pi/2, 2\pi)\). En la siguiente tabla calculamos estos signos.

\begin{longtable}[]{@{}
  >{\raggedright\arraybackslash}p{(\linewidth - 6\tabcolsep) * \real{0.2500}}
  >{\raggedright\arraybackslash}p{(\linewidth - 6\tabcolsep) * \real{0.2500}}
  >{\raggedright\arraybackslash}p{(\linewidth - 6\tabcolsep) * \real{0.2500}}
  >{\raggedright\arraybackslash}p{(\linewidth - 6\tabcolsep) * \real{0.2500}}@{}}
\toprule\noalign{}
\begin{minipage}[b]{\linewidth}\raggedright
Intervalo
\end{minipage} & \begin{minipage}[b]{\linewidth}\raggedright
\(x_0\)
\end{minipage} & \begin{minipage}[b]{\linewidth}\raggedright
\(f'(x_0)\)
\end{minipage} & \begin{minipage}[b]{\linewidth}\raggedright
\(f\)
\end{minipage} \\
\midrule\noalign{}
\endhead
\bottomrule\noalign{}
\endlastfoot
\(\displaystyle \left(0, \frac{\pi}{2}\right)\) &
\(\displaystyle\frac{\pi}{4}\) & \(-\) & cóncava \\
\(\displaystyle \left(\frac{\pi}{2}, \frac{3\pi}{2}\right)\) & \(\pi\) &
\(+\) & convexa \\
\(\displaystyle \left(\frac{3\pi}{2}, 2\pi\right)\) &
\(\displaystyle\frac{7\pi}{4}\) & \(-\) & cóncava \\
\end{longtable}

Entonces \(f\) resulta convexa en \((\pi/2,3\pi/2)\) y cóncava en
\((0,\pi/2)\) y en \((3\pi/2,2\pi)\) como consecuencia de la
\hyperref[teo-prueba-concavidad]{prueba de la concavidad}. El signo de
\(f''\) cambia en \(x=\pi/2\) y en \(x=3\pi/2\), ambos puntos de
continuidad de \(f\), por lo que \(P=(\pi/2,f(\pi/2))=(\pi/2,0)\) y
\(Q=(3\pi/2, f(3\pi/2))=(3\pi/2,0)\) son los dos puntos de inflexión de
la gráfica.

\textbf{(H) Trazar la gráfica} Con lo calculado anteriormente trazamos
la gráfica aproximada de \(f\).

\includegraphics[width=0.5\linewidth,height=\textheight,keepaspectratio]{index_files/mediabag/imagenes_cap4/4.5-funcion4.pdf}

\end{tcolorbox}

Ejemplo

Graficar \(y=\ln(4-x^2)\).

\begin{tcolorbox}[enhanced jigsaw, opacityback=0, left=2mm, opacitybacktitle=0.6, toprule=.15mm, bottomrule=.15mm, colbacktitle=quarto-callout-tip-color!10!white, titlerule=0mm, bottomtitle=1mm, title=\textcolor{quarto-callout-tip-color}{\faLightbulb}\hspace{0.5em}{Solución}, toptitle=1mm, colframe=quarto-callout-tip-color-frame, colback=white, breakable, rightrule=.15mm, arc=.35mm, leftrule=.75mm, coltitle=black]

\textbf{(A) Dominio} Debido a que el logaritmo de cero y de números
negativos no existe, planteamos que \(4-x^2>0\), lo que nos lleva a
\(|x|<2\), o también \(-2<x<2\). En consecuencia, \(D_f=(-2,2)\).

\textbf{(B) Intersecciones}

\begin{itemize}
\tightlist
\item
  Para encontrar intersecciones con el eje \(x\) hacemos
\end{itemize}

\[
f(x)=0, \quad \text{ es decir,} \quad \ln(4-x^2)=0,
\]

que nos conduce a \(x^2=3\) y entonces \(x=\pm \sqrt{3}\).

\begin{itemize}
\tightlist
\item
  Para calcular la intersección con el eje \(y\) evaluamos
\end{itemize}

\[
f(0)=\ln(4-0^2)=\ln 4\approx 1.386.
\]

Entonces la gráfica interseca al eje \(y\) en \(y=\ln 4\).

\textbf{(C) Simetría y periodicidad}

Comenzamos calculando \(f(-x)\)

\[
f(-x)=\ln(4-(-x)^2)=\ln(4-x^2)=f(x),
\]

para todo \(x\) en el dominio de \(f\), de donde resulta que \(f\) es
par.

En este caso \(f\) no resulta periódica.

\textbf{(D) Asíntotas}

\begin{itemize}
\tightlist
\item
  Dado que \(f\) es continua en su dominio, analizamos los puntos
  extremos del intervalo.
\end{itemize}

Observar que tanto cuando \(x\to 2^-\) como cuando \(x\to -2^+\),
\(4-x^2\to 0\) con valores positivos. Como
\(\lim_{t\to 0^+} ln(t)=-\infty\), concluimos que

\[
\lim_{x\to -2^+} f(x)=-\infty \quad \text{ y también } \quad \lim_{x\to 2^-} f(x)=-\infty.
\]

Es decir, \(x=-2\) y \(x=2\) son asíntotas verticales de la gráfica de
\(f\).

\begin{itemize}
\tightlist
\item
  La función no presenta asíntotas horizontales, ya que \(f\) sólamente
  está definida en \((-2,2)\).
\end{itemize}

\textbf{(E) Intervalos de crecimiento y decrecimiento}

Calculamos \(f'\) utilizando la derivada del logaritmo natural y la
regla de la cadena

\[
f'(x)=-\frac{2x}{4-x^2},
\]

de donde podemos ver que \(f'(x)=0\) si y sólo si \(x=0\). Por lo tanto
\(f\) tiene un único punto crítico, \(x=0\). Utilizando la
\hyperref[teo-propiedad-signoconst]{propiedad del signo constante},
\(f'\) tendrá signo constante en los intervalos \((-2, 0)\) y en
\((0,2)\). La siguiente tabla muestra el signo de \(f'\) en cada uno.

\begin{longtable}[]{@{}llll@{}}
\toprule\noalign{}
Intervalo & \(x_0\) & \(f'(x_0)\) & \(f\) \\
\midrule\noalign{}
\endhead
\bottomrule\noalign{}
\endlastfoot
\((-2,0)\) & \(-1\) & \(+\) & crece \\
\((0,2)\) & 1 & \(-\) & decrece \\
\end{longtable}

Utilizando la \hyperref[teo-prueba-cd]{prueba creciente/decreciente}
concluimos que \(f\) crece en \((-2,0)\) y decrece en \((0,2)\).

\textbf{(F) Valores extremos locales}

En la tabla anterior vemos que \(f'\) cambia de positiva a negativa en
\(x=0\), por lo que \(f\) presenta un máximo local en este punto (por la
\hyperref[teo-prueba-primerader]{prueba de la primera derivada}) y su
valor es \(f(0)=\ln 4\).

\textbf{(G) Concavidad y puntos de inflexión}

Calculamos la derivada segunda de \(f\) usando la regla del cociente

\[
f''(x)=-\frac{2(4-x^2)-2x(-2x)}{(4-x^2)^2}=-\frac{2x^2+8}{(4-x^2)^2}.
\]

Podemos observar que \(f''(x)<0\) en todo su dominio, es decir, en todo
el intervalo \((-2,2)\). Por la \hyperref[teo-prueba-concavidad]{prueba
de la concavidad} \(f\) resulta cóncava en todo su dominio, y en
consecuencia no hay puntos de inflexión.

\textbf{(H) Trazar la gráfica} Con todo lo calculado anteriormente
trazamos la gráfica aproximada de \(f\).

\includegraphics[width=0.45\linewidth,height=\textheight,keepaspectratio]{index_files/mediabag/imagenes_cap4/4.5-funcion5.pdf}

\end{tcolorbox}

\hyperref[seccion_4.5]{↑ Volver al inicio de la sección}

\section{4.7. Problemas de optimización}\label{seccion_4.7}

En esta sección aplicaremos los métodos aprendidos para hallar extremos
de una función a problemas concretos. La idea es optimizar cierta
cantidad, estableciendo una función que debe maximizarse (como
utilidades o ganancias) o minimizarse (como un costo, una distancia o un
tiempo de recorrido). Veamos algunos ejemplos.

Ejemplo

Un granjero tiene 2400 pies de cerca y desea cercar un campo rectangular
que limita con un río recto. No necesita cercar a lo largo del río.
¿Cuáles son las dimensiones del campo que tiene el área más grande?

\begin{center}
\includegraphics[width=0.3\linewidth,height=\textheight,keepaspectratio]{index_files/mediabag/imagenes_cap4/4.7-campo-rectangular.pdf}
\end{center}

\begin{tcolorbox}[enhanced jigsaw, opacityback=0, left=2mm, opacitybacktitle=0.6, toprule=.15mm, bottomrule=.15mm, colbacktitle=quarto-callout-tip-color!10!white, titlerule=0mm, bottomtitle=1mm, title=\textcolor{quarto-callout-tip-color}{\faLightbulb}\hspace{0.5em}{Solución}, toptitle=1mm, colframe=quarto-callout-tip-color-frame, colback=white, breakable, rightrule=.15mm, arc=.35mm, leftrule=.75mm, coltitle=black]

El primer paso es definir adecuadamente la función a optimizar,
especificando las unidades de medida y el intervalo al que pertenece la
variable independiente. Si llamamos \(x\) e \(y\) a las longitudes
medidas en pies de los laterales del campo y del fondo, respectivamente,
entonces el área encerrada es \[
A=xy.
\]

Sin embargo, sabemos que el perímetro de la región resultante debe sewr
de 2400 pies, con lo cual

\[
2x+y=2400, \quad \text{ o bien} \quad y=2400-2x.
\]

Ahora podemos expresar el área como una función de \(x\)

\[
A(x)=x(2400-2x)=2400x-2x^2,
\]

donde \(0\leq x\leq 1200\). Esta restricción aparece al imponer las
condiciones de que tanto \(x\) como \(A(x)\) sean no negativas.

Como \(A\) es una función continua en el intervalo \([0,1200]\), el
\hyperref[teo-valores-extremos]{teorema de los valores extremos} asegura
que alcanza tanto un máximo como un mínimo absoluto en el intervalo.

Buscamos entonces los puntos críticos de \(A\). Tenemos que

\[
A'(x)=2400-4x,
\]

y \(A'(x)=0\) si y sólo si \(x=600\). Ahora evaluamos \(A\) en este
punto, y en los extremos del intervalo

\[
A(600)=2400\cdot 600-2(600)^2=720000
\]

\[
 A(0)=2400\cdot 0 - 2\cdot 0^2=0, \quad A(1200)=2400\cdot 1200- 2(1200)^2=0.
\]

De esta manera, el máximo valor posible para el área es \(720000\)
pies2, y se logra cuando \(x=600\) pies. Luego,
\(y=2400-2\cdot 600=1200\) pies. Por lo que el sector que encierra área
máxima debe tener \(600\) pies en los laterales y \(1200\) pies en el
fondo.

\end{tcolorbox}

Ejemplo

Se va a fabricar una lata para que contenga 1 litro de aceite. Hallar
las dimensiones que minimizarán el costo del metal para fabricar la
lata.

\begin{center}
\includegraphics[width=0.25\linewidth,height=\textheight,keepaspectratio]{index_files/mediabag/imagenes_cap4/4.7-lata.pdf}
\end{center}

\begin{tcolorbox}[enhanced jigsaw, opacityback=0, left=2mm, opacitybacktitle=0.6, toprule=.15mm, bottomrule=.15mm, colbacktitle=quarto-callout-tip-color!10!white, titlerule=0mm, bottomtitle=1mm, title=\textcolor{quarto-callout-tip-color}{\faLightbulb}\hspace{0.5em}{Solución}, toptitle=1mm, colframe=quarto-callout-tip-color-frame, colback=white, breakable, rightrule=.15mm, arc=.35mm, leftrule=.75mm, coltitle=black]

Como queremos minimizar el costo de fabricación, debemos minimizar el
área de la superficie de la lata. Si llamamos \(r\) y \(h\) al radio de
la base y la altura, respectivamente, medidas en cm, el área de la
superficie es

\[
A=A_\ell+A_b=2\pi rh+2\pi r^2.
\]

La condición sobre la capacidad de la lata nos da

\[
1000 =\pi r^2 h,
\]

y despejando \(h\) en términos de \(r\)

\[
h=\frac{1000}{\pi r^2}.
\]

Usando esta relación en la fórmula para el área resulta

\[
A(r)=2\pi r\frac{1000}{\pi r^2}+2\pi r^2=\frac{2000}{r}+2\pi r^2
\]

con \(r\in(0,+\infty)\).

Buscamos los puntos críticos de \(A\), calculando la derivada.

\[
A'(r)=-\frac{2000}{r^2}+4\pi r
\]

con lo cual \(A'(r)=0\) si y sólo si \(r=\sqrt[3]{500/\pi}=r_0\). Para
determinar la naturaleza de este punto crítico, podemos emplear la
\hyperref[teo-prueba-segundader]{prueba de la segunda derivada}.

\[
A''(r)=\frac{4000}{r^3}+4\pi,
\]

de donde \(A''(r_0)>0\) y entonces la función \(A\) alcanza un mínimo
local en \(r_0\), cuyo valor es

\[
A(r_0)=\frac{2000}{r_0}+2\pi r_0^2=\frac{2000}{\sqrt[3]{500/\pi}}+2\pi (500/\pi)^{2/3}=\frac{3000}{\sqrt[3]{500/\pi}}\, \mathrm{cm}^2 \approx 553.58 \mathrm{cm}^2.
\]

Para asegurarnos de que este valor mínimo local es también absoluto,
estudiamos los límites de \(A\) cuando \(r\) se aproxima a cero o a
infinito

\[
\lim_{r\to 0^+} A(r)=\lim_{r\to 0^+} \frac{2000}{r}+2\pi r^2=\infty
\]

y

\[
\lim_{r\to \infty} A(r)=\lim_{r\to \infty} \frac{2000}{r}+2\pi r^2=\infty,
\]

con lo cual el mínimo encontrado es absoluto. Las dimensiones de la lata
que minimizan el área de la superficien (y, por lo tanto, el costo de
fabricación) son

\[
r=r_0=\sqrt[3](500/pi)\approx 5.419 \,\mathrm{cm} \quad \text{ y } \quad h=\frac{1000}{\pi r_0^2}=2r_0=2\sqrt[3](500/pi)\approx 10.838 \,\mathrm{cm}.
\]

\end{tcolorbox}

\begin{tcolorbox}[enhanced jigsaw, opacityback=0, left=2mm, opacitybacktitle=0.6, toprule=.15mm, bottomrule=.15mm, colbacktitle=quarto-callout-caution-color!10!white, titlerule=0mm, bottomtitle=1mm, title=\textcolor{quarto-callout-caution-color}{\faFire}\hspace{0.5em}{Importante}, toptitle=1mm, colframe=quarto-callout-caution-color-frame, colback=white, breakable, rightrule=.15mm, arc=.35mm, leftrule=.75mm, coltitle=black]

Podemos utilizar otro método para determinar la naturaleza del punto
crítico hallado. Como \(A'\) es continua en \((0,+\infty)\) y se anula
en \(r_0\), la \hyperref[teo-prueba-signoconst]{propiedad del signo
constante} establece que tendrá signo constante en \((0,r_0)\) y en
\((r_0,+\infty)\). En la siguiente tabla tenemos el signo de \(A'\) en
cada uno.

\begin{longtable}[]{@{}llll@{}}
\toprule\noalign{}
Intervalo & \(r\) & \(A'(r)\) & \(A\) \\
\midrule\noalign{}
\endhead
\bottomrule\noalign{}
\endlastfoot
\((0,r_0)\) & \(1\) & \(-\) & decrece \\
\((r_0,+\infty)\) & \(6\) & \(+\) & crece \\
\end{longtable}

Por la \hyperref[teo-prueba-cd]{prueba creciente/decreciente} sabemos
que \(A\) decrece en \((0,r_0)\) y crece en \((r_0,+\infty)\), por lo
que en \(r_0\) alcanza su único mínimo local, que también es absoluto.

\end{tcolorbox}

Podemos dar una versión de la \hyperref[teo-prueba-primerader]{prueba de
la primera derivada} para extremos absolutos.

\phantomsection\label{teo-primerader-global}
Prueba de la primera derivada para valores extremos absolutos

Supongamos que \(c\) es un punto crítico de una función continua
definida sobre un intervalo \(I\).

\begin{itemize}
\item
  Si \(f'(x)>0\) para todo \(x<c\) con \(x\in I\) y \(f'(x)<0\) para
  todo \(x>c\), \(x\in I\) entonces \(f(c)\) es el valor máximo absoluto
  de \(f\) en \(I\).
\item
  Si \(f'(x)<0\) para todo \(x<c\) con \(x\in I\) y \(f'(x)>0\) para
  \(x>c\), \(x\in I\) entonces \(f(c)\) es el valor mínimo absoluto de
  \(f\) en \(I\).
\end{itemize}

Ejemplo

Encontrar el punto sobre la parábola \(y^2=2x\) más cercano al punto
\((1,4)\).

\begin{center}
\includegraphics[width=0.35\linewidth,height=\textheight,keepaspectratio]{index_files/mediabag/imagenes_cap4/4.7-parabola.pdf}
\end{center}

\begin{tcolorbox}[enhanced jigsaw, opacityback=0, left=2mm, opacitybacktitle=0.6, toprule=.15mm, bottomrule=.15mm, colbacktitle=quarto-callout-tip-color!10!white, titlerule=0mm, bottomtitle=1mm, title=\textcolor{quarto-callout-tip-color}{\faLightbulb}\hspace{0.5em}{Solución}, toptitle=1mm, colframe=quarto-callout-tip-color-frame, colback=white, breakable, rightrule=.15mm, arc=.35mm, leftrule=.75mm, coltitle=black]

Si \((x,y)\) es cualquier punto sobre la parábola, entonces queremos que
la distancia desde \((x,y)\) a \((1,4)\), es decir

\[
d=\sqrt{(x-1)^2+(y-4)^2}
\]

sea lo más chica posible. Como \(y^2=2x\) por estar el punto en la
parábola, también podemos escribir \(x=y^2/2\). Con esta relación
expresamos a la distancia en términos de \(y\)

\[
d(y)=\sqrt{\left(\frac{1}{2}y^2-1\right)^2+(y-4)^2}.
\]

Dado que \(d(y)\geq 0\) para todo \(y\), \(d\) es mínima si y sólo si
\(d^2\) es mínima. Para evitar trabajar con la raíz cuadrada,
consideramos

\[
f(y)=d^2(y)=\left(\frac{1}{2}y^2-1\right)^2+(y-4)^2,
\]

para \(y\in \mathbb{R}\). Buscamos los puntos críticos, hallando \(f'\).

\[
f'(y)=2\left(\frac{1}{2}y^2-1\right)y+2(y-4)=y^3-8,
\]

con lo que \(f'(y)=0\) si y sólo si \(y=2\). Como \(f\) presenta un
único punto crítico en \(\mathbb{R}\), hacemos la tabla

\begin{longtable}[]{@{}llll@{}}
\toprule\noalign{}
Intervalo & \(y_0\) & \(f'(y_0)\) & \(f\) \\
\midrule\noalign{}
\endhead
\bottomrule\noalign{}
\endlastfoot
\((-\infty,2)\) & \(0\) & \(-\) & decrece \\
\((2,+\infty)\) & \(3\) & \(+\) & crece \\
\end{longtable}

Utilizando la \hyperref[teo-primerader-global]{prueba de la primera
derivada para valores extremos absolutos} concluimos que \(f\) presenta
un mínimo absoluto en \(y=2\), cuyo valor es

\[
f(2)=\left(\frac{1}{2}2^2-1\right)^2+(2-4)^2=5,
\]

con lo que la distancia mínima es \(d=\sqrt{5}\) (recordar que
minimizamos \(d^2\)). El punto sobre la parábola donde se alcanza esta
distancia mínima es \((2^2/2,2)=(2,2)\).

\end{tcolorbox}

Ejemplo

Un hombre está en un punto \(A\) sobre una de las riberas de un río
recto que tiene 3 km de ancho y desea llegar hasta el punto \(B\), 8 km
corriente abajo en la ribera opuesta, tan rápido como le sea posible.
Podría remar en su bote, cruzar directamente el río hasta el punto \(C\)
y correr hasta \(B\), o podría remar hasta \(B\) o, en última instancia,
remar hasta algún punto \(D\), entre \(C\) y \(B\), y luego correr hasta
\(B\). Si puede remar a 6 km/h y correr a 8 km/h, ¿dónde debe
desembarcar para llegar a \(B\) tan pronto como sea posible? (Suponer
que la rapidez del agua es insignificante comparada con la rapidez a la
que rema el hombre.)

\begin{center}
\includegraphics[width=0.4\linewidth,height=\textheight,keepaspectratio]{index_files/mediabag/imagenes_cap4/4.7-rio.pdf}
\end{center}

\begin{tcolorbox}[enhanced jigsaw, opacityback=0, left=2mm, opacitybacktitle=0.6, toprule=.15mm, bottomrule=.15mm, colbacktitle=quarto-callout-tip-color!10!white, titlerule=0mm, bottomtitle=1mm, title=\textcolor{quarto-callout-tip-color}{\faLightbulb}\hspace{0.5em}{Solución}, toptitle=1mm, colframe=quarto-callout-tip-color-frame, colback=white, breakable, rightrule=.15mm, arc=.35mm, leftrule=.75mm, coltitle=black]

Supongamos que el hombre rema desde \(A\) hasta un punto \(D\) ubicado
entre \(C\) y \(B\) (pudiendo, tal vez, coincidir con alguno de éstos).
Sea \(x\) la distancia en km desde el punto \(C\) al punto \(D\). Si
\(d_1(x)\) es la distancia en km desde \(A\) hasta \(D\), utilizando el
teorema de Pitágoras resulta

\[
d_1^2(x)=x^2+3^2, \quad \text{ o bien }\quad d_1(x)=\sqrt{9+x^2}.
\]

Entonces el hombre rema \(d_1(x)\) km desde \(A\) hasta \(D\), y luego
corre \(d_2(x)=8-x\) kilómetros desde \(D\) hasta \(B\). Podemos
calcular el tiempo total dividiendo cada distancia por la velocidad
correspondiente. Si llamamos \(T(x)\) tal tiempo en horas, en función de
\(x\) obtenemos

\[
T(x)=\frac{d_1(x)}{6}+\frac{d_2(x)}{8}=\frac{\sqrt{9+x^2}}{6}+\frac{8-x}{8},
\]

donde \(x\in [0,8]\).

Para llegar al punto \(B\) lo más rapido posible, es preciso minimizar
el tiempo de recorrido.

Como estamos en el caso de una función continua en un intervalo cerrado,
buscamos los puntos críticos en \((0,8)\) y comparamos los valores de la
función en estos puntos con los valores que toma en los extremos del
intervalo.

\[
T'(x)=\frac{1}{6}\frac{1}{2\sqrt{9+x^2}}2x+\frac{1}{8}(-1)=\frac{x}{6\sqrt{9+x^2}}-\frac{1}{8}
\]

que se anula cuando \(x=9\sqrt{7}/7\approx 3.4\) km. Dado que

\[
T(9\sqrt{7}/7)=\frac{\sqrt{9+\frac{81\cdot 7}{49}}}{6}+\frac{8-\frac{9\sqrt{7}}{7}}{8}=1+\frac{\sqrt{7}}{8}\approx 1.33
\]

y

\[
T(0)=\frac{\sqrt{9+0^2}}{6}+\frac{8-0}{8}=\frac{1}{2}+1=1.5 \quad \text{ y }\quad T(8)=\frac{\sqrt{9+8^2}}{6}+\frac{8-8}{8}=\frac{\sqrt{73}}{6}\approx 1.42, 
\]

el mínimo tiempo posible se consigue cuando \(x=9\sqrt{7}/7\) km.

\end{tcolorbox}

Ejemplo

Encontrar el área del rectángulo más grande que puede inscribirse en una
semicircunferencia de radio \(r\).

\begin{center}
\includegraphics[width=0.4\linewidth,height=\textheight,keepaspectratio]{index_files/mediabag/imagenes_cap4/4.7-semicircunferencia.pdf}
\end{center}

\begin{tcolorbox}[enhanced jigsaw, opacityback=0, left=2mm, opacitybacktitle=0.6, toprule=.15mm, bottomrule=.15mm, colbacktitle=quarto-callout-tip-color!10!white, titlerule=0mm, bottomtitle=1mm, title=\textcolor{quarto-callout-tip-color}{\faLightbulb}\hspace{0.5em}{Solución}, toptitle=1mm, colframe=quarto-callout-tip-color-frame, colback=white, breakable, rightrule=.15mm, arc=.35mm, leftrule=.75mm, coltitle=black]

Por simplicidad podemos ubicar la semicircunferencia centrada en el
origen, abarcando el primer y el segundo cuadrantes. Sean \(x\) e \(y\)
las coordenadas del vértice superior derecho del rectángulo que toca a
la circunferencia (es decir, \(x\) e \(y\) son positivos). Entonces el
rectángulo tiene altura \(y\) y base de longitud \(2x\). Su área resulta
entonces

\[
A=2xy.
\]

Como el punto \((x,y)\) también está ubicado sobre la
semicircunferencia, se cumple que \(x^2+y^2=r^2\), o también
\(y=\sqrt{r^2-x^2}\). Por lo tanto, el área como función de \(x\)
resulta ser

\[
A(x)=2x\sqrt{r^2-x^2},
\]

con \(0\leq x\leq r\). Como \(A\) es continua en el intervalo cerrado
dado, sabemos que alcanza extremos absolutos. La derivada de \(A\) es

\[
A'(x)=2\sqrt{r^2-x^2}+2x\frac{1}{2\sqrt{r^2-x^2}}(-2x)=\frac{2r^2-4x^2}{\sqrt{r^2-x^2}}.
\]

El único punto crítico en el intervalo \((0,r)\) es \(x=r\sqrt{2}/2\).
Evaluamos entonces

\[
A(r\sqrt{2}/2)=2\frac{\sqrt{2}}{2}r\sqrt{r^2-\frac{r^2}{2}}=r^2
\]

y también en los extremos del intervalo

\[
A(0)=2\cdot 0\sqrt{r^2-0^2}=0 \quad \text{ y }\quad A(r)=2r\sqrt{r^2-r^2}=0.
\]

Con lo cual, el área máxima es \(r^2\) y se alcanza cuando
\(x=r\sqrt{2}/2\) e \(y=\sqrt{r^2-\frac{r^2}{2}}=r\sqrt{2}/2\). Es
decir, el rectángulo tiene altura igual a la mitad de su base.

\end{tcolorbox}

\hyperref[seccion_4.7]{↑ Volver al inicio de la sección}

\bookmarksetup{startatroot}

\chapter{5. Integrales}\label{integrales}

\section{5.1. Áreas y distancias}\label{seccion_5.1}

Comenzaremos este capítulo motivando la definición del concepto de
integral. Con la derivada, le dábamos sentido al cociente de cantidades
infinitesimales, y para ello calculábamos el límite de cocientes
incrementales. Con la integral daremos sentido formal a ``sumar''
cantidades infinitesimales.

\subsection{Área bajo la curva}\label{uxe1rea-bajo-la-curva}

Supongamos que queremos calcular el área de una región plana. Muchas
veces disponemos de una fórmula para el cálculo en término de sus
dimensiones, sobre todo cuando la región es de una forma sencilla.

\begin{center}
\includegraphics[width=0.65\linewidth,height=\textheight,keepaspectratio]{index_files/mediabag/imagenes_cap5/5.1-area-figuras.pdf}
\end{center}

Ahora quisiéramos calcular el área de una figura plana un poco más
general. Supongamos que dicha región \(S\) está comprendida entre el eje
\(x\) y el gráfico de una función no negativa y continua \(y = f(x)\),
entre \(x = a\) y \(x = b\).

\begin{center}
\includegraphics[width=0.5\linewidth,height=\textheight,keepaspectratio]{index_files/mediabag/imagenes_cap5/5.1-area-de-S.pdf}
\end{center}

\begin{equation}\phantomsection\label{eq-regionS}{
S=\{x\in\mathbb{R}^2: 0\leq y\leq f(x), \quad a\leq x\leq b\}.
}\end{equation}

Sin embargo, como no tenemos definido qué es el área de una región de
este tipo, comenzaremos obteniendo una aproximación del área de \(S\)
utilizando rectángulos, cuyas áreas conocemos y podemos calcular con
facilidad.

Ejemplo

Utilizar rectángulos para estimar el área debajo de la parábola
\(y = x^2\), desde \(x=0\) hasta \(x=1\).

\begin{center}
\includegraphics[width=0.45\linewidth,height=\textheight,keepaspectratio]{index_files/mediabag/imagenes_cap5/5.1-parabola.pdf}
\end{center}

\begin{tcolorbox}[enhanced jigsaw, opacityback=0, left=2mm, opacitybacktitle=0.6, toprule=.15mm, bottomrule=.15mm, colbacktitle=quarto-callout-tip-color!10!white, titlerule=0mm, bottomtitle=1mm, title=\textcolor{quarto-callout-tip-color}{\faLightbulb}\hspace{0.5em}{Solución}, toptitle=1mm, colframe=quarto-callout-tip-color-frame, colback=white, breakable, rightrule=.15mm, arc=.35mm, leftrule=.75mm, coltitle=black]

Si \(A\) denota el área de \(S\), entonces es claro que \(0<A<1\), dado
que la región \(S\) puede encerrarse en un cuadrado de lados de longitud
\(1\) cuya área es, por supuesto, \(1\). Ahora dividimos \(S\) en cuatro
franjas \(S_1\), \(S_2\), \(S_3\) y \(S_4\), trazando cuatro líneas
verticales en \(x=1/4\), \(x=1/2\) y \(x=3/4\), respectivamente.

\includegraphics[width=0.45\linewidth,height=\textheight,keepaspectratio]{index_files/mediabag/imagenes_cap5/5.1-parabola-franjas.pdf}

Ahora bien, podemos aproximar el área de cada franja por la del
rectángulo que tiene la misma base que la franja, y la altura es la
misma que la del lado derecho de la correspondiente franja. Es decir,
las alturas de los rectángulos son los valores que toma la función en
los extremos derechos de cada subintervalo en los que quedó dividido el
\([0,1]\). Si denotamos con \(R_4\) a la suma de las áreas de estos
rectángulos, resulta

\[
R_4=\frac{1}{4} \left(\frac{1}{4}\right)^2 + \frac{1}{4} \left(\frac{1}{2}\right)^2 + \frac{1}{4} \left(\frac{3}{4}\right)^2 + \frac{1}{4} 1^2= \frac{15}{32} = 0.46875.
\]

\includegraphics[width=0.6\linewidth,height=\textheight,keepaspectratio]{index_files/mediabag/imagenes_cap5/5.1-parabola-izqder.pdf}

Como vemos en el dibujo, ésto nos da una aproximación de \(A\) por
exceso. Otra manera posible para aproximar \(A\) es sumar las áreas de
los rectángulos que tienen la misma base que cada franja, pero la altura
igual a la del lado izquierdo de cada una. En este caso las alturas
corresponden a los valores de \(f\) en los extremos izquierdos de cada
subintervalo en los que dividimos a \([0,1]\). SI denotamos con \(L_4\)
a la suma de éstas áreas, entonces

\[
L_4=\frac{1}{4} 0^2 + \frac{1}{4} \left(\frac{1}{4}\right)^2 + \frac{1}{4} \left(\frac{1}{2}\right)^2 + \frac{1}{4} \left(\frac{3}{4}\right)^2 = \frac{7}{32} = 0.21875,
\]

resultando en este caso una aproximación por defecto. Es decir

\[
0.21875<A<0.46875.
\]

Este procedimiento puede repetirse considerando una mayor cantidad de
franjas. En la siguiente tabla se muestran los valores calculados en
distintos casos.

\begin{longtable}[]{@{}lll@{}}
\toprule\noalign{}
\(n\) franjas & \(L_n\) & \(R_n\) \\
\midrule\noalign{}
\endhead
\bottomrule\noalign{}
\endlastfoot
10 & 0.2850000 & 0.3850000 \\
20 & 0.3087500 & 0.3587500 \\
30 & 0.3168519 & 0.3501852 \\
50 & 0.3234000 & 0.3434000 \\
100 & 0.3283500 & 0.3383500 \\
1000 & 0.3328335 & 0.3338335 \\
\end{longtable}

\includegraphics[width=0.7\linewidth,height=\textheight,keepaspectratio]{index_files/mediabag/imagenes_cap5/5.1-parabola-rectangulos.pdf}

Todo parece indicar que los valores de \(L_n\) y \(R_n\) se aproximan a
\(1/3\) cuando \(n\) crece.

\end{tcolorbox}

Ejemplo

Para la región \(S\) del ejemplo anterior, demostrar que la suma de las
áreas de los rectángulos superiores de aproximación tiende a \(1/3\). Es
decir, comprobar que
\(\displaystyle \lim_{n\to \infty} R_n=\frac{1}{3}\).

\begin{tcolorbox}[enhanced jigsaw, opacityback=0, left=2mm, opacitybacktitle=0.6, toprule=.15mm, bottomrule=.15mm, colbacktitle=quarto-callout-tip-color!10!white, titlerule=0mm, bottomtitle=1mm, title=\textcolor{quarto-callout-tip-color}{\faLightbulb}\hspace{0.5em}{Solución}, toptitle=1mm, colframe=quarto-callout-tip-color-frame, colback=white, breakable, rightrule=.15mm, arc=.35mm, leftrule=.75mm, coltitle=black]

Fijado un número \(n\in \mathbb{N}\), dividimos el intervalo \([0,1]\)
en \(n\) subintervalos de longitud \(\frac{1}{n}\). Cada rectángulo de
aproximación tiene entonces base \(1/n\) y altura \(f(k/n)\), donde
\(1\leq k\leq n\). De esta manera

\[
\begin{aligned}
R_n
&= \frac{1}{n}\!\left(\frac{1}{n}\right)^2
   + \frac{1}{n}\!\left(\frac{2}{n}\right)^2
   + \frac{1}{n}\!\left(\frac{3}{n}\right)^2
   + \ldots
   + \frac{1}{n}\!\left(\frac{n-1}{n}\right)^2
   + \frac{1}{n}1^2 \\
   \\
&= \sum_{k=1}^n \frac{1}{n}\!\left(\frac{k}{n}\right)^2\\
\\
&=\frac{1}{n^3}\sum_{k=1}^n k^2.
\end{aligned}
\]

Teniendo en cuenta que

\[
\sum_{k=1}^{n} k^2 = \frac{n(n+1)(2n+1)}{6}
\]

resulta

\[
R_n = \frac{1}{n^3} \sum_{k=1}^{n} k^2 =  \frac{(n+1)(2n+1)}{6n^2}.
\]

Por lo tanto

\[
\begin{aligned}
\lim_{n\rightarrow \infty} R_n &= \lim_{n\rightarrow \infty} \frac{(n+1)(2n+1)}{6n^2}\\
&= \lim_{n\rightarrow \infty} \frac{1}{6}\frac{(n+1)}{n}\frac{(2n+1)}{n}\\
\\
&= \frac{1}{6}\cdot 1 \cdot 2 = \frac{1}{3}.
\end{aligned}
\]

Si repetimos el proceso considerando los rectángulos a la izquierda,
obtenemos que

\[
\begin{aligned}
L_n &= \frac{1}{n} 0^2 + \frac{1}{n}\left(\frac{1}{n}\right)^2 + \frac{1}{n}\left(\frac{2}{n}\right)^2 + \frac{1}{n}\left(\frac{3}{n}\right)^2 + \ldots + \frac{1}{n}\left(\frac{n-1}{n}\right)^2 \\
\\
&= \sum_{k=1}^{n-1} \frac{1}{n}\!\left(\frac{k}{n}\right)^2=\frac{(n-1)n(2n-1)}{6n^3}=\frac{1}{6}\frac{n-1}{n}\frac{2n-1}{n}.
\end{aligned}
\]

Y entonces

\[
\begin{aligned}
\lim_{n\rightarrow \infty} L_n &= \lim_{n\rightarrow \infty} \frac{1}{6}\frac{n-1}{n}\frac{2n-1}{n}\\
&= \frac{1}{6}\cdot 1 \cdot 2 = \frac{1}{3}.
\end{aligned}
\]

Con lo cual \(\lim_{n\to \infty} L_n=1/3\).

\end{tcolorbox}

Los ejemplos anteriores nos dan una idea de cómo definir el área de una
region general como la dada en (\ref{eq-regionS}).

En efecto, dividimos la región \(S\) en \(n\) franjas de igual ancho
\(S_i\), \(1\leq i\leq n\). Como el ancho del intervalo es \(b-a\), cada
franja tendrá ancho \(\displaystyle\Delta x= \frac{b-a}{n}\). Estas
franjas dividen al intervalos en los subintervalos

\[
[x_0, x_1], \quad [x_1, x_2], \quad \dots \quad [x_{n-1}, x_n],
\]

donde \(x_0=a\), \(x_n=b\) y, en general

\[
x_i=a+i\Delta x=a+i\frac{b-a}{n}, \quad 0\leq i\leq n.
\]

Podemos aproximar el área de cada franja \(S_i\) por la del rectángulo
con base \(\Delta x\) y altura igual a la imagen por \(f\) del extremo
derecho del intervalo \([x_{i-1}, x_i]\). Sumando todas las áreas de
estas figuras obtenemos la aproximación dada por rectángulos superiores
como

\[
R_n=f(x_1)\Delta x+f(x_2)\Delta x+\dots+f(x_n)\Delta x=\sum_{i=1}^n f(x_i)\Delta x.
\]

\begin{center}
\includegraphics[width=0.6\linewidth,height=\textheight,keepaspectratio]{index_files/mediabag/imagenes_cap5/5.1-rectangulos-der.pdf}
\end{center}

Esta aproximación parece mejorar conforme tomamos \(n\) cada vez más
grande, como vemos en la siguiente figura.

\begin{center}
\includegraphics[width=0.7\linewidth,height=\textheight,keepaspectratio]{index_files/mediabag/imagenes_cap5/5.1-rectangulos-der2.pdf}
\end{center}

\begin{tcolorbox}[enhanced jigsaw, opacityback=0, left=2mm, opacitybacktitle=0.6, toprule=.15mm, bottomrule=.15mm, colbacktitle=quarto-callout-note-color!10!white, titlerule=0mm, bottomtitle=1mm, title=\textcolor{quarto-callout-note-color}{\faInfo}\hspace{0.5em}{Definición (Área)}, toptitle=1mm, colframe=quarto-callout-note-color-frame, colback=white, breakable, rightrule=.15mm, arc=.35mm, leftrule=.75mm, coltitle=black]

El \textbf{área} \(A\) de la región \(S\) dada por la expresión

\[
S=\{x\in\mathbb{R}^2: 0\leq y\leq f(x),\quad a\leq x\leq b\},
\]

donde \(f\) es continua y no negativa en \([a,b]\), se define como

\[
A=\lim_{n\to \infty} R_n=\lim_{n\to \infty} \sum_{i=1}^n f(x_i)\Delta x,
\]

siendo \(\Delta x=(b-a)/n\) y \(x_i=a+i\Delta x\) para
\(0\leq i\leq n\).

\end{tcolorbox}

Es posible demostrar que el límite de arriba siempre existe cuando \(f\)
es continua en \([a,b]\). Más aún, si consideramos rectángulos cuyas
alturas sean las imágenes por \(f\) de los extremos izquierdos de cada
intervalo obtenemos el mismo resultado. Es decir

\[
A=\lim_{n\to \infty} L_n=\lim_{n\to \infty} \sum_{i=0}^{n-1} f(x_i)\Delta x.
\]

Más aún, el valor del límite no cambia si en cada subintervalo
\([x_{i-1},x_i]\) elegimos un punto cualquiera \(x_i^*\) del intervalo.
Estos puntos se llaman \textbf{puntos muestra} y la expresión del área
en este caso se escribe

\[
A=\lim_{n\to \infty} \sum_{i=1}^{n} f(x_i^*)\Delta x.
\]

\subsection{El problema de la
distancia}\label{el-problema-de-la-distancia}

Supongamos que intentamos conocer la distancia recorrida por un objeto
que se mueve con una velocidad \(v\) dada. Por ejemplo, si un automóvil
viaja a \(30\) km/h durante media hora, en este período recorre

\[
30 \,\frac{\mathrm{km}}{\mathrm{h}} \cdot 0.5\, \mathrm{h}  = 15\,\mathrm{km}. 
\]

Observar que si graficamos la velocidad en función del tiempo, esto
equivale a calcular el área debajo del gráfico de \(v\) entre \(t=0\) y
\(t=0.5\).

En general, supongamos que el objeto se mueve con velocidad \(v=f(t)\),
donde \(a\leq t\leq b\) y \(f(t)\geq 0\), de modo que el objeto siempre
se mueve en la dirección positiva.

\begin{center}
\includegraphics[width=0.7\linewidth,height=\textheight,keepaspectratio]{index_files/mediabag/imagenes_cap5/5.1-velocidad.pdf}
\end{center}

Dado \(n\in\mathbb{N}\), subdividimos el intervalo \([a,b]\) en \(n\)
subintervalos \([t_{i-1},t_i]\), con \(1\leq i\leq n\) de igual longitud
\(\Delta t=(b-a)/n\) y aproximamos la distancia recorrida en cada
subintervalo por el producto \(f(t_{i-1})\Delta t\). La distancia total
recorrida aproximada en el intervalo \([a,b]\) será

\[
f(t_0)\Delta t+f(t_1)\Delta t+\dots+f(t_{n-1})\Delta t=\sum_{i=0}^{n-1}f(t_{i})\Delta t.
\]

Si en lugar de estimar la velocidad en los extremos izquierdos de cada
subintervalo lo hacemos con los derechos, obtenemos la aproximación

\[
f(t_1)\Delta t+f(t_2)\Delta t+\dots+f(t_{n})\Delta t=\sum_{i=1}^{n}f(t_{i})\Delta t.
\]

La distancia recorrida \textbf{exacta} \(d\) la encontramos tomando el
límite de cualquiera de las expresiones anteriores cuando
\(n\to\infty\), es decir

\[
d=\lim_{n\to \infty} \sum_{i=0}^{n-1}f(t_{i})\Delta t=\lim_{n\to \infty} \sum_{i=1}^{n}f(t_{i})\Delta t.
\]

Es decir, podemos interpretar la distancia total recorrida por el objeto
en el intervalo de tiempo \([a,b]\) como el área debajo del gráfico de
la función velocidad del mismo, entre \(t=a\) y \(t=b\).

\begin{center}
\includegraphics[width=0.6\linewidth,height=\textheight,keepaspectratio]{index_files/mediabag/imagenes_cap5/5.1-velocidad-area.pdf}
\end{center}

\hyperref[seccion_5.1]{↑ Volver al inicio de la sección}

\section{5.2. La integral definida}\label{seccion_5.2}

Hemos visto que cuando intentamos calcular el área de una región como la
(\ref{eq-regionS}) o la distancia recorrida por un objeto sabiendo su
función velocidad, aparece naturalmente un límite de la forma

\[
\lim_{n\to \infty} \sum_{i=1}^{n} f(x_i^*)\Delta x.
\]

\begin{tcolorbox}[enhanced jigsaw, opacityback=0, left=2mm, opacitybacktitle=0.6, toprule=.15mm, bottomrule=.15mm, colbacktitle=quarto-callout-note-color!10!white, titlerule=0mm, bottomtitle=1mm, title=\textcolor{quarto-callout-note-color}{\faInfo}\hspace{0.5em}{Definición (Integral definida)}, toptitle=1mm, colframe=quarto-callout-note-color-frame, colback=white, breakable, rightrule=.15mm, arc=.35mm, leftrule=.75mm, coltitle=black]

Si \(f\) es una función continua definida para \(a\leq x\leq b\),
dividimos el intervalo \([a,b]\) en \(n\) subintervalos de igual ancho
\(\Delta x=(b-a)/n\). Si \(x_i=a+i\Delta x\), para \(0\leq i\leq n\) y
elegimos un punto muestra \(x_i^*\) en cada subintervalo
\([x_{i-1},x_i]\) para \(1\leq i\leq n\), entonces la \textbf{integral
definida de \(f\) desde \(a\) hasta \(b\)}, es

\[
\int_a^b f(x)\,dx=\lim_{n\to \infty} \sum_{i=1}^n f(x_i^*)\Delta x,
\]

siempre que este límite exista. Cuando es el caso, decimos que \(f\) es
\textbf{integrable} sobre \([a,b]\).

\end{tcolorbox}

Algunas observaciones:

\begin{itemize}
\item
  El símbolo \(\int\) se llama \textbf{símbolo integral}, y fue
  introducido por Leibniz. El mismo recuerda a una ``S'' alargada, dado
  que la integral es un límite de sumas.
\item
  En la expresión \(\displaystyle \int_a^b f(x)\,dx\), \(f(x)\) se llama
  \textbf{integrando}, \(a\) y \(b\) son los \textbf{límites superior} e
  \textbf{inferior} de integración, respectivamente. El símbolo \(dx\)
  no representa nada por sí mismo, sino que indica que la variable
  respecto a la que integramos es \(x\).
\item
  Si existe la integral definida \(\displaystyle \int_a^b f(x)\,dx\),
  entonces es un número, que no depende de \(x\). En este sentido, la
  variable de integración es ``muda'' y cualquiera de estas expresiones
  \[
  \int_a^b f(x)\,dx \qquad \int_a^b f(t)\,dt \qquad \int_a^b f(r)\,dr 
  \] puede usarse para denotar a ese número.
\item
  La suma \(\displaystyle \sum_{i=1}^n f(x_i^*)\Delta x\) se conoce como
  \textbf{suma de Riemann}.
\end{itemize}

Cuando \(f\) es una función positiva, la suma de Riemann puede
interpretarse como la suma de las áreas de los rectángulos de
aproximación.

\begin{center}
\includegraphics[width=0.7\linewidth,height=\textheight,keepaspectratio]{index_files/mediabag/imagenes_cap5/5.2-rectangulos-medios.pdf}
\end{center}

Al comparar las definiciones de integral definida con la de área dada en
la \hyperref[seccion_5.1]{sección} anterior, podemos interpretar la
integral definida \(\displaystyle \int_a^b f(x)\,dx\) como el área
debajo del gráfico de \(y=f(x)\), entre \(x=a\) y \(x=b\).

Si \(f\) toma valores tanto positivos como negativos, entonces la suma
de Riemann es la suma de las áreas de los rectángulos que se encuentran
arriba del eje \(x\) y los negativos de las áreas de los rectángulos que
están debajo del eje \(x\).

\begin{center}
\includegraphics[width=0.5\linewidth,height=\textheight,keepaspectratio]{index_files/mediabag/imagenes_cap5/5.2-rectangulos-masmenos.pdf}
\end{center}

Una integral definida puede interpretarse entonces como un \textbf{área
neta}, es decir, una diferencia de áreas

\[
\int_a^b f(x)\,dx=A_1-A_2,
\]

siendo \(A_1\) el área de la región arriba del eje \(x\) y debajo de la
gráfica de \(f\), y \(A_2\) corresponde a la región debajo del eje \(x\)
y arriba de la gráfica de \(f\).

\begin{center}
\includegraphics[width=0.5\linewidth,height=\textheight,keepaspectratio]{index_files/mediabag/imagenes_cap5/5.2-grafica-masmenos.pdf}
\end{center}

Teorema

Si \(f\) es continua en \([a, b]\), o si \(f\) tiene únicamente un
número finito de saltos discontinuos en el intervalo, entonces \(f\) es
integrable en \([a, b]\). Es decir, la integral definida

\[
\int_a^b f(x)\,dx
\] existe.

Si \(f\) es integrable en \([a,b]\), entonces el límite

\[
\lim_{n\to \infty} \sum_{i=1}^n f(x_i*)\Delta x
\]

existe independientemente de la elección de los puntos muestra
\(x_i^*\). En muchas ocasiones solemos considerar \(x_i^*=x_i\), es
decir, como el extremo derecho de cada subintervalo, para simplificar
los cálculos. Esto nos permite enunciar el siguiente resultado.

\phantomsection\label{teo-puntosder}
Teorema

Si \(f\) es integrable en \([a, b]\), entonces

\[
\int_a^b f(x)\,dx=\lim_{n\to \infty} \sum_{i=1}^n f(x_i)\Delta x,
\]

donde

\[
\Delta x = \frac{b-a}{n} \quad \text{ y }\quad x_i=a+i\Delta x, \quad 0\leq i\leq n.
\]

Las siguientes fórmulas resultan de utilidad para el cálculo de
integrales por definición.

\phantomsection\label{form-sumas}
\[
\begin{aligned}
\sum_{i=1}^n i &= \frac{n(n+1)}{2} \quad \text{ (suma de números consecutivos) }\\
\\
\sum_{i=1}^n i^2 &= \frac{n(n+1)(2n+1)}{6} \quad \text{ (suma de cuadrados de números consecutivos) }\\
\\
\sum_{i=1}^n i^3 &= \left[\frac{n(n+1)}{2}\right]^2 \quad \text{ (suma de cubos de números consecutivos) }
\end{aligned}
\]

Ejemplo

Evaluar \(\displaystyle \int_0^3 (x^3-6x)\,dx\) utilizando la definición
de integral.

\begin{tcolorbox}[enhanced jigsaw, opacityback=0, left=2mm, opacitybacktitle=0.6, toprule=.15mm, bottomrule=.15mm, colbacktitle=quarto-callout-tip-color!10!white, titlerule=0mm, bottomtitle=1mm, title=\textcolor{quarto-callout-tip-color}{\faLightbulb}\hspace{0.5em}{Solución}, toptitle=1mm, colframe=quarto-callout-tip-color-frame, colback=white, breakable, rightrule=.15mm, arc=.35mm, leftrule=.75mm, coltitle=black]

Dado \(n\in\mathbb{N}\), definimos \(\Delta x=(3-0)/n=3/n\) y
\(x_i=0+i\Delta x=3i/n\), para \(1\leq i\leq n\).

Calculamos primero

\[
\begin{aligned}
R_n=\sum_{i=1}^n f(x_i)\Delta x=\sum_{i=1}^n \left[\left(\frac{3i}{n}\right)^3-6\frac{3i}{n}\right]\frac{3}{n}&=\frac{81}{n^4}\sum_{i=1}^n i^3-\frac{54}{n^2}\sum_{i=1}^n i\\
\\
&=\frac{81}{n^4} \left[\frac{n(n+1)}{2}\right]^2-\frac{54}{n^2}\frac{n(n+1)}{2}\\
\\
&=\frac{81}{4} \left(\frac{n+1}{n}\right)^2-27\frac{n+1}{n},
\end{aligned}
\]

donde hemos usado los resultados de la \hyperref[form-sumas]{fórmula
para las sumas} de arriba. Por el \hyperref[teo-puntosder]{teorema}
anterior, tenemos que

\[
\int_0^3 (x^3-6x)\,dx=\lim_{n\to \infty} \sum_{i=1}^n f(x_i)\Delta x= \lim_{n\to \infty} \left[\frac{81}{4} \left(\frac{n+1}{n}\right)^2-27\frac{n+1}{n}\right]=\frac{81}{4} \cdot 1^2-27\cdot 1=-\frac{81}{12}.
\]

\begin{quote}
{[}Gráfico{]}
\end{quote}

\end{tcolorbox}

Ejemplo

Evaluar las siguientes integrales, interpretándolas como un área.

\begin{enumerate}
\def\labelenumi{\arabic{enumi}.}
\item
  \(\displaystyle \int_0^1 \sqrt{1-x^2}\,dx\).
\item
  \(\displaystyle \int_0^3 (x-1)\,dx\).
\end{enumerate}

\begin{tcolorbox}[enhanced jigsaw, opacityback=0, left=2mm, opacitybacktitle=0.6, toprule=.15mm, bottomrule=.15mm, colbacktitle=quarto-callout-tip-color!10!white, titlerule=0mm, bottomtitle=1mm, title=\textcolor{quarto-callout-tip-color}{\faLightbulb}\hspace{0.5em}{Solución}, toptitle=1mm, colframe=quarto-callout-tip-color-frame, colback=white, breakable, rightrule=.15mm, arc=.35mm, leftrule=.75mm, coltitle=black]

Para la primera escribimos \(y=\sqrt{1-x^2}\), de donde se tiene la
relación \(x^2+y^2=1\). Como \(y\geq 0\), entonces el gráfico de \(f\)
es la mitad superior de la circunferencia con centro en \((0,0)\) y
radio 1. Como la integral es entre 0 y 1, el área debajo del gráfico y
sobre el eje \(x\) corresponde al de un cuarto de la circunferencia. Por
lo tanto

\[
\int_0^1 \sqrt{1-x^2}\,dx = \frac{1}{4} \pi \cdot 1^2 =\frac{\pi}{4}.
\]

\includegraphics[width=0.45\linewidth,height=\textheight,keepaspectratio]{index_files/mediabag/imagenes_cap5/5.2-ejemplo-circunferencia.pdf}

Para la segunda función, observemos que \(y(x)=x-1\) es no positiva para
\(0\leq x\leq 1\) y no negativa para \(1\leq x\leq 3\). Entonces la
integral resulta ser una diferencia de áreas, la del triángulo por
arriba del eje \(x\) menos la del que se encuentra por debajo. Es decir

\[
\int_0^3 (x-1)\,dx = A_1-A_2=\frac{2\cdot 2}{2}-\frac{1\cdot 1}{2}=2-\frac{1}{2}=\frac{3}{2}. 
\]

\includegraphics[width=0.4\linewidth,height=\textheight,keepaspectratio]{index_files/mediabag/imagenes_cap5/5.2-ejemplo-recta.pdf}

\end{tcolorbox}

\subsection{Propiedades de la integral
definida}\label{propiedades-de-la-integral-definida}

Al definir la integral \(\int_a^b f(x)\,dx\) hemos supuesto que \(a<b\),
pues trabajamos con el intervalo \([a,b]\). Sin embargo, las sumas de
Riemann pueden calcularse aún cuando \(a>b\), y en tal caso \(\Delta x\)
cambia de signo. Por este motivo, resulta que

\[
\int_b^a f(x)\,dx=-\int_a^b f(x)\,dx.
\]

Por otro lado, si \(a=b\) entonces \(\Delta x=0\) y en consecuencia

\[
\int_a^a f(x)\,dx=0.
\]

En el siguiente teorema enunciamos otras propiedades adicionales de la
integral.

\phantomsection\label{teo-propiedades-integraldef}
Propiedades de la integral definida

Sean \(f\) y \(g\) funciones continuas y \(c\) un número real. Entonces

\begin{enumerate}
\def\labelenumi{\arabic{enumi}.}
\item
  \(\displaystyle \int_a^b c\,dx =c(b-a)\).
\item
  \(\displaystyle \int_a^b (f(x)+g(x))\,dx =\int_a^b f(x)\,dx+\int_a^b g(x)\,dx\).
\item
  \(\displaystyle \int_a^b (cf(x))\,dx=c\int_a^b f(x)\,dx\).
\item
  \(\displaystyle \int_a^b (f(x)-g(x))\,dx =\int_a^b f(x)\,dx-\int_a^b g(x)\,dx\).
\item
  \(\displaystyle \int_a^b f(x)\,dx  =\int_a^c f(x)\,dx+\int_c^b f(x)\,dx\),
  para cualesquiera \(a,b\) y \(c\) en \(\mathbb{R}\).
\end{enumerate}

Ejemplo

Evaluar la integral \(\displaystyle \int_0^1 (4+3x^2)\,dx\) utilizando
propiedades.

\begin{tcolorbox}[enhanced jigsaw, opacityback=0, left=2mm, opacitybacktitle=0.6, toprule=.15mm, bottomrule=.15mm, colbacktitle=quarto-callout-tip-color!10!white, titlerule=0mm, bottomtitle=1mm, title=\textcolor{quarto-callout-tip-color}{\faLightbulb}\hspace{0.5em}{Solución}, toptitle=1mm, colframe=quarto-callout-tip-color-frame, colback=white, breakable, rightrule=.15mm, arc=.35mm, leftrule=.75mm, coltitle=black]

Utilizando las \hyperref[teo-propiedades-integraldef]{propiedades} (2),
(1) y (3) de la integral definida, obtenemos

\[
\int_0^1 (4+3x^2)\,dx = \int_0^1 4\,dx+\int_0^1 (3x^2)\,dx=4(1-0)+3\int_0^1 x^2\,dx=3+3\cdot \frac{1}{3}=4, 
\]

donde hemos usado también que
\(\displaystyle \int_0^1 x^2\,dx=\frac{1}{3}\), resultado que obtuvimos
en el Ejemplo 1 de la \hyperref[seccion_5.1]{Sección 5.1}.

\end{tcolorbox}

A continuación veremos algunas propiedades de comparación de integrales
definidas.

\phantomsection\label{teo-comparacion-integraldef}
Propiedades de comparación de la integral definida

Sean \(f\) y \(g\) funciones integrables en \([a,b]\), con \(a\leq b\).

\begin{enumerate}
\def\labelenumi{\arabic{enumi}.}
\item
  Si \(f(x)\geq 0\) en \([a,b]\), entonces
  \(\displaystyle\int_a^b f(x)\,dx\geq 0\).
\item
  Si \(f(x)\geq g(x)\) en \([a,b]\), entonces
  \(\displaystyle \int_a^b f(x)\,dx \geq \int_a^b g(x)\,dx\).
\item
  Si \(m\leq f(x)\leq M\) en \([a,b]\), entonces
  \(\displaystyle m(b-a)\leq \int_a^b f(x)\,dx \leq M(b-a)\).
\end{enumerate}

Ejemplo

Estimar el valor de la integral \(\displaystyle \int_0^1 e^{-x^2}\,dx\)
utilizando propiedades de comparación.

\begin{tcolorbox}[enhanced jigsaw, opacityback=0, left=2mm, opacitybacktitle=0.6, toprule=.15mm, bottomrule=.15mm, colbacktitle=quarto-callout-tip-color!10!white, titlerule=0mm, bottomtitle=1mm, title=\textcolor{quarto-callout-tip-color}{\faLightbulb}\hspace{0.5em}{Solución}, toptitle=1mm, colframe=quarto-callout-tip-color-frame, colback=white, breakable, rightrule=.15mm, arc=.35mm, leftrule=.75mm, coltitle=black]

Observemos que \(f(x)=e^{-x^2}\) es función decreciente en \([0,1]\).
Por lo tanto su máximo absoluto se alcanza en el extremo izquierdo
\(x=0\) y su mínimo absoluto en el extremo derecho, \(x=1\). De esta
manera

\[
e^{-1}\leq e^{-x^2}\leq 1 \quad \text{ para todo } x\in [0,1].
\]

Utilizando la \hyperref[teo-comparacion-integraldef]{propiedad} (3) de
comparación resulta

\[
e^{-1}(1-0)\leq \int_0^1 e^{-x^2}\,dx\leq 1(1-0).
\]

Dado que \(e^{-1}\approx 0.3679\), obtenemos la estimación

\[
0.367\leq \int_0^1 e^{-x^2}\,dx\leq 1.
\]

\end{tcolorbox}

\hyperref[seccion_5.2]{↑ Volver al inicio de la sección}

\section{5.3. El teorema fundamental del cálculo}\label{seccion_5.3}

En esta sección veremos cómo se relacionan dos problemas en apariencia
muy distintos: el de encontrar la recta tangente a una curva dada y el
de calcular el área de cierta región general en el plano.

Comencemos considerando una función continua \(f\) en un intervalo
\([a,b]\), y definimos la función

\[
g(x)=\int_a^x f(t)\,dt.
\]

Notemos que si \(a\leq x\leq b\) está fijo, entonces
\(\displaystyle \int_a^x f(t)\,dt\) es un número, que depende del valor
que tome \(x\). Si dejamos que \(x\) se mueva sobre todo el intervalo,
obtenemos la función \(g\). Notemos también que hemos utilizado \(t\)
como variable de integración, para no confundir con \(x\) que es el
límite superior en la integral.

En el caso particular en que \(f\) sea no negativa en el intervalo
\([a,b]\), \(g(x)\) da el área de la región debajo del gráfico de \(f\),
entre \(a\) y \(x\).

\begin{center}
\includegraphics[width=0.5\linewidth,height=\textheight,keepaspectratio]{index_files/mediabag/imagenes_cap5/5.3-area-hasta-x.pdf}
\end{center}

Consideremos ahora un ejemplo concreto. Si \(f(x)=x\) sobre el intervalo
\([0,1]\), entonces podemos verificar usando la definición de integral
que

\[
g(x)=\int_0^x t\,dt= \frac{x^2-0^2}{2}=\frac{1}{2}x^2.
\]

Entonces notamos que \(g'(x)=x=f(x)\), es decir, la función \(g\) tiene
como derivada a la función \(f\) de la que partimos. En este caso
decimos que \(g\) es una \textbf{antiderivada} de \(f\).

\begin{tcolorbox}[enhanced jigsaw, opacityback=0, left=2mm, opacitybacktitle=0.6, toprule=.15mm, bottomrule=.15mm, colbacktitle=quarto-callout-note-color!10!white, titlerule=0mm, bottomtitle=1mm, title=\textcolor{quarto-callout-note-color}{\faInfo}\hspace{0.5em}{Definición (Antiderivada o primitiva)}, toptitle=1mm, colframe=quarto-callout-note-color-frame, colback=white, breakable, rightrule=.15mm, arc=.35mm, leftrule=.75mm, coltitle=black]

Dada una función \(f\), decimos que la función \(F\) es una
\textbf{antiderivada} o \textbf{primitiva} de \(f\) en un intervalo
\(I\) si \(F'(x)=f(x)\) para todo \(x\) del intervalo.

\end{tcolorbox}

Notemos que si \(F\) es una antiderivada de una función \(f\), entonces
\(G(x)=F(x)+C\) también es una antiderivada de \(f\) para cualquier
constante \(C\), pues

\[
G'(x)=(F(x)+C)'=F'(x)+C'=f(x)+0=f(x).
\]

Si \(F\) es cualquier antiderivada de \(f\) en el intervalo \(I\),
entonces \(F(x)+C\) se denomina \textbf{antiderivada general} de \(f\)
en \(I\).

Volviendo al ejemplo anterior, ¿es una casualidad que \(g'=f\)? La
respuesta es no, y esta importante relación forma parte del teorema
fundamental del cálculo que enunciaremos a continuación.

\phantomsection\label{teo-TFC1}
Teorema fundamental del cálculo - Parte I

Si \(f\) es continua en \([a,b]\), entonces la función \(g\) definida
por \[
g(x) = \int_a^x f(t) \, dt \quad \text{ para } a \leq x \leq b,
\] es continua en \([a,b]\), derivable en \((a,b)\), y además
\(g'(x) = f(x)\) para todo \(x\in (a,b)\).

Con la notación de Leibniz de los diferenciales, podemos reescribir la
igualdad \(g'(x)=f(x)\) como

\[
\frac{d}{dx}\left(\int_a^x f(t)\,dt\right)=f(x).
\]

Observar que la primera parte del teorema fundamental del cálculo puede
ser reescrita de la siguiente manera: si \(f\) es continua en \([a,b]\),
entonces la función

\[
g(x)=\int_a^x f(t)\,dt
\]

resulta ser una antiderivada de \(f\) en \((a,b)\). Como consecuencia,
inferimos que \textbf{toda función continua tiene una antiderivada},
hecho que puede resultar útil en muchas ocasiones.

En la siguiente tabla podemos encontrar las antiderivadas de algunas
funciones básicas que aparecen con frecuencia en la práctica.

\phantomsection\label{tabla-antiderivadas}
\begin{longtable}[]{@{}ll@{}}
\toprule\noalign{}
Función \(f\) & Antiderivada \(F\) \\
\midrule\noalign{}
\endhead
\bottomrule\noalign{}
\endlastfoot
\(k\) & \(kx+C\) \\
\(x^n\), \(n\neq -1\) & \(\displaystyle\frac{x^{n+1}}{n+1}+C\) \\
\(\displaystyle\frac{1}{x}\) & \(\ln|x|+C\) \\
\(e^x\) & \(e^x+C\) \\
\(a^x\), \(a>0\) & \(\displaystyle\frac{a^x}{\ln a} + C\) \\
\(\sin x\) & \(-\cos x + C\) \\
\(\cos x\) & \(\sin x + C\) \\
\(\sec^2 x\) & \(\tan x + C\) \\
\(\csc^2 x\) & \(-\cot x + C\) \\
\(\sec x \tan x\) & \(\sec x + C\) \\
\(\csc x \cot x\) & \(-\csc x + C\) \\
\(\displaystyle \frac{1}{\sqrt{1 - x^2}}\) & \(\arcsin x + C\) \\
\(\displaystyle\frac{-1}{\sqrt{1 - x^2}}\) & \(\arccos x + C\) \\
\(\displaystyle \frac{1}{1 + x^2}\) & \(\arctan x + C\) \\
\end{longtable}

Ejemplo

Calcular la derivada de la función
\(\displaystyle g(x)=\int_0^x \sqrt{1+t^2}\,dt\).

\begin{tcolorbox}[enhanced jigsaw, opacityback=0, left=2mm, opacitybacktitle=0.6, toprule=.15mm, bottomrule=.15mm, colbacktitle=quarto-callout-tip-color!10!white, titlerule=0mm, bottomtitle=1mm, title=\textcolor{quarto-callout-tip-color}{\faLightbulb}\hspace{0.5em}{Solución}, toptitle=1mm, colframe=quarto-callout-tip-color-frame, colback=white, breakable, rightrule=.15mm, arc=.35mm, leftrule=.75mm, coltitle=black]

Dado que \(f(t)=\sqrt{1+t^2}\) es continua en \(\mathbb{R}\), en
particular lo es en cualquier intervalo cerrado, y utilizando la primera
parte del \hyperref[teo-TFC1]{teorema fundamental del cálculo} resulta

\[
\frac{dg}{dx}(x)=\frac{d}{dx}\left(\int_0^x \sqrt{1+t^2}\,dt\right)=f(x)=\sqrt{1+x^2}.
\]

\end{tcolorbox}

\phantomsection\label{teo-TFC2}
Teorema fundamental del cálculo - Parte II

Si \(f\) es continua en \([a,b]\), entonces

\[
\int_a^b f(x)\,dx=F(b)-F(a),
\]

donde \(F\) es una antiderivada de \(f\) en \((a,b)\), es decir,
\(F'(x)=f(x)\) para \(x\in (a,b)\).

Esta segunda parte del teorema fundamental del cálculo establece que si
se conoce una antiderivada \(F\) de \(f\) en \((a,b)\), entonces la
integral definida \(\displaystyle \int_a^b f(x)\,dx\) puede calcularse
simplemente calculando la diferencia de los valores de la función \(F\)
en los extremos del intervalo. A este teorema también se lo conoce
comúnmente como \textbf{regla de Barrow}. En muchas ocasiones
escribiremos

\[
F(x)\Big{]}_a^{b}=F(x)\Big{]}_{x=a}^{x=b}
\]

para denotar a la diferencia \(F(b)-F(a)\).

Ejemplo

Evaluar la integral \(\displaystyle \int_1^3 e^x\,dx\).

\begin{tcolorbox}[enhanced jigsaw, opacityback=0, left=2mm, opacitybacktitle=0.6, toprule=.15mm, bottomrule=.15mm, colbacktitle=quarto-callout-tip-color!10!white, titlerule=0mm, bottomtitle=1mm, title=\textcolor{quarto-callout-tip-color}{\faLightbulb}\hspace{0.5em}{Solución}, toptitle=1mm, colframe=quarto-callout-tip-color-frame, colback=white, breakable, rightrule=.15mm, arc=.35mm, leftrule=.75mm, coltitle=black]

Sabemos que \(f(x)=e^x\) es continua en \(\mathbb{R}\), con lo cual es
continua en \([1,3]\). Además, como \((e^x)'=e^x\), la misma \(f\) es
una antiderivada de sí misma. Utilizando la segunda parte del
\hyperref[teo-TFC2]{teorema fundamental del cálculo} obtenemos

\[
\int_1^3 e^x\,dx=e^x \Big{]}_{1}^{3} = e^3-e^1=e^3-e.
\]

\end{tcolorbox}

Ejemplo

Encontrar la derivada de la función
\(\displaystyle \int_1^{x^4} \sec t\,dt\).

\begin{tcolorbox}[enhanced jigsaw, opacityback=0, left=2mm, opacitybacktitle=0.6, toprule=.15mm, bottomrule=.15mm, colbacktitle=quarto-callout-tip-color!10!white, titlerule=0mm, bottomtitle=1mm, title=\textcolor{quarto-callout-tip-color}{\faLightbulb}\hspace{0.5em}{Solución}, toptitle=1mm, colframe=quarto-callout-tip-color-frame, colback=white, breakable, rightrule=.15mm, arc=.35mm, leftrule=.75mm, coltitle=black]

Observar que \(f(t)=\sec t\) es continua en todo \(\mathbb{R}\) excepto
donde \(\cos t=0\), es decir, \(t=(2n+1)\pi/2\), con \(n\in\mathbb{Z}\).
Como \(1\in (-\pi/2,\pi/2)\), vamos a suponer que \(x^4\) también se
encuentra en este intervalo, para poder integrar una función continua.
Es decir, supongamos que \(x\in (-\sqrt[4]{\pi/2},\sqrt[4]{\pi/2})\).
Por la primera parte del teorema fundamental del cálculo sabemos que
existe \(F\) una antiderivada de \(f\). Aplicando la segunda parte

\[
\int_1^{x^4} \sec t\,dt=F(t)\Big{]}_{1}^{x^4}=F(x^4)-F(1).
\]

Ahora bien, como queremos calcular la derivada de esta integral,
utilizamos la regla de la cadena y el hecho de que \(F'=f\) para
escribir

\[
\frac{d}{dx}\left(\int_1^{x^4} \sec t\,dt\right)=\frac{d}{dx}\left(F(x^4)-F(1)\right)=F'(x^4) (x^4)'-0=4x^3 \,\sec(x^4).
\]

\end{tcolorbox}

Ejemplo

Calcular \(\displaystyle \int_3^{6} \frac{1}{x}\,dx\).

\begin{tcolorbox}[enhanced jigsaw, opacityback=0, left=2mm, opacitybacktitle=0.6, toprule=.15mm, bottomrule=.15mm, colbacktitle=quarto-callout-tip-color!10!white, titlerule=0mm, bottomtitle=1mm, title=\textcolor{quarto-callout-tip-color}{\faLightbulb}\hspace{0.5em}{Solución}, toptitle=1mm, colframe=quarto-callout-tip-color-frame, colback=white, breakable, rightrule=.15mm, arc=.35mm, leftrule=.75mm, coltitle=black]

Dado que \(f(x)\) es continua en \([3,6]\), podemos aplicar la segunda
parte del \hyperref[teo-TFC2]{teorema fundamental del cálculo}. Según
vimos en la \hyperref[tabla-antiderivadas]{tabla de antiderivadas}, una
antiderivada para \(f\) es \(F(x)=\ln |x|\), con lo cual

\[
\int_3^{6} \frac{1}{x}\,dx= \ln x\Big{]}_{3}^{6}=\ln 6-\ln 3=\ln 2.
\]

\end{tcolorbox}

Ejemplo

Calcular el área debajo de la curva \(f(x)=\cos x\), desde \(x=0\) hasta
\(x=b\), con \(0\leq b\leq \pi/2\).

\begin{tcolorbox}[enhanced jigsaw, opacityback=0, left=2mm, opacitybacktitle=0.6, toprule=.15mm, bottomrule=.15mm, colbacktitle=quarto-callout-tip-color!10!white, titlerule=0mm, bottomtitle=1mm, title=\textcolor{quarto-callout-tip-color}{\faLightbulb}\hspace{0.5em}{Solución}, toptitle=1mm, colframe=quarto-callout-tip-color-frame, colback=white, breakable, rightrule=.15mm, arc=.35mm, leftrule=.75mm, coltitle=black]

Sea \(0\leq b\leq \pi/2\). Dado que \(f(x)\geq 0\) en \([0,\pi/2]\), el
área pedida está dada por la integral definida

\[
\int_0^b \cos x\,dx.
\]

Una \hyperref[tabla-antiderivadas]{antiderivada} para \(f\) es
\(F(x)=\operatorname{sen} x\), con lo cual

\[
\int_0^b \cos x\,dx=\operatorname{sen} x\Big{]}_{0}^b=\operatorname{sen} b-\operatorname{sen} 0=\operatorname{sen} b.
\]

\end{tcolorbox}

Ejemplo

¿Por qué es incorrecto el siguiente cálculo?

\[
\int_{-1}^3 \frac{1}{x^2}\,dx = \left(-\frac{1}{x}\right)\Bigg{]}_{-1}^{3}=-\frac{1}{3}-\left(-\frac{1}{-1}\right)=-\frac{4}{3}.
\]

\begin{tcolorbox}[enhanced jigsaw, opacityback=0, left=2mm, opacitybacktitle=0.6, toprule=.15mm, bottomrule=.15mm, colbacktitle=quarto-callout-tip-color!10!white, titlerule=0mm, bottomtitle=1mm, title=\textcolor{quarto-callout-tip-color}{\faLightbulb}\hspace{0.5em}{Solución}, toptitle=1mm, colframe=quarto-callout-tip-color-frame, colback=white, breakable, rightrule=.15mm, arc=.35mm, leftrule=.75mm, coltitle=black]

Es incorrecto aplicar la segunda parte del \hyperref[teo-TFC2]{teorema
fundamental del cálculo}, puesto que la función \(f(x)=\frac{1}{x^2}\)
no es continua en el intervalo \([-1,3]\), ya que tiene un salto
infinito en \(x=0\). En este caso, la integral se llama
\textbf{impropia} y este tipo de situaciones lo estudiaremos más
adelante.

\end{tcolorbox}

\hyperref[seccion_5.3]{↑ Volver al inicio de la sección}

\section{5.4. Integrales indefinidas y el teorema del cambio
total}\label{seccion_5.4}

\subsection{Integrales indefinidas}\label{integrales-indefinidas}

Recordando las dos partes del teorema fundamental del cálculo, la
primera de ellas establece que toda función continua \(f\) tiene una
antiderivada dada por la expresión

\[
\int_a^x f(t)\,dt,
\]

mientras que la segunda asegura que el valor de la integral definida

\[
\int_a^b f(x)\,dx
\]

puede encontrarse haciendo \(F(b)-F(a)\), donde \(F\) es una
antiderivada de \(f\). Ahora vamos a introducir una notación para las
antiderivadas o primitivas de una función dada. Escribiremos
\(\displaystyle \int f(x)\,dx\) para simbolizar a una antiderivada de
\(f\), y la llamaremos \textbf{integral indefinida}. Es decir

\[
\int f(x)\,dx=F(x) \quad \text{ significa que }\quad F'(x)=f(x).
\]

Por ejemplo,

\[
\int x^2\,dx =\frac{1}{3}x^3+C
\] siendo \(C\) cualquier constante real, dado que al sumar cualquier
constante a una antiderivada de \(f\) obtenemos otra antiderivada (pues
la derivada de las constantes es cero).

\begin{tcolorbox}[enhanced jigsaw, opacityback=0, left=2mm, opacitybacktitle=0.6, toprule=.15mm, bottomrule=.15mm, colbacktitle=quarto-callout-caution-color!10!white, titlerule=0mm, bottomtitle=1mm, title=\textcolor{quarto-callout-caution-color}{\faFire}\hspace{0.5em}{Importante}, toptitle=1mm, colframe=quarto-callout-caution-color-frame, colback=white, breakable, rightrule=.15mm, arc=.35mm, leftrule=.75mm, coltitle=black]

No confundir una integral definida con una indefinida. Por un lado,
tenemos que \(\int_a^b f(x)\,dx\) es un número, mientras que
\(\int f(x)\,dx\) es una familia de funciones de la forma \(F(x)+C\),
siendo \(F'=f\). La relación entre ambas la da la segunda parte del
\hyperref[teo-TFC2]{teorema fundamental del cálculo}

\[
\int_a^b f(x)\,dx =\left(\int f(x)\,dx\right)\Bigg{]}_a^b.
\]

\end{tcolorbox}

Con esta nueva notación, podemos reescribir la tabla de antiderivadas de
la sección anterior como sigue.

\phantomsection\label{form-int-indefinidas}
\begin{enumerate}
\def\labelenumi{\arabic{enumi}.}
\item
  \(\displaystyle \int kf(x)\,dx=k\int f(x)\, dx\).
\item
  \(\displaystyle \int (f(x)+g(x))\,dx=\int f(x)\, dx+\int g(x)\,dx\).
\end{enumerate}

\begin{enumerate}
\def\labelenumi{\arabic{enumi}.}
\setcounter{enumi}{2}
\item
  \(\displaystyle \int k\,dx=kx+C\).
\item
  \(\displaystyle \int x^n\,dx=\frac{x^{n+1}}{n+1}+C\), \(n\neq -1\).
\item
  \(\displaystyle \int e^x\,dx=e^x+C.\)
\item
  \(\displaystyle \int a^x\,dx=\frac{a^x}{\ln a}+C\), \(a>0, a\neq 1\).
\item
  \(\displaystyle \int \operatorname{sen} x\,dx=-\cos x+C\).
\item
  \(\displaystyle \int \cos x\,dx=\operatorname{sen} x+C\).
\item
  \(\displaystyle \int \sec^2 x\,dx=\tan x+C\).
\end{enumerate}

\begin{enumerate}
\def\labelenumi{\arabic{enumi}.}
\setcounter{enumi}{9}
\item
  \(\displaystyle \int \csc^2 x\,dx=-\cot x+C\).
\item
  \(\displaystyle \int \sec x\tan x\,dx=\sec x+C\).
\item
  \(\displaystyle \int \csc x\cot x\,dx=-\csc x+C\).
\item
  \(\displaystyle \int \frac{1}{\sqrt{1-x^2}}\,dx=\operatorname{arcsen} x+C\).
\item
  \(\displaystyle \int \frac{-1}{\sqrt{1-x^2}}\,dx=\arccos x+C\).
\item
  \(\displaystyle \int \frac{1}{1+x^2}\,dx=\arctan x+C\).
\end{enumerate}

Cuando escribamos la expresión

\[
\int f(x)\,dx=F(x)+C
\]

recordemos que ésta siempre es válida en algún intervalo \(I\), que por
lo general omitimos.

Ejemplo

Evaluar
\(\displaystyle \int \frac{\cos \theta}{\operatorname{sen}^2 \theta}\,d\theta\).

\begin{tcolorbox}[enhanced jigsaw, opacityback=0, left=2mm, opacitybacktitle=0.6, toprule=.15mm, bottomrule=.15mm, colbacktitle=quarto-callout-tip-color!10!white, titlerule=0mm, bottomtitle=1mm, title=\textcolor{quarto-callout-tip-color}{\faLightbulb}\hspace{0.5em}{Solución}, toptitle=1mm, colframe=quarto-callout-tip-color-frame, colback=white, breakable, rightrule=.15mm, arc=.35mm, leftrule=.75mm, coltitle=black]

Primero escribimos

\[
\frac{\cos \theta}{\operatorname{sen}^2 \theta}=\frac{1}{\operatorname{sen} \theta} \frac{\cos \theta}{\operatorname{sen} \theta}=\csc \theta \cot \theta.
\]

Aplicando la \hyperref[form-int-indefinidas]{fórmula} (12) de
antiderivadas obtenemos

\[
\int \csc \theta \cot \theta\,d\theta=-\csc \theta + C.
\]

\end{tcolorbox}

Ejemplo

Calcular \(\displaystyle \int_0^3 (x^3-6x)\,dx\).

\begin{tcolorbox}[enhanced jigsaw, opacityback=0, left=2mm, opacitybacktitle=0.6, toprule=.15mm, bottomrule=.15mm, colbacktitle=quarto-callout-tip-color!10!white, titlerule=0mm, bottomtitle=1mm, title=\textcolor{quarto-callout-tip-color}{\faLightbulb}\hspace{0.5em}{Solución}, toptitle=1mm, colframe=quarto-callout-tip-color-frame, colback=white, breakable, rightrule=.15mm, arc=.35mm, leftrule=.75mm, coltitle=black]

Aplicamos la segunda parte del \hyperref[teo-TFC2]{teorema fundamental
del cálculo} y las fórmulas (2), (3) y (4) de la
\hyperref[form-int-indefinidas]{tabla de antiderivadas} para obtener

\[
\int_0^3 (x^3-6x)\,dx= \left(\frac{x^4}{4}-6\frac{x^2}{2}\right)\Bigg{]}_0^3=\frac{3^4}{4}-3\cdot 3^2 -\left(\frac{0^4}{4}-3\cdot 0^2\right)=\frac{81}{4}-27=-\frac{81}{12}.
\]

Notar que a este mismo valor habíamos llegado utilizando la definición
de integral, en el Ejemplo 3 de la \hyperref[seccion_5.2]{Sección 5.2}.

\end{tcolorbox}

Ejemplo

Evaluar la integral
\(\displaystyle \int_1^9 \frac{2t^2+t^2\sqrt{t}-1}{t^2}\,dt\).

\begin{tcolorbox}[enhanced jigsaw, opacityback=0, left=2mm, opacitybacktitle=0.6, toprule=.15mm, bottomrule=.15mm, colbacktitle=quarto-callout-tip-color!10!white, titlerule=0mm, bottomtitle=1mm, title=\textcolor{quarto-callout-tip-color}{\faLightbulb}\hspace{0.5em}{Solución}, toptitle=1mm, colframe=quarto-callout-tip-color-frame, colback=white, breakable, rightrule=.15mm, arc=.35mm, leftrule=.75mm, coltitle=black]

Primero notemos que

\[
\frac{2t^2+t^2\sqrt{t}-1}{t^2}=2+t^{1/2}-t^{-2}.
\]

Aplicamos las fórmulas (1), (2), (3) y (4) de la
\hyperref[form-int-indefinidas]{tabla de antiderivadas} para obtener

\[
\begin{aligned}
\int_1^9  \frac{2t^2+t^2\sqrt{t}-1}{t^2}\,dt=\int_1^9 (2+t^{1/2}-t^{-2})\,dt &=\left(2t+\frac{2}{3}t^{3/2}+\frac{1}{t}\right)\Bigg{]}_1^9\\
\\
&=\left(2\cdot 9+\frac{2}{3}9^{3/2}+\frac{1}{9}\right)-\left(2\cdot 1+\frac{2}{3}1^{3/2}+\frac{1}{1}\right)\\
\\
&=\frac{325}{9}-\frac{11}{3}=\frac{292}{9}.
\end{aligned}
\]

\end{tcolorbox}

\subsection{Aplicaciones}\label{aplicaciones}

Recordemos nuevamente que la segunda parte del teorema fundamental del
cálculo establece que si \(f\) es continua en \([a,b]\), entonces

\[
\int_a^bf(x)\,dx=F(b)-F(a),
\]

siendo \(F\) cualquier antiderivada de \(f\), es decir, cualquier
función tal que \(F'=f\). Pero entonces la igualdad de arriba puede
reescribirse como

\[
\int_a^bF'(x)\,dx=F(b)-F(a).
\]

Sabemos que \(F'\) representa la relación de cambio de \(y=F(x)\) con
respecto a \(x\), mientras que la diferencia \(F(b)-F(a)\) representa el
cambio total de la función \(F\) en el intervalo \([a,b]\). En
definitiva, la igualdad de arriba nos está diciendo que si integramos la
relación de cambio de una función sobre un intervalo \([a,b]\),
obtenemos el cambio total de la misma en el intervalo.

El siguiente teorema es una reformulación de la segunda parte del
teorema fundamental del cálculo en este contexto.

Teorema del cambio total

La integral de la relación de cambio de una función, es el cambio total
de la función. Es decir,

\[
\int_a^bF'(x)\,dx=F(b)-F(a).
\]

Si un objeto se mueve en línea recta con función de posición \(s(t)\),
entonces su velocidad es \(v(t)=s'(t)\) y en consecuencia

\[
\int_{t_1}^{t_2} v(t)\,dt=\int_{t_1}^{t_2} s'(t)\,dt=s(t_2)-s(t_1)
\]

es el \textbf{cambio de la posición} o \textbf{desplazamiento} del
objeto durante el periodo de tiempo desde \(t_1\) hasta \(t_2\).

Si queremos calcular la distancia recorrida durante el intervalo
\([t_1,t_2]\) hay que considerar los intervalos donde \(v(t)\geq 0\) (la
partícula se mueve hacia la derecha) y también aquellos donde
\(v(t)\leq 0\) (la partícula se mueve hacia la izquierda).

En ambos casos la distancia se calcula al integrar \(|v(t)|\), la
magnitud de la rapidez. Por lo tanto, la \textbf{distancia total
recorrida} es

\[
d=\int_{t_1}^{t_2} |v(t)|\,dt.
\]

En la siguiente figura vemos cómo interpretar el desplazamiento y la
distancia recorrida en términos de las áreas debajo de una curva de
velocidad.

\begin{center}
\includegraphics[width=0.5\linewidth,height=\textheight,keepaspectratio]{index_files/mediabag/imagenes_cap5/5.4-despl-y-dist.pdf}
\end{center}

El desplazamiento total en el intervalo \([t_1,t_2]\) viene dado por
\(\displaystyle\int_{t_1}^{t_2} f(t)\,dt=-A_1+A_2-A_3\), y la distancia
total recorrida es
\(\displaystyle \int_{t_1}^{t_2} |f(t)|\,dt=A_1+A_2+A_3\).

Ejemplo

Una partícula se mueve a lo largo de una recta de modo que su velocidad
en metros por segundo en el instante \(t\) es \(v(t)=t^2-t-6\).

\begin{itemize}
\item
  Encontrar el desplazamiento \(\Delta\) de la partícula durante el
  periodo \(1\leq t\leq 4\).
\item
  Hallar la distancia \(d\) recorrida durante este período.
\end{itemize}

\begin{tcolorbox}[enhanced jigsaw, opacityback=0, left=2mm, opacitybacktitle=0.6, toprule=.15mm, bottomrule=.15mm, colbacktitle=quarto-callout-tip-color!10!white, titlerule=0mm, bottomtitle=1mm, title=\textcolor{quarto-callout-tip-color}{\faLightbulb}\hspace{0.5em}{Solución}, toptitle=1mm, colframe=quarto-callout-tip-color-frame, colback=white, breakable, rightrule=.15mm, arc=.35mm, leftrule=.75mm, coltitle=black]

\begin{itemize}
\tightlist
\item
  Según vimos, el desplazamiento lo encontramos integrando la función
  velocidad.
\end{itemize}

\[
\Delta=\int_1^4 v(t)\,dt=\int_1^4 (t^2-t-6)\,dt =\left(\frac{t^3}{3}-\frac{t^2}{2}-6t\right)\Bigg{]}_1^4=-\frac{9}{2}.
\]

\begin{itemize}
\tightlist
\item
  Para hallar la distancia total, primero factorizamos
  \(t^2-t-6=(t-3)(t+2)\), con lo cual \(v\) es negativa para
  \(1\leq t<3\) y \(v(t)\geq 0\) en el intervalo \([3,4]\). Por lo tanto
\end{itemize}

\[
d=\int_1^4 |v(t)|\,dt=\int_1^3 (-v(t))\,dt+\int_3^4 v(t)\,dt=\int_1^3 (6+t-t^2)\,dt+\int_3^4 (t^2-t-6)\,dt=\frac{61}{6}.
\]

Por lo tanto, desde el primer al cuarto segundo de recorrido, la
partícula se desplazó 4.5 m hacia la izquierda, y recorrió un total de
aproximadamente 10.17 m.

\end{tcolorbox}

\hyperref[seccion_5.4]{↑ Volver al inicio de la sección}

\section{5.5. La regla de la sustitución}\label{seccion_5.5}

En esta sección veremos una forma de calcular integrales de muchas
funciones que se escapan de los casos que aparecen en las tablas vistas
hasta el momento.

Si \(F\) es una primitiva o antiderivada de \(f\) y \(g\) es una función
derivable, entonces la composición \(F\circ g\) es una primitiva de
\(f(g(x))g'(x)\). En efecto, utilizando la regla de la cadena

\[
\frac{d}{dx}\left(F(g(x))\right)=\frac{dF}{dx}(g(x))\frac{dg}{dx}(x)=F'(g(x))
g'(x)=f(g(x))g'(x).
\]

En términos de integrales indefinidas, la igualdad de arriba puede
expresarse como

\[
\int f(g(x))g'(x)\,dx=F(g(x))+C.
\]

Si llamamos \(u=g(x)\) entonces tenemos que
\(\displaystyle \frac{du}{dx}(x)=g'(x)\) y entonces podemos reescribir
lo anterior como

\[
\int f(u)\frac{du}{dx}\,dx=F(u)+C=\int f(u)\,du.
\]

En definitiva, si tenemos una integral de la forma
\(\int f(g(x))g'(x)\,dx\) podemos transformarla en la integral
\(\int f(u)\,du\) haciendo la \textbf{sustitución} o \textbf{cambio de
variable} \(u=g(x)\).

\phantomsection\label{teo-regla-sustitucion}
Regla de sustitución

Si \(u=g(x)\) es una función derivable con rango \(I\) y \(f\) es
continua sobre \(I\), entonces

\[
\int f(g(x))g'(x)\,dx=\int f(u)\,du
\]

Por ejemplo, si queremos calcular

\[
\int 2x \sqrt{1+x^2}\, dx
\]

escribimos \(f(x)=\sqrt{x}\) y \(g(x)=1+x^2\), con lo que la integral
anterior resulta

\[
\int f(g(x))g'(x)\, dx.
\]

En vistas de la \hyperref[teo-regla-sustitucion]{regla de sustitución},
si \(u=g(x)=1+x^2\) podemos escribir

\[
\begin{aligned}
\int 2x \sqrt{1+x^2}\, dx=\int f(g(x))g'(x)\, dx=\int f(u)\,du=\int \sqrt{u}\,du&=\frac{2}{3}u^{3/2}+C\\
\\
&=\frac{2}{3}(1+x^2)^{3/2}+C.
\end{aligned}
\]

Una manera común de denotar esta sustitución en la práctica es la
siguiente: si \(u=1+x^2\), entonces \(du=g'(x)\,dx=2x\,dx\), entonces
sustituyendo directamente en la integral definida resulta

\[
\int 2x \sqrt{1+x^2}\, dx=\int \sqrt{u}\,du.
\]

Ejemplo

Calcular la integral indefinida \[
\int x^3 \cos(x^4+2)\, dx.
\]

\begin{tcolorbox}[enhanced jigsaw, opacityback=0, left=2mm, opacitybacktitle=0.6, toprule=.15mm, bottomrule=.15mm, colbacktitle=quarto-callout-tip-color!10!white, titlerule=0mm, bottomtitle=1mm, title=\textcolor{quarto-callout-tip-color}{\faLightbulb}\hspace{0.5em}{Solución}, toptitle=1mm, colframe=quarto-callout-tip-color-frame, colback=white, breakable, rightrule=.15mm, arc=.35mm, leftrule=.75mm, coltitle=black]

Haciendo la sustitución \(u = x^4 + 2\) resulta que \(du=4x^3\,dx\), o
equivalentemente \(\displaystyle  x^3\,dx = \tfrac{1}{4}\,du\). De esta
manera

\[
\int x^3 \cos(x^4+2)\, dx = \tfrac{1}{4} \int \cos(u)\,du=\operatorname{sen} u+C=\operatorname{sen}(x^4+2)+C.
\]

\end{tcolorbox}

Ejemplo

Calcular \(\displaystyle \int \frac{x}{\sqrt{1-4x^2}}\,dx\).

\begin{tcolorbox}[enhanced jigsaw, opacityback=0, left=2mm, opacitybacktitle=0.6, toprule=.15mm, bottomrule=.15mm, colbacktitle=quarto-callout-tip-color!10!white, titlerule=0mm, bottomtitle=1mm, title=\textcolor{quarto-callout-tip-color}{\faLightbulb}\hspace{0.5em}{Solución}, toptitle=1mm, colframe=quarto-callout-tip-color-frame, colback=white, breakable, rightrule=.15mm, arc=.35mm, leftrule=.75mm, coltitle=black]

En este caso parece razonable sustituir el polinomio dentro de la raíz
por una nueva variable, ya que su derivada será de la forma \(ax\), y
aparece en el numerador. Escribiendo \(u=1-4x^2\) tenemos que
\(du=-8x\,dx\). Esta expresión no coincide exactamente con el numerador,
pero podemos hacerla aparecer multiplicando y dividiendo por \(-8\). En
definitiva

\[
\begin{aligned}
\int \frac{x}{\sqrt{1-4x^2}}\,dx=-\frac{1}{8}\int \frac{-8x}{\sqrt{1-4x^2}}\,dx=-\frac{1}{8}\int \frac{du}{\sqrt{u}}=-\frac{1}{8} \int u^{-1/2}\,du& =-\frac{1}{8} 2u^{1/2}+C\\
\\
&=-\frac{1}{4}\sqrt{1-4x^2}+C.
\end{aligned}
\]

\end{tcolorbox}

Ejemplo

Evaluar \(\displaystyle \int \sqrt{1+x^2}\,x^5\,dx\).

\begin{tcolorbox}[enhanced jigsaw, opacityback=0, left=2mm, opacitybacktitle=0.6, toprule=.15mm, bottomrule=.15mm, colbacktitle=quarto-callout-tip-color!10!white, titlerule=0mm, bottomtitle=1mm, title=\textcolor{quarto-callout-tip-color}{\faLightbulb}\hspace{0.5em}{Solución}, toptitle=1mm, colframe=quarto-callout-tip-color-frame, colback=white, breakable, rightrule=.15mm, arc=.35mm, leftrule=.75mm, coltitle=black]

Si llamamos \(u=1+x^2\), resulta que \(du=2x\,dx\). Como el factor que
acompaña a la raíz es \(x^5\), reescribimos
\(x^5=x^4\cdot x=(x^2)^2\cdot x\), con lo cual

\[
\begin{aligned}
\int \sqrt{1+x^2}\,x^5\,dx=\int\sqrt{1+x^2}\,(x^2)^2\cdot x \,dx
&=\frac{1}{2}\int \sqrt{u} (u-1)^2\,du\\
\\
&=\frac{1}{2}\int \sqrt{u}(u^2-2u+1)\,du\\
\\
&=\frac{1}{2}\int(u^{5/2}-2u^{3/2}+u^{1/2})\,du\\
\\
&=\frac{1}{2}\left(\frac{2}{7}u^{7/2}-\frac{4}{5}u^{5/2}+\frac{2}{3}u^{3/2}\right)+C\\
\\
&=\frac{1}{7}(1+x^2)^{7/2}-\frac{2}{5}(1+x^2)^{5/2}+\frac{1}{3}(1+x^2)^{3/2}+C.
\end{aligned}
\]

\end{tcolorbox}

Ejemplo

Calcular \(\displaystyle \int \tan x\,dx\).

\begin{tcolorbox}[enhanced jigsaw, opacityback=0, left=2mm, opacitybacktitle=0.6, toprule=.15mm, bottomrule=.15mm, colbacktitle=quarto-callout-tip-color!10!white, titlerule=0mm, bottomtitle=1mm, title=\textcolor{quarto-callout-tip-color}{\faLightbulb}\hspace{0.5em}{Solución}, toptitle=1mm, colframe=quarto-callout-tip-color-frame, colback=white, breakable, rightrule=.15mm, arc=.35mm, leftrule=.75mm, coltitle=black]

Si escribimos
\(\displaystyle \tan x=\frac{\operatorname{sen} x}{\cos x}\), entonces
podemos hacer la sustitución \(u=\cos x\), con lo cual
\(du=-\operatorname{sen} x\,dx\). Por lo tanto

\[
\int \tan x\,dx=\int \frac{\operatorname{sen} x}{\cos x}\,dx =-\int \frac{du}{u}=-\ln|u|+C=-\ln|\cos x|+C=\ln|\sec x|+C.
\]

\end{tcolorbox}

\subsection{Sustitución para integrales
definidas}\label{sustituciuxf3n-para-integrales-definidas}

Cuando trabajemos con integrales definidas también es posible aplicar la
regla de sustitución. Por ejemplo, si quisiéramos calcular

\[
\int_0^1 2x\sqrt{1+x^2}\,dx,
\]

hemos visto que

\[
\int 2x\sqrt{1+x^2}\,dx=\frac{2}{3}(1+x^2)^{3/2}+C
\]

por lo que utilizando el teorema fundamental del cálculo, resultaría

\[
\int_0^1 2x\sqrt{1+x^2}\,dx=\left(\frac{2}{3}(1+x^2)^{3/2}\right)\Bigg{]}_{0}^1=\frac{2}{3}(1+1^2)^{3/2}-\frac{2}{3}(1+0^2)^{3/2}=\frac{2}{3}(2\sqrt{2}-1).
\]

Sin embargo, si queremos hacer todo en un sólo cálculo, es preciso
\textbf{cambiar} los extremos de integración en función de la
sustitución hecha, como establece la siguiente regla.

\phantomsection\label{teo-sustitucion-definida}
Regla de sustitución para integrales definidas

Si \(g'\) es continua en \([a,b]\) y \(f\) es continua sobre el rango de
\(g\), entonces

\[
\int_a^b f(g(x))g'(x)\,dx=\int_{g(a)}^{g(b)} f(u)\,du.
\]

Si calculamos ahora \(\displaystyle \int_0^1 2x\sqrt{1+x^2}\,dx\) usando
esta regla, como \(u=g(x)=1+x^2\), tenemos que \(g(0)=1\) y \(g(1)=2\)

\[
\int_0^1 2x\sqrt{1+x^2}\,dx=\int_1^2 \sqrt{u}\,du =\frac{2}{3}u^{3/2}\Big{]}_1^2=\frac{2}{3}\left(2^{3/2}-1^{3/2}\right)=\frac{2}{3}\left(2\sqrt{2}-1\right).
\]

Ejemplo

Evaluar \(\displaystyle \int_0^4 \sqrt{2x+1}\,dx\).

\begin{tcolorbox}[enhanced jigsaw, opacityback=0, left=2mm, opacitybacktitle=0.6, toprule=.15mm, bottomrule=.15mm, colbacktitle=quarto-callout-tip-color!10!white, titlerule=0mm, bottomtitle=1mm, title=\textcolor{quarto-callout-tip-color}{\faLightbulb}\hspace{0.5em}{Solución}, toptitle=1mm, colframe=quarto-callout-tip-color-frame, colback=white, breakable, rightrule=.15mm, arc=.35mm, leftrule=.75mm, coltitle=black]

Hacemos el cambio de variable \(u=2x+1\), con lo cual \(du=2\,dx\).
Notar también que \(u=1\) cuando \(x=0\) y \(u=9\) cuando \(x=4\).
Utilizando la \hyperref[teo-sustitucion-definida]{regla de sustitución
para integrales definidas} obtenemos

\[
\int_0^4 \sqrt{2x+1}\,dx=\frac{1}{2}\int_1^9\sqrt{u}\,du=\frac{1}{2}\left(\frac{2}{3}u^{3/2}\right)\Bigg{]}_1^9=\frac{1}{3}(27-1)=\frac{26}{3}. 
\]

\end{tcolorbox}

\hyperref[seccion_5.5]{↑ Volver al inicio de la sección}

\section{7.1. Integración por partes}\label{seccion_7.1}

En la sección anterior explotamos la regla de la cadena para calcular
integrales utilizando un cambio de variable. En esta sección
estudiaremos otro método de integración, basado en la regla del producto
para la derivación. Recordar que si \(f\) y \(g\) son funciones
derivables, entonces

\[
(f(x)g(x))' = f'(x)g(x) + f(x)g'(x).
\]

Esta última igualdad nos dice que \(f(x)g(x)\) es una antiderivada para
la función \(f'(x)g(x) + f(x)g'(x)\) o, en término de integrales
indefinidas

\[
\int (f'(x)g(x) + f(x)g'(x))\,dx =f(x)g(x)+C.
\]

Como la integral se puede distribuir en la suma, también podemos
escribir

\[
\int f'(x)g(x)\,dx + \int f(x)g'(x)\,dx =f(x)g(x)+C, \quad \text{ o de forma equivalente } 
\]

\phantomsection\label{form-partes1}
\[
\int f(x)g'(x)\,dx = f(x)g(x) - \int f'(x)g(x)\,dx.
\]

Esta expresión se conoce como \textbf{fórmula de integración por
partes}. De manera alternativa, si denotamos \(u=f(x)\) y \(v=g(x)\)
entonces formalmente podemos escribir \(du=f'(x)\,dx\),
\(dv=g'(x)\,dx\), y así

\phantomsection\label{form-partes2}
\[
\int udv = uv - \int vdu.
\]

Este método de integración resulta útil cuando tenemos un producto de
dos funciones, siendo una de ellas la derivada de otra función adecuada.

Ejemplo

Calcular \(\displaystyle \int x\operatorname{sen} x\,dx\).

\begin{tcolorbox}[enhanced jigsaw, opacityback=0, left=2mm, opacitybacktitle=0.6, toprule=.15mm, bottomrule=.15mm, colbacktitle=quarto-callout-tip-color!10!white, titlerule=0mm, bottomtitle=1mm, title=\textcolor{quarto-callout-tip-color}{\faLightbulb}\hspace{0.5em}{Solución}, toptitle=1mm, colframe=quarto-callout-tip-color-frame, colback=white, breakable, rightrule=.15mm, arc=.35mm, leftrule=.75mm, coltitle=black]

\begin{itemize}
\item
  Vamos a proceder primero utilizando la \hyperref[form-partes1]{primera
  fórmula} vista de integración por partes. Denotamos \(f(x)=x\) y
  \(g'(x)=\operatorname{sen} x\). Por lo tanto, \(f'(x)=1\) y
  \(g(x)=-\cos x\) . De hecho, \(g\) será cualquier antiderivada de
  \(g'\): podemos escribir \[
  g(x)=\int g'(x)\,dx=\int \operatorname{sen} x\,dx=-\cos x+C,
  \] y cualquier valor fijo de \(C\) servirá. Por comodidad, en general
  tomaremos \(C=0\).

  Aplicando la fórmula de integración por partes obtenemos \[
  \begin{aligned}
  \int x\operatorname{sen} x\,dx=\int f(x)g'(x)\,dx&=f(x)g(x)-\int f'(x)g(x)\,dx\\
  \\
  &=-x\cos x-\int 1\cdot (-\cos x)\,dx=-x\cos x+\int \cos x\,dx\\
  &=-x\cos x+\operatorname{sen } x + C.
  \end{aligned}
  \]
\item
  Utilicemos ahora la \hyperref[form-partes2]{segunda fórmula} de
  integración por partes. Si llamamos \(u=x\) y
  \(dv=\operatorname{sen} x\,dx\), entonces tenemos que \(du=dx\) y
  \(v=-\cos x\), con lo cual
\end{itemize}

\[
\int u\,dv=uv-\int v\,du=-x\cos x-\int (-\cos x)\,dx=-x\cos x+\operatorname{sen} x+C.
\]

\end{tcolorbox}

Ejemplo

Encontrar \(\displaystyle \int \ln x\,dx\).

\begin{tcolorbox}[enhanced jigsaw, opacityback=0, left=2mm, opacitybacktitle=0.6, toprule=.15mm, bottomrule=.15mm, colbacktitle=quarto-callout-tip-color!10!white, titlerule=0mm, bottomtitle=1mm, title=\textcolor{quarto-callout-tip-color}{\faLightbulb}\hspace{0.5em}{Solución}, toptitle=1mm, colframe=quarto-callout-tip-color-frame, colback=white, breakable, rightrule=.15mm, arc=.35mm, leftrule=.75mm, coltitle=black]

En este caso elegimos \(u=\ln x\), con lo que \(dv\) será simplemente
\(dx\). Con esta elección, \(\displaystyle du=\frac{1}{x}\,dx\) y
\(v=x\). Por lo tanto

\[
\int \ln x\,dx=x\ln x-\int x\frac{1}{x}\,dx=x\ln x-\int\,dx=x\ln x-x+C=x(\ln x-1)+C.
\]

\end{tcolorbox}

Ejemplo

Hallar \(\displaystyle \int t^2e^t\,dt\).

\begin{tcolorbox}[enhanced jigsaw, opacityback=0, left=2mm, opacitybacktitle=0.6, toprule=.15mm, bottomrule=.15mm, colbacktitle=quarto-callout-tip-color!10!white, titlerule=0mm, bottomtitle=1mm, title=\textcolor{quarto-callout-tip-color}{\faLightbulb}\hspace{0.5em}{Solución}, toptitle=1mm, colframe=quarto-callout-tip-color-frame, colback=white, breakable, rightrule=.15mm, arc=.35mm, leftrule=.75mm, coltitle=black]

En este caso tenemos un producto de un polinomio por una exponencial.
Como las derivadas de un polinomio bajan su grado en uno y la integral
lo aumenta en uno, es conveniente elegir \(u=t^2\) y \(dv=e^t\,dt\), ya
que para el caso de la exponencial la derivada o la integral producirán
como resultado la misma función. Entonces obtenemos \(du=2t\,dt\) y
\(v=e^t\), y usando la \hyperref[form-partes2]{fórmula de integración
por partes} \begin{equation}\phantomsection\label{eq-ejemplo-polexp1}{
\int t^2e^t\,dt=t^2e^t-2\int te^t\,dt.
}\end{equation}

Vemos que aplicar el proceso de integrar por partes nos lleva a la
integral \(\int te^t\,dt\). Calculémosla antes de continuar. Volviendo a
usar la fórmula con \(u=t\) y \(dv=e^t\,dt\), de donde \(du=dt\) y
\(v=e^t\) resulta

\[
\int t e^t\,dt=te^t-\int e^t\,du=te^t-e^t+C.
\]

Volviendo a (\ref{eq-ejemplo-polexp1}) obtenemos

\[
\int t^2e^t\,dt=t^2e^t-2\left(te^t-e^t+C\right)=t^2e^t-2te^t+2e^t-2C=e^t(t^2-2t+2)+C',
\]

donde \(C'=-2C\).

\end{tcolorbox}

Ejemplo

Evaluar la integral \(\displaystyle \int e^x\operatorname{sen} x\,dx\).

\begin{tcolorbox}[enhanced jigsaw, opacityback=0, left=2mm, opacitybacktitle=0.6, toprule=.15mm, bottomrule=.15mm, colbacktitle=quarto-callout-tip-color!10!white, titlerule=0mm, bottomtitle=1mm, title=\textcolor{quarto-callout-tip-color}{\faLightbulb}\hspace{0.5em}{Solución}, toptitle=1mm, colframe=quarto-callout-tip-color-frame, colback=white, breakable, rightrule=.15mm, arc=.35mm, leftrule=.75mm, coltitle=black]

En este caso parece que no hay una elección mejor que otra, ya que tanto
\(e^x\) como \(\operatorname{sen} x\) se pueden derivar e integrar
fácilmente, produciendo funciones similares. Elijamos \(u=e^x\) y
\(dv=\operatorname{sen} x\),dx. Así, \(du=e^x\,dx\) y \(v=-\cos x\).
Aplicando la fórmula de integración por partes

\begin{equation}\phantomsection\label{eq-ejemplo-exptrig1}{
\int e^x\operatorname{sen} x\,dx=-e^x\cos x+\int \cos x e^x\,dx.
}\end{equation}

Para continuar, debemos calcular la integral \(\int \cos x\, e^x\,dx\).
Para ello, volvemos a aplicar la fórmula de partes con \(u=e^x\) y
\(dv=\cos x\,dx\), resultando \(du=e^x\,dx\) y
\(v=\operatorname{sen} x\). Entonces

\begin{equation}\phantomsection\label{eq-ejemplo-exptrig2}{
\int \cos x\, e^x\,dx=e^x\operatorname{sen} x-\int \operatorname{sen} x\,e^x\,dx.
}\end{equation}

Pareciera ser que este procedimiento no dará resultado, pues intentando
calcular la integral recaemos en otra, que a su vez depende de la
primera. Sin embargo, si llamamos \(I=\int e^x\operatorname{sen} x\,dx\)
reemplazando la ecuación (\ref{eq-ejemplo-exptrig2}) en la
(\ref{eq-ejemplo-exptrig1}) obtenemos

\[
I=-e^x\cos x+e^x \operatorname{sen} x-I,
\]

de donde

\[
2I=-e^x\cos x+e^x \operatorname{sen} x.
\]

De esta última expresión podemos despejar \(I\), y agregando la
constante arbitraria \(C\) resulta

\[
I=\frac{1}{2}\left(-e^x\cos x+e^x \operatorname{sen} x\right)+C.
\]

Por lo tanto

\[
\int e^x\operatorname{sen} x\,dx=\frac{1}{2}e^x\left(\operatorname{sen} x-\cos x\right)+C.
\]

\end{tcolorbox}

\subsection{Integración por partes en integrales
definidas}\label{integraciuxf3n-por-partes-en-integrales-definidas}

Si combinamos la \hyperref[form-partes1]{fórmula de integración por
partes} con la regla de Barrow, obtenemos una manera de calcular
integrales definidas con este método.

\phantomsection\label{form-partes-definida}
\[
\int_a^b f(x)g'(x)\,dx = f(b)g(b) - f(a)g(a) - \int_a^b f'(x)g(x)\,dx.
\]

Ejemplo

Calcular el valor de la integral
\(\displaystyle \int_0^1 \arctan x\,dx\).

\begin{tcolorbox}[enhanced jigsaw, opacityback=0, left=2mm, opacitybacktitle=0.6, toprule=.15mm, bottomrule=.15mm, colbacktitle=quarto-callout-tip-color!10!white, titlerule=0mm, bottomtitle=1mm, title=\textcolor{quarto-callout-tip-color}{\faLightbulb}\hspace{0.5em}{Solución}, toptitle=1mm, colframe=quarto-callout-tip-color-frame, colback=white, breakable, rightrule=.15mm, arc=.35mm, leftrule=.75mm, coltitle=black]

Como es más dificil integrar la función arcotangente que derivarla,
elegimos \(u=\arctan x\) y \(dv=dx\). De esta manera resultan

\[
du=\frac{1}{1+x^2}\,dx \quad \text{ y }\quad v=x.
\]

Aplicando la \hyperref[form-partes-definida]{fórmula de integración por
partes para integrales definidas} resulta

\begin{equation}\phantomsection\label{eq-ejemplo-arctan}{
\int_0^1 \arctan x\,dx=\left(x\arctan x\right)\Big{]}_{0}^1-\int_0^1 \frac{x}{1+x^2}\,dx.
}\end{equation}

Antes de terminar, calculemos la integral definida resultante utilizando
el método de sustitución. Si llamamos \(t=1+x^2\), entonces
\(dt=2x\,dx\) y los nuevos extremos de integración son \(1\) y \(2\),
respectivamente. Así

\[
\int_0^1 \frac{x}{1+x^2}\,dx=\frac{1}{2}\int_1^2 \frac{du}{u}=\frac{1}{2}\left(\ln |u|\right)\Big{]}_1^2=\frac{1}{2}(\ln 2-\ln 1)=\frac{\ln 2}{2}. 
\]

Volviendo a la ecuación (\ref{eq-ejemplo-arctan}) obtenemos

\[
\int_0^1 \arctan x\,dx=1\cdot \arctan 1-0\cdot \arctan 0-\frac{\ln 2}{2}=\frac{\pi}{4}-\frac{\ln 2}{2}=\frac{\pi-\ln 4}{4}.
\]

\end{tcolorbox}

\hyperref[seccion_7.1]{↑ Volver al inicio de la sección}

\bookmarksetup{startatroot}

\chapter{6. Aplicaciones de la
integración}\label{aplicaciones-de-la-integraciuxf3n}

\section{6.1. Áreas entre curvas}\label{seccion_6.1}

Consideremos la región \(S\) del plano que se ubica entre los gráficos
de dos funciones continuas \(y=f(x)\) e \(y=g(x)\), y entre las rectas
verticales \(x=a\) y \(x=b\), donde \(f\) y \(g\) son funciones
continuas tales que \(0 \leq g(x)\leq f(x)\) para toda \(x\) en
\([a,b]\).

\begin{center}
\includegraphics[width=0.55\linewidth,height=\textheight,keepaspectratio]{index_files/mediabag/imagenes_cap6/6.1-area-entre-curvas.pdf}
\end{center}

El área deseada se puede pensar como la diferencia entre el área debajo
del gráfico de \(f\) y la que está debajo del gráfico de \(g\).

\begin{center}
\includegraphics[width=0.7\linewidth,height=\textheight,keepaspectratio]{index_files/mediabag/imagenes_cap6/6.1-areas-separadas.pdf}
\end{center}

Entonces el área de la región es

\[
A=A_1-A_2=\int_a^b f(x)\,dx-\int_a^b g(x)\,dx=\int_a^b (f(x)-g(x))\,dx.
\]

Esta idea se puede extender fácilmente al caso más general, donde
sólamente sabemos que \(g \leq f\) en \([a, b]\). Podemos encontrar una
constante \(C\) adecuada de manera que los gráficos de \(y=g(x)+C\) y de
\(y=f(x)+C\) estén ambos por encima del eje \(x\), es decir, tal que
\(0\leq g(x)+C\leq f(x)+C\).

\begin{center}
\includegraphics[width=0.8\linewidth,height=\textheight,keepaspectratio]{index_files/mediabag/imagenes_cap6/6.1-areas-desplazadas.pdf}
\end{center}

Si llamamos \(F(x)=f(x)+C\) y \(G(x)=g(x)+C\), entonces en el intervalo
\([a,b]\) se cumple que \(0\leq G(x)\leq F(x)\). Aplicando la fórmula
anterior para el área \(A'\) comprendida entre los gráficos de \(F\) y
de \(G\) resulta

\[
A'=\int_a^b (F(x)-G(x))\,dx=\int_a^b (f(x)+C-(g(x)+C))\,dx=\int_a^b (f(x)-g(x))\,dx,
\]

y el área \(A'\) entre los gráficos de \(F\) y \(G\), para
\(a\leq x\leq b\), es la misma que el área \(A\) entre los gráficos de
\(f\) y \(g\) ya que al sumar la constante \(C\) sólo hemos desplazado
la región \(C\) unidades verticalmente hacia arriba.

De esta manera, tenemos la siguiente fórmula para el cálculo de áreas
determinadas por el gráfico de funciones continuas.

Si \(f\) y \(g\) son funciones continuas tales que \(g(x)\leq f(x)\)
para todo \(x\in [a,b]\), entonces el área \(A\) de la región limitada
por los gráficos de \(f\) y \(g\), y lateralmente por las rectas \(x=a\)
y \(x=b\) es

\[
A=\int_a^b (f(x)-g(x))\,dx.
\]

En particular, si \(g(x)=0\) en todo el intervalo \([a,b]\) y
\(f(x)\geq g(x)=0\) en el intervalo, entonces el área es

\[
A=\int_a^b (f(x)-g(x))\,dx=\int_a^b (f(x)-0)\,dx=\int_a^b f(x)\,dx.
\]

Es decir, como caso particular obtenemos que si \(f\) es continua y no
negativa en el intervalo \([a,b]\), la integral definida
\(\int_a^b f(x)\,dx\) nos da el área bajo la curva de \(y=f(x)\) y sobre
el eje \(x\), entre \(a\) y \(b\).

Ejemplo

Encontrar el área de la región delimitada por las parábolas \(y=x^2\) e
\(y=2x-x^2\).

\begin{tcolorbox}[enhanced jigsaw, opacityback=0, left=2mm, opacitybacktitle=0.6, toprule=.15mm, bottomrule=.15mm, colbacktitle=quarto-callout-tip-color!10!white, titlerule=0mm, bottomtitle=1mm, title=\textcolor{quarto-callout-tip-color}{\faLightbulb}\hspace{0.5em}{Solución}, toptitle=1mm, colframe=quarto-callout-tip-color-frame, colback=white, breakable, rightrule=.15mm, arc=.35mm, leftrule=.75mm, coltitle=black]

Primero buscamos los puntos de intersección entre las curvas. Igualando
las ordenadas obtenemos

\[
x^2=2x-x^2 \quad \text{ o también } \quad 2x(x-1)=0,
\]

de donde resulta \(x=0\) o \(x=1\). Entonces los puntos de corte entre
ambas son \((0,0)\) y \((1,1)\). Con ayuda del gráfico de las funciones
podemos ver que \(x^2\leq 2x-x^2\) en el intervalo \([0,1]\).

\includegraphics[width=0.4\linewidth,height=\textheight,keepaspectratio]{index_files/mediabag/imagenes_cap6/6.1-ejemplo-parabolas.pdf}

Por lo tanto, el área de la región es

\[ 
A=\int_0^1 \left[\left(2x-x^2\right)-x^2\right]\,dx =\int_0^1 \left(2x-2x^2\right)\,dx =\left(x^2-\frac{2}{3}x^3\right)\Bigg{]}_0^1=\frac{1}{3}.
\]

\end{tcolorbox}

\subsection{Área entre curvas (más
general)}\label{uxe1rea-entre-curvas-muxe1s-general}

Si tenemos que calcular el área entre dos curvas correspondientes a los
gráficos de dos funciones continuas \(f\) y \(g\), pero
\(f(x)\geq g(x)\) para algunos valores de \(x\) y \(f(x)\leq g(x)\) para
otros, entonces es preciso dividir el intervalo \([a,b]\) en
subintervalos de modo que en aquellos donde ocurra que \(f(x)\geq g(x)\)
integraremos la diferencia \(f(x)-g(x)\), y en los restantes, donde
\(f(x)\leq g(x)\), integraremos la expresión \(g(x)-f(x)\).

\begin{center}
\includegraphics[width=0.6\linewidth,height=\textheight,keepaspectratio]{index_files/mediabag/imagenes_cap6/6.1-funciones-cruzadas.pdf}
\end{center}

Por ejemplo, para calcular el área entre las curvas \(y=f(x)\) e
\(y=g(x)\) de la figura, observamos que \(f(x)\geq g(x)\) en \([a,c]\) y
en \([d,b]\), y \(g(x)\geq f(x)\) en \([c,d]\). Por lo tanto el área
\(A\) entre las curvas es

\[
A=\int_{a}^c(f(x)-g(x))\,dx+\int_{c}^d (g(x)-f(x))\,dx+\int_{d}^b (f(x)-g(x))\,dx.
\]

Recordando la definición de valor absoluto, tenemos que

\[
|f(x)-g(x)|=
\begin{cases}
f(x)-g(x) & \text{ si } f(x)\geq g(x),\\
g(x)-f(x) & \text{ si } f(x)< g(x).
\end{cases}
\]

En definitiva, podemos calcular el área entre las curvas de \(f\) y
\(g\) utilizando la siguiente fórmula.

\phantomsection\label{form-areageneral}
Si \(f\) y \(g\) son continuas en \([a,b]\), el área encerrada por las
curvas \(y=f(x)\) e \(y=g(x)\) es

\[
A=\int_a^b |f(x)-g(x)|\,dx.
\]

Ejemplo

Calcular el área encerrada por las curvas \(y=\operatorname{sen} x\),
\(y=\cos x\), \(x=0\) y \(x=\pi/2\).

\begin{tcolorbox}[enhanced jigsaw, opacityback=0, left=2mm, opacitybacktitle=0.6, toprule=.15mm, bottomrule=.15mm, colbacktitle=quarto-callout-tip-color!10!white, titlerule=0mm, bottomtitle=1mm, title=\textcolor{quarto-callout-tip-color}{\faLightbulb}\hspace{0.5em}{Solución}, toptitle=1mm, colframe=quarto-callout-tip-color-frame, colback=white, breakable, rightrule=.15mm, arc=.35mm, leftrule=.75mm, coltitle=black]

Primero buscamos los puntos de intersección entre ambas. Tenemos que

\[
\operatorname{sen} x=\cos x \quad \text{ equivale a } \quad \tan x=1,
\]

con lo cual \(x=\pi/4\). Graficando las funciones

\includegraphics[width=0.4\linewidth,height=\textheight,keepaspectratio]{index_files/mediabag/imagenes_cap6/6.1-ejemplo-seno-coseno.pdf}

podemos notar que \(\cos x\geq \operatorname{sen} x\) en \([0,\pi/4]\) y
\(\cos x\leq \operatorname{sen} x\) en \([\pi/4,\pi/2]\).

Utilizando la \hyperref[form-areageneral]{fórmula para el área entre
curvas} anterior, el área \(A\) es

\[
\begin{aligned}
A=\int_0^{\pi/2} |\operatorname{sen} x-\cos x|\,dx&=\int_0^{\pi/4} (\cos x-\operatorname{sen} x)\,dx+\int_{\pi/4}^{\pi/2} (\operatorname{sen} x-\cos x)\,dx\\
&=(\operatorname{sen} x+\cos x)\Big{]}_0^{\pi/4}+(-\cos x-\operatorname{sen} x)\Big{]}_{\pi/4}^{\pi/2}\\
&=2\sqrt{2}-2.
\end{aligned}
\]

\end{tcolorbox}

\subsection{\texorpdfstring{Áreas integrando respecto a
\(y\)}{Áreas integrando respecto a y}}\label{uxe1reas-integrando-respecto-a-y}

En algunos casos es conveniente expresar a \(x\) como una función de
\(y\).

\begin{center}
\includegraphics[width=0.4\linewidth,height=\textheight,keepaspectratio]{index_files/mediabag/imagenes_cap6/6.1-area-vertical.pdf}
\end{center}

Si la región está delimitada por los gráficos de las funciones continuas
\(x=f(y)\) y \(x=g(y)\) en un intervalo \([c,d]\), y se cumple que
\(f(y)\geq g(y)\) para todo \(y\) de este intervalo, entonces el área
\(A\) es

\phantomsection\label{form-area-respectoay}
\[
A=\int_c^d (f(y)-g(y))\,dy.
\]

Ejemplo

Hallar el área comprendida entre la recta \(y=x-1\) y la parábola
\(y^2=2x+6\).

\begin{tcolorbox}[enhanced jigsaw, opacityback=0, left=2mm, opacitybacktitle=0.6, toprule=.15mm, bottomrule=.15mm, colbacktitle=quarto-callout-tip-color!10!white, titlerule=0mm, bottomtitle=1mm, title=\textcolor{quarto-callout-tip-color}{\faLightbulb}\hspace{0.5em}{Solución}, toptitle=1mm, colframe=quarto-callout-tip-color-frame, colback=white, breakable, rightrule=.15mm, arc=.35mm, leftrule=.75mm, coltitle=black]

Buscamos primero los puntos de intersección. Despejando \(x\) en
términos de \(y\) en ambas funciones e igualando resulta

\[
y+1=\frac{y^2-6}{2}, \quad \text{ o también }\quad y^2-2y-8=0,
\]

de donde obtenemos \(y=-2\) e \(y=4\). Por lo tanto, los puntos de corte
son \((-1,-2)\) y \((5,4)\).

\includegraphics[width=0.5\linewidth,height=\textheight,keepaspectratio]{index_files/mediabag/imagenes_cap6/6.1-ejemplo-recta-parabola.pdf}

En el gráfico podemos notar que es más conveniente integrar con respecto
a \(y\), para \(-2\leq y\leq 4\), y este intervalo resulta
\(y+1\geq (y^2-6)/2\). Entonces el área es

\[
A=\int_{-2}^4 \left[(y+1)-\frac{1}{2}(y^2-6)\right]\,dy=\int_{-2}^4 \left(-\frac{1}{2}y^2+y+4\right)\,dy=\left(-\frac{y^3}{6}+\frac{y^2}{2}+4y\right)\Bigg{]}_{-2}^4=18.
\]

\end{tcolorbox}

\hyperref[seccion_6.1]{↑ Volver al inicio de la sección}

\section{6.2. Volúmenes}\label{seccion_6.2}

En esta sección veremos cómo utilizar integrales definidas para calcular
volúmenes de ciertos sólidos. Para empezar, consideremos una región
plana \(\Omega\).

\begin{center}
\includegraphics[width=0.6\linewidth,height=\textheight,keepaspectratio]{index_files/mediabag/imagenes_cap6/6.2-cilindros.pdf}
\end{center}

Un \textbf{cilindro recto con sección transversal \(\Omega\)} es un
sólido obtenido de trasladar la región \(\Omega\) a lo largo de un eje
perpendicular a ella.

Si \(A\) es el área de \(\Omega\) y el cilindro está formado por la
traslación de esta región a lo largo de una distancia \(h\), entonces el
volumen del cilindro es \(V=A\cdot h\).

Cuando \(\Omega\) es un círculo o un rectángulo, el sólido obtenido
resulta un cilindro circular recto y un prisma de base rectangular,
respectivamente.

\begin{center}
\includegraphics[width=0.6\linewidth,height=\textheight,keepaspectratio]{index_files/mediabag/imagenes_cap6/6.2-cilindro-prisma.pdf}
\end{center}

Si queremos calcular el volumen de un sólido \(S\) que no es un
cilindro, podemos comenzar cortando a \(S\) en rebanadas, cuyos
respectivos volúmenes puedan aproximarse por el volumen de un cilindro.

\begin{center}
\includegraphics[width=0.6\linewidth,height=\textheight,keepaspectratio]{index_files/mediabag/imagenes_cap6/6.2-volumen-rebanada.pdf}
\end{center}

\begin{center}
\includegraphics[width=0.6\linewidth,height=\textheight,keepaspectratio]{index_files/mediabag/imagenes_cap6/6.2-rebanadas.pdf}
\end{center}

Podemos hacer una estimación del volumen de \(S\) sumando los volúmenes
de los cilindros. El valor del volumen exacto de \(S\) se obtendrá
tomando límite cuando la cantidad de rebanadas en las que cortamos
tiende a infinito.

Por ejemplo, si quisiéremos aproximar el volumen de una esfera
tendríamos la siguiente situación

\begin{center}
\includegraphics[width=0.7\linewidth,height=\textheight,keepaspectratio]{index_files/mediabag/imagenes_cap6/6.2-vol-esfera-aprox.pdf}
\end{center}

Comenzamos cortando a \(S\) con un plano perpendicular y obteniendo una
región plana que se denomina \textbf{sección transversal} de \(S\). Si
\(A(x)\) denota el área de la sección transversal de \(S\) en un plano
perpendicular al eje \(x\) y que pasa por el punto \(x\), esta función
variará cuando \(x\) se incrementa desde \(a\) hasta \(b\).

\begin{center}
\includegraphics[width=0.6\linewidth,height=\textheight,keepaspectratio]{index_files/mediabag/imagenes_cap6/6.2-seccion-transversal.pdf}
\end{center}

Sea \(A(x)\) el área de la sección transversal de \(S\) en un plano
\(P_x\) perpendicular al eje \(x\) y que pasa por el punto \(x\), donde
\(a\leq x\leq b\).

Si dividimos el intervalo \([a,b]\) en en puntos con
\(\Delta x = \frac{b-a}{n}\) y \(x_i = a + i \Delta x\), con
\(0\leq i\leq n\), el volúmen de la rebanada entre los planos
\(P_{x_i}\) y \(P_{x_{i+1}}\) será aproximadamente el volúmen de un
cilindro con base \(A(x_{i})\) y altura \(\Delta x\).

Podemos decir que el volumen \(V\) es aproximadamente la suma \[
V\approx \sum_{i=0}^n A(x_{i})\,\Delta x
\]

y el valor exacto lo obtenemos tomando límite para \(n\to\infty\), es
decir,

\[
V=\lim_{n\to \infty} \sum_{i=0}^n A(x_{i})\,\Delta x=\int_a^b A(x)\,dx.
\]

En particular, cuando \(A(x)\) es una función continua el límite
anterior existe. Esto nos motiva a dar la siguiente definición.

\begin{tcolorbox}[enhanced jigsaw, opacityback=0, left=2mm, opacitybacktitle=0.6, toprule=.15mm, bottomrule=.15mm, colbacktitle=quarto-callout-note-color!10!white, titlerule=0mm, bottomtitle=1mm, title=\textcolor{quarto-callout-note-color}{\faInfo}\hspace{0.5em}{Definición (Volumen de un sólido)}, toptitle=1mm, colframe=quarto-callout-note-color-frame, colback=white, breakable, rightrule=.15mm, arc=.35mm, leftrule=.75mm, coltitle=black]

Sea \(S\) un sólido que se encuentra entre \(x=a\) y \(x=b\), y sea
\(A(x)\) el área de la sección transversal de \(S\) a través de \(x\) y
perpendicular al eje \(x\). Si \(A\) es continua en \([a,b]\) entonces
el volumen del sólido es

\[
V=\int_a^b A(x)\,dx.
\]

\end{tcolorbox}

Ejemplo

Hallar el volumen de una pirámide cuya altura es \(h\) y su base es un
cuadrado de lado \(L\).

\begin{center}
\includegraphics[width=0.7\linewidth,height=\textheight,keepaspectratio]{index_files/mediabag/imagenes_cap6/6.2-ejemplo-piramide.pdf}
\end{center}

\begin{tcolorbox}[enhanced jigsaw, opacityback=0, left=2mm, opacitybacktitle=0.6, toprule=.15mm, bottomrule=.15mm, colbacktitle=quarto-callout-tip-color!10!white, titlerule=0mm, bottomtitle=1mm, title=\textcolor{quarto-callout-tip-color}{\faLightbulb}\hspace{0.5em}{Solución}, toptitle=1mm, colframe=quarto-callout-tip-color-frame, colback=white, breakable, rightrule=.15mm, arc=.35mm, leftrule=.75mm, coltitle=black]

Ubicamos la pirámide de modo que su eje sea el eje de abscisas y su
cúspide esté ubicada sobre el origen de coordenadas, de modo que debemos
encontrar la función de área transversal \(A(x)\) para
\(0\leq x\leq h\). Observando la figura del enunciado, podemos decir que
\(A(x)=\ell^2\), con lo cual bastará expresar a \(\ell\) en términos de
\(x\). Para ello, por semejanza de triángulos podemos escribir

\[
\frac{h}{\frac{L}{2}}=\frac{x}{\ell}, \quad \text{ de donde }\quad \ell=\frac{Lx}{2h}.
\]

De esta manera

\[
A(x)=(2\ell)^2=\frac{L^2x^2}{h^2}.
\]

Esta función es continua en el intervalo \([0,h]\), por lo que el
volumen \(V\) de la pirámide es

\[
V=\int_0^h \frac{L^2x^2}{h^2}\,dx =\frac{L^2}{h^2} \left(\frac{x^3}{3}\right)\Bigg{]}_0^h=\frac{1}{3}L^2h.
\]

\end{tcolorbox}

\subsection{Sólidos de revolución}\label{suxf3lidos-de-revoluciuxf3n}

Supongamos que \(f\) es no negativa y continua en \([a, b]\). Si giramos
el gráfico de \(f\) alrededor del eje \(x\) obtenemos una superficie que
determina un sólido que llamaremos \textbf{sólido de revolución}.

\begin{center}
\includegraphics[width=0.7\linewidth,height=\textheight,keepaspectratio]{index_files/mediabag/imagenes_cap6/6.2-solidos-revolucion.pdf}
\end{center}

En este caso, para cada \(x\) en \([a,b]\) la sección transversal es un
círculo de radio \(f(x)\). Entonces la función de área transversal es

\[
A(x) = \pi[f(x)]^2,
\]

para \(a \leq x \leq b\). Como \(f\) es continua, \(A\) resulta continua
en \([a, b]\) y por lo visto anteriormente el volumen lo obtenemos
integrando esta función

\[
V=\pi\int_a^b \left[f(x)\right]^2\,dx.
\]

Ejemplo

Demostrar que el volumen de una esfera de radio \(r\) es
\(\displaystyle V=\frac{4}{3}\pi r^3\).

\begin{tcolorbox}[enhanced jigsaw, opacityback=0, left=2mm, opacitybacktitle=0.6, toprule=.15mm, bottomrule=.15mm, colbacktitle=quarto-callout-tip-color!10!white, titlerule=0mm, bottomtitle=1mm, title=\textcolor{quarto-callout-tip-color}{\faLightbulb}\hspace{0.5em}{Solución}, toptitle=1mm, colframe=quarto-callout-tip-color-frame, colback=white, breakable, rightrule=.15mm, arc=.35mm, leftrule=.75mm, coltitle=black]

La esfera es un sólido de revolución, obtenido a partir de rotar una
circunferencia de radio \(r\). Consideremos entonces la circunferencia
con centro en el origen y radio \(r\), dada por la ecuación
\(x^2+y^2=r^2\).

\includegraphics[width=0.4\linewidth,height=\textheight,keepaspectratio]{index_files/mediabag/imagenes_cap6/6.2-circunferencia.pdf}

Pero como esta curva no es el gráfico de una función, despejamos \(y\) y
nos quedamos con la semicircunferencia superior, cuya ecuación es
\(y=\sqrt{r^2-x^2}\), para \(-r\leq x\leq r\).

\includegraphics[width=0.4\linewidth,height=\textheight,keepaspectratio]{index_files/mediabag/imagenes_cap6/6.2-semicircunferencia.pdf}

La función de área transversal en este caso es

\[
A(x)=\pi (y(x))^2=\pi \left(\sqrt{r^2-x^2}\right)^2=\pi (r^2-x^2).
\]

Y el volumen \(V\) es

\[
V=\pi \int_{-r}^r (r^2-x^2)\,dx=\pi \left(r^2\,x-\frac{x^3}{3}\right)\Bigg{]}_{-r}^r=\pi\left[\left(r^3-\frac{r^3}{3}\right)-\left(-r^3+\frac{r^3}{3}\right)\right]=\frac{4}{3}\pi r^3.
\]

\end{tcolorbox}

Ejemplo

Determinar el volumen del sólido que se obtiene al girar la región bajo
la curva \(y=\sqrt{x}\) con respecto al eje \(x\), desde \(x=0\) hasta
\(x=1\).

\begin{tcolorbox}[enhanced jigsaw, opacityback=0, left=2mm, opacitybacktitle=0.6, toprule=.15mm, bottomrule=.15mm, colbacktitle=quarto-callout-tip-color!10!white, titlerule=0mm, bottomtitle=1mm, title=\textcolor{quarto-callout-tip-color}{\faLightbulb}\hspace{0.5em}{Solución}, toptitle=1mm, colframe=quarto-callout-tip-color-frame, colback=white, breakable, rightrule=.15mm, arc=.35mm, leftrule=.75mm, coltitle=black]

La función de área transversal es \[
A(x)=\pi(\sqrt{x})^2=\pi x.
\]

\includegraphics[width=0.35\linewidth,height=\textheight,keepaspectratio]{index_files/mediabag/imagenes_cap6/6.2-solidos-raiz.pdf}

Entonces el volumen resulta

\[
V=\pi \int_0^1 x\,dx =\pi \left(\frac{x^2}{2}\right)\Bigg{]}_0^1=\frac{\pi}{2}.
\]

\end{tcolorbox}

Ejemplo

Calcular el volumen del sólido generado al rotar la región definida por
\(y=x^3\), \(y=8\) y \(x=0\) con respecto al eje \(y\).

\begin{tcolorbox}[enhanced jigsaw, opacityback=0, left=2mm, opacitybacktitle=0.6, toprule=.15mm, bottomrule=.15mm, colbacktitle=quarto-callout-tip-color!10!white, titlerule=0mm, bottomtitle=1mm, title=\textcolor{quarto-callout-tip-color}{\faLightbulb}\hspace{0.5em}{Solución}, toptitle=1mm, colframe=quarto-callout-tip-color-frame, colback=white, breakable, rightrule=.15mm, arc=.35mm, leftrule=.75mm, coltitle=black]

En este caso rotamos alrededor del eje \(y\), con lo que debemos
expresar \(x\) en términos de \(y\).

\includegraphics[width=0.3\linewidth,height=\textheight,keepaspectratio]{index_files/mediabag/imagenes_cap6/6.2-solidos-cubica.pdf}

Como \(y=x^3\), resulta \(x=y^{1/3}\) y en consecuencia la función de
área transversal es

\[
A(y)=\pi (x(y))^2=\pi y^{2/3}.
\]

El volumen del sólido es

\[
V=\pi \int_0^8 y^{2/3}\,dy= \pi \left(\frac{3}{5}y^{5/3}\right)\Bigg{]}_0^8 =\frac{96}{5}\pi.
\]

\end{tcolorbox}

Ejemplo

Hallar el volumen del sólido generado al rotar la región delimitada por
\(y=x\) e \(y=x^2\), con \(0\leq x\leq 1\), alrededor de la recta
\(x=-1\).

\begin{tcolorbox}[enhanced jigsaw, opacityback=0, left=2mm, opacitybacktitle=0.6, toprule=.15mm, bottomrule=.15mm, colbacktitle=quarto-callout-tip-color!10!white, titlerule=0mm, bottomtitle=1mm, title=\textcolor{quarto-callout-tip-color}{\faLightbulb}\hspace{0.5em}{Solución}, toptitle=1mm, colframe=quarto-callout-tip-color-frame, colback=white, breakable, rightrule=.15mm, arc=.35mm, leftrule=.75mm, coltitle=black]

Como la rotación es con respecto a un eje vertical, debemos expresar las
funciones en términos de la variable \(y\). Notemos además que en este
caso se trata de una región delimitada por dos funciones, por lo tanto
las secciones transversales en cada \(y\) corresponden a arandelas, con
radio interior \(r(y)\) y radio exterior \(R(y)\).

\includegraphics[width=0.4\linewidth,height=\textheight,keepaspectratio]{index_files/mediabag/imagenes_cap6/6.2-solidos-recta-parabola.pdf}

Las funciones dadas son \(x=y\) y \(x=\sqrt{y}\), y los radios los
calculamos tomando las respectivas distancias desde el eje de rotación a
cada curva, es decir

\[
R(y)=\sqrt{y}-(-1)=\sqrt{y}+1\quad \text{ y } \quad   r(y)=y-(-1)=y+1.
\]

De esta manera, para cada \(0\leq y\leq 1\)

\[
A(y)=\pi(R^2(y)-r^2(y))=\pi \left[(\sqrt{y}+1)^2-(y+1)^2\right]=\pi\left[-y^2-y+2y^{1/2}\right].
\]

El volumen del sólido es

\[
V=\int_0^1 A(y)\,dy =\pi \int_0^1 \left(-y^2-y+2y^{1/2}\right)\,dy =\pi \left(-\frac{y^3}{3}-\frac{y^2}{2}+\frac{4}{3}y^{3/2}\right)\Bigg{]}_0^1=\frac{\pi}{2}.
\]

\end{tcolorbox}

\hyperref[seccion_6.2]{↑ Volver al inicio de la sección}

\section{6.4. Trabajo}\label{seccion_6.4}

En esta sección veremos cómo aplicar integrales a determinadas
situaciones físicas. Consideremos un objeto de masa \(m\) que se
desplaza en línea recta con función de posición \(s(t)\), por lo tanto
la fuerza \(F\) sobre el objeto está definida por la segunda ley de
Newton como

\[
F(t) = \text{masa} \times \text{aceleración} = m\frac{d^2}{dt^2}s.
\]

Recordemos unidades: la masa se mide en kilogramos (kg), el
desplazamiento en me- tros (m), el tiempo en segundos (seg) y la fuerza
en newtons ( \(\text{N} = \text{kg}\cdot \text{m} /\text{s}^2\) ). Es
decir, una fuerza de 1 N que actúa en una masa de 1 kg produce una
aceleración de 1 \(\text{m}/\text{seg}^2\).

Si la aceleración del objeto es constante, la fuerza también, y el
\textbf{trabajo} realizado se define como el producto de la fuerza por
la distacia \(d\) recorrida por el objeto, esto es \[
W = F \cdot d.
\]

Si medimos la fuerza en newtons y la distancia en metros, la unidad de
trabajo es un newton-metro, llamado joule (J).

Ejemplo

¿Qué tanto trabajo se realiza al levantar un libro de 1.2 kg desde el
suelo y colocarlo en un escritorio que tiene 0.7 m de altura?

\begin{tcolorbox}[enhanced jigsaw, opacityback=0, left=2mm, opacitybacktitle=0.6, toprule=.15mm, bottomrule=.15mm, colbacktitle=quarto-callout-tip-color!10!white, titlerule=0mm, bottomtitle=1mm, title=\textcolor{quarto-callout-tip-color}{\faLightbulb}\hspace{0.5em}{Solución}, toptitle=1mm, colframe=quarto-callout-tip-color-frame, colback=white, breakable, rightrule=.15mm, arc=.35mm, leftrule=.75mm, coltitle=black]

Considerando que la aceleración de la gravedad es \(9.8\)
\(\text{m}/\text{seg}^2\), y que la fuerza ejercida es igual y opuesta a
la que ejerce la gravedad, tenemos que la fuerza es \[
F = 1.2\, \text{kg} \cdot 9.8 \,\text{m}/\text{seg}^2 = 11.76\, \text{N}
\] por lo tanto, el trabajo realizado es \[
W = 11.76\, \text{N} \cdot 0.7 \, \text{m} = 8.232\, \text{J}.
\]

\end{tcolorbox}

Lo que vimos describe un fenómeno en el que la fuerza es constante. ¿Qué
sucede si la fuerza es variable?

Supongamos que el objeto se desplaza a lo largo del eje \(x\) en la
dirección positiva, desde \(x=a\) hasta \(x=b\), y en cada punto \(x\)
entre \(a\) y \(b\) actúa sobre el objeto una fuerza \(f(x)\), donde
\(f\) es una función continua.

Dividamos el intervalo \([a,b]\) como siempre, en \(n\) subintervalitos
\([x_{i-1},x_{i}]\) de longitud \(\Delta x\). Si el intervalo
\([x_{i-1},x_{i}]\) es suficientemente pequeño, como \(f\) es continua,
puede considerarse casi constante en el mismo. Por lo tanto, el trabajo
realizado será aproximadamente \[
W_i \approx f(x_i)\Delta x
\]

Es claro que el trabajo total \(W\) puede considerarse como la suma de
todos los trabajos realizados en cada subintervalo \([x_{i-1},x_{i}]\),
es decir \[
W = \sum_{i=1}^{n} W_i \approx \sum_{i=1}^{n} f(x_i)\Delta x.
\]

El valor exacto del trabajo se obtiene tomando el límite de las sumas de
arriba cuando \(n\to\infty\). Cuando la función \(f\) es continua, este
límite existe y nos da la integral definida de la función de fuerza en
el intervalo dado.

\begin{tcolorbox}[enhanced jigsaw, opacityback=0, left=2mm, opacitybacktitle=0.6, toprule=.15mm, bottomrule=.15mm, colbacktitle=quarto-callout-note-color!10!white, titlerule=0mm, bottomtitle=1mm, title=\textcolor{quarto-callout-note-color}{\faInfo}\hspace{0.5em}{Definición (Trabajo)}, toptitle=1mm, colframe=quarto-callout-note-color-frame, colback=white, breakable, rightrule=.15mm, arc=.35mm, leftrule=.75mm, coltitle=black]

El \textbf{trabajo} realizado por un objeto que se mueve en un segmento
representado por el intervalo \([a,b]\) y donde se ejerce una fuerza
\(f(x)\) para cada \(x\) del intervalo, es la integral \[
W = \int_a^b f(x)dx.
\]

\end{tcolorbox}

Ejemplo

Cuando una partícula se ubica a \(x\) metros del origen, una fuerza de
\(x^2+2x\) newtons actúa sobre ella. ¿Cuánto trabajo se efectúa al
moverla desde \(x=1\) hasta \(x=3\)?

\begin{tcolorbox}[enhanced jigsaw, opacityback=0, left=2mm, opacitybacktitle=0.6, toprule=.15mm, bottomrule=.15mm, colbacktitle=quarto-callout-tip-color!10!white, titlerule=0mm, bottomtitle=1mm, title=\textcolor{quarto-callout-tip-color}{\faLightbulb}\hspace{0.5em}{Solución}, toptitle=1mm, colframe=quarto-callout-tip-color-frame, colback=white, breakable, rightrule=.15mm, arc=.35mm, leftrule=.75mm, coltitle=black]

Calculamos la integral

\[
W = \int_1^3 (x^2+2x)\,dx = \left(\frac{x^3}{3} + x^2\right)\Bigg{]}_1^3 = \frac{50}{3}.
\]

El trabajo realizado es de \(\displaystyle \frac{50}{3}\) joules.

\end{tcolorbox}

Ejemplo

Una fuerza de 40 N se requiere para detener un resorte que está estirado
desde su longitud natural de 10 cm a una longitud de 15 cm. ¿Cuánto
trabajo se hace al estirar el resorte de 15 a 18 cm?

\begin{center}
\includegraphics[width=0.4\linewidth,height=\textheight,keepaspectratio]{index_files/mediabag/imagenes_cap6/6.4-resortes.pdf}
\end{center}

\begin{tcolorbox}[enhanced jigsaw, opacityback=0, left=2mm, opacitybacktitle=0.6, toprule=.15mm, bottomrule=.15mm, colbacktitle=quarto-callout-tip-color!10!white, titlerule=0mm, bottomtitle=1mm, title=\textcolor{quarto-callout-tip-color}{\faLightbulb}\hspace{0.5em}{Solución}, toptitle=1mm, colframe=quarto-callout-tip-color-frame, colback=white, breakable, rightrule=.15mm, arc=.35mm, leftrule=.75mm, coltitle=black]

De acuerdo con la \textbf{ley de Hooke}, la fuerza que se requiere para
mantener el resorte estirado \(x\) metros más allá de su longitud
natural es \(f(x)=kx\), donde \(k\) es la constante asociada al resorte.

Cuando el resorte pasa de los 10 cm a los 15 cm, la longitud de
estiramiento es 5 cm=0.05 m, y como la fuerza requerida para producir
este estiramiento es de 40 N, tenemos que \(f(0.05)=40\), de donde
podemos deducir que \[
0.05 k=40 \quad \text{ o bien } \quad k=800.
\]

Calculamos ahora el trabajo

\[
W=\int_{0.05}^{0.08} 800 x\,dx=800\frac{x^2}{2}\Big{]}_{0.05}^{0.08}=1.56.
\]

El trabajo realizado para estirar el resorte de los 15 cm a los 18 cm es
de 1.56 J.

\end{tcolorbox}

\hyperref[seccion_6.4]{↑ Volver al inicio de la sección}

\section{6.5. Valor promedio de una
función}\label{valor-promedio-de-una-funciuxf3n}

Cuando hablamos de \textbf{promedio} de cierta cantidad de números en la
vida cotidiana, entendemos que sumamos todos los valores y dividimos por
la cantidad de elementos que sumamos: comúnmente esto es lo que hacemos
al calcular el promedio de calificaciones en una carrera o el promedio
de edades en un grupo de personas, entre otros.

Si tenemos \(n\) valores \(y_1, y_2, \dots, y_n\), el valor promedio es
\[
\overline{y}=\frac{y_1 + y_2 + \dots + y_n}{n}=\frac{\sum_{i=1}^n y_i}{n}=\frac{1}{n}\sum_{i=1}^n y_i.
\]

Pero ¿qué pasa cuando hablamos de calcular el promedio sobre una
cantidad infinita de elementos? Esto sucedería si intentásemos promediar
los valores que toma una función continua \(f\) en todos los puntos
\(x\) de un intervalo \([a,b]\). Situaciones como estas podrían ocurrir
en la práctica si, por ejemplo, queremos conocer la temperatura promedio
de un proceso o la velocidad promedio durante un viaje, eventos en los
que contaríamos con una cantidad infinita de mediciones dadas por el
gráfico de una función adecuada.

Si tenemos una función \(f\) definida en un intervalo \([a,b]\), ¿cómo
podríamos darle sentido a esto? Para ello, dividimos el intervalo
\([a,b]\) tomando puntos \(x_i = a + i \Delta x\), con \(0\leq i\leq n\)
y \(\Delta x = \frac{b-a}{n}\) y promediamos de la forma usual los
números \(f(x_i)\), ya que son finitos. Es decir, calculamos

\[
\frac{\sum_{i=1}^n f(x_i)}{n},
\]

que también podemos escribir de la siguiente manera

\[
\frac{\sum_{i=1}^n f(x_i)}{n} = \frac{\sum_{i=1}^n f(x_i)}{\frac{b-a}{\Delta x}} = \frac{1}{b-a}\sum_{i=1}^n f(x_i) \Delta x.
\]

Conforme \(n\) crece, estas sumas nos van a dar el promedio de una
cantidad muy grande de puntos, correspondientes a valores ubicados sobre
el gráfico de \(f\) y muy próximos entre sí. Cuando \(f\) es continua,
el límite de estas expresiones para \(n\to \infty\) existe. Más aún,
tenemos que

\[
\lim_{n\to \infty} \frac{1}{b-a}\sum_{i=1}^n f(x_i)\Delta x=\frac{1}{b-a}\int_a^b f(x)\,dx,
\]

lo que nos conduce a la siguiente definición.

\begin{tcolorbox}[enhanced jigsaw, opacityback=0, left=2mm, opacitybacktitle=0.6, toprule=.15mm, bottomrule=.15mm, colbacktitle=quarto-callout-note-color!10!white, titlerule=0mm, bottomtitle=1mm, title=\textcolor{quarto-callout-note-color}{\faInfo}\hspace{0.5em}{Definición (Valor promedio de una función en un intervalo)}, toptitle=1mm, colframe=quarto-callout-note-color-frame, colback=white, breakable, rightrule=.15mm, arc=.35mm, leftrule=.75mm, coltitle=black]

Si \(f\) es continua en \([a,b]\) entonces su \textbf{valor promedio} o
directamente su \textbf{promedio} en este intervalo es

\[
\overline{f}_{[a,b]}=\frac{1}{b-a}\int_a^b f(x)\,dx.
\]

\end{tcolorbox}

Ejemplo

Determinar el valor promedio de \(f(x)=1+x^2\) en el intervalo
\([-1,2]\).

\begin{tcolorbox}[enhanced jigsaw, opacityback=0, left=2mm, opacitybacktitle=0.6, toprule=.15mm, bottomrule=.15mm, colbacktitle=quarto-callout-tip-color!10!white, titlerule=0mm, bottomtitle=1mm, title=\textcolor{quarto-callout-tip-color}{\faLightbulb}\hspace{0.5em}{Solución}, toptitle=1mm, colframe=quarto-callout-tip-color-frame, colback=white, breakable, rightrule=.15mm, arc=.35mm, leftrule=.75mm, coltitle=black]

En virtud de la definición anterior tenemos que

\[
\overline{f}_{[-1,2]}=\frac{1}{2-(-1)}\int_{-1}^2 (1+x^2)\,dx=\frac{1}{3}\left(x+\frac{x^3}{3}\right)\Bigg{]}_{-1}^2=2.
\]

\end{tcolorbox}

En el ejemplo anterior podemos ver que el promedio de \(f\) nos dio
\(2\), y \(f(1)=2\), es decir, existe un valor de \(x\) en el intervalo
dado donde la imagen de \(x\) es exactamente su valor promedio sobre el
intervalo. Esto ocurre siempre cuando \(f\) es continua, como lo
establece el siguiente resultado.

\phantomsection\label{teo-valor-medio-integral}
Teorema del valor medio para integrales

Si \(f\) es continua en \([a,b]\), entonces existe un número \(c\) en
\([a,b]\) tal que \(f(c)=\overline{f}_[a,b]\), es decir, de modo que se
verifica la igualdad

\[
f(c) = \frac{1}{b-a}\int_a^b f(x)dx.
\]

Una interpretación geométrica del teorema del valor medio para
integrales es que, para funciones positivas \(f\), hay un número \(c\)
tal que el rectángulo con base \([a,b]\) y altura \(f(c)\) tiene la
misma área que la región bajo la gráfica de \(f\).

\begin{center}
\includegraphics[width=0.4\linewidth,height=\textheight,keepaspectratio]{index_files/mediabag/imagenes_cap6/6.5-teorema-valor-medio.pdf}
\end{center}

Ejemplo

Encontrar todos los puntos \(c\) sobre el intervalo \([-1,2]\) tales que
\(\overline{f}_{[-1,2]}=f(c)\), siendo \(f(x)=1+x^2\).

\begin{tcolorbox}[enhanced jigsaw, opacityback=0, left=2mm, opacitybacktitle=0.6, toprule=.15mm, bottomrule=.15mm, colbacktitle=quarto-callout-tip-color!10!white, titlerule=0mm, bottomtitle=1mm, title=\textcolor{quarto-callout-tip-color}{\faLightbulb}\hspace{0.5em}{Solución}, toptitle=1mm, colframe=quarto-callout-tip-color-frame, colback=white, breakable, rightrule=.15mm, arc=.35mm, leftrule=.75mm, coltitle=black]

Ya hemos encontrado \(c=1\) por inspección. De todos modos, procederemos
de forma general. El \hyperref[teo-valor-medio-integral]{teorema del
valor medio integral} nos asegura que existe al menos un valor \(c\) en
\([-1,2]\) ya que \(f\) es continua sobre este intervalo. En el ejemplo
anterior hemos calculado que

\[
\overline{f}_{[a,b]}=2,
\]

por lo que ahora buscamos números \(c\) tales que \(f(c)=2\), es decir

\[
1+c^2=2, \quad \text{ con lo cual } \quad c^2=1,
\]

de donde deducimos que \(c=\pm 1\). Entonces los valores de \(c\) dentro
del intervalo donde \(f\) toma su valor promedio son \(c=-1\) y \(c=1\).

\end{tcolorbox}

\hyperref[seccion_6.5]{↑ Volver al inicio de la sección}

\section{7.8. Integrales impropias}\label{seccion_7.8}

Hasta ahora hemos trabajado con integrales definidas en las que aparece
una función \(f\) en un intervalo cerrado \([a,b]\) y supusimos que
\(f\) es continua a trozos sin discontinuidades infinitas.

Para ciertos modelos es necesario considerar el caso de intervalos
infinitos o funciones que tengan singularidades, esto es, que tiendan a
infinito en algún punto del intervalo.

Una integral se llama \textbf{impropia} cuando el intervalo de
integración es infinito o cuando el integrando tiene una singularidad
(no está acotado) en el intervalo de integración. A continuación
distinguiremos entre estos dos tipos de comportamiento.

\subsection{Tipo I: intervalos
infinitos}\label{tipo-i-intervalos-infinitos}

Consideremos la región \(S\) bajo la curva
\(y=\displaystyle\frac{1}{x^2}\), por encima del eje \(x\) y a la
derecha de la recta \(x=1\).

\begin{center}
\includegraphics[width=0.4\linewidth,height=\textheight,keepaspectratio]{index_files/mediabag/imagenes_cap6/7.8-region-S.pdf}
\end{center}

Para \(t>1\) podemos calcular la integral \[
\int_1^t \frac{dx}{x^2} = \left(-\frac{1}{x}\right)\Bigg{]}_1^t = 1 - \frac{1}{t},
\]

que nos da el área de la subregión \(S_t\) que se ecuentra a la
izquierda de \(x=t\).

\begin{center}
\includegraphics[width=0.4\linewidth,height=\textheight,keepaspectratio]{index_files/mediabag/imagenes_cap6/7.8-region-St.pdf}
\end{center}

Observermos que \[
\lim_{t\to\infty} \left(1 - \frac{1}{t}\right) = 1
\]

y con esto podemos decir que el área de \(S\) es igual a 1, ya que las
regiones \(S_t\) tienen un área que converge a 1 cuando \(t\to \infty\).
Esto motiva a la siguiente definición.

\begin{tcolorbox}[enhanced jigsaw, opacityback=0, left=2mm, opacitybacktitle=0.6, toprule=.15mm, bottomrule=.15mm, colbacktitle=quarto-callout-note-color!10!white, titlerule=0mm, bottomtitle=1mm, title=\textcolor{quarto-callout-note-color}{\faInfo}\hspace{0.5em}{Definición (Integral impropia de tipo I)}, toptitle=1mm, colframe=quarto-callout-note-color-frame, colback=white, breakable, rightrule=.15mm, arc=.35mm, leftrule=.75mm, coltitle=black]

\begin{itemize}
\item
  Si la integral \(\displaystyle \int_a^{t} f(x)\,dx\) existe para todo
  número \(t\geq a\), definimos \[
  \int_a^{\infty} f(x)\,dx = \lim_{t\to\infty} \int_a^{t} f(x)\,dx.
  \]
\item
  Si la integral \(\displaystyle \int_{t}^{b} f(x)\,dx\) existe para
  todo número \(t\leq b\), definimos \[
  \int_{-\infty}^b f(x)\,dx = \lim_{t\to-\infty} \int_{t}^b f(x)\,dx.
  \]
\end{itemize}

Las integrales impropias se llaman \textbf{convergentes} si el límite
correspondiente existe y \textbf{divergentes} si el límite no existe.

\begin{itemize}
\tightlist
\item
  Si tanto \(\displaystyle \int_{-\infty}^a f(x)\,dx\) como
  \(\displaystyle \int_a^{-\infty} f(x)\,dx\) existen para algún número
  \(a\in\mathbb{R}\), definimos \[
  \int_{-\infty}^{\infty} f(x)\,dx = \int_{-\infty}^a f(x)\,dx + \int_a^{\infty} f(x)\,dx.
  \]
\end{itemize}

\end{tcolorbox}

Estas integrales impropias pueden interpretarse como áreas cuando la
función \(f\) es no negativa. Por ejemplo, para el primer caso de la
definición anterior, si \(f(x)\geq 0\) y la integral

\[
\int_a^\infty f(x)\,dx
\]

converge, entonces define el área de la región \(S\) del plano dada por

\[
S=\{(x,y)\in \mathbb{R}^2: 0\leq y\leq f(x), x\geq a\}.
\]

\begin{center}
\includegraphics[width=0.6\linewidth,height=\textheight,keepaspectratio]{index_files/mediabag/imagenes_cap6/7.8-region-S-general.pdf}
\end{center}

Ejemplo

Determinar si la integral
\(\displaystyle \int_{1}^{\infty} \frac{1}{x}\,dx\) es convergente o
divergente.

\begin{tcolorbox}[enhanced jigsaw, opacityback=0, left=2mm, opacitybacktitle=0.6, toprule=.15mm, bottomrule=.15mm, colbacktitle=quarto-callout-tip-color!10!white, titlerule=0mm, bottomtitle=1mm, title=\textcolor{quarto-callout-tip-color}{\faLightbulb}\hspace{0.5em}{Solución}, toptitle=1mm, colframe=quarto-callout-tip-color-frame, colback=white, breakable, rightrule=.15mm, arc=.35mm, leftrule=.75mm, coltitle=black]

Según la definición, \[
\begin{aligned}
\int_{1}^{\infty} \frac{1}{x}\,dx 
&= \lim_{t\to\infty} \int_{1}^{t} \frac{1}{x}\,dx = \lim_{t\to\infty} \left(\ln|x|\right)\bigg{]}_{1}^{t} \\
\\
&= \lim_{t\to\infty} (\ln t - \ln 1) = \lim_{t\to\infty} \ln t = \infty.
\end{aligned}
\]

Como el límite no existe por infinitud, la integral impropia
\(\displaystyle \int_{1}^{\infty} \frac{1}{x}\,dx\) es divergente.

\end{tcolorbox}

Ejemplo

Evaluar la integral
\(\displaystyle \int_{-\infty}^{\infty} \frac{1}{1+x^2}\,dx\).

\begin{tcolorbox}[enhanced jigsaw, opacityback=0, left=2mm, opacitybacktitle=0.6, toprule=.15mm, bottomrule=.15mm, colbacktitle=quarto-callout-tip-color!10!white, titlerule=0mm, bottomtitle=1mm, title=\textcolor{quarto-callout-tip-color}{\faLightbulb}\hspace{0.5em}{Solución}, toptitle=1mm, colframe=quarto-callout-tip-color-frame, colback=white, breakable, rightrule=.15mm, arc=.35mm, leftrule=.75mm, coltitle=black]

Comencemos estudiando las integrales
\(\displaystyle\int_a^\infty f(x)\,dx\) y
\(\displaystyle\int_{-\infty}^a f(x)\,dx\), para algún valor fijo de
\(a\).

Si elegimos \(a=0\), para \(t>0\) tenemos que

\[
\int_0^t \frac{1}{1+x^2}\,dx=(\arctan x)\big{]}_0^t=\arctan t-\arctan 0 =\arctan t.
\]

Por lo tanto

\[
\int_0^\infty f(x)\, dx=\lim_{t\to \infty} \int_0^t \frac{1}{1+x^2}\,dx =\lim_{t\to \infty} \arctan t =\frac{\pi}{2}.
\]

De manera similar, para \(t<0\) podemos escribir

\[
\int_t^0 \frac{1}{1+x^2}\,dx=(\arctan x)\big{]}_t^0=\arctan 0-\arctan t =-\arctan t.
\]

Con lo cual

\[
\int_{-\infty}^0 f(x)\, dx=\lim_{t\to -\infty} \int_t^0 \frac{1}{1+x^2}\,dx =\lim_{t\to -\infty} -\arctan t =\frac{\pi}{2}.
\]

Como ambas integrales impropias son convergentes, concluimos que

\[
\int_{-\infty}^\infty \frac{1}{1+x^2}\,dx=\int_{-\infty}^0 \frac{1}{1+x^2}\,dx+\int_{0}^\infty \frac{1}{1+x^2}\,dx=\frac{\pi}{2}+\frac{\pi}{2}=\pi.
\]

\end{tcolorbox}

Ejemplo

¿Para qué valores de \(p\) es convergente la integral
\(\displaystyle \int_{1}^{\infty} \frac{1}{x^{p}}\,dx \, ?\)

\begin{tcolorbox}[enhanced jigsaw, opacityback=0, left=2mm, opacitybacktitle=0.6, toprule=.15mm, bottomrule=.15mm, colbacktitle=quarto-callout-tip-color!10!white, titlerule=0mm, bottomtitle=1mm, title=\textcolor{quarto-callout-tip-color}{\faLightbulb}\hspace{0.5em}{Solución}, toptitle=1mm, colframe=quarto-callout-tip-color-frame, colback=white, breakable, rightrule=.15mm, arc=.35mm, leftrule=.75mm, coltitle=black]

Cuando \(p=1\) la integral es divergente, según lo visto en el Ejemplo
14. Supongamos entonces que \(p \ne 1\).

Para \(t>1\) tenemos que

\[
\int_1^t \frac{1}{x^p}\,dx=\left(\frac{x^{-p+1}}{-p+1}\right)\Bigg{]}_1^t=\frac{1}{1-p}\left(t^{1-p}-1\right).
\]

Cuando \(p>1\) obtenemos

\[
\begin{aligned}
\int_{1}^{\infty} \frac{1}{x^{p}}\,dx
&= \lim_{t \to \infty} \frac{1}{p - 1}\,\left(1-\frac{1}{t^{p-1}}\right)\\
\\
&=\frac{1}{p-1},
\end{aligned}
\]

dado que \(\lim_{t\to \infty} \frac{1}{t^{p-1}}=0\) por ser \(p>1\).

Por otro lado, si \(p < 1\), entonces

\[
\begin{aligned}
\int_{1}^{\infty} \frac{1}{x^{p}}\,dx
&= \lim_{t \to \infty} \frac{1}{1 - p}\,\left(t^{1-p}-1\right)\\
\\
&=\infty,
\end{aligned}
\]

ya que \(t^{1-p}\to \infty\) por ser \(p<1\). En este caso, la integral
es divergente.

Concluimos entonces que la integral impropia es convergente sólamente
cuando \(p>1\), y en este caso

\[
\int_{1}^{\infty} \frac{1}{x^{p}}\,dx=\frac{1}{p-1}.
\]

\end{tcolorbox}

\subsection{Tipo II: integrandos
discontinuos}\label{tipo-ii-integrandos-discontinuos}

Supongamos ahora que \(f\) es una función continua y definida en un
intervalo finito \([a,b)\), pero tiene una asíntota vertical en \(x=b\).
Es claro que si \(a<t<b\), la integral \(\displaystyle \int_a^t f(x)dx\)
existe. ¿Qué significado podemos darle, entonces, a la integral
\(\displaystyle\int_a^b f(x)dx\)?

Para este tipo de situaciones será de utilidad la siguiente definición.

\begin{tcolorbox}[enhanced jigsaw, opacityback=0, left=2mm, opacitybacktitle=0.6, toprule=.15mm, bottomrule=.15mm, colbacktitle=quarto-callout-note-color!10!white, titlerule=0mm, bottomtitle=1mm, title=\textcolor{quarto-callout-note-color}{\faInfo}\hspace{0.5em}{Definición (Integral impropia de tipo II)}, toptitle=1mm, colframe=quarto-callout-note-color-frame, colback=white, breakable, rightrule=.15mm, arc=.35mm, leftrule=.75mm, coltitle=black]

\begin{itemize}
\tightlist
\item
  Si \(f\) es continua en \([a,b)\) y discontinua en \(b\), entonces\\
  \[
  \int_a^b f(x)\,dx \;=\; \lim_{t\to b^-}\int_a^t f(x)\,dx,
  \]\\
  si este límite existe.
\end{itemize}

\includegraphics[width=0.4\linewidth,height=\textheight,keepaspectratio]{index_files/mediabag/imagenes_cap6/7.8-discontinua-derecha.pdf}

\begin{itemize}
\tightlist
\item
  Si \(f\) es continua en \((a,b]\) y discontinua en \(a\), entonces\\
  \[
  \int_a^b f(x)\,dx \;=\; \lim_{t\to a^+}\int_t^b f(x)\,dx,
  \]\\
  siempre que este límite exista.
\end{itemize}

\includegraphics[width=0.4\linewidth,height=\textheight,keepaspectratio]{index_files/mediabag/imagenes_cap6/7.8-discontinua-izquierda.pdf}

\begin{itemize}
\tightlist
\item
  Si \(f\) tiene una discontinuidad en \(c\), donde \(a<c<b\), y tanto
  \(\displaystyle\int_a^c f(x)\,dx\) como
  \(\displaystyle\int_c^b f(x)\,dx\) son convergentes, definimos\\
  \[
  \int_a^b f(x)\,dx \;=\; \int_a^c f(x)\,dx \;+\; \int_c^b f(x)\,dx.
  \]
\end{itemize}

\includegraphics[width=0.4\linewidth,height=\textheight,keepaspectratio]{index_files/mediabag/imagenes_cap6/7.8-discontinua-medio.pdf}

Como en las integrales impropias de tipo I, diremos que la integral
impropia \(\displaystyle \int_a^b f(x)\,dx\) es \textbf{convergente}
cuando los límites correspondientes existen, y \textbf{divergente} en
caso contrario.

\end{tcolorbox}

Ejemplo

Determinar para qué número \(p\), la integral \[
\int_{0}^{1} \frac{1}{x^{p}}\,dx
\]

es convergente o divergente.

\begin{tcolorbox}[enhanced jigsaw, opacityback=0, left=2mm, opacitybacktitle=0.6, toprule=.15mm, bottomrule=.15mm, colbacktitle=quarto-callout-tip-color!10!white, titlerule=0mm, bottomtitle=1mm, title=\textcolor{quarto-callout-tip-color}{\faLightbulb}\hspace{0.5em}{Solución}, toptitle=1mm, colframe=quarto-callout-tip-color-frame, colback=white, breakable, rightrule=.15mm, arc=.35mm, leftrule=.75mm, coltitle=black]

Notemos primero que si \(p\leq 0\) entonces la función
\(\displaystyle f(x)=\frac{1}{x^p}\) es continua en \([0,1]\) y, por lo
tanto, integrable. Para estos valores de \(p\) no estamos ante una
integral impropia. Para \(p>0\) tenemos una integral impropia de tipo II
ya que \(f\) tiene una asíntota vertical en \(x=0\), el extremo
izquierdo del intervalo. Fijado \(p\neq 1\), considerando \(0<t<1\)
tenemos que

\[
\int_t^1 \frac{1}{x^p}\,dx=\left(\frac{x^{1-p}}{1-p}\right)\Bigg{]}_t^1=\frac{1}{1-p}\left(1-t^{1-p}\right).
\]

Luego, por la definición de este tipo de integrales impropias, si
\(0<p<1\) obtenemos

\[
\int_0^1 \frac{1}{x^p}\,dx=\lim_{t\to 0^+} \int_t^1 \frac{1}{x^p}\,dx=\lim_{t\to 0^+} \frac{1}{1-p}\left(1-t^{1-p}\right)=\frac{1}{1-p},
\] ya que \(t^{1-p}\to 0\) cuando \(t\to 0^+\) por ser \(1-p>0\).

Cuando \(p>1\) resulta

\[
\int_0^1 \frac{1}{x^p}\,dx=\lim_{t\to 0^+} \int_t^1 \frac{1}{x^p}\,dx=\lim_{t\to 0^+} \frac{1}{p-1}\left(\frac{1}{t^{p-1}}-1\right)=\infty,
\]

ya que \(\displaystyle \frac{1}{t^{p-1}}\to \infty\) cuando
\(t\to 0^+\), puesto que \(p-1>0\).

Finalmente, si \(p=1\) obtenemos

\[
\int_0^1 \frac{1}{x}\,dx=\lim_{t\to 0^+} \int_t^1 \frac{1}{x}\,dx=\lim_{t\to 0^+} (\ln |x|)\Big{]}_t^1=\lim_{t\to 0^+} (-\ln t)=\infty,
\]

y la integral es divergente en este caso también.

En definitiva, la integral dada es convergente cuando \(p<1\) y
divergente si \(p\geq 1\).

\end{tcolorbox}

Ejemplo

Decidir si la integral impropia
\(\displaystyle \int_2^5 \frac{1}{\sqrt{x-2}}\) converge o diverge.

\begin{tcolorbox}[enhanced jigsaw, opacityback=0, left=2mm, opacitybacktitle=0.6, toprule=.15mm, bottomrule=.15mm, colbacktitle=quarto-callout-tip-color!10!white, titlerule=0mm, bottomtitle=1mm, title=\textcolor{quarto-callout-tip-color}{\faLightbulb}\hspace{0.5em}{Solución}, toptitle=1mm, colframe=quarto-callout-tip-color-frame, colback=white, breakable, rightrule=.15mm, arc=.35mm, leftrule=.75mm, coltitle=black]

Notemos que se trata de una integral impropia de tipo II, ya que la
función tiene una asíntota vertical en \(x=2\). Si \(2<t<5\) entonces

\[
\int_t^5 \frac{1}{\sqrt{x-2}}\,dx=\int_{t-2}^3 \frac{1}{\sqrt{u}}\,du=\left(2\sqrt{u}\right)\Big{]}_{t-2}^3=2\sqrt{3}-2\sqrt{t-2},
\]

donde hemos hecho la sustitución \(u=x-2\).

Luego

\[
\int_2^5 \frac{1}{\sqrt{x-2}}\,dx=\lim_{t\to 2^+} \int_t^5 \frac{1}{\sqrt{x-2}}\,dx=\lim_{t\to 2^+} \left(2\sqrt{3}-2\sqrt{t-2}\right)=2\sqrt{3},
\]

por lo que la integral dada es convergente.

\end{tcolorbox}

Ejemplo

Evaluar, si es posible, la integral impropia \[
\int_{0}^{3}\frac{1}{x-1}\,dx.
\]

\begin{tcolorbox}[enhanced jigsaw, opacityback=0, left=2mm, opacitybacktitle=0.6, toprule=.15mm, bottomrule=.15mm, colbacktitle=quarto-callout-tip-color!10!white, titlerule=0mm, bottomtitle=1mm, title=\textcolor{quarto-callout-tip-color}{\faLightbulb}\hspace{0.5em}{Solución}, toptitle=1mm, colframe=quarto-callout-tip-color-frame, colback=white, breakable, rightrule=.15mm, arc=.35mm, leftrule=.75mm, coltitle=black]

La recta \(x=1\) es una asíntota vertical del integrando y está dentro
del intervalo \([0,3]\). Por la definición de integral impropia con
discontinuidad en \(x=1\), debemos dividir la integral en dos partes. La
integral impropia será convergente si ambas integrales impropias

\[
\int_{0}^{1}\frac{1}{x-1}\,dx \quad \text{ y }\quad \int_{1}^{3}\frac{1}{x-1}\,dx
\]

son convergentes. Comencemos analizando la primera. Para \(0<t<1\),
haciendo el cambio de variable \(u=x-1\) tenemos que

\[
\int_{0}^{t}\frac{1}{x-1}\,dx=\int_{-1}^{t-1}\frac{1}{u}\,du=\left(\ln |u|\right)\Big{]}_{-1}^{t-1}=\ln (1-t)-\ln 1=\ln(1-t).
\]

Entonces

\[
\int_{0}^{1}\frac{1}{x-1}\,dx=\lim_{t\to 0^+}\int_{0}^{t}\frac{1}{x-1}\,dx=\lim_{t\to 1^-} \ln (1-t)=\lim_{z\to 0^+}\ln z=-\infty,
\]

con lo cual la primera de las integrales es divergente. Con esto, ya
podems concluir que la integral dada es divergente, sin necesidad de
evaluar la segunda integral impropia.

\end{tcolorbox}

\begin{tcolorbox}[enhanced jigsaw, opacityback=0, left=2mm, opacitybacktitle=0.6, toprule=.15mm, bottomrule=.15mm, colbacktitle=quarto-callout-caution-color!10!white, titlerule=0mm, bottomtitle=1mm, title=\textcolor{quarto-callout-caution-color}{\faFire}\hspace{0.5em}{Importante}, toptitle=1mm, colframe=quarto-callout-caution-color-frame, colback=white, breakable, rightrule=.15mm, arc=.35mm, leftrule=.75mm, coltitle=black]

Si en el ejemplo anterior no hubiéramos notado la asíntota vertical en
\(x=1\) y, en cambio, hubiéramos confundido la integral con una integral
definida ordinaria, podríamos haber hecho un cálculo erróneo usando la
regla de Barrow \[
\int_{0}^{3}\frac{1}{x^{2}-1}\,dx
\stackrel{\text{(mal)}}{=}
\ln|x-1|\big{]}_{0}^{3} = \ln 2 - \ln 1 = \ln 2,
\] obteniendo así un número finito. Este procedimiento es incorrecto,
pues no se cumplen las hipótesis del teorema fundamental del cálculo ya
que \(f\) tiene una discontinuidad infinita y por lo tanto no puede
nunca ser continua a trozos en \([0,3]\).

De aquí en adelante, cada vez que aparezca el símbolo
\(\displaystyle \int_{a}^{b} f(x)\,dx\) debemos decidir, observando el
comportamiento de \(f\) en \([a,b]\), si se trata de una integral
definida ordinaria o de una integral impropia que requiere límites.

\end{tcolorbox}

Ejemplo

Evaluar \(\displaystyle \int_{0}^{1}\ln x\,dx\).

\begin{tcolorbox}[enhanced jigsaw, opacityback=0, left=2mm, opacitybacktitle=0.6, toprule=.15mm, bottomrule=.15mm, colbacktitle=quarto-callout-tip-color!10!white, titlerule=0mm, bottomtitle=1mm, title=\textcolor{quarto-callout-tip-color}{\faLightbulb}\hspace{0.5em}{Solución}, toptitle=1mm, colframe=quarto-callout-tip-color-frame, colback=white, breakable, rightrule=.15mm, arc=.35mm, leftrule=.75mm, coltitle=black]

La función \(y=\ln x\) tiene una asíntota vertical en \(x=0\), con lo
cual se trata de una integral impropia de tipo II. Por lo tanto

\[
\int_0^1 \ln x\,dx =\lim_{t\to 0^+} \int_t^1 \ln x\,dx.
\]

En el Ejemplo 20 de la \hyperref[seccion_7.1]{Sección 7.1} obtuvimos que

\[
\int \ln x\,dx=x(\ln x-1)+C,
\]

con lo cual

\[
\int_0^1 \ln x\,dx =\lim_{t\to 0^+} \left(x(\ln x -1)\right)\Big{]}_{t}^1=\lim_{t\to 0^+} [1(\ln 1-1)-t(\ln t-1)]=\lim_{t\to 0^+} [-1-t\ln t+t]=-1.
\]

Entonces la integral impropia es convergente.

\end{tcolorbox}

\subsection{Prueba de comparación para integrales
impropias}\label{prueba-de-comparaciuxf3n-para-integrales-impropias}

Muchas veces en la práctica nos encontramos con integrales impropias muy
difíciles de evaluar con los métodos aprendidos. Sin embargo, si
sólamente nos interesa conocer si son convergentes o divergentes, es
posible compararlas con otras más sencillas. Enunciamos el criterio de
comparación en el siguiente teorema.

\phantomsection\label{teo-comparacion}
Teorema de comparación para integrales impropias

Supongamos que \(f\) y \(g\) son funciones continuas en \([a,\infty)\)
para algún \(a\in\mathbb{R}\) y que \(f(x)\geq g(x)\geq 0\) para todo
\(x\geq a\).

\begin{itemize}
\item
  Si \(\displaystyle \int_a^\infty f(x)\,dx\) es convergente, entonces
  \(\displaystyle \int_a^\infty g(x)\,dx\) también.
\item
  Si \(\displaystyle \int_a^\infty g(x)\,dx\) es divergente, entonces
  \(\displaystyle \int_a^\infty f(x)\,dx\) también.
\end{itemize}

Ejemplo

Mostrar que la integral impropia
\(\displaystyle \int_{0}^{\infty}e^{-x^2}\,dx\) es convergente.

\begin{tcolorbox}[enhanced jigsaw, opacityback=0, left=2mm, opacitybacktitle=0.6, toprule=.15mm, bottomrule=.15mm, colbacktitle=quarto-callout-tip-color!10!white, titlerule=0mm, bottomtitle=1mm, title=\textcolor{quarto-callout-tip-color}{\faLightbulb}\hspace{0.5em}{Solución}, toptitle=1mm, colframe=quarto-callout-tip-color-frame, colback=white, breakable, rightrule=.15mm, arc=.35mm, leftrule=.75mm, coltitle=black]

No podemos calcular la integral de manera directa, ya que
\(g(x)=e^{-x^2}\) no tiene una antiderivada expresable en término de
funciones elementales. Sin embargo, cuando \(x\geq 1\) tenemos que
\(x\leq x^2\), o equivalentemente \(-x^2\leq -x\). Esto implica que
\(e^{-x^2}\leq e^{-x}\) en el intervalo \(I=[1,\infty)\). Si llamamos
\(f(x)=e^{-x}\), entonces tenemos que \(f(x)\geq g(x)\geq 0\) en \(I\).

Ahora bien,

\[
\int_1^\infty f(x)\,dx=\lim_{t\to \infty} \int_1^t e^{-x}\,dx =\lim_{t\to \infty} \left(-e^{-x}\right)\Big{]}_1^t=\lim_{t\to \infty} \left(-e^{-t}+1\right)=1
\]

con lo que la integral impropia \(\displaystyle \int_1^\infty f(x)\,dx\)
es convergente. Utilizando el \hyperref[teo-comparacion]{teorema de
comparación para integrales impropias} concluimos que
\(\displaystyle \int_1^{\infty} g(x)\,dx\) es convergente.

Finalmente

\[
\int_0^\infty e^{-x^2}\,dx=\int_0^1 e^{-x^2}\,dx+\int_1^\infty e^{-x^2}\,dx
\]

es convergente por ser la suma de una integral definida de \(g\) en
\([0,1]\) y una integral impropia convergente.

\end{tcolorbox}

Ejemplo

Mostrar que la integral impropia
\(\displaystyle \int_{1}^{\infty}\frac{1+e^{-x}}{x}\,dx\) es divergente.

\begin{tcolorbox}[enhanced jigsaw, opacityback=0, left=2mm, opacitybacktitle=0.6, toprule=.15mm, bottomrule=.15mm, colbacktitle=quarto-callout-tip-color!10!white, titlerule=0mm, bottomtitle=1mm, title=\textcolor{quarto-callout-tip-color}{\faLightbulb}\hspace{0.5em}{Solución}, toptitle=1mm, colframe=quarto-callout-tip-color-frame, colback=white, breakable, rightrule=.15mm, arc=.35mm, leftrule=.75mm, coltitle=black]

Nuevamente estamos ante un caso donde el cálculo directo de la integral
es complicado. Pero para conocer sólo su carácter de convergente o
divergente basta con comparar esta integral con otra más sencilla. De
hecho, observemos que para cualquier \(x\geq 1\)

\[
f(x)=\frac{1+e^{-x}}{x}\geq \frac{1}{x}=g(x),
\]

ya que \(e^{-x}\)\textgreater0 siempre.

Ahora bien, hemos visto que

\[
\int_1^\infty \frac{1}{x}\,dx=\infty,
\]

es decir, la integral impropia \(\displaystyle \int_1^\infty g(x)\,dx\)
es divergente. Además hemos visto que \(0\leq g(x)\leq f(x)\) para todo
\(x\in[1,\infty)\). Por el \hyperref[teo-comparacion]{teorema de
comparación para integrales impropias} concluimos que la integral

\[
\int_{1}^{\infty}\frac{1+e^{-x}}{x}\,dx
\]

es divergente.

\end{tcolorbox}

\hyperref[seccion_7.8]{↑ Volver al inicio de la sección}


\backmatter


\end{document}
